% !TeX root = ../Book3.tex
% 纲要标题：# 第二编　计算、维数、平坦性与完备化
% 纲要位置：工作记录/计划纲要.md，第 2034 行。


\BookPart{计算、维数、平坦性与完备化}{b3:part:02}
\BookPartGuide
第一编已经建立理想、局部环及整扩张的基本关系。现在要进一步回答几个可以实际检验的问题：怎样计算理想，几个方程能限制多少维，局部代数随基底改变时是否保持结构，以及有限阶近似能恢复哪些信息。本编第8—14章分别发展这些问题所需的计算、维数、增长、平坦性、完备化和微分方法。

\ProseParagraphBreak
\chapref{b3:ch:08}从单项式序和多元除法开始，在本卷完整证明 Gröbner 基的基础理论，再讨论消去、理想运算、参数例外和零维商代数。\secref{b3:sec:0805}用乘法矩阵求出相容的坐标并解释重数，\secref{b3:sec:0806}以结式提供另一种消去方法。学习时应保存每次约化的多项式等式，尤其区分“同一理想”“同一零点集”和“相同的投影闭包”。

\ProseParagraphBreak
\chapref{b3:ch:09}用素理想链定义维数，证明主理想定理及广义高度定理，再建立有限型代数的维数、参数系和嵌入维数。\chapref{b3:ch:10}把这些结论用于一维正规整环，证明 DVR 的刻画和 Dedekind 整环的理想分解，形成一条研究一维局部结构的算术与曲线支线。赋值记录局部零点阶，理想类群则记录局部主理想未必能由一个全局元素生成的现象。第8、9、11—14章更连续地联系理想计算、维数增长、基变换、完备化与微分；学习这些内容不以第10章的全部算术例子为前提。

\ProseParagraphBreak
\chapref{b3:ch:11}从分次片的维数或长度建立 Hilbert 级数及多项式，最终用 Hilbert–Samuel 增长\footnote{萨缪尔（Pierre Samuel，1921-2009），法国数学家，研究交换代数及其在代数几何中的应用。Hilbert 已在本卷前文介绍。}刻画局部维数。计数中必须始终区分单层长度与累计长度：正维情形下，它们的多项式次数相差一。证明先在原环上完成，不借用后面完备化保持维数的结论。

\ProseParagraphBreak
\chapref{b3:ch:12}在本卷重建所需的平坦判据，然后讨论有限展示平坦模、局部自由性、Fitting 理想\footnote{菲廷（Hans Fitting，1906-1938），德国数学家，研究群论，并以行列式理想描述模的展示。}、泛自由性和忠实平坦性。Fitting 理想记录局部最少生成元数的变化，还能检测某些纤维维数看不到的幂零关系；不能仅因纤维维数恒定就判定平坦。末尾的\secref{b3:sec:1206}说明带余循环相容条件的有限展示模下降。

\ProseParagraphBreak
\chapref{b3:ch:13}通过 Rees 环\footnote{里斯（David Rees，1918-2013），威尔士数学家，研究半群与交换代数。}和 Artin–Rees 引理处理滤过，再构造完备化。第11.4节在证明中临时把滤过各层合成分次对象，用有限生成性得到诱导滤过的粗估计；这里才系统命名并研究 Rees 环，证明精确的 Artin–Rees 等式，并解释两种滤过拓扑的关系。随后按正合性、有限生成模的张量表示、Noether 性、平坦性及维数保持依次展开。形式幂级数来自所有阶的相容数据，一个有限截断通常不能代替它。\chapref{b3:ch:14}则从导子的乘法法则构造 Kähler 微分，以基本正合列计算余切空间与 Jacobian 矩阵\footnote{雅可比（Carl Gustav Jacob Jacobi，1804-1851），德国数学家，研究椭圆函数、行列式和力学。}，并将可分性与微分联系起来，区分形式光滑、形式非分歧、形式平展，以及光滑、非分歧、平展的条件。

\begin{center}
\begin{tikzcd}[column sep=large,row sep=large]
\begin{gathered}\text{第8章}\\\text{Gröbner 基与计算}\end{gathered}
&\begin{gathered}\text{第9章}\\\text{维数与高度}\end{gathered}\arrow[r]\arrow[d]
&\begin{gathered}\text{第10章}\\\text{DVR 与 Dedekind 环}\end{gathered}\\
&\begin{gathered}\text{第11章}\\\text{Hilbert 增长}\end{gathered}\arrow[d]
&\begin{gathered}\text{第12章}\\\text{平坦性与基变换}\end{gathered}\arrow[dl]\arrow[d]\\
&\begin{gathered}\text{第13章}\\\text{滤过与完备化}\end{gathered}
&\begin{gathered}\text{第14章}\\\text{微分与切空间}\end{gathered}
\end{tikzcd}
\end{center}

图中只画本编关键证明的直接依赖。第8章的计算方法还可帮助理解第11章的单项式计数和第12章泛自由性的首项论证，后二者所需的分次或模论版本会在使用处说明。第14章前半可在第一编后单独阅读；讨论平展时才需要第12章的平坦性。不必先完成第10章的全部算术例子再学习增长或微分。

\ProseParagraphBreak
本编默认读者已掌握环、模、张量积和有限维线性代数的基本运算；关键的局部判据、首项计算和有限关系论证均在本卷给出。维数、赋值、Hilbert 理论及完备化可配合松村英之的《Commutative Ring Theory》\cite[Chapters 3--5]{Matsumura1989CommutativeRingTheory}研读，各章还列出相应定理的出处。下一编将研究局部环的深度、正则序列和同调刻画，并在使用处引入复形与消解。

\ProseParagraphBreak
读过本编，应能为理想计算保留可核验的等式，用高度和 Hilbert 增长判断方程留下的维数，并选择平坦判据、完备化或微分来检验局部结构。对于随参数变化的方程，还应分清纤维的点数、维数与幂零关系，判断哪些性质在基变换中保持，并区别单个有限截断与全部相容近似的逆极限提供的信息。

% !TeX root = ../../Book3.tex
% 纲要标题：## 第8章　Gröbner 基与消去
% 纲要位置：工作记录/计划纲要.md，第 2036 行。
% 纲要细目（写作范围）：
% > 第一卷附录F提供先行入门。本章按本卷的独立阅读需要完整建立单项式序、多元除法、Dickson有限性、Gröbner基与正规式、Buchberger判据及算法、约化基唯一性和消去定理，不以跨卷回引代替基础证明；继而发展理想运算、几何消去与零维计算。高效实现与参数化算法转附录B，局部／形式标准基转附录F；这些进阶附录不构成后续结构理论的未说明前提。



\chapter{Gröbner 基与消去}
\label{b3:ch:08}
\BookChapterNumber{b3:ch:08}

\begin{chapterabstract}
本章以单项式序和多元除法建立 Gröbner 基、正规式及 Buchberger 算法，证明理想成员与消去关系的可计算判据。理想交、商和饱和由同一工具处理；零维商代数的标准单项式、乘法矩阵与结式则服务于多项式方程求解。
\end{chapterabstract}

\begin{learninggoals}
\begin{itemize}
\item 正确选择单项式序，完成多元除法，并解释一般余式与正规式的区别。
\item 写出最高项抵消的有界表示，证明并执行 Buchberger 算法。
\item 通过扩展变量计算消去、交、理想商和饱和，同时保留成员表示；在不变环算例中分别证明生成性和关系的完备性。
\item 用标准单项式和乘法矩阵计算有限维商代数，辨认参数例外及非约化结构。
\item 从 Sylvester 矩阵计算结式与判别式，并在参数特殊化后核查候选解。
\end{itemize}
\end{learninggoals}

设 \(k\) 为域，\(I=(x^2-y,y^2-1)\subset k[x,y]\)。由
\[
x^4-1=(x^2+y)(x^2-y)+(y^2-1)
\]
可立即核验 \(x^4-1\in I\cap k[x]\)。这是一份成员证书，却还不能只凭它断言 \(I\cap k[x]\) 恰由 \(x^4-1\) 生成：也许存在次数更低、只含 \(x\) 的关系。消去不仅要找到一个关系，还要证明没有遗漏。

\ProseParagraphBreak

要系统回答这个问题，必须为多项式的首项规定次序，使除法和余式可重复执行；然后再挑出足够好的理想生成元，让余式为零真正等价于理想成员资格。这样的生成元就是 Gröbner 基。Buchberger 的过程\footnote{布赫贝格尔（Bruno Buchberger，1942-），奥地利数学家，提出 Gröbner 基算法。}既负责补足原生成元遗漏的首项，也让计算停止的理由可以核验。消去定理把首项控制转化为投影、交与商的计算方法。

\ProseParagraphBreak

当理想的商代数为有限维时，计算又有了另一种组织方式：标准单项式的剩余类给出向量空间基，乘以各坐标则成为一组交换矩阵。多项式之间的关系因而可以用线性相关性检验，扩至代数闭域后，方程的根也能从这些矩阵的共同特征信息恢复。\secref{b3:sec:0805}沿着这项联系讨论换序与单变量表示，\secref{b3:sec:0806}再用结式消去变量。计算中出现的系数分母及参数例外，则在\secref{b3:sec:0804}用精确等式逐项控制。

\ProseParagraphBreak
本章需要多项式、理想、商环及有限维线性代数；投影的几何解释还使用本卷的零点定理。单项式有限性和 Gröbner 基的基础证明都在这里展开。手算例子同时保留当前多项式及其由原生成元给出的表示，使计算结果附有可核验的成员证书。

\ProseParagraphBreak
零维部分可配合 Sturmfels 的《Solving Systems of Polynomial Equations》\cite[Sections 2.1--2.3, pp. 13--19]{Sturmfels2002PolynomialEquations}阅读；其中以标准单项式、局部重数及乘法矩阵组织计算。本章先建立这些计算所需的代数依据，附录B再讨论高效算法与实现，附录F另行处理局部标准基。


% !TeX root = ../../Book3.tex
% 纲要标题：### 8.1 全局单项式序与多元除法
% 纲要位置：工作记录/计划纲要.md，第 2141 行。
% 纲要细目（写作范围）：
% * 定义并验证乘法相容的良序及三种常用项序，以指数比较说明两项公理的作用
% * 建立首项、首项理想和标准单项式，完整证明多元除法及首项界；结合计算说明选序的影响
% * 说明余式依赖性，完整证明单项式成员判据、Dickson有限性、Gröbner基存在及正规式唯一性
% * 本章的除法和终止证明依赖全局良序；局部序与混合序转附录 F 简介，不作为后续正文所需的基础

\section{全局单项式序与多元除法}
\label{b3:sec:0801}

理想的有限生成性只保证存在有限描述，并不立即给出成员判定的方法。一元多项式的带余除法提示我们先处理最高项；多元情形没有唯一的“最高次数项”，必须先选择一个与乘法相容的次序。以下固定域 \(k\)、正整数 \(n\) 和 \(R=k[x_1,\ldots,x_n]\)。代数结论适用于任意域；称某过程为可执行算法时，还要求系数域支持精确四则运算及零判定。
\subsection{字典序、分次字典序与分次逆字典序}
\label{b3:subsec:0801:01}

对 \(\alpha=(\alpha_1,\ldots,\alpha_n)\in\mathbb N^n\)，记
\[
x^\alpha=x_1^{\alpha_1}\cdots x_n^{\alpha_n},
\qquad |\alpha|=\alpha_1+\cdots+\alpha_n.
\]
这里 \(\mathbb N\) 含零。\(x^\alpha\) 称为单项式，\(|\alpha|\) 为其全次数；非零系数与单项式的乘积称为项。单项式整除等价于各坐标指数分别不大于。

\begin{definition}\label{b3:def:global-monomial-order}
设 \(k\) 为域，\(R=k[x_1,\ldots,x_n]\)，其中 \(1\le n<\infty\)。给定 \(R\) 的单项式集合上的全序 \(\prec\)。如果每个非空单项式子集都有最小元，且对 \(\alpha,\beta,\gamma\in\mathbb N^n\) 有
\[
x^\alpha\prec x^\beta
\quad\Longrightarrow\quad
x^{\alpha+\gamma}\prec x^{\beta+\gamma}
\quad(\gamma\in\mathbb N^n),
\]
则称 \(\prec\) 为\term[danxiangshi-xu]{单项式序}[monomial order]，本章也称全局单项式序。良序性与乘法相容性是两个分别需要核验的条件。
\end{definition}

良序性用于排除约化时首单项式的无限下降，乘法相容性则保证相乘以后仍可控制首项；下面首单项式的乘法公式正依赖后一条件。

单位单项式 \(1\) 必为最小元。否则若 \(m\prec1\)，相乘给出 \(1\succ m\succ m^2\succ\cdots\)，违反良序。反之，一个与乘法相容且以一为最小元的全序，在有限个变量的交换单项式上也必为良序；本节稍后用 Dickson 引理\footnote{迪克森（Leonard Eugene Dickson，1874-1954），美国数学家，研究群论与数论，著有《数论史》。}说明这个结论。

\ProseParagraphBreak
固定变量优先顺序 \(x_1\succ\cdots\succ x_n\)，三种常用次序如下。
\begin{enumerate}
\item \term[zidian-xu]{字典序}[lexicographic order]：在 \(\alpha-\beta\) 的第一个非零坐标处，正数对应较大的单项式。
\item \term[fenci-zidian-xu]{分次字典序}[graded lexicographic order]：先比较全次数，全次数相同时再按字典序比较。
\item \term[fenci-ni-zidian-xu]{分次逆字典序}[graded reverse lexicographic order]：先比较全次数；相同时，在 \(\alpha-\beta\) 的最后一个非零坐标处，负数对应较大的单项式。
\end{enumerate}
字典序的非空指数集有最小元：依次在剩余指数中最小化第一、第二直至第 \(n\) 个坐标。后两种次序先选最小全次数，该次数下只有有限个单项式，再取最小者。给两指数同时加上 \(\gamma\) 不改变差向量，也不改变全次数之差，所以三种次序都与乘法相容。

\ProseParagraphBreak
例如按变量顺序 \(x\succ y\succ z\)，\(x^2z\) 与 \(xy^2\) 的指数分别为 \((2,0,1)\) 与 \((1,2,0)\)，全次数同为三。两指数之差为 \((1,-2,1)\)：字典序看第一个非零坐标，故 \(x^2z\succ xy^2\)；分次字典序因全次数相同，也得到这个结果。分次逆字典序看最后一个非零坐标，\(xy^2\) 的 \(z\) 指数较小，反而有 \(xy^2\succ x^2z\)。字典序还把 \(x\) 排在每个 \(y^N\) 前面，而与 \(N\) 多大无关。所谓“逆”不是把一个良序整体倒过来，而只是在全次数相等时改变最后一个不同指数的比较方向。后文的消去需要特定次序，不能把这三种次序互换而不检查结论。

\subsection{首项、标准单项式与多元除法}
\label{b3:subsec:0801:02}

\begin{definition}\label{b3:def:leading-data}
设 \(k\) 为域，\(R=k[x_1,\ldots,x_n]\)，并固定单项式序 \(\prec\)。将 \(0\ne f\in R\) 按此序写成有限个项之和。最大的单项式记为 \(\operatorname{LM}(f)\)，其系数记为 \(\operatorname{LC}(f)\)，二者的乘积记为 \(\operatorname{LT}(f)\)，分别称为\term[shou-danxiangshi]{首单项式}[leading monomial]、\term[shouxiang-xishu]{首项系数}[leading coefficient]和\term[shouxiang]{首项}[leading term]。零多项式不定义首项。
\end{definition}
\symidx{\(\operatorname{LM}(f),\operatorname{LC}(f),\operatorname{LT}(f)\)：首单项式、首项系数与首项}

若 \(f,g\ne0\)，乘法相容性保证
\[
\operatorname{LM}(fg)=\operatorname{LM}(f)\operatorname{LM}(g),
\qquad
\operatorname{LC}(fg)=\operatorname{LC}(f)\operatorname{LC}(g).
\]
因为除了两个最高项的乘积，其余乘积至少有一因子较小，故严格小于最高乘积，不可能抵消它。这里用到了系数域中两个非零数的乘积仍非零。

多元除法允许用几个除数依次消去首项。可能有不止一个 \(\operatorname{LM}(g_i)\) 整除当前首单项式，选择哪个先做便可能影响余式，因此以下将除数列为有序组。一元单除数除法的余式唯一性，并不自动延续到这样的生成组。

\begin{theorem}[多元除法]\label{b3:thm:multivariate-division}
设 \(k\) 为域，\(R=k[x_1,\ldots,x_n]\)，并固定全局单项式序 \(\prec\)。给定 \(R\) 中的有序组 \(G=(g_1,\ldots,g_s)\)，其中各 \(g_i\ne0\)，每个 \(f\in R\) 都能在有限步内写成
\begin{equation}\label{b3:eq:multivariate-division}
f=\sum_{i=1}^s q_i g_i+r,
\end{equation}
其中 \(r\) 的每个单项式都不被任何 \(\operatorname{LM}(g_i)\) 整除。若 \(f\ne0\)，所有非零乘积 \(q_i g_i\) 以及非零余式 \(r\) 的首单项式均不大于 \(\operatorname{LM}(f)\)。
\end{theorem}
\begin{proof}
令 \(h=f\)，各 \(q_i=0\)，\(r=0\)。若 \(h\ne0\)，查找第一个满足 \(\operatorname{LM}(g_i)\mid\operatorname{LM}(h)\) 的下标。若存在，置
\[
u=\frac{\operatorname{LT}(h)}{\operatorname{LT}(g_i)},
\qquad q_i\longleftarrow q_i+u,\qquad h\longleftarrow h-ug_i.
\]
这里 \(u\) 是一个项，系数除法在域 \(k\) 中进行。若不存在，就把 \(\operatorname{LT}(h)\) 加入 \(r\)，再从 \(h\) 中减去它。每步均保持恒等式
\[
f=\sum_iq_ig_i+r+h.
\]
新的非零待处理式的首单项式严格下降，良序性排除无限执行。这里下降的是指定单项式序下的首单项式，全次数却可能增加：在字典序 \(x\succ y\) 下，用 \(x-y^2\) 约化 \(x\)，剩下 \(y^2\)，其全次数从一增至二，但 \(y^2\prec x\)。最终 \(h=0\)，得到除法等式。移入余式的项不能被任何 \(\operatorname{LM}(g_i)\) 整除，且后来的项更小，不会改变这一性质。

每步加入的乘积 \(ug_i\) 的首单项式等于当时 \(\operatorname{LM}(h)\)，因而不超过初始 \(\operatorname{LM}(f)\)。对同一个 \(i\) 把这些乘子相加，也不会产生更大的项；移入余式的项同样受这个界限制。这证明最后的首项界。\(f=0\) 时全部输出为零；空除数组时逐项移入余式，输出 \(r=f\)。
\end{proof}

\begin{algorithm}[htbp]
\caption{有序多元除法}
\label{b3:alg:ordered-multivariate-division}
\begin{algorithmsteps}{可精确计算且可判零的域 \(k\)、可有效比较的全局单项式序，以及 \(f\in k[x_1,\ldots,x_n]\) 和有限有序非零多项式组 \(G=(g_1,\ldots,g_s)\)。}{多项式 \(q_1,\ldots,q_s,r\)，满足 \(f=\sum_iq_ig_i+r\)，且 \(r\) 的每个单项式都不被任何 \(\operatorname{LM}(g_i)\) 整除。}
\State 置 \(h\gets f\)、\(r\gets0\)，并对每个 \(i\) 置 \(q_i\gets0\)。
\While{\(h\ne0\)}
\State 查找最小的 \(i\)，使 \(\operatorname{LM}(g_i)\mid\operatorname{LM}(h)\)。
\If{找到这样的 \(i\)}
\State 置 \(u\gets\operatorname{LT}(h)/\operatorname{LT}(g_i)\)，\(q_i\gets q_i+u\)，\(h\gets h-ug_i\)。
\Else
\State 置 \(r\gets r+\operatorname{LT}(h)\)，\(h\gets h-\operatorname{LT}(h)\)。
\EndIf
\EndWhile
\State \Return \((q_1,\ldots,q_s,r)\)。
\end{algorithmsteps}
\end{algorithm}

\cref{b3:alg:ordered-multivariate-division}同时产生余式和一个理想成员的线性组合表示。但是，非零余式目前只说明这次除法没有化为零，尚不能断言 \(f\notin(G)\)。要使这种逆向判断成立，必须对生成组增加条件。

\begin{definition}\label{b3:def:groebner-basis}
设 \(k\) 为域，\(R=k[x_1,\ldots,x_n]\)，并固定全局单项式序 \(\prec\)。理想 \(I\subseteq R\) 的\term[shouxiang-lixiang]{首项理想}[initial ideal]为
\[
\operatorname{in}_{\prec}(I)
=\bigl(\operatorname{LM}(f)\mid 0\ne f\in I\bigr).
\]
给定有限组 \(G\subseteq I\setminus\{0\}\)。如果
\[
\operatorname{in}_{\prec}(I)
=\bigl(\operatorname{LM}(g)\mid g\in G\bigr),
\]
则称 \(G\) 为 \(I\) 的 \term[Groebner-ji]{Gröbner 基}[Gröbner basis]。不属于 \(\operatorname{in}_{\prec}(I)\) 的单项式称为\term[biaozhun-danxiangshi]{标准单项式}[standard monomial]。
\end{definition}
\symidx{\(\operatorname{in}_{\prec}(I)\)：理想关于指定单项式序的首项理想}

首项理想是一个单项式理想，它把原理想中可能出现的首单项式转化为指数的整除问题。Gröbner 基的条件要求，由全部理想成员的首单项式生成的这个理想，已经由有限个 \(\operatorname{LM}(g)\) 生成。首项理想也可由首项生成，因为每个非零首系数在域中可逆。定义先要求 \(G\subseteq I\)，尚未把“生成 \(I\)”作为额外公理；下一小节证明这项生成性自动成立。零理想以空组为 Gröbner 基，单位理想以 \(\{1\}\) 为 Gröbner 基。

\ProseParagraphBreak
单项式整除在指数平面上有直接的形状：\(x^ay^b\) 整除 \(x^iy^j\)，当且仅当 \(i\ge a,j\ge b\)。例如 \(J=(x^2,xy,y^3)\) 覆盖三个生成指数右上方的格点；图中的灰点属于 \(J\)，蓝点是剩下的标准单项式。
\begin{center}
\begin{tikzpicture}[scale=0.85]
\fill[gray!15] (1.75,-0.2) rectangle (4.2,4.2);
\fill[gray!15] (0.75,0.75) rectangle (1.75,4.2);
\fill[gray!15] (-0.2,2.75) rectangle (0.75,4.2);
\draw[->] (-0.3,0)--(4.6,0) node[right] {\(i\)};
\draw[->] (0,-0.3)--(0,4.6) node[above] {\(j\)};
\foreach \i in {0,...,4} {\foreach \j in {0,...,4} {\fill[gray] (\i,\j) circle (1.5pt);}}
\foreach \i/\j in {0/0,1/0,0/1,0/2} {\fill[blue!70!black] (\i,\j) circle (2.5pt);}
\draw[thick] (-0.2,2.5)--(0.5,2.5)--(0.5,0.5)--(1.5,0.5)--(1.5,-0.2);
\node[below left] at (0,0) {\(1\)};
\node[below] at (1,0) {\(x\)};
\node[left] at (0,1) {\(y\)};
\node[left] at (0,2) {\(y^2\)};
\node[below] at (2,0) {\(x^2\)};
\node[above right] at (1,1) {\(xy\)};
\node[left] at (0,3) {\(y^3\)};
\node at (3.2,3.3) {属于 \(J\)};
\end{tikzpicture}
\end{center}
因此 \(k[x,y]/J\) 的基为 \(1,x,y,y^2\)。向右上方封闭对应乘以单项式。

\subsection{余式的顺序依赖与 Dickson 引理}
\label{b3:subsec:0801:03}

在 \(k[x,y]\) 中取字典序 \(x\succ y\)，令
\[
g_1=xy-1,\qquad g_2=y^2-1,\qquad f=xy^2.
\]
先用 \(g_1\) 得到 \(f=yg_1+y\)，余式为 \(y\)；先用 \(g_2\) 得到 \(f=xg_2+x\)，余式为 \(x\)。两个余式的项都不能被 \(xy,y^2\) 整除，却不相等。进一步，\(x-y=yg_1-xg_2\) 属于原理想，按这组生成元除法却留下非零余式。问题在生成组，而不在除法恒等式。

图中的阶梯拐点对应单项式理想的极小生成元：一个拐点的指数不能逐坐标大于等于另一个拐点，否则它已被后者覆盖。Dickson 引理说明，有限维指数格子中不可能需要无限多个这样互不支配的拐点。

\begin{lemma}[Dickson]\label{b3:lem:dickson}
设 \(n\) 为有限正整数。对任意 \(E\subseteq\mathbb N^n\)，存在有限子集 \(E_0\subseteq E\)，使每个 \(\alpha\in E\) 都逐坐标不小于某个 \(\beta\in E_0\)。因此，对任意域 \(k\)，环 \(k[x_1,\ldots,x_n]\) 的每个单项式理想都由有限个单项式生成，单项式理想的升链必稳定。
\end{lemma}
\begin{proof}
先说明每个无限自然数序列都含无限非下降子序列。如果某个值出现无限次，就取常值子序列；否则每个有界区间中只有有限项，可以逐次在尾部选出更大的值。对 \(\mathbb N^n\) 中的无限序列先按第一坐标抽取，再按第二坐标抽取，直至第 \(n\) 个坐标，便得到所有坐标均非下降的无限子序列。特别地，原序列中存在 \(i<j\)，使第 \(i\) 项逐坐标不大于第 \(j\) 项。

若不存在所需有限 \(E_0\)，先选 \(\alpha^{(1)}\in E\)，然后总可选取一个不逐坐标大于等于任何已选指数的新指数。这样得到无限序列，其中没有上一段所说的一对，矛盾。\(E=\varnothing\) 时取空集即可。

由单项式生成的理想，作为向量空间恰由这些单项式的全部倍数生成：展开多项式系数给出一个包含，倍数本身属于理想给出反向包含。不同单项式线性无关，因此单项式属于该理想，当且仅当它被某个生成单项式整除。把指数结论用于所有生成单项式，便得有限生成。对单项式理想升链取并，选并理想的有限组生成元；它们同时落在某个链成员中，此后各成员都等于并理想，故升链稳定。
\end{proof}

这一论证称为 \term[Dickson-yinli]{Dickson 引理}[Dickson's lemma]。它只使用自然数指数的有限维性质，不预先借用多项式环的 Hilbert 基定理。

\begin{theorem}\label{b3:thm:groebner-existence-normal-form}
设 \(k\) 为域，\(R=k[x_1,\ldots,x_n]\)，并固定全局单项式序。每个理想 \(I\subseteq R\) 都存在有限 Gröbner 基。若 \(G\) 为其 Gröbner 基，则 \(I=(G)\)；对每个 \(f\in R\)，按 \(G\) 除法所得的余式是只含标准单项式、与 \(f\) 模 \(I\) 同余的唯一多项式。特别地，
\[
f\in I\quad\Longleftrightarrow\quad\operatorname{NF}_G(f)=0,
\]
其中 \(\operatorname{NF}_G(f)\) 称为 \(f\) 的\term[zhenggui-shi]{正规式}[normal form]。
\end{theorem}
\begin{proof}
由\cref{b3:lem:dickson}，可从全部 \(\operatorname{LM}(f)\)、\(0\ne f\in I\)，中选出有限个单项式，生成首项理想。各选一个产生它的 \(g_i\in I\)，便得有限 Gröbner 基。

若 \(f\in I\)，按该基除法得到的余式 \(r=f-\sum_iq_ig_i\) 仍在 \(I\) 中。若 \(r\ne0\)，它的首单项式属于首项理想，却又不被任何 \(\operatorname{LM}(g_i)\) 整除，与\cref{b3:lem:dickson}证明中的单项式成员判据矛盾。因此 \(r=0\)，并得到 \(I=(G)\)。

对一般 \(f\)，除法余式只含标准单项式且与 \(f\) 同余。若还有一个这样的 \(r'\)，则非零的 \(r-r'\in I\) 也只含标准单项式，同样矛盾，故 \(r=r'\)。这证明唯一性。理想成员的余式已证为零；反向，余式为零的除法等式直接给出 \(f\in(G)=I\)。
\end{proof}
\symidx{\(\operatorname{NF}_G(f)\)：关于 Gröbner 基的正规式}

这个存在性证明从整个理想的首单项式中选择生成元，并各取一个产生它的多项式；它没有给出怎样从已知有限生成组找到这些多项式的计算过程。下一节的 Buchberger 算法才从有限输入出发，用四则运算和单项式比较构造 Gröbner 基。

\subsection{全局良序与算法终止的适用范围}
\label{b3:subsec:0801:04}

良序性不是一个可以在证明末尾略去的形式要求。若把一元单项式按 \(1\succ x\succ x^2\succ\cdots\) 排列，则 \(1-x\) 的首项为一。用它约化一，会依次留下 \(x,x^2,x^3,\ldots\)，永不终止。这类局部次序有其用途，但需要其他终止机制，不能直接使用本节的除法算法。

\ProseParagraphBreak
在本章的有限变量交换情形下，若一个乘法相容的全序以一为最小元，Dickson 引理反而能推出其良序性。否则有严格下降序列 \(m_1\succ m_2\succ\cdots\)；从中取 \(i<j\) 且 \(m_i\mid m_j\)。写 \(m_j=m_i u\)，因 \(1\preceq u\)，乘法相容性使 \(m_i\preceq m_j\)，矛盾。此论证依赖指数属于 \(\mathbb N^n\)，不能直接用于非交换字或无限多个变量。

% !TeX root = ../../Book3.tex
% 纲要标题：### 8.2 Buchberger 理论
% 纲要位置：工作记录/计划纲要.md，第 2148 行。
% 纲要细目（写作范围）：
% * 从首项对齐构造 S 多项式，以最大首单项式最小的表示和有界表示完整证明 Buchberger 判据
% * 给出 Buchberger 算法及完整终止性、正确性证明，以 Dickson 有限性控制首项理想升链
% * 构造约化 Gröbner 基，并完整证明固定基域和单项式序下的唯一性
% * 用手算例题与零余式表示证书解释基础算法，完整书写当前应用的推导

\section{Buchberger 理论}
\label{b3:sec:0802}

Gröbner 基控制理想中全部多项式的首项，定义本身包含无限多个检验。有限算法的关键在于理解最高项怎样抵消：两个生成元的首项抵消由一个 S 多项式记录，而多个最高项同时抵消可以分解为这样的两两抵消。本节将这项观察写成首项下降的严格论证，再据此建立算法。仍固定 \(R=k[x_1,\ldots,x_n]\) 及全局单项式序。
\subsection{S 多项式、Buchberger 判据与算法正确性}
\label{b3:subsec:0802:01}

在 \(k[x,y]\) 的字典序 \(x\succ y\) 下，取 \(f=x^2-y\)、\(g=xy-1\)。首单项式 \(x^2\) 与 \(xy\) 的最小公倍式为 \(x^2y\)，所以分别乘以 \(y\) 与 \(x\)，就能使两个首项相同，再相减得到
\[
yf-xg=y(x^2-y)-x(xy-1)=x-y^2.
\]
原来的最高单项式 \(x^2y\) 已经消失。对于非首一多项式，还需除以各自的首系数，使这两个最高项的系数也相同。

\begin{definition}\label{b3:def:s-polynomial}
设 \(k\) 为域，\(R=k[x_1,\ldots,x_n]\)，并固定全局单项式序。对非零 \(f,g\in R\)，置
\[
L=\operatorname{lcm}\bigl(\operatorname{LM}(f),\operatorname{LM}(g)\bigr).
\]
称
\begin{equation}\label{b3:eq:s-polynomial}
S(f,g)=
\frac{L}{\operatorname{LM}(f)}\frac{f}{\operatorname{LC}(f)}
-\frac{L}{\operatorname{LM}(g)}\frac{g}{\operatorname{LC}(g)}
\end{equation}
为 \(f,g\) 的 \term[S-duoxiangshi]{\(S\) 多项式}[S-polynomial]。单项式最小公倍式的各个指数取两指数的较大者。
\end{definition}
\symidx{\(S(f,g)\)：多项式 \(f,g\) 的 S-多项式}

两项的首项都归一化成 \(L\)，故它们相减后或者为零，或者首单项式严格小于 \(L\)。仅有 \(S(f,g)\in(f,g)\) 没有检验意义；重要的是能否用原生成元把它表示出来，同时让表示中的每个乘积也严格小于 \(L\)。

\begin{theorem}[Buchberger 判据]\label{b3:thm:buchberger-criterion}
设 \(k\) 为域，\(R=k[x_1,\ldots,x_n]\)，并固定全局单项式序。设 \(G=(g_1,\ldots,g_s)\) 由 \(R\) 中的非零多项式组成，\(I=(G)\)。以下条件等价：
\begin{equivalences}
\item\label{b3:item:buchberger-basis} \(G\) 是 \(I\) 的 Gröbner 基；
\item\label{b3:item:buchberger-zero} 每个 \(S(g_i,g_j)\)、\(i<j\)，按 \(G\) 除法的余式为零；
\item\label{b3:item:buchberger-bound} 对每个 \(i<j\)，存在表示
\begin{equation}\label{b3:eq:s-bounded-representation}
S(g_i,g_j)=\sum_{\ell=1}^s h_{ij\ell}g_\ell,
\qquad
\operatorname{LM}(h_{ij\ell}g_\ell)\prec L_{ij}
\quad(h_{ij\ell}\ne0),
\end{equation}
其中 \(L_{ij}=\operatorname{lcm}\bigl(\operatorname{LM}(g_i),\operatorname{LM}(g_j)\bigr)\)。
\end{equivalences}
空组对应零理想，后两个条件此时为空条件。
\end{theorem}
\begin{proof}
若\itemref{b3:item:buchberger-basis}成立，S 多项式属于 \(I\)，由\cref{b3:thm:groebner-existence-normal-form}余式为零，得到\itemref{b3:item:buchberger-zero}。再用多元除法的首项界，每个非零乘积的首单项式不超过 S 多项式的首单项式，而后者严格小于 \(L_{ij}\)，便得\itemref{b3:item:buchberger-bound}。S 多项式为零时取所有系数为零。

反过来，假定有这些有界表示。给定 \(0\ne f\in I\)，要证明某个 \(\operatorname{LM}(g_i)\) 整除 \(\operatorname{LM}(f)\)，我们希望找到一个生成元表示，使最大的乘积首单项式恰为 \(\operatorname{LM}(f)\)。若这个最大首单项式比 \(\operatorname{LM}(f)\) 更大，它的系数必在求和时抵消；S 多项式的有界表示正用来重新组织这些抵消。记 \(\widetilde g_i=g_i/\operatorname{LC}(g_i)\)，将生成元表示的系数展开为项，可写成
\[
f=\sum_{\nu=1}^t a_\nu m_\nu\widetilde g_{i_\nu},
\qquad a_\nu\in k^\times,\quad m_\nu\;\text{为单项式}.
\]
在全部这样的表示中，选择
\[
L_{\max}=\max_\nu\bigl\{m_\nu\operatorname{LM}(g_{i_\nu})\bigr\}
\]
最小的一种，良序性保证它存在。相加不会产生更大的项，故 \(\operatorname{LM}(f)\preceq L_{\max}\)。如果相等，某个 \(\operatorname{LM}(g_i)\) 整除 \(\operatorname{LM}(f)\)，已经得到所需结论。

否则令 \(J\) 为达到 \(L_{\max}\) 的那些指标。最高项系数抵消给出 \(\sum_{\nu\in J}a_\nu=0\)。固定 \(\nu_0\in J\)，于是
\[
\sum_{\nu\in J}a_\nu m_\nu\widetilde g_{i_\nu}
=\sum_{\nu\in J\setminus\{\nu_0\}}
a_\nu\bigl(m_\nu\widetilde g_{i_\nu}
-m_{\nu_0}\widetilde g_{i_{\nu_0}}\bigr).
\]
括号中若两个生成元下标相同，则乘子也相同，差为零；否则，令它们首单项式的最小公倍式为 \(L\)。二者都整除 \(L_{\max}\)，故 \(L\mid L_{\max}\)，括号中的差恰为 \((L_{\max}/L)S(g_{i_\nu},g_{i_{\nu_0}})\)。下标反序只改变符号。按假设，这可以用各 \(g_\ell\) 表示，且每个非零乘积的首单项式严格小于 \(L_{\max}\)。系数展开成项后这个界仍保持。

替换全部达到 \(L_{\max}\) 的项，其他原项本来就小于 \(L_{\max}\)，便得到最大首单项式更小的表示，与所选 \(L_{\max}\) 的最小性矛盾。因此必有 \(\operatorname{LM}(f)=L_{\max}\)。每个非零理想成员的首单项式都被某个 \(\operatorname{LM}(g_i)\) 整除，结合 \(G\subset I\) 即得首项理想相等，证明\itemref{b3:item:buchberger-basis}。
\end{proof}

这称为\term[Buchberger-panju]{Buchberger 判据}[Buchberger criterion]。证明中的第三种表述不仅解释正确性，还使我们能够保存每次计算的具体依据：以后即使加入新生成元，旧的有界表示仍然成立，不必假定不同除法顺序给出相同的中间余式。

\subsection{首项理想升链与算法终止性}
\label{b3:subsec:0802:02}

如果一个 S 多项式留下非零余式，余式仍在原理想中，但它的首单项式未被当前生成元控制。把余式加入即可改进生成组；同时必须登记它与原有全部生成元组成的新多项式对。

\begin{algorithm}[htbp]
\caption{Buchberger}
\label{b3:alg:buchberger}
\begin{algorithmsteps}{有限组 \(F=(f_1,\ldots,f_m)\subseteq R\)，可有效比较的全局单项式序，以及支持精确四则运算和零判定的系数域 \(k\)。}{\(I=(F)\) 的一个有限 Gröbner 基 \(G\)。}
\State 删除零多项式，将其余元素首一化，记为 \(G=(g_1,\ldots,g_s)\)。
\State 建立待处理集合 \(P=\bigl\{(i,j)\mid1\le i<j\le s\bigr\}\)。
\While{\(P\ne\varnothing\)}
\State 取出并删除一对 \((i,j)\)，按当前 \(G\) 对 \(S(g_i,g_j)\) 除法，得到余式 \(r\)。
\If{\(r\ne0\)}
\State 把 \(r/\operatorname{LC}(r)\) 加到 \(G\) 末尾，记新下标为 \(s+1\)。
\State 将 \((1,s+1),\ldots,(s,s+1)\) 加入 \(P\)，再令 \(s\) 增加一。
\EndIf
\EndWhile
\State \Return \(G\)。
\end{algorithmsteps}
\end{algorithm}

\begin{theorem}\label{b3:thm:buchberger-algorithm}
设域 \(k\) 支持精确四则运算与零判定，输入为 \(k[x_1,\ldots,x_n]\) 中的有限生成组及一个可有效比较的全局单项式序。\Cref{b3:alg:buchberger}在有限步后终止，输出确为原理想的 Gröbner 基。算法还可同时记录每个新生成元关于原输入的多项式线性组合。
\end{theorem}
\begin{proof}
每轮中的多元除法由\cref{b3:thm:multivariate-division}终止，依据是待处理首单项式在全局良序中严格下降。该轮有等式 \(S(g_i,g_j)=\sum_\ell q_\ell g_\ell+r\)，所以余式仍在当前理想中，添加它不改变理想。非零余式的首单项式不被已有首单项式整除，故每次添加使它们生成的单项式理想严格增大；\cref{b3:lem:dickson}的升链稳定性排除无限次添加。这控制了生成组的增长。每次添加只引入有限个新对，每轮删除一个待处理对；不再添加之后，待处理集合已经有限，余下的对也会在有限轮内处理完，因此整个算法终止。

处理某对时，若余式为零，多元除法给出\eqref{b3:eq:s-bounded-representation}。若余式非零，添加首一 \(g_{\rm new}=r/\operatorname{LC}(r)\) 后，有
\[
S(g_i,g_j)=\sum_\ell q_\ell g_\ell+
\operatorname{LC}(r)g_{\rm new}.
\]
所有非零乘积的首单项式都不超过原 S 多项式的首单项式，因而严格小于 \(L_{ij}\)。算法保留已有生成元，这些有界表示在后续步骤中仍成立。终止时最终组的每一对都已处理，由\cref{b3:thm:buchberger-criterion}即得 Gröbner 性。

最初首一化仅除以非零常数。已有生成元关于原输入的表示时，将它们代入 S 多项式和除法等式，就得到余式的表示，再除以首系数得到新生成元的表示。因此可以与算法同步，逐步记录全部关系。
\end{proof}

还可以省去部分 S 多项式的实际除法。例如首一 \(f=a+u,g=b+v\)，其中 \(a,b\) 为首单项式；若 \(a,b\) 互素，则
\[
S(f,g)=bf-ag=ug-vf.
\]
这里两个非零乘积的首单项式都严格小于 \(ab\)。因此这对已经具有判据所需的有界表示。这个\term[Buchberger-chengji-panju]{Buchberger 乘积判据}[Buchberger product criterion]只略去有数学依据的对；单凭两个多项式看起来简单，不能跳过检验。

\subsection{约化 Gröbner 基的唯一性}
\label{b3:subsec:0802:03}

\begin{definition}\label{b3:def:reduced-groebner}
设 \(k\) 为域，\(R=k[x_1,\ldots,x_n]\)，并固定全局单项式序。给定理想 \(I\subseteq R\) 的 Gröbner 基 \(G\)。若 \(G\) 的各元素首一，任意两个不同元素的首单项式互不整除，而且每个元素除首单项式以外的全部单项式都是 \(I\) 的标准单项式，则称 \(G\) 为\term[yuehua-Groebner-ji]{约化 Gröbner 基}[reduced Gröbner basis]。零理想的约化基为空，单位理想的约化基为 \(\{1\}\)。
\end{definition}

这两项限制分别去除首项生成组中的整除冗余，并把每个基元素的尾部完全化为标准单项式。

\begin{theorem}\label{b3:thm:reduced-groebner-unique}
设 \(k\) 为域，并在 \(R=k[x_1,\ldots,x_n]\) 上固定全局单项式序。每个理想 \(I\subseteq R\) 都有唯一的约化 Gröbner 基。由任意 Gröbner 基可以在有限步内得到它。
\end{theorem}
\begin{proof}
先在一组 Gröbner 基的首单项式中只保留整除关系下的极小者，相同单项式只保留一个对应多项式。保留者仍属于原理想，且首单项式生成原首项理想，因而仍为 Gröbner 基。将它们首一化，记为 \(G=\{g_1,\ldots,g_s\}\)，并置 \(m_i=\operatorname{LM}(g_i)=\operatorname{LT}(g_i)\)。对每个 \(i\)，用 \(G\setminus\{g_i\}\) 对整个尾部 \(g_i-m_i\) 作完整的多元除法，得到
\[
g_i-m_i=\sum_{j\ne i}q_{ij}g_j+r_i.
\]
以 \(g_i'=m_i+r_i=g_i-\sum_{j\ne i}q_{ij}g_j\) 替换 \(g_i\)。除法的首项界保证 \(r_i\) 的每个单项式都小于 \(m_i\)，所以首项不变；余式条件保证这些单项式不被任何 \(m_j\)（\(j\ne i\)）整除。它们也不能被 \(m_i\) 整除，因为 \(m_i\) 的真倍数严格大于 \(m_i\)。全部 \(g_i'\) 仍在 \(I\) 中，且其首单项式仍生成原首项理想，故它们构成约化 Gröbner 基。空基的情形无需除法。

两个约化基的首单项式都恰为首项理想中整除关系的极小单项式。这组单项式只由理想决定：任一理想单项式被其中一个整除，而极小者不能被另一个严格更小的理想单项式整除。按相同首单项式配对两个基元素 \(g,h\)。二者首一，故首项相消；所有余下项均为标准单项式。又 \(g-h\in I\)，由\cref{b3:thm:groebner-existence-normal-form}只能有 \(g-h=0\)。逐对相等，证明唯一性。
\end{proof}

唯一的是固定次序下的约化基，而不是任意 Gröbner 基。增加一个已在理想中的冗余多项式通常不会破坏 Gröbner 基性质；改变次序却可能同时改变首项理想、标准单项式和约化基。两种变化应分别考察。

\subsection{手算例题与零余式的线性组合表示}
\label{b3:subsec:0802:04}

在 \(\mathbb Q[x,y]\) 中取字典序 \(x\succ y\)，以
\[
f_1=x^2-y,\qquad f_2=xy-1
\]
为输入。前面的首项对齐已经给出 \(h=yf_1-xf_2=x-y^2\)。\(h\) 的首单项式 \(x\) 不被 \(x^2,xy\) 整除，所以要加入。随后 \(f_2\) 与 \(h\) 的 S 多项式为
\[
p=f_2-yh=y^3-1.
\]
这又提供未被控制的纯 \(y\) 首单项式。我们可以不再保留冗余的原生成元，而用以下双向等式核验最终理想：
\begin{equation}\label{b3:eq:groebner-cubic-certificates}
\begin{aligned}
h&=yf_1-xf_2,&
p&=-y^2f_1+(1+xy)f_2,\\
f_1&=(x+y^2)h+yp,&
f_2&=yh+p.
\end{aligned}
\end{equation}
因此 \(I=(f_1,f_2)=(h,p)\)。对最后一组，只需检验一对：
\[
S(h,p)=y^3h-xp=x-y^5=h-y^2p.
\]
先减去 \(h\)，再消去 \(-y^2p\)，余式为零。判据表明 \(\{x-y^2,y^3-1\}\) 是 Gröbner 基，而且已经约化。所有等式在任意系数域中仍成立；若进一步数几何点，才须另看 \(y^3-1\) 是否有重根。

\ProseParagraphBreak
换成字典序 \(y\succ x\)，同一理想的约化基为
\[
\{y-x^2,\ x^3-1\}.
\]
因为 \(y-x^2=-f_1\)、\(x^3-1=xf_1+f_2\)，反向有 \(f_2=x(y-x^2)+(x^3-1)\)；首单项式为 \(y,x^3\)，互素，乘积判据给出 Gröbner 性。两种次序分别先消去 \(x\) 与 \(y\)，商代数虽然相同，适合展示的坐标却不同。

\ProseParagraphBreak
等式\eqref{b3:eq:groebner-cubic-certificates}还有一个实际作用：检查者只需展开有限个多项式等式，再检查一个 S 多项式，就能确认最终答案。只报告“软件返回这组基”，或只验证原多项式对新组的余式为零，都少了一半理想包含的依据。

% !TeX root = ../../Book3.tex
% 纲要标题：### 8.3 消去与理想运算
% 本轮补充的纲要细目：
% * 对特征不为二的平面同时变号作用，分别证明不变环的生成性与全部关系，再把核写成消去理想
% * 区分候选不变量所生成的子代数与整个不变环，说明特征二时生成结论的变化
% #### 8.3.1 消去定理与投影的 Zariski 闭包
% #### 8.3.2 理想成员、交、理想商与饱和
% #### 8.3.3 隐式化、参数消去与 Rabinowitsch 技巧
% #### 8.3.4 平面同时变号的不变环
% #### 8.3.5 理想等式与乘子表示的核验
% 纲要位置：工作记录/计划纲要.md，第 2054 行。
% 纲要细目（写作范围）：
% * 消去定理与投影的 Zariski 闭包；不要把消去所得闭包无条件当作投影像
% * 理想成员、交、理想商与饱和的计算
% * 隐式化、参数消去及 Rabinowitsch 技巧的有效版本
% * 平面同时变号的不变环：先证明生成性，再求全部关系
% * 保留理想等式与乘子表示，以便独立核验


\section{消去与理想运算}
\label{b3:sec:0803}

消去变量的目的，是把一个含参数或辅助变量的问题转成原坐标中的理想问题。但消去所得方程只记录闭条件；被排除的分母和不能提升的点，未必能由这些方程单独恢复。本节在计算理想时保留等式，在解释零点时另外核对基域、闭包和分母条件。


\subsection{消去定理与投影的 Zariski 闭包}
\label{b3:subsec:0803:01}

设 \(k\) 为域，\(R=k[x_1,\ldots,x_r,y_1,\ldots,y_s]\)，\(S=k[y_1,\ldots,y_s]\)。对理想 \(I\subseteq R\)，交 \(I\cap S\) 称为对前 \(r\) 个变量的\term[xiaoqu-lixiang]{消去理想}[elimination ideal]。取字典序，且全部 \(x\) 变量先于 \(y\) 变量。它的关键性质是：只含 \(y\) 的单项式小于任何含 \(x\) 的单项式。

\begin{theorem}[消去定理]\label{b3:thm:groebner-elimination}
设 \(k\) 为域，\(I\subseteq R=k[x_1,\ldots,x_r,y_1,\ldots,y_s]\) 为理想，\(S=k[y_1,\ldots,y_s]\)。对全部 \(x\) 变量先于 \(y\) 变量的字典序，若 \(G\) 是 \(I\) 的 Gröbner 基，则 \(G\cap S\) 是 \(I\cap S\) 的 Gröbner 基。
\end{theorem}
\begin{proof}
取 \(0\ne f\in I\cap S\)。因 \(G\) 为 Gröbner 基，某个 \(g\in G\) 的首单项式整除 \(\operatorname{LM}(f)\)。后者只含 \(y\)，故 \(\operatorname{LM}(g)\) 也只含 \(y\)。如果 \(g\) 中另有含 \(x\) 的项，该项在当前次序下将大于 \(\operatorname{LM}(g)\)，矛盾。因此 \(g\in G\cap S\)。于是 \(I\cap S\) 中任一非零元素的首单项式都被 \(G\cap S\) 的某个首单项式整除；这些基元素本来就在交理想中，故首项理想相等。由\cref{b3:thm:groebner-existence-normal-form}，它们也生成整个交理想。
\end{proof}

这个结果称为\term[xiaoqu-dingli]{消去定理}[elimination theorem]。证明同样适用于满足上述变量分块比较性质的全局序；对任意次序则没有这种保证。

\begin{theorem}[投影的闭包]\label{b3:thm:elimination-projection-closure}
设 \(k\) 代数闭，\(R=k[x_1,\ldots,x_r,y_1,\ldots,y_s]\)，\(S=k[y_1,\ldots,y_s]\)，\(I\subseteq R\) 为理想，\(\pi:k^{r+s}\to k^s\) 为忘掉 \(x\) 坐标的投影。则
\[
\overline{\pi\bigl(Z_k(I)\bigr)}^{\mathrm{Zar}}=Z_k(I\cap S).
\]
此外，投影像的消失理想为 \(\sqrt I\cap S=\sqrt{I\cap S}\)。
\end{theorem}
\begin{proof}
对 \(f\in S\)，它在 \(\pi\bigl(Z_k(I)\bigr)\) 上为零，当且仅当把它看成不含 \(x\) 的多项式后，在 \(Z_k(I)\) 上为零。由\cref{b3:thm:strong-nullstellensatz}，后者等价于 \(f\in\sqrt I\)。因此消失理想为 \(\sqrt I\cap S\)。对 \(f\in S\)，某个 \(f^N\) 在 \(I\) 中当且仅当在 \(I\cap S\) 中，故这也等于 \(\sqrt{I\cap S}\)。任一集合的 Zariski 闭包是其消失理想的零点集，结论随之成立。
\end{proof}

例如 \(I=(xy-1)\subseteq k[x,y]\) 满足 \(I\cap k[y]=0\)：商环同构于 \(k[y,y^{-1}]\)，所以 \(k[y]\) 映入商环为单射。投影像是 \(k\setminus\{0\}\)，闭包却是整个直线。闭包定理没有断言每个闭包中的点都能补出被消去的坐标。

\ProseParagraphBreak
代数闭条件也有实际作用。在 \(\mathbb R[x,y]\) 中取 \(I=(x^2+1)\)，其消去理想为零，实零点集却为空。代数的消去计算仍正确，关于实点投影闭包的上述公式则不适用；实零点的消去需要额外的有序域理论。

\subsection{理想成员、交、理想商与饱和}
\label{b3:subsec:0803:02}

成员判定由\cref{b3:thm:groebner-existence-normal-form}完成。下面把其他理想运算也化成可以消去的理想。暂设 \(I=(f_1,\ldots,f_a)\)、\(J=(g_1,\ldots,g_b)\) 都在 \(R=k[x_1,\ldots,x_n]\) 中。

\begin{proposition}\label{b3:prop:ideal-intersection-elimination}
设 \(k\) 为域，\(R=k[x_1,\ldots,x_n]\)，\(I=(f_1,\ldots,f_a)\)、\(J=(g_1,\ldots,g_b)\) 为 \(R\) 的理想。在 \(R[u]\) 中扩张这两个理想后，有
\[
I\cap J=\bigl(uI+(1-u)J\bigr)\cap R.
\]
右侧可对生成组 \(uf_i,(1-u)g_j\) 取消去 \(u\) 的 Gröbner 基求得。
\end{proposition}
\begin{proof}
若 \(h\in I\cap J\)，写 \(h=uh+(1-u)h\)，便得右侧包含它。反过来，若 \(h\in R\) 有表示
\(h=u\sum A_if_i+(1-u)\sum B_jg_j\)，置 \(u=1\) 得 \(h\in I\)，置 \(u=0\) 得 \(h\in J\)。两次代入也把计算所得的表示分别变成原环中的成员表示。
\end{proof}

\begin{proposition}\label{b3:prop:colon-saturation-elimination}
设 \(k\) 为域，\(R=k[x_1,\ldots,x_n]\)，\(I\subseteq R\) 为理想。若 \(0\ne g\in R\)，并已求得 \(I\cap(g)=(h_1,\ldots,h_c)\)，则各 \(h_i\) 被 \(g\) 整除，且
\[
(I:g)=(h_1/g,\ldots,h_c/g).
\]
对任意 \(g\in R\)，记 \((I:g^\infty)=\bigcup_{N\ge0}(I:g^N)\)，则
\[
(I:g^\infty)=\bigl(IR[u]+(1-ug)\bigr)\cap R.
\]
对理想 \(J=(g_1,\ldots,g_b)\subseteq R\)，记
\[
(I:J^\infty)=\bigcup_{N\ge0}(I:J^N),\qquad
(I:J^N)=\{h\in R\mid J^Nh\subseteq I\},
\]
其中 \(J^0=R\)。则有
\[
(I:J)=\bigcap_{j=1}^b(I:g_j),\qquad
(I:J^\infty)=\bigcap_{j=1}^b(I:g_j^\infty).
\]
若 \(J=0\)，后两式的左端均为 \(R\)；用空生成组表示 \(J\) 时，右端按空交等于 \(R\) 的约定解释。
\end{proposition}
\begin{proof}
对第一式，若 \(ag\in I\)，它属于 \(I\cap(g)\)，写成 \(ag=\sum c_i h_i\)，在整环中消去非零 \(g\) 得 \(a=\sum c_i(h_i/g)\)。反向直接乘以 \(g\) 即得。

若 \(g=0\)，饱和与所述消去理想都等于 \(R\)。以下设 \(g\ne0\)。若 \(g^Nh\in I\)，在 \(R[u]\) 中有恒等式
\[
h=u^Ng^Nh+(1-ug)h\sum_{j=0}^{N-1}(ug)^j,
\]
故 \(h\) 在所述消去理想中，\(N=0\) 时直接使用 \(h\in I\)。反向，对一个成员表示代入 \(u=g^{-1}\)，得到 \(h/1\in IR_g\)；有限清分母给出某个 \(g^Nh\in I\)。

理想商的交式直接来自逐个生成元检验。对饱和，\(J=0\) 时两端均为 \(R\)，以下设 \(J\ne0\)。若 \(J^Nh\subseteq I\)，则 \(g_j^Nh\in I\)，给出一个包含。反向，设分别有 \(g_j^{N_j}h\in I\)，可增大到 \(N_j\ge1\)。取 \(N=1+\sum_j(N_j-1)\)。\(J^N\) 的每个生成单项式至少含某个 \(g_j^{N_j}\)，故乘以 \(h\) 后都在 \(I\) 中。这证明另一包含。
\end{proof}

例如 \(I=(x^2,xy)\subseteq k[x,y]\)，直接计算得
\[
(I:x)=(x,y),\qquad (I:y)=(x),\qquad
(I:x^\infty)=R,\qquad (I:y^\infty)=(x).
\]
第一式由 \(xh=x(ax+by)\) 消去 \(x\) 得到；第二式中 \(yh\in(x)\) 迫使 \(h\in(x)\)，反向由 \(xy\in I\) 得到。因 \(x^2\in I\)，关于 \(x\) 的饱和为全环；关于 \(y\) 的饱和保持 \((x)\)。于是
\(\bigl(I:(x,y)^\infty\bigr)=(x)\)。若误把多生成元饱和写成关于乘积 \(xy\) 的饱和，就会错误地得到全环。

\subsection{隐式化、参数消去与 Rabinowitsch 技巧}
\label{b3:subsec:0803:03}

\begin{proposition}[参数映射的隐式理想]\label{b3:prop:parametric-implicit-ideal}
设 \(k\) 为域，\(t=(t_1,\ldots,t_m)\) 为代数无关元，\(0\ne q\in k[t]\)，\(p_1,\ldots,p_n\in k[t]\)。定义
\[
\phi:k[x_1,\ldots,x_n]\longrightarrow k[t,q^{-1}],
\qquad x_i\longmapsto p_i/q.
\]
则
\[
\ker\phi=
(qx_1-p_1,\ldots,qx_n-p_n,uq-1)\cap k[x_1,\ldots,x_n],
\]
右边的理想先在 \(k[t,x,u]\) 中生成，再消去 \(t,u\)。
\end{proposition}
\begin{proof}
令右边取交前的理想为 \(H\)。商环中 \(q\) 可逆，逆元为 \(u\)，并有 \(x_i=up_i\)。所以把 \(t_j\) 送到自身、\(u\) 送到 \(q^{-1}\)、\(x_i\) 送到 \(p_i/q\)，给出商环到 \(k[t,q^{-1}]\) 的同构；逆映射由局部化泛性质将 \(q^{-1}\) 送回 \(\bar u\)。限制到 \(k[x]\)，核就是 \(H\cap k[x]\)，得到公式。
\end{proof}

例如有理参数化
\[
x=\frac{t-1}{t},\qquad y=\frac{t}{t-1},\qquad t(t-1)\ne0
\]
满足 \(xy=1\)。它的隐式理想恰为 \((xy-1)\)：商 \(k[x,y]/(xy-1)\cong k[x,x^{-1}]\) 映入 \(k(t)\) 时，\(x\) 的像 \(1-t^{-1}\) 是超越元，所以映射单射。在代数闭域上，参数像是双曲线删去 \((1,1)\)；因为其余点可由 \(t=(1-x)^{-1}\) 恢复。消去得到整条双曲线，原来的分母条件仍须另记。

\ProseParagraphBreak
Rabinowitsch 技巧还给出一个有效的根理想成员判据。设 \(k\) 为域，\(R=k[x_1,\ldots,x_n]\)，\(I=(f_1,\ldots,f_m)\subseteq R\) 为理想。对 \(f\in R\)，
\begin{equation}\label{b3:eq:effective-rabinowitsch}
f\in\sqrt I\quad\Longleftrightarrow\quad
1\in IR[u]+(1-uf).
\end{equation}
这是\cref{b3:prop:colon-saturation-elimination}中“关于 \(f\) 的饱和为全环”的情形，也可直接使用其中的恒等式。若 \(f=0\)，右边的理想含 \(1\)，而 \(f^1=0\in I\) 就是所需证书。若 \(f\ne0\)，右边可以用 Gröbner 基检验；结果为真时，保存
\(1=\sum A_i(u,x)f_i+B(u,x)(1-uf)\)，代入 \(u=f^{-1}\) 再清分母，便产生某个实际指数 \(N\ge1\) 和等式 \(f^N=\sum c_i f_i\)。这个判据不要求底域代数闭；只有把它解释为“在全部几何零点上消失”时才调用零点定理。

\subsection{平面同时变号的不变环}
\label{b3:subsec:0803:sign-invariants}

将 \(x,y\) 同时换成 \(-x,-y\) 后，\(x^2,xy,y^2\) 都保持不变，而且满足 \(x^2y^2=(xy)^2\)。这提示我们用三个新变量记录这些多项式，再消去 \(x,y\) 求出它们之间的全部关系。要由此描述所有保持不变的多项式，还须证明这三个元素确实生成相应子代数。

\begin{definition}\label{b3:def:invariant-subring}
设 \(k\) 为域，\(A\) 为交换 \(k\) 代数，群 \(G\) 通过 \(k\) 代数自同构作用于 \(A\)。这意味着每个 \(g\in G\) 给定一个保持加法、乘法和 \(k\) 中元素的双射 \(a\mapsto g\cdot a\)，满足
\[
1\cdot a=a,\qquad (gh)\cdot a=g\cdot(h\cdot a).
\]
由所有满足 \(g\cdot a=a\)、\(g\in G\) 的元素组成的子代数
\[
A^G=\{\,a\in A\mid g\cdot a=a\text{ 对所有 }g\in G\,\}
\]
称为这一作用的\term[bubian-huan]{不变环}[invariant ring]。
\symidx{\(A^G\)：群作用的不变环}
\end{definition}

它确为子代数：若 \(a,b\) 都固定，则每个自同构也固定 \(a+b\)、\(ab\) 及所有标量。现在取 \(A=k[x,y]\)、\(C_2=\{1,s\}\)，令
\[
s\cdot f(x,y)=f(-x,-y).
\]
代入两次恢复原多项式，故这给出一个作用。下面不假定 \(k\) 代数闭，也不把多项式等式替换为在 \(k^2\) 中逐点取值相等；后者在有限域上不足以判定多项式。

\begin{proposition}\label{b3:prop:sign-invariant-presentation}
设 \(k\) 为特征不为二的域，\(C_2\) 在 \(k[x,y]\) 上通过同时变号作用。则
\[
k[x,y]^{C_2}=k[x^2,xy,y^2].
\]
把 \(u,v,w\) 分别送到 \(x^2,xy,y^2\) 的同态诱导同构
\begin{equation}\label{b3:eq:sign-invariant-presentation}
k[u,v,w]/(uw-v^2)\ \cong\ k[x,y]^{C_2}.
\end{equation}
\end{proposition}
\begin{proof}
对 \(f=\sum_{i,j}a_{ij}x^iy^j\) 比较系数，有
\[
s\cdot f=f
\quad\Longleftrightarrow\quad
\bigl((-1)^{i+j}-1\bigr)a_{ij}=0\quad\text{对所有 }i,j.
\]
因 \(2\ne0\)，这等价于所有奇数全次数的系数为零。当 \(i+j\) 为偶数时，\(i,j\) 或同时为偶数，或同时为奇数，相应单项式分别写为
\[
x^{2a}y^{2c}=(x^2)^a(y^2)^c,\qquad
x^{2a+1}y^{2c+1}=(xy)(x^2)^a(y^2)^c.
\]
所以全部不变多项式由 \(x^2,xy,y^2\) 生成。这三个元素本身固定，反向包含也成立。

记上述满射为 \(\phi:k[u,v,w]\to k[x,y]^{C_2}\)。恒等式 \(x^2y^2=(xy)^2\) 说明 \(uw-v^2\in\ker\phi\)，但尚不能据此断定核只有这一条生成关系。把任意 \(F\in k[u,v,w]\) 看作系数在 \(k[u,w]\) 中的 \(v\) 多项式，对首一多项式 \(v^2-uw\) 作除法，得到
\[
F=Q(v^2-uw)+A(u,w)+v\cdot B(u,w).
\]
若 \(\phi(F)=0\)，则
\[
A(x^2,y^2)+xyB(x^2,y^2)=0.
\]
第一项只含两个指数都为偶数的单项式，第二项只含两个指数都为奇数的单项式。两组不相交，而且每组内部不同的 \(u^aw^c\) 也给出不同的单项式。因此所有系数都为零，即 \(A=B=0\)。这证明 \(\ker\phi=(uw-v^2)\)，从而得到所述同构。
\end{proof}

这个核也正是一个消去理想。应用\cref{b3:prop:parametric-implicit-ideal}的多项式情形，取参数 \(x,y\) 和分母 \(q=1\)，得到
\[
(u-x^2,v-xy,w-y^2)\cap k[u,v,w]=(uw-v^2),
\]
其中左侧的理想先在 \(k[x,y,u,v,w]\) 中生成。命题中的除法与单项式独立性直接证明了这个等式；也可以用消去序计算左侧。若只猜出三个不变多项式而未证明它们生成整个不变环，同样的消去计算便只描述这三个元素所生成的子代数。

\ProseParagraphBreak
特征条件用在了排除奇数次单项式的步骤。若 \(\operatorname{char}k=2\)，同时变号就是恒等作用，不变环为整个 \(k[x,y]\)；子代数 \(k[x^2,xy,y^2]\) 不包含 \(x\)。然而三个给定生成元之间的核仍为 \((uw-v^2)\)，因为上面的首一除法与单项式独立性在任意特征下都成立。生成性与关系的完备性，确实需要分别检查。

\subsection{理想等式与乘子表示的核验}
\label{b3:subsec:0803:04}
\label{b3:subsec:0803:certificate-verification}

消去算法输出的一组生成元，应当带着它们属于原理想的表示。若原生成组为 \(F=(f_1,\ldots,f_m)\)，计算时可把每个新多项式连同一个向量一起保存：
\[
g_j=\sum_{i=1}^m a_{ij}f_i.
\]
加减多项式、乘单项式及首一化时，对表示向量执行同一运算，因而始终保持等式。若最终还验证每个 \(f_i\) 对 \(G=(g_j)\) 的余式为零，就得到反向表示 \(f_i=\sum_j b_{ji}g_j\)。两组等式分别证明 \((G)\subseteq(F)\) 与 \((F)\subseteq(G)\)；Gröbner 性则由有界 S 多项式表示单独证明，不能只凭理想相等推出。

\ProseParagraphBreak
对交理想，\cref{b3:prop:ideal-intersection-elimination}的表示在 \(u=0,1\) 的两次取值给出两边的成员证明。对单个元素的饱和，清分母后得到 \(g^Nh\in I\)。若 \(J=(g_1,\ldots,g_b)\)、\(b\ge1\)，各生成元分别给出 \(g_j^{N_j}h\in I\)、\(N_j\ge1\)，则按\cref{b3:prop:colon-saturation-elimination}的证明取 \(N=1+\sum_j(N_j-1)\)，才可保证 \(J^Nh\subseteq I\)。例如在 \(k[x,y]\) 中，\(I=(x^2,y^2)\)、\(J=(x,y)\) 时，\(x^2,y^2\in I\)，但 \(xy\notin I\)，故 \(J^2\not\subseteq I\)，而 \(J^3\subseteq I\)。

\ProseParagraphBreak
对于消去理想的完备性，只检验“所得多项式确实消去变量”还不够，还需确认使用了消去序且原扩展理想的 Gröbner 基已经算完。

\ProseParagraphBreak
例如要证明 \((xy,xz)\cap(y,z)=(xy,xz)\)，只须观察两生成元都在 \((y,z)\) 中，原理想本来已包含于后者；不必为每个简单理想关系都重算一次 Gröbner 基。算法提供一般办法，具体结构往往给出更短的证明。无论采用哪种方法，都应保留足以恢复两侧包含的等式或论证。

\ProseParagraphBreak
\cref{b3:tab:elimination-constructions}汇集了本节的辅助变量构造。交理想中，\(u\) 用来分别取值；饱和与有理参数消去中，它用来使指定元素可逆。根理想成员判定则把“某个幂属于理想”转成辅助理想是否含 \(1\) 的问题。

\begin{table}[htbp]
\centering
\caption[消去构造的比较]{消去构造的比较。设 \(k\) 为域，\(R=k[x_1,\ldots,x_n]\)，\(I,J\subseteq R\) 为理想，\(f,g\in R\)，\(u\) 为辅助变量。最后一行另取 \(t=(t_1,\ldots,t_m)\)、\(p_1,\ldots,p_n,q\in k[t]\) 且 \(q\ne0\)。}
\label{b3:tab:elimination-constructions}
\begin{tabularx}{\textwidth}{@{}>{\raggedright\arraybackslash}p{0.20\textwidth}>{\raggedright\arraybackslash}X>{\raggedright\arraybackslash}p{0.29\textwidth}@{}}
\toprule
所求对象 & 辅助理想 \(H\) & 提取方式 \\
\midrule
交理想 \(I\cap J\) & \(uIR[u]+(1-u)JR[u]\) & \(H\cap R\) \\
单元素饱和 \(I:g^\infty\) & \(IR[u]+(1-ug)\) & \(H\cap R\) \\
根理想成员 \(f\in\sqrt I\) & \(IR[u]+(1-uf)\) & 判定 \(1\in H\)，并恢复 \(f^N\in I\)、\(N\ge1\) \\
有理参数映射的核 & \((qx_i-p_i\ (1\le i\le n),\,uq-1)\)，生成于 \(k[t,x,u]\) & \(H\cap R=\ker\phi\)，其中 \(\phi:R\to k[t,q^{-1}]\)、\(x_i\mapsto p_i/q\) \\
\bottomrule
\end{tabularx}
\end{table}

% !TeX root = ../../Book3.tex
% 纲要标题：### 8.4 可验证的计算与系数域
% 纲要位置：工作记录/计划纲要.md，第 2061 行。
% 纲要细目（写作范围）：
% * 有理数、有限域与有理函数域上的精确运算
% * 输入、输出、正确性证书及可重复计算的记录格式
% * 终止性与实际复杂度的区别；系数膨胀和单项式序的影响
% * 数值近似不能替代理想成员与理想相等的精确证明


\section{可验证的计算与系数域}
\label{b3:sec:0804}


同一组符号式放在不同的系数环中，可能给出不同的理想。精确计算除了记录多项式，还须说明允许哪些除法、参数取值是否使分母为零，以及怎样独立检验结果。本节用具体算例说明这些要求，并区分算法终止、计算规模和数值近似。

\subsection{有理数、有限域与有理函数域上的精确运算}
\label{b3:subsec:0804:01}

前面各算法只用到系数的加、减、乘、除及零判定。因此，给定一个抽象域还不够：实际执行时，必须说明它的元素怎样表示、这些运算怎样完成。\(\mathbb Q\) 中的系数可用互素整数对表示，每次运算后约分；\(\mathbb F_p\) 中的系数用模 \(p\) 的剩余类表示，除法由扩展 Euclid 算法求逆。对于有限扩域 \(\mathbb F_p[z]/(h)\)，还须给定不可约多项式 \(h\)，并把系数约化为次数小于 \(\deg h\) 的多项式。

\ProseParagraphBreak
在 \(k(a_1,\ldots,a_s)\) 上计算时，参数 \(a_i\) 被视为代数无关元。系数是多项式之比，零判定归结为分子是否为零多项式；其分母在有理函数域中可逆，却未必在参数的每个具体取值处可逆。例如 \((ax-1)\subset k(a)[x]\) 的约化 Gröbner 基为 \(\{x-a^{-1}\}\)。当 \(a=c\ne0\) 时它可以直接专门化；当 \(a=0\) 时原生成元成为 \(-1\)，所得理想是单位理想，原公式则没有意义。

\begin{proposition}\label{b3:prop:groebner-specialization-open}
设 \(k\) 为域，\(a_1,\ldots,a_s\) 为 \(k\) 上的代数无关元，\(A=k[a_1,\ldots,a_s]\)、\(K=\operatorname{Frac}(A)\)，并固定 \(K[x_1,\ldots,x_n]\) 中的全局单项式序。若 \(F\subset A[x_1,\ldots,x_n]\) 为有限集，\(G\) 是 \((F)K[x_1,\ldots,x_n]\) 的首一 Gröbner 基，则存在非零 \(d\in A\)，使 \(G\subset A_d[x_1,\ldots,x_n]\)，且对任意域扩张 \(k\subset L\) 和任意满足 \(d(c)\ne0\) 的 \(c\in L^s\)，专门化后的 \(G(c)\) 是 \(\bigl(F(c)\bigr)\) 的 Gröbner 基，其首单项式与 \(G\) 相同。
\end{proposition}
\begin{proof}
记录 \(G\) 的全部系数、两组生成元互相表示的系数，以及每个 \(S(g_i,g_j)\) 约化为零所得的有界表示系数。这些都是有限多个有理函数。取其所有非零分母的乘积为 \(d\)，便可把全部恒等式写在 \(A_d[x_1,\ldots,x_n]\) 中。

当 \(d(c)\ne0\) 时，代入给出良定义的环同态 \(A_d\to L\)。两组互相表示的等式保证 \(\bigl(G(c)\bigr)=\bigl(F(c)\bigr)\)。每个 \(g_i\) 首一，故专门化不改变其首单项式；其余系数即使变为零，也只会删去较低的项。因此，各 \(S\)-多项式专门化后仍具有小于相应最小公倍数的有界表示。由\cref{b3:thm:buchberger-criterion}，\(G(c)\) 是所求 Gröbner 基。
\end{proof}

这一命题只说明存在一个可以沿用计算结果的非空开集。对 \(d(c)=0\) 的参数必须另算，不能把它们当作无解。例如 \(F=\{ax,y-x\}\) 在 \(k(a)[x,y]\) 中生成 \((x,y)\)，而在 \(a=0\) 时生成 \((y-x)\)，零点从一点变成一条直线。按参数分情形的系统方法留待附录B。

\subsection{计算步骤与结果验证}
\label{b3:subsec:0804:02}

一份可重复的计算至少应写出系数域、变量及其顺序、单项式序、原生成元，以及输出多项式的精确系数。若使用有限扩域，应写出定义该域的不可约多项式；若使用有理函数域，应区分参数和多项式变量。例如，\(k(a)[x]\) 中可用 \(a\) 除，而 \(k[a,x]\) 中一般不可以，这会改变理想成员问题本身。

\ProseParagraphBreak
验证 Gröbner 基时可把任务分成三个有限检查。给原生成元 \(F=(f_i)\) 和候选基 \(G=(g_j)\)，先检查多项式恒等式
\[
g_j=\sum_i a_{ij}f_i,\qquad f_i=\sum_j b_{ji}g_j,
\]
证明两者生成同一理想；再检查每一对 \(g_i,g_j\) 的 \(S\)-多项式有\cref{b3:thm:buchberger-criterion}要求的有界表示，或按已指定的除法次序确实得到零余式。前两项不能由第三项取代：\(\{x\}\) 自身是 Gröbner 基，却不是 \((x^2)\) 的 Gröbner 基。

\ProseParagraphBreak
理想成员的证书可以更短。证明 \(f\in(F)\)，只需给出并核验 \(f=\sum a_if_i\)，无须重新验证一次完整 Gröbner 基计算；证明 \(f\notin(F)\)，则可以先验证 \(G\) 的 Gröbner 基资格，再展示 \(f\) 的非零正规形。保存计算过程中的表示向量，正是为了在结果出来后提取这些较短的恒等式。它们能由另一套程序或人工逐项检查，与发现结果所用的实现相互独立。

\subsection{复杂度、系数膨胀与单项式序}
\label{b3:subsec:0804:03}

Buchberger 算法的终止性来自首单项式理想不可能形成无限严格升链。这个论证没有给出“只需少量步骤”的保证，也没有限制途中多项式的次数、项数和系数长度。即使最后得到很短的答案，中间表达式也可能很大。

\ProseParagraphBreak
消去本身便可能提高次数。在 \(k[x_1,\ldots,x_n]\) 中考虑
\[
I=(x_2-x_1^2,x_3-x_2^2,\ldots,x_n-x_{n-1}^2).
\]
逐次代入表明商环同构于 \(k[x_1]\)，且
\[
I\cap k[x_1,x_n]=(x_n-x_1^{2^{n-1}}).
\]
后一等式也可直接检验：以 \(x_n-x_1^{2^{n-1}}\) 作一元首一除法，余式只含 \(x_1\)，代入后为零便迫使余式为零。原方程全是二次式，消去后的关系却有指数增长的次数。这是具体输出变大的例子，并不声称所有单项式序或所有表示方法都具有相同复杂度。

\ProseParagraphBreak
系数也可能增长。对方程 \(x_1=2\)、\(x_{i+1}=x_i^2\) 逐次消去，得到整数 \(x_n=2^{2^{n-1}}\)，其二进制长度为 \(2^{n-1}+1\)。有理数运算虽然精确，存储分子、分母和做整数运算的成本仍须计算。单项式序还会改变中间基及最后基的形状：字典序适合消去，但不总是最快；先用另一种全局序得到商代数，再在线性代数中换序，是下一节要解释的做法。附录B再讨论具体实现，正文只以已经证明终止的算法承担存在性与有效性。

\subsection{数值近似与精确理想判定}
\label{b3:subsec:0804:04}

在 \(\mathbb R[x]\) 中，\(I_\varepsilon=(x,x-\varepsilon)\) 当 \(\varepsilon\ne0\) 时等于整个环，因为
\[
1=\varepsilon^{-1}x-\varepsilon^{-1}(x-\varepsilon).
\]
当 \(\varepsilon=0\) 时却等于 \((x)\)。若把足够小的非零系数统一舍为零，便会把一个无公共零点的系统误判为有零点的系统。这里出错的不是近似精度还不够高，而是“系数等于零”这个精确问题没有被近似数据决定。

\ProseParagraphBreak
同样，在 \(\mathbb C[x]\) 中 \(x^2\) 与 \(x^2-\varepsilon\) 的系数可以任意接近，但前者有一个二重根，后者在 \(\varepsilon\ne0\) 时有两个不同的根。数值计算适合近似根的位置，精确理想运算则要判断恒等式、理想相等和重数。两者可以相互帮助：例如先猜测有理系数关系，再用精确运算验证；但未经验证的近似零余式不能作为理想成员关系的证明。

% !TeX root = ../../Book3.tex
% 纲要标题：### 8.5 零维方程组与商代数
% 纲要位置：工作记录/计划纲要.md，第 2068 行。
% 纲要细目（写作范围）：
% * 标准单项式给出的商代数基与乘法矩阵
% * 有限维商代数的维数、重数与约化性；不能把向量空间维数一概当作不同几何点数
% * FGLM 换序的动机及零维假设；具体算法转附录 B
% * 三角形式与有理单变量表示简介；额外假设和退化分支逐项说明
% * 乘法矩阵的共同特征值与约化性在正文证明；完整求解算例和非约化对照
% ---


\section{零维方程组与商代数}
\label{b3:sec:0805}


有限维商代数把多项式关系转成有限个矩阵的关系。标准单项式给出坐标，乘法矩阵保存全部代数运算；它们既能辅助求解，也能记录不同零点个数所不能反映的幂零结构。本节同时说明单变量表示需要哪些附加条件。

\subsection{商代数的标准单项式基与乘法矩阵}
\label{b3:subsec:0805:01}

设 \(R=k[x_1,\ldots,x_n]\)、\(I\subsetneq R\)，并固定一个全局单项式序。正规形把每个剩余类写成标准单项式的线性组合。如果标准单项式只有有限多个，所有乘法也就能由有限个矩阵表示；求解方程的问题由此获得另一种计算方法。

\begin{proposition}\label{b3:prop:standard-basis-finite}
设 \(k\) 为域，\(R=k[x_1,\ldots,x_n]\)，\(I\subsetneq R\)，并固定全局单项式序。不属于 \(\operatorname{in}(I)\) 的标准单项式的剩余类构成 \(A=R/I\) 的 \(k\) 基。下列条件等价：
\begin{enumerate}
\item \(\dim_k A<\infty\)；
\item 标准单项式只有有限多个；
\item 对每个 \(i\)，存在 \(d_i\ge1\) 使 \(x_i^{d_i}\in\operatorname{in}(I)\)。
\end{enumerate}
\end{proposition}
\begin{proof}
用 Gröbner 基除法得到的余式只含标准单项式，故它们生成商空间。若其非零线性组合属于 \(I\)，首单项式便同时属于且不属于 \(\operatorname{in}(I)\)，矛盾。

由基的性质，前两项等价。若对某个 \(i\) 没有这样的幂，则 \(1,x_i,x_i^2,\ldots\) 都是标准单项式，违反有限性。反之，第三项使任何标准单项式的 \(x_i\) 指数均小于 \(d_i\)，故总数至多 \(d_1\cdots d_n\)。
\end{proof}

本节把有限维商代数称为零维商代数。它与素理想链意义的零维一致：有限维代数是 Artin 环，素理想都极大；反之，有限型 \(k\) 代数是 Noether 环，若所有素理想都极大，则由\cref{b3:thm:noether-dimension-zero}和\cref{b3:thm:artin-finite-length}知它有有限合成列，每个简单因子是某个极大理想的剩余域。\Cref{b3:cor:finite-type-residue-fields}保证这些剩余域在 \(k\) 上有限维，逐个相加便得 \(\dim_k A<\infty\)。第9章将系统定义和讨论 Krull 维数。

\ProseParagraphBreak
设标准单项式基为 \(B=(b_1,\ldots,b_D)\)。对 \(a\in A\)，记乘法映射 \(m_a:A\to A\)、\(v\mapsto av\) 在此基下的矩阵为 \(T_a\)，称为\term[chengfa-juzhen]{乘法矩阵}[multiplication matrix]。其第 \(j\) 列就是 \(ab_j\) 的正规形系数列。因此 \(T_aT_b=T_{ab}\)，各 \(T_{x_i}\) 两两交换，且对任意 \(f\in R\) 有
\[
T_{\bar f}=f(T_{x_1},\ldots,T_{x_n}).
\]
映射 \(a\mapsto T_a\) 是单射，因为 \(m_a(1)=a\)。特别地，\(f\in I\) 当且仅当这个矩阵为零；\(a\) 为单位当且仅当 \(T_a\) 可逆，后者的逆作用于 \(1\) 就给出 \(a^{-1}\)。

\ProseParagraphBreak
对 \(I=(x-y^2,y^3-1)\)，取基 \((1,y,y^2)\)，则
\[
T_y=\begin{pmatrix}0&0&1\\1&0&0\\0&1&0\end{pmatrix},
\qquad T_x=T_y^2=\begin{pmatrix}0&1&0\\0&0&1\\1&0&0\end{pmatrix}.
\]
因此 \(x^2=y\)、\(x^3=1\) 都可以直接从矩阵乘法读出。Sturmfels 在\cite[Section 2.3, pp. 17--19]{Sturmfels2002PolynomialEquations}中进一步讨论了乘法矩阵与方程求解的联系。

\begin{theorem}[乘法矩阵与共同特征值]\label{b3:thm:multiplication-joint-eigenvalues}
设 \(k\) 为代数闭域，\(I\subsetneq k[x_1,\ldots,x_n]\)，且 \(A=k[x_1,\ldots,x_n]/I\) 的 \(k\) 维数为有限正整数 \(D\)。选定 \(A\) 的一组 \(k\) 基，记由 \(\overline{x_i}\) 乘法给出的矩阵为 \(T_i\)（\(1\le i\le n\)）。点 \(a=(a_1,\ldots,a_n)\in k^n\) 是 \(I\) 的零点，当且仅当存在 \(0\ne v\in k^D\)，使
\[
T_i v=a_i v\qquad(1\le i\le n).
\]
此外，这些矩阵能同时对角化，当且仅当 \(A\) 约化。
\end{theorem}
\begin{proof}
若有这样的 \(v\)，任取 \(f\in I\)，乘法表示给出 \(f(T_1,\ldots,T_n)=0\)，所以 \(f(a)\cdot v=0\)。因 \(v\ne0\)，有 \(f(a)=0\)，故 \(a\) 是零点。

反过来，取零点 \(a\) 对应的极大理想 \(\mathfrak m_a\)。由\cref{b3:thm:artin-decomposition}，\(A\) 是各局部 Artin 因子的直积。在因子 \(A_{\mathfrak m_a}\) 中，极大理想 \(\mathfrak n\) 幂零；若 \(\mathfrak n=0\)，取 \(v=1\)，否则在最后一个非零幂 \(\mathfrak n^s\) 中取非零 \(v\)。两种情形都有 \(\mathfrak n v=0\)。因 \(x_i-a_i\in\mathfrak n\)，得到 \(x_iv=a_iv\)。将它放在这个因子中，其余因子取零，就得到 \(A\) 中的共同特征向量。

若 \(A\) 约化，局部因子的极大理想既幂零又只能为零。其剩余域为 \(k\)，因为它是代数闭域的有限扩张。因此 \(A\cong k^D\)，各因子的单位向量给出全部乘法矩阵的共同对角基。反之，若坐标乘法矩阵同时对角化，则每个 \(T_b\) 都是它们的多项式，也在同一基下对角化。若 \(b\) 幂零，\(T_b\) 是幂零对角矩阵，故为零；再用 \(m_b(1)=b\) 得 \(b=0\)。
\end{proof}

上述结论可对照\cite[Theorem 2.6 and Corollary 2.7, pp. 17--18]{Sturmfels2002PolynomialEquations}。同一个非零向量必须同时满足全部特征方程；分别列出各矩阵的特征值，再任意拼成坐标，会加入本来不存在的零点。

\subsection{商代数的维数、重数与约化性}
\label{b3:subsec:0805:02}

商空间的维数计入了剩余域的次数及幂零元所造成的重数。为看清这两项贡献，先在任意域 \(k\) 上写出 Artin 分解
\[
A\simeq\prod_{i=1}^r A_{\mathfrak m_i},
\qquad L_i=A/\mathfrak m_i.
\]
由\cref{b3:thm:artin-decomposition}，这些正是全部局部因子。设 \(\ell_i\) 为 \(A_{\mathfrak m_i}\) 在自身上作为模的长度。局部环的简单模只有 \(L_i\)，沿合成列计算 \(k\) 维数得到
\begin{equation}\label{b3:eq:zero-dimensional-length}
\dim_k A=\sum_{i=1}^r\ell_i[L_i:k].
\end{equation}
因子 \(\ell_i\) 称为这一闭点的重数。当 \(k\) 代数闭时，\(L_i=k\)，故 \(D=\sum_i\ell_i\)；此时不同点数等于 \(D\) 当且仅当每个局部因子都是 \(k\)，也就是 \(A\) 约化。

\ProseParagraphBreak
这个等价可以不借助任何求根公式来证明。局部 Artin 环的极大理想是幂零理想，所以约化时它只能为零，局部因子成为域；反之，域的有限直积没有非零幂零元。若 \(k\) 不代数闭，则即使 \(A\) 约化，各 \([L_i:k]\) 仍可能大于一。例如 \(\mathbb R[x]/(x^2+1)\) 维数为二，有一个闭点而没有实坐标点。

\ProseParagraphBreak
在前述三维代数 \(k[x,y]/(x-y^2,y^3-1)\) 中，若 \(k\) 代数闭且特征不为三，\(y^3-1\) 有三个不同的根，代数为 \(k^3\)。若特征为三，则 \(y^3-1=(y-1)^3\)，代数同构于 \(k[\varepsilon]/(\varepsilon^3)\)，只剩一个重数为三的点。这两种情形有完全相同的标准单项式数和乘法矩阵公式，却有不同的约化性。

\ProseParagraphBreak
基域扩张也须留意可分性。令 \(k=\mathbb F_p(s)\)，则 \(A=k[t]/(t^p-s)\) 是维数为 \(p\) 的域：\(s\) 不是 \(k\) 中的 \(p\) 次幂，其不可约性已在\cref{b3:exm:inseparable-geometric-thickening}中通过最小多项式系数证明。但扩张到含 \(s^{1/p}\) 的域后，关系变成 \((t-s^{1/p})^p=0\)。所以原代数约化，不代表基域扩张后仍约化；谈几何点的数目时，不能略去这一差别。

\subsection{FGLM 换序与零维假设}
\label{b3:subsec:0805:03}

同一个零维理想在两种单项式序下的标准单项式可能不同，却给出同一个有限维向量空间。已知一种序下的 Gröbner 基后，可以用其正规形把任意单项式表示成 \(k^D\) 中的列向量，再按照新的单项式序寻找线性相关关系。这就是 FGLM 换序方法的基本思想。

\ProseParagraphBreak
具体地，已有若干线性无关的单项式剩余类时，考察它们乘各变量所得的候选单项式。若一个候选 \(m\) 的向量在先前较小单项式的线性生成空间中，写成
\[
\bar m=\sum_j c_j\bar b_j,
\]
便得到首单项式为 \(m\) 的关系 \(m-\sum_jc_jb_j\in I\)；其倍数无需再作为基候选。若向量独立，则把它加入新的基。所有向量均可由乘法矩阵计算。候选的组织方式、完整算法和终止性证明放在附录B；这里先说明为什么新问题能转为有限维线性代数。

\ProseParagraphBreak
例如旧序 \(x>y\) 的字典序给出基 \(\{x-y^2,y^3-1\}\)。改用 \(y>x\) 后，在旧商空间中 \(1,x,x^2\) 依次对应 \(1,y^2,y\)，因而独立；随后 \(x^3=1\)、\(y=x^2\) 给出新基 \(\{y-x^2,x^3-1\}\)。这次换序只须三维空间中的计算，不需要从原生成元重新完成所有 \(S\)-多项式约化。

\ProseParagraphBreak
零维条件提供的是有限的 \(D\) 和有限个 \(D\times D\) 乘法矩阵。对 \(k[x,y]/(y)\simeq k[x]\)，单项式 \(1,x,x^2,\ldots\) 永远线性无关，不能用有限维乘法矩阵表示这个商环。正维理想仍有有限 Gröbner 基，但那是另一种有限性。

\subsection{三角形式、有理单变量表示与退化分支}
\label{b3:subsec:0805:04}

若能找到一个元素 \(t\in A\) 使 \(1,t,\ldots,t^{D-1}\) 成为 \(k\) 基，那么每个坐标都能写成 \(x_i=q_i(t)\)，且 \(t\) 满足一个首一 \(D\) 次多项式 \(p\)。于是
\[
A\simeq k[T]/\bigl(p(T)\bigr),\qquad x_i=q_i(T).
\]
这是一种三角描述：先解一个一元方程，再恢复全部坐标。它需要 \(A\) 是单生成代数，不能仅从有限维性推出。

\begin{proposition}\label{b3:prop:finite-reduced-shape}
设 \(k\) 为代数闭域，\(R=k[x_1,\ldots,x_n]\)，\(I\subsetneq R\) 是根理想，且 \(A=R/I\) 的 \(k\) 维数为有限正整数 \(D\)。存在一个线性形式 \(t=\sum_{i=1}^n c_ix_i\)，在 \(I\) 的所有零点上取值两两不同。对任意这样的 \(t\)，存在无重根的首一 \(D\) 次多项式 \(p(T)\) 及次数小于 \(D\) 的 \(q_i(T)\)，使
\[
A\simeq k[T]/(p(T)),\qquad \overline{x_i}\longleftrightarrow q_i(T),\quad \bar t\longleftrightarrow T.
\]
\end{proposition}
\begin{proof}
由式\eqref{b3:eq:zero-dimensional-length}及约化性，\(A\simeq k^D\)，各因子对应点 \(P_1,\ldots,P_D\)。对每对不同点，等式 \(\sum c_i(P_{ji}-P_{hi})=0\) 定义系数空间中的一个真超平面。无限域上的有限个真超平面不能覆盖整个空间：这些非零线性多项式之积非零，而非零多项式不在无限域的全部点上消失，后一事实对变量数归纳即可。故可以选取避开这些超平面的 \(c\)。

令 \(\lambda_j=t(P_j)\)，取 \(p(T)=\prod_j(T-\lambda_j)\)。Lagrange 插值给出次数小于 \(D\) 的 \(q_i\)，满足 \(q_i(\lambda_j)=P_{ji}\)。映射 \(k[T]\to k^D\)、\(f\mapsto\bigl(f(\lambda_j)\bigr)_j\) 满射，核为 \((p)\)，并把 \(q_i\) 送到第 \(i\) 个坐标函数，这就给出所需同构。
\end{proof}

也可把坐标写为 \(x_i=q_i(T)/q_0(T)\)。若 \(\gcd(q_0,p)=1\)，Euclid 算法给出 \(uq_0+vp=1\)，从而 \(q_0\) 在 \(k[T]/(p)\) 中可逆；这种写法称为\term[youli-danbianliang-biaoshi]{有理单变量表示}[rational univariate representation]，没有增加或丢失根。如果 \(q_0\) 与 \(p\) 有公因子，分母为零的根不能直接代入；在无重根情形应先按 \(\gcd(p,q_0)\) 与其互素补因子分支，在含重根情形则必须保留相应幂次，另行计算局部因子。

\ProseParagraphBreak
最后列出两种不能忽略的情形。有限域上的线性形式未必分离全部有理点，例如 \(\mathbb F_2^2\) 的四个点不可能通过一个 \(\mathbb F_2\) 值的函数取得四个不同值。含幂零元的代数也未必单生成：\(k[x,y]/(x,y)^2\) 维数为三，每个元素都形如 \(c+n\)、\(n^2=0\)，其生成的子代数最多二维。一般零维计算因此要允许域扩张、多个三角分支或更一般的有限代数表示；不能把一个漂亮的单变量形式当作无条件存在的输出。

\subsection{从乘法矩阵恢复解与重数}
\label{b3:subsec:0805:05}

现在把上述步骤用于一个完整方程组。工作系数域取 \(\mathbb Q\)，要求在 \(\overline{\mathbb Q}\) 中求出
\[
x-y^2=0,\qquad y^3-y=0.
\]
在字典序 \(x\succ y\) 下，两个首单项式为 \(x,y^3\)，互素；其唯一的 S 多项式约化为零，故原生成组就是 Gröbner 基。标准单项式为 \(1,y,y^2\)，从而商代数维数为三。用关系 \(y^3=y\) 按列计算乘法，得到
\[
T_y=\begin{pmatrix}0&0&0\\1&0&1\\0&1&0\end{pmatrix},
\qquad
T_x=T_y^2=\begin{pmatrix}0&0&0\\0&1&0\\1&0&1\end{pmatrix}.
\]
\(T_y\) 的特征多项式为 \(T(T-1)(T+1)\)，三个根不同。分别取特征向量
\[
v_0=(1,0,-1)^{\mathsf T},\quad
v_1=(0,1,1)^{\mathsf T},\quad
v_{-1}=(0,1,-1)^{\mathsf T}.
\]
它们在 \(T_x\) 下的特征值依次为 \(0,1,1\)。由\cref{b3:thm:multiplication-joint-eigenvalues}，全部零点正是
\[
(0,0),\quad(1,1),\quad(1,-1).
\]
各点代回两个方程都为零；共同对角化又证明商代数约化，所以这三个点的重数均为一。若只把 \(T_x\) 的特征值零与 \(T_y\) 的特征值一拼接，所得 \((0,1)\) 就不能通过代回检查。

\ProseParagraphBreak
把第二个方程改为 \(y^2(y-1)=0\)，标准单项式仍是 \(1,y,y^2\)，但这次
\[
T_y=\begin{pmatrix}0&0&0\\1&0&0\\0&1&1\end{pmatrix},
\qquad \det(TI-T_y)=T^2(T-1).
\]
由 \(x=y^2\)，不同零点只有 \((0,0)\) 与 \((1,1)\)。要解释重数，使用互素因子给出的中国剩余分解
\[
\mathbb Q[y]/\bigl(y^2(y-1)\bigr)
\cong \mathbb Q[y]/(y^2)\times\mathbb Q[y]/(y-1).
\]
两因子的维数分别为二与一，因此这两个点的重数分别为二与一。\(y^2-y\) 在商中非零而平方为零，正说明第一个因子还保存了单个点所看不见的幂零结构。

\ProseParagraphBreak
这两次计算使用相同的三维线性空间，但答案包含不同的信息。乘法矩阵用于恢复相容的坐标，局部因子用于解释重数；最后的代回和维数比较各有作用。对于不能精确写出根的实例，可以保留一元多项式及坐标多项式作为精确答案；实根的计数与隔离见\appref{b3:app:G}。

% !TeX root = ../../Book3.tex
% 纲要标题：### 8.6 结式与消去计算
% * Sylvester 矩阵、公共因子判据与结式的 Bézout 表示
% * 结式与乘法行列式、范数及首一多项式判别式
% * 参数特殊化时的次数下降及例外分支
% * 消去、回代和全部解的精确核验；一般多变量及稀疏结式暂不展开
% ---

\section{结式与消去计算}
\label{b3:sec:0806}

消去一个变量时，有时不必先计算整个理想的 Gröbner 基。两个一元多项式是否在系数域的代数闭包中有公共根，可以由一个行列式判断；把其系数看成其余变量的多项式，就得到消去关系。本节建立这一方法，并说明消去所得的候选解为何仍须回代。这一矩阵构造及公共因子、消去性质可参见 Cox、Little 与 O'Shea 的《Using Algebraic Geometry》\cite[Chapter 3, \S 1, pp. 71--72]{CoxLittleOShea1998UsingAG}。

\subsection{Sylvester 矩阵与公共因子}
\label{b3:subsec:0806:01}

设 \(k\) 为域，\(f=a_mX^m+\cdots+a_0\)、\(g=b_nX^n+\cdots+b_0\) 属于 \(k[X]\)，其中 \(a_m b_n\ne0\)、\(m,n>0\)。记 \(k[X]_{<d}\) 为次数小于 \(d\) 的多项式空间。寻找次数受限的关系 \(uf+vg=0\)，就是研究线性映射
\[
\Phi_{f,g}:k[X]_{<n}\oplus k[X]_{<m}
\longrightarrow k[X]_{<m+n},\qquad (u,v)\longmapsto uf+vg.
\]
两侧维数相同，所以非零核可以由行列式检测。

\begin{definition}\label{b3:def:resultant}
设 \(k\) 为域，\(f,g\in k[X]\) 的次数分别为正整数 \(m,n\)。考虑 \(k\) 线性映射
\[
\Phi_{f,g}:k[X]_{<n}\oplus k[X]_{<m}\longrightarrow k[X]_{<m+n},
\qquad (u,v)\longmapsto uf+vg,
\]
其中 \(k[X]_{<d}\) 表示次数小于 \(d\) 的多项式组成的空间。在各空间中按幂次从高到低取基，并先列第一个直和因子的基。所得矩阵的列依次为
\[
X^{n-1}f,\ldots,f,\ X^{m-1}g,\ldots,g
\]
在 \((X^{m+n-1},\ldots,1)\) 下的系数列，称为 \(f,g\) 的\term[Sylvester-juzhen]{Sylvester 矩阵}[Sylvester matrix]，\footnote{西尔维斯特（James Joseph Sylvester，1814-1897），英国数学家，研究矩阵、行列式与不变量理论。}记为 \(S(f,g)\)。其行列式
\[
\operatorname{Res}_X(f,g)\coloneqq\det S(f,g)
\]
称为关于 \(X\) 的\term[jieshi]{结式}[resultant]。
\end{definition}

例如 \(f=aX+b\)、\(g=cX+d\) 时，结式为 \(ad-bc\)。交换两个多项式，相当于交换长分别为 \(n,m\) 的两组列，所以
\(\operatorname{Res}_X(g,f)=(-1)^{mn}\operatorname{Res}_X(f,g)\)。固定次数后，结式是系数的整系数多项式；因此同一矩阵也可在任意交换系数环上计算。

\begin{proposition}\label{b3:prop:resultant-common-factor}
设 \(k\) 为域，\(f,g\in k[X]\) 分别具有正次数 \(m,n\)。则 \(\operatorname{Res}_X(f,g)=0\)，当且仅当 \(f,g\) 有正次数公因子，也当且仅当它们在 \(\overline k\) 中有公共根。
\end{proposition}
\begin{proof}
若有正次数公因子 \(h\)，取 \(u=g/h\)、\(v=-f/h\)，便得 \(uf+vg=0\)，且 \(\deg u<n\)、\(\deg v<m\)。所以 \(\Phi_{f,g}\) 有非零核，行列式为零。

若 \(f,g\) 互素而 \(uf+vg=0\)，则 \(f\mid vg\)。由 Bézout 等式\footnote{贝祖（Étienne Bézout，1730-1783），法国数学家，研究代数方程与消去理论。}与互素性得 \(f\mid v\)；次数限制迫使 \(v=0\)，进而 \(u=0\)。所以映射单射，行列式非零。正次数公因子在代数闭包中有根；反过来，若互素，多项式 Bézout 等式 \(af+bg=1\) 排除任何扩域中的公共根。
\end{proof}

\begin{proposition}\label{b3:prop:resultant-bezout}
设 \(R\) 为交换含幺环，\(f,g\in R[X]\) 的次数分别不超过固定的正整数 \(m,n\)，按这两个次数补零构造 \(S(f,g)\)，并记其行列式为 \(D\)。存在 \(u,v\in R[X]\)，满足 \(\deg u<n\)、\(\deg v<m\) 及 \(uf+vg=D\)。这里 \(u,v\) 的系数均为 \(f,g\) 系数的整系数多项式。
\end{proposition}
\begin{proof}
设 \(e\) 是目标空间中常数一的系数列，即最后一个标准基向量。伴随矩阵恒等式给出
\(S\operatorname{adj}(S)e=De\)。把向量 \(\operatorname{adj}(S)e\) 的前 \(n\) 项和后 \(m\) 项解释为 \(u,v\) 的系数，即得等式。伴随矩阵各项是子式，故还得到最后的系数说明。论证没有除法，也不要求 \(R\) 为域。
\end{proof}

\subsection{结式、乘法行列式与判别式}
\label{b3:subsec:0806:02}

\begin{proposition}\label{b3:prop:resultant-norm}
设 \(k\) 为域，\(f\in k[X]\) 首一且次数为 \(m>0\)，\(g\in k[X]\)。令 \(\bar g\) 为 \(g\) 在有限维 \(k\)-代数 \(A=k[X]/(f)\) 中的像，记乘 \(\bar g\) 的矩阵为 \(T_g\)。则
\[
\operatorname{Res}_X(f,g)=\det T_g
=\prod_{i=1}^m g(\alpha_i),
\]
其中 \(\alpha_1,\ldots,\alpha_m\) 是 \(f\) 在 \(\overline k\) 中按重数排列的根；对常数 \(g=c\)，约定结式为 \(c^m\)。若 \(f\) 的首项系数为任意 \(a_m\ne0\)，且 \(\deg g=n>0\)，则相应根公式为 \(a_m^n\prod_i g(\alpha_i)\)。
\end{proposition}
\begin{proof}
先设 \(g\) 次数为 \(n>0\)，且 \(f\) 首一。目标空间的有序基
\[
X^{n-1}f,\ldots,f,\ X^{m-1},\ldots,1
\]
到原幂基的变换矩阵为对角元全一的三角矩阵。关于新基，\(\Phi_{f,g}\) 的矩阵分块为 \(\left(\begin{smallmatrix}1_n&Q\\0&T_g\end{smallmatrix}\right)\)：后 \(m\) 个坐标正是除以 \(f\) 的余式。因此两边行列式相等。

变基到 \(\overline k\) 不改变行列式。置 \(h_j=\prod_{i=1}^j(X-\alpha_i)\)，则商代数中有理想滤过 \((1)\supset(h_1)\supset\cdots\supset(h_m)=0\)。每个 \((h_{j-1})/(h_j)\) 一维，且乘 \(X\) 的作用为 \(\alpha_j\)。适应该滤过取基，乘 \(g\) 的矩阵为三角矩阵，对角元为 \(g(\alpha_j)\)，从而得到乘积公式。这也包含重根情形。常数情形的乘法矩阵为 \(c1_m\)。最后，将一般的 \(f\) 换成 \(f/a_m\)，Sylvester 矩阵的前 \(n\) 列各提出 \(a_m\)，便得一般公式。
\end{proof}

这里的行列式是有限维代数中乘法的范数；即使 \(A\) 有幂零元也有意义。当 \(f\) 不可约时，它就是第一卷域扩张范数的计算。因而同一个数既检测公共根，也检测 \(\bar g\) 在 \(A\) 中是否为单位。

\begin{corollary}\label{b3:cor:resultant-discriminant}
设 \(k\) 为域，\(f\in k[X]\) 首一、次数为 \(m>0\)，其在代数闭包 \(\overline k\) 中的根按重数记为 \(\alpha_1,\ldots,\alpha_m\)。定义
\(\operatorname{disc}(f)=\prod_{i<j}(\alpha_i-\alpha_j)^2\)。则
\[
\operatorname{disc}(f)=(-1)^{m(m-1)/2}\operatorname{Res}_X(f,f').
\]
其中 \(f'=0\) 时结式取零。判别式属于 \(k\)，并且非零当且仅当 \(f\) 无重根。
\end{corollary}
\begin{proof}
若有重根，\(f,f'\) 有公共根，等式两侧均为零。若无重根，则 \(f'(\alpha_i)=\prod_{j\ne i}(\alpha_i-\alpha_j)\)；相乘后，每一对指标贡献一个负号与一个平方。用\cref{b3:prop:resultant-norm}即得公式。右端由 \(k\) 中系数的行列式组成，故属于 \(k\)。\(m=1\) 时空乘积与 \(\operatorname{Res}(f,1)\) 都为一。
\end{proof}

例如在 \(k[X]/(X^2+aX+b)\) 的基 \((1,X)\) 下，
\[
T_X=\begin{pmatrix}0&-b\\1&-a\end{pmatrix},\qquad
T_{f'}=2T_X+a1_2=\begin{pmatrix}a&-2b\\2&-a\end{pmatrix}.
\]
其行列式为 \(4b-a^2\)，故判别式为 \(a^2-4b\)。这个计算在特征二时仍有效，此时无重根条件化为 \(a\ne0\)。

\subsection{参数特殊化与消去条件}
\label{b3:subsec:0806:03}

设 \(f,g\in k[t_1,\ldots,t_s][X]\)，按其关于 \(X\) 的次数构造结式 \(D(t)\)。由\cref{b3:prop:resultant-bezout}，任何共同零点 \((t,x)\) 都满足 \(D(t)=0\)。另一方面，只有在特殊化后两个首项系数仍非零时，才能直接用\cref{b3:prop:resultant-common-factor}从 \(D(t)=0\) 推出存在公共根。

\ProseParagraphBreak
例如在 \(\mathbb Q[t,X]\) 中取 \(f=tX+1\)、\(g=tX+2\)，则结式为 \(t\)。候选参数 \(t=0\) 使两个次数都下降，所得常数方程 \(1=0\)、\(2=0\) 无解。实际上 \(g-f=1\)，原理想是单位理想。结式为零在这里反映了固定尺寸矩阵的退化，不能证明仿射方程存在解。

\ProseParagraphBreak
因此参数计算须分开处理首项系数非零的部分与其零点部分。前一部分使用公共根判据；后一部分把参数条件代入原生成元，重算实际次数和公共因子。若某个多项式变成零多项式，只剩另一个方程；若出现非零常数，则该分支无解。这些操作也解释了为什么消去关系通常是必要条件，而非不经检查的完整解集。

\subsection{消去、回代与全部解的验证}
\label{b3:subsec:0806:04}

在 \(\mathbb Q[x,y]\) 中考虑
\[
f=xy-1=0,\qquad g=x^2+y^2-5=0.
\]
将它们视为关于 \(x\) 的一次与二次多项式，得到
\[
\operatorname{Res}_x(f,g)
=\det\begin{pmatrix}y&0&1\\-1&y&0\\0&-1&y^2-5\end{pmatrix}
=y^4-5y^2+1.
\]
这一结果还能由恒等式直接核验：
\[
y^2g-(xy+1)f=y^4-5y^2+1.
\]
结式的任何根 \(y\) 都非零，故可由第一个方程唯一恢复 \(x=1/y\)。代回第二个方程得到 \(g=(y^4-5y^2+1)/y^2=0\)，说明每个候选确实能提升，不存在遗漏或多余分支。

\ProseParagraphBreak
令 \(z=y^2\)，则 \(z^2-5z+1=0\)，从而
\[
y=\pm\sqrt{\dfrac{5+\sqrt{21}}{2}}
\quad\text{或}\quad
y=\pm\sqrt{\dfrac{5-\sqrt{21}}{2}},\qquad x=1/y.
\]
两个根号内的数都为正，故四个解均为实点，且彼此不同。若只要求精确代数表示，保留 \(y^4-5y^2+1=0\)、\(x=1/y\) 已足够。

\ProseParagraphBreak
本例还可比较两种消去方法。由结式关系得到
\[
(f,g)=(x+y^3-5y,\ y^4-5y^2+1).
\]
确实，置 \(h=x+y^3-5y\)、\(r=y^4-5y^2+1\)，有 \(h=yg-xf\)、\(r=y^2g-(xy+1)f\)，反向有 \(f=yh-r\)、\(g=xh-(y^2-5)f\)。右端是字典序 \(x\succ y\) 的 Gröbner 基，标准单项式为 \(1,y,y^2,y^3\)。所以商代数维数为四，与已找出的四个不同点一致。结式给出了消去多项式，双向生成元表示和标准单项式则进一步核验了原理想的完整结构。


\begin{chaptersummary}
全局单项式序提供严格下降，Dickson 引理提供有限性。二者结合，既保证多元除法和补充生成元的过程终止，也使理想的全部首项能够由有限个首项控制。Buchberger 判据把最高项抵消归结到有限对的有界表示，约化基则为固定次序下的理想给出唯一描述。

\ProseParagraphBreak
消去把投影和理想运算化为扩展多项式环中的计算。几何投影通常只得到闭包；参数分母、特殊取值和不能提升的点仍须另行检查。同时变号的不变环则展示了另一种用法：先由作用找齐不变多项式，再以核理想描述生成元之间的全部关系。两向生成元表示证明理想相等，S 多项式的约化证明 Gröbner 性；这两项核验不能互相代替。

\ProseParagraphBreak
零维商代数把问题带入有限维线性代数。标准单项式数计算代数维数，局部因子的长度记录重数，乘法矩阵的共同特征信息恢复相容的坐标。结式则以行列式提供消去关系；参数导致次数下降时，还须重新检查原方程。后续维数与 Hilbert 函数将把有限计数推广为多项式增长；本章的精确首项方法也将用于泛自由性的证明。
\end{chaptersummary}

\BookExercises

\begin{exercise}\label{b3:ex:08-three-orders}
对变量顺序 \(x\succ y\succ z\)，分别按字典序、分次字典序和分次逆字典序将 \(x^4,x^2z,xy^2,yz^3\) 从大到小排列。
\end{exercise}
\begin{solution}
字典序先看 \(x\) 的指数，得到 \(x^4\succ x^2z\succ xy^2\succ yz^3\)。分次字典序先分成四次与三次两组，再看字典序，得到 \(x^4\succ yz^3\succ x^2z\succ xy^2\)。分次逆字典序中，三次两项的最后一个不同坐标是 \(z\) 的指数，较小者较大，故得到 \(x^4\succ yz^3\succ xy^2\succ x^2z\)。
\end{solution}

\begin{exercise}\label{b3:ex:08-weighted-order}
设 \(n\ge1\) 为整数，给定正整数权 \(w_1,\ldots,w_n\)，先比较 \(\sum_iw_i\alpha_i\)，相同再按字典序比较。证明这给出全局单项式序。若允许负权，这个证明在哪一步失效？
\end{exercise}
\begin{solution}
权重值相加，差向量不变，所以乘法相容。任意非空指数集先有最小权重；因各 \(w_i>0\)，不超过某个固定权重的指数只有有限个，从中按字典序取最小者便得良序。若某个 \(w_i<0\)，则 \(1\succ x_i\succ x_i^2\succ\cdots\)，良序性失效。这不是只须调整同权时的比较即可修复的问题。
\end{solution}

\begin{exercise}\label{b3:ex:08-monomial-minimal}
设 \(k\) 为域。在 \(k[x,y]\) 中求 \(J=(x^3y,x^2y^2,xy^3,x^4y)\) 的极小单项式生成组，并判断 \(x^2y^3,x^3\) 是否属于 \(J\)。再说明无限变量环中的理想 \((x_1,x_2,\ldots)\) 不存在有限生成组。
\end{exercise}
\begin{solution}
\(x^4y\) 被 \(x^3y\) 整除，其余三项互不整除，故极小组为 \(x^3y,x^2y^2,xy^3\)。\(x^2y^3\) 被后二者整除，属于 \(J\)；\(x^3\) 不含 \(y\)，不属于 \(J\)。无限变量环中，若上述理想由有限个多项式生成，这些多项式只涉及有限个变量。它们都无常数项，把涉及的变量全置零后，生成理想映为零，但某个未涉及的 \(x_j\) 仍非零，矛盾。
\end{solution}

\begin{exercise}\label{b3:ex:08-division-choice}
设 \(k\) 为域。在 \(k[x,y]\) 的字典序 \(x\succ y\) 下，用 \(g_1=xy-1,g_2=y^2-y\) 的两种排列分别除 \(f=xy^2\)。找出应补入的关系，并给出原理想的约化 Gröbner 基。
\end{exercise}
\begin{solution}
先用 \(g_1\) 得 \(f=yg_1+y\)，余式为 \(y\)。先用 \(g_2\)，得到 \(f=xg_2+xy=xg_2+g_1+1\)，余式为一。差 \(y-1=(1-y)g_1+xg_2\) 在原理想中。又
\[
x-1=g_1-x(y-1),\qquad
g_1=y(x-1)+(y-1),\qquad g_2=y(y-1).
\]
所以原理想是 \((x-1,y-1)\)。其首单项式 \(x,y\) 互素，两多项式首一且尾项均为常数，乘积判据说明它们构成约化 Gröbner 基。此时 \(f\) 的正规式是一。
\end{solution}

\begin{exercise}\label{b3:ex:08-nonreduced-buchberger}
设 \(k\) 为域。在 \(k[x,y]\) 的字典序 \(x\succ y\) 下计算 \(I=(x^2-y,xy)\) 的约化 Gröbner 基和标准单项式，并在商环中计算 \(x^3,y^2\)。
\end{exercise}
\begin{solution}
初始 S 多项式为 \(y(x^2-y)-x(xy)=-y^2\)，故加入 \(y^2\)。三对 S 多项式分别为 \(-y^2\)、\(-y^3\)、零，均约化为零。因此
\[
G=\{x^2-y,xy,y^2\},\qquad \mathcal B=\{1,x,y\}.
\]
各尾项不可再约化，所以该基约化。商中 \(x^3=x(x^2-y)+xy=0\)，且 \(y^2=0\)，但 \(x^2=y\ne0\)，后一点由标准单项式独立性保证。
\end{solution}

\begin{exercise}\label{b3:ex:08-order-change-nilpotent}
设 \(k\) 为域。对 \(k[x,y]\) 中的理想 \(I=(x^2-y,xy)\)，改用字典序 \(y\succ x\)，求约化基，并写出两组标准单项式基之间的变换。
\end{exercise}
\begin{solution}
令 \(h=y-x^2\)。有 \(h=-(x^2-y)\)、\(x^3=x(x^2-y)+xy\)，反向 \(xy=xh+x^3\)，故 \(I=(h,x^3)\)。首单项式 \(y,x^3\) 互素，得到约化基 \(\{y-x^2,x^3\}\)。新标准单项式基是 \((1,x,x^2)\)，旧基为 \((1,x,y)\)，在商中第三项相同，所以按此顺序的换基矩阵是单位矩阵；它们所用的多项式代表却不同。
\end{solution}

\begin{exercise}\label{b3:ex:08-product-criterion-use}
设 \(k\) 为域。在 \(k[x,y]\) 的字典序 \(x\succ y\) 下，证明 \(\{x^2+y,y^2+1\}\) 是 Gröbner 基，并求 \(x^4+y^2\) 的正规式。结果随 \(\operatorname{char}k=2\) 有何变化？
\end{exercise}
\begin{solution}
首单项式 \(x^2,y^2\) 互素，乘积判据适用。商中 \(x^2=-y\)、\(y^2=-1\)，故 \(x^4+y^2=2y^2=-2\)，常数已为标准余式。特征二时正规式为零，原多项式属于理想；其他特征下不属于。两种情况下所给基仍是 Gröbner 基。
\end{solution}

\begin{exercise}\label{b3:ex:08-ideal-equality-insufficient}
设 \(k\) 为域，在 \(k[x,y]\) 中说明 \(\{x^2+y,xy+1\}\) 在字典序 \(x\succ y\) 下不是 Gröbner 基，尽管它当然生成自己的理想。给出一个属于该理想却对这组除法留下非零余式的多项式。
\end{exercise}
\begin{solution}
取 S 多项式 \(y(x^2+y)-x(xy+1)=y^2-x\)。其首单项式为 \(x\)，不能被 \(x^2\) 或 \(xy\) 整除，其他项 \(y^2\) 也不能，故余式就是非零的 \(y^2-x\)。这违反 Gröbner 基的成员判据。理想相等只检查生成关系，不能取代首项控制。
\end{solution}

\begin{exercise}\label{b3:ex:08-triangular-normal}
设 \(k\) 为域，\(p(y)\in k[y]\) 首一且次数 \(d\ge1\)，\(h(y)\in k[y]\)。在字典序 \(x\succ y\) 下，证明 \(G=\bigl\{x-h(y),p(y)\bigr\}\) 是 Gröbner 基，并写出它约化的充要条件及任意 \(f\) 的正规式求法。
\end{exercise}
\begin{solution}
首单项式为 \(x,y^d\)，互素，故为 Gröbner 基。第一元素的尾项不可被 \(y^d\) 整除恰当且仅当 \(\deg h<d\)，零多项式也允许；第二元素只有较低的 \(y\) 次幂，所以这就是约化的充要条件。差 \(f(x,y)-f\bigl(h(y),y\bigr)\) 被 \(x-h(y)\) 整除，这由每个 \(x^j-h(y)^j\) 的差幂公式得到。先代入 \(x=h(y)\)，再除以 \(p\)，余式只含 \(1,y,\ldots,y^{d-1}\)，因此为唯一正规式。
\end{solution}

\begin{exercise}\label{b3:ex:08-cusp-elimination}
设 \(k\) 为域。对参数 \(x=t^2,y=t^3\)，求 \(k[t,x,y]\) 中 \((x-t^2,y-t^3)\) 与 \(k[x,y]\) 的交，并说明全部几何点能否提升成参数点。
\end{exercise}
\begin{solution}
代入同态的核就是所求交。显然 \(y^2-x^3\) 在核中；对它按 \(y\) 作首一除法，任意余式形如 \(a(x)+y\cdot b(x)\)。代入后成为 \(a(t^2)+t^3\cdot b(t^2)\)，前者只有偶次幂，后者只有奇次幂，且各组内部指数互异，故只有 \(a=b=0\) 才能为零。因此交为 \((y^2-x^3)\)。对任一扩域中的点 \((a,b)\)，若 \(a\ne0\)，取 \(t=b/a\)，由 \(b^2=a^3\) 得 \(t^2=a,t^3=b\)；若 \(a=0\)，则 \(b=0\)，取 \(t=0\)。这里的提升来自具体公式，并非一般闭包定理。
\end{solution}

\begin{exercise}\label{b3:ex:08-projection-missing}
设 \(k\) 代数闭，\(X=Z_k(u-x,v-xy)\subseteq k^4\)。求向 \((u,v)\) 平面的投影像、消去理想及其零点集，并描述不能提升成 \(X\) 中点的那些平面点。
\end{exercise}
\begin{solution}
若 \(u\ne0\)，取 \(x=u,y=v/u\) 即可提升；若 \(u=0\)，则 \(x=0\)，只能有 \(v=0\)。因此像为 \(\{u\ne0\}\cup\bigl\{(0,0)\bigr\}\)。商环同构于 \(k[x,y]\)，其中 \(u\mapsto x,v\mapsto xy\)。不同单项式 \(u^iv^j\) 映成不同的 \(x^{i+j}y^j\)，故限制同态 \(k[u,v]\to k[x,y]\) 单射，消去理想为零。其零点集及投影闭包都是整个平面；遗漏的点恰为 \(u=0,v\ne0\)，并非某个完整不可约分量。
\end{solution}

\begin{exercise}\label{b3:ex:08-intersection-example}
设 \(k\) 为域。在 \(k[x,y]\) 中求 \((x^2,y)\cap(x,y^2)\)，并给出每个输出生成元属于两理想的表示。
\end{exercise}
\begin{solution}
单项式属于交，当且仅当同时被第一组和第二组中的某项整除，因而被某一对生成元的最小公倍式整除。四个最小公倍式为 \(x^2,x^2y^2,xy,y^2\)，删去冗余得交为 \((x^2,xy,y^2)\)。分别有 \(x^2=1\cdot x^2=x\cdot x\)、\(xy=x\cdot y=y\cdot x\)、\(y^2=y\cdot y=1\cdot y^2\)，即为两侧的成员表示。单项式成员判据同时证明这些生成元已经穷尽交理想。
\end{solution}

\begin{exercise}\label{b3:ex:08-colon-example}
设 \(k\) 为域，\(I=(x^3,x^2y,y^3)\subset k[x,y]\)。计算 \((I:x^2)\)、\((I:y)\) 及 \(\bigl(I:(x,y)\bigr)\)。
\end{exercise}
\begin{solution}
单项式理想除以单项式 \(m\) 的理想商由 \(u/\gcd(u,m)\) 生成，其中 \(u\) 遍历原生成元；逐个指数比较即可验证。于是
\[
(I:x^2)=(x,y),\quad (I:y)=(x^3,x^2,y^2)=(x^2,y^2).
\]
同理 \((I:x)=(x^2,xy,y^3)\)。与 \((x^2,y^2)\) 取交，使用两组最小公倍式，删去冗余后得到
\[
\bigl(I:(x,y)\bigr)=(x^2,xy^2,y^3).
\]
\end{solution}

\begin{exercise}\label{b3:ex:08-saturation-components}
设 \(k\) 为域，\(I=(xy)\subset k[x,y]\)。求 \((I:x^\infty)\)、\((I:y^\infty)\)、\(\bigl(I:(x,y)^\infty\bigr)\) 和 \(\bigl(I:(xy)^\infty\bigr)\)，并解释后两者为什么不同。
\end{exercise}
\begin{solution}
由于 \(x,y\) 互素，某个 \(x^Nh\) 被 \(xy\) 整除恰等价于 \(y\mid h\)，所以第一项为 \((y)\)；第二项同理为 \((x)\)。多生成元饱和的交公式给出第三项为 \((x)\cap(y)=(xy)\)。因 \(xy\in I\)，第四项含一，等于全环。第三项要求一个足够高幂理想中的全部乘子有效，第四项只允许乘积 \(xy\) 的幂，量词和乘法集合不同。
\end{solution}

\begin{exercise}\label{b3:ex:08-radical-certificate}
设 \(k\) 为域，\(I=(x^2,y^3)\subseteq k[x,y]\)。给出一个适用于任意特征的恒等式，证明 \(x+y\in\sqrt I\)，并据此写出 \(1\in I\,k[x,y,u]+\bigl(1-u(x+y)\bigr)\) 的证书。
\end{exercise}
\begin{solution}
展开得
\[
(x+y)^4=x^2(x^2+4xy+6y^2)+y^3(4x+y).
\]
这个整数系数恒等式在任意特征下都成立。记 \(f=x+y\)，则
\[
1=u^4f^4+(1-uf)(1+uf+u^2f^2+u^3f^3).
\]
将第一式代入第二式的 \(f^4\)，便得到一关于 \(x^2,y^3,1-uf\) 的显式表示。证书不依赖求出任何零点。
\end{solution}

\begin{exercise}\label{b3:ex:08-rational-denominator}
设 \(k\) 为代数闭域。在 \(k[t,x,u]\) 中令 \(H=(tx-1,ut-1)\)。计算 \(H\cap k[x]\)，并比较参数 \(x=1/t\)、\(t\ne0\) 的像与其闭包。为什么不能只消去方程 \(tx-1\) 后宣称每个点都可取参数？
\end{exercise}
\begin{solution}
商环由 \(t,t^{-1}\) 生成，且 \(x=u=t^{-1}\)，同构于 \(k[t,t^{-1}]\)。限制映射 \(k[x]\to k[t,t^{-1}]\) 单射，所以交为零。在代数闭域上，参数像为 \(k^\times\)，闭包为整条直线；零不能提升。消去计算保留了所有多项式关系，却不把“参数分母非零”转成普通闭方程。即使加入辅助变量保证分母可逆，投影后仍可能漏点。
\end{solution}

\begin{exercise}\label{b3:ex:08-parameter-split}
设 \(k\) 为域。将 \(I=(ax-y,y^2-y)\) 视为 \(k(a)[x,y]\) 中的理想，求字典序 \(x\succ y\) 的约化 Gröbner 基。讨论专门化 \(a=c\) 后的商代数及其维数，并说明为什么 \(a=0\) 要单独处理。
\end{exercise}
\begin{solution}
在有理函数域中，约化基为 \(\{x-a^{-1}y,y^2-y\}\)，两首单项式 \(x,y^2\) 互素，尾项均不被其他首单项式整除。对 \(c\ne0\)，可直接代入，商代数同构于 \(k[y]/(y^2-y)\simeq k\times k\)，维数二，两个点为 \((0,0),(c^{-1},1)\)。当 \(c=0\) 时，原生成组成为 \(-y,y^2-y\)，理想为 \((y)\)，商为 \(k[x]\)，维数无限。原基含分母 \(a\)，在零处没有定义；单纯删去其分母会错误保留泛点处才成立的条件。
\end{solution}

\begin{exercise}\label{b3:ex:08-field-coefficients}
在 \(\mathbb F_2[\alpha]=\mathbb F_2[z]/(z^2+z+1)\) 上取 \(I=(x-\alpha y,y^2-1)\)，按字典序 \(x\succ y\) 求约化基、标准单项式基及 \(x^2\) 的正规式；这个商代数约化吗？
\end{exercise}
\begin{solution}
给定两生成元首一且首单项式互素，尾项均不被其他首单项式整除，已构成约化基；标准单项式基为 \(1,y\)。由 \(\alpha^2=\alpha+1\) 得 \(x^2=\alpha^2y^2=\alpha+1\)。但特征二下 \(y^2-1=(y-1)^2\)，商中的 \(y-1\) 非零而平方为零，故商不约化。系数域本身的不可约定义多项式并不保证另一个商代数约化。
\end{solution}

\begin{exercise}\label{b3:ex:08-modular-exception}
在 \(\mathbb Q[x]\) 中和 \(\mathbb F_2[x]\) 中分别计算由整系数多项式 \(2x,x+1\) 生成的理想。说明一次模素数计算为何不能自动证明有理数域中的理想相等。
\end{exercise}
\begin{solution}
在 \(\mathbb Q[x]\) 中，
\[
1=(x+1)-\dfrac{1}{2}(2x),
\]
故理想为全环。在 \(\mathbb F_2[x]\) 中第一生成元变成零，理想为 \((x+1)\)。原证书的分母二在该特征下不可逆，原首系数也消失。因此模素数计算可以提供线索，但要回到有理数域验证精确表示，并排除坏素数。
\end{solution}

\begin{exercise}\label{b3:ex:08-certificate-directions}
设 \(k\) 为域。在 \(k[x]\) 中分别用 \(F=\{x^2\},G=\{x\}\) 和 \(F=\{x\},G=\{x^2\}\)，说明只核验原生成元对候选组的零余式，或者只核验候选组属于原理想，都不能证明两组理想相等。
\end{exercise}
\begin{solution}
第一组有 \(x^2=x\cdot x\)，故 \(F\) 对 \(G\) 的余式为零，但 \(x\notin(x^2)\)，仅得到 \((F)\subsetneq(G)\)。第二组中 \(x^2\in(x)\)，候选元素确属原理想，但 \(x\notin(x^2)\)，仅得到 \((G)\subsetneq(F)\)。两个候选单元素组本身都是 Gröbner 基；仍须同时验证两向理想包含。
\end{solution}

\begin{exercise}\label{b3:ex:08-elimination-growth}
设 \(k\) 为域，\(n\) 为不小于二的整数，在 \(k[x_1,\ldots,x_n]\) 中取 \(I=(x_{i+1}-x_i^2\mid1\le i<n)\)。求 \(I\cap k[x_1,x_n]\)，并给出关系 \(x_n-x_1^{2^{n-1}}\) 属于 \(I\) 的递推证明。
\end{exercise}
\begin{solution}
记 \(h_j=x_j-x_1^{2^{j-1}}\)。有 \(h_1=0\)，且
\[
h_{j+1}=(x_{j+1}-x_j^2)
+\left(x_j+x_1^{2^{j-1}}\right)h_j.
\]
递推得到 \(h_n\in I\)。另一方面，商环通过逐次代入同构于 \(k[x_1]\)。任一 \(f(x_1,x_n)\) 对首一 \(h_n\) 除法，余式只含 \(x_1\)；若 \(f\in I\)，该余式在商环中为零，只能是零多项式。所以交理想恰为 \((h_n)\)。输入二次关系的数目随 \(n\) 线性增长，输出关系的次数却为 \(2^{n-1}\)。
\end{solution}

\begin{exercise}\label{b3:ex:08-near-zero-not-zero}
令 \(I_\varepsilon=(x^2-\varepsilon,xy,y^2)\subseteq\mathbb R[x,y]\)。比较 \(\varepsilon=0\)、\(\varepsilon>0\)、\(\varepsilon<0\) 时商代数的维数、约化性及实零点数，说明小的系数扰动可以改变哪些性质。
\end{exercise}
\begin{solution}
\(\varepsilon=0\) 时理想为 \((x,y)^2\)，商有基 \(1,x,y\)，维数三，非约化，实零点只有原点。若 \(\varepsilon\ne0\)，商中 \(x^2=\varepsilon\) 使 \(x\) 可逆，再由 \(xy=0\) 得 \(y=0\)；反向 \(y=\varepsilon^{-1}\bigl(x(xy)-y(x^2-\varepsilon)\bigr)\) 给出理想中的显式表示。因此商为 \(\mathbb R[x]/(x^2-\varepsilon)\)，维数二且约化。正参数时有两个实根，商为 \(\mathbb R\times\mathbb R\)；负参数时无实根，商为一个同构于 \(\mathbb C\) 的域。任意小的参数区间都包含这三种不同情形，近似系数不能单独决定上述离散数据。
\end{solution}

\begin{exercise}\label{b3:ex:08-multiplication-nilpotent}
设 \(k\) 为域。对 \(A=k[x,y]/(x^2-y,xy)\)，用基 \((1,x,y)\) 写出 \(T_x,T_y\)，求 \(T_x\) 的最小多项式，并计算 \((1+x)^{-1}\)。
\end{exercise}
\begin{solution}
由\cref{b3:ex:08-nonreduced-buchberger}，
\[
T_x=\begin{pmatrix}0&0&0\\1&0&0\\0&1&0\end{pmatrix},
\qquad T_y=T_x^2=\begin{pmatrix}0&0&0\\0&0&0\\1&0&0\end{pmatrix}.
\]
\(T_x^3=0\)、\(T_x^2\ne0\)，所以最小多项式为 \(T^3\)。商中 \((1+x)(1-x+y)=1\)，因为 \(x^2=y\)、\(xy=0\)。故逆为 \(1-x+y\)，任意特征下均成立。
\end{solution}

\begin{exercise}\label{b3:ex:08-double-points-length}
设 \(k\) 为域。对 \(A=k[x,y]/\bigl(x-y,y^2(y-1)^2\bigr)\)，求其标准单项式数、局部因子及各点重数。为什么底域已经含有全部根仍不能把维数当作不同点数？
\end{exercise}
\begin{solution}
字典序 \(x\succ y\) 的三角基给出标准单项式 \(1,y,y^2,y^3\)，所以维数四。因 \(y^2\) 与 \((y-1)^2\) 互素，中国剩余定理给出
\[
A\simeq k[y]/(y^2)\times k[y]/\bigl((y-1)^2\bigr).
\]
两个局部因子均长二，分别对应 \((0,0)\)、\((1,1)\)，故每个点重数二。非零幂零元保留在局部因子中，底域是否含根不能将它们消去。零与一在任意特征下都不同，因而这里不需要排除特征二。
\end{solution}

\begin{exercise}\label{b3:ex:08-joint-eigenvalues}
设 \(k\) 代数闭，\(A=k[x_1,\ldots,x_n]/I\) 非零且有限维，\(T_i\) 为坐标乘法矩阵。对零点 \(a\)，证明求值泛函 \(\varepsilon_a\in A^*\) 是共同左特征向量，即 \(\varepsilon_aT_i=a_i\varepsilon_a\)。证明对应 \(a\) 的共同左特征空间总为一维，即使 \(A\) 非约化也如此。
\end{exercise}
\begin{solution}
因 \(I\) 中的多项式在 \(a\) 消失，求值在商上良定义，且 \(\varepsilon_a(x_ib)=a_i\varepsilon_a(b)\)，所以给出共同左特征向量。若 \(\ell T_i=a_i\ell\)，则对任意多项式 \(f\)，有 \(\ell(f\cdot1)=f(a)\ell(1)\)。所有商元素均由多项式表示，故 \(\ell=\ell(1)\varepsilon_a\)。非零 \(\ell\) 必有 \(\ell(1)\ne0\)，共同左特征空间因此恰为 \(k\varepsilon_a\)。
\end{solution}

\begin{exercise}\label{b3:ex:08-simultaneous-diagonalization}
设 \(k\) 代数闭，\(A=k[x,y]/(x,y)^2\)。求坐标乘法矩阵及其对应 \((0,0)\) 的共同右特征空间，并比较该空间的维数、不同零点数及局部因子长度。说明共同特征向量的个数为何不能直接代替重数。
\end{exercise}
\begin{solution}
取基 \((1,x,y)\)，则 \(T_x\) 只有第 \((2,1)\) 项为一，\(T_y\) 只有第 \((3,1)\) 项为一，其余项均为零。两个核的交为 \(kx\oplus ky\)，维数二。不同零点只有原点一个，而商代数的长度为三，因为它是剩余域为 \(k\) 的三维局部代数。二、一、三依次是两核交的维数、零点数和商代数长度。由\cref{b3:thm:multiplication-joint-eigenvalues}，这些非零幂零乘法矩阵也不可能同时对角化。
\end{solution}

\begin{exercise}\label{b3:ex:08-fglm-small}
设 \(k\) 为域。已知 \(I=(x-y^2,y^3-y)\) 的字典序 \(x\succ y\) Gröbner 基。用商空间基 \(1,y,y^2\) 找出字典序 \(y\succ x\) 的约化基，并写出新标准单项式。
\end{exercise}
\begin{solution}
商中 \(x=y^2\)，故 \(x^2=x\)、\(xy=y\)、\(y^2=x\)。因此候选基为
\[
G'=\{y^2-x,\ xy-y,\ x^2-x\}.
\]
第一关系是原生成元的负数；第二为 \(y(x-y^2)+(y^3-y)\)；第三由 \(x^2-x=(x+y^2-1)(x-y^2)+y(y^3-y)\) 给出。反向 \(x-y^2=-(y^2-x)\)，\(y^3-y=y(y^2-x)+(xy-y)\)，故理想相等。三对 S 多项式依次可写为 \(y^2-x^2=(y^2-x)-(x^2-x)\)、\(xy^2-x^3=x(y^2-x)-x(x^2-x)\)、零；在该序下均可约化为零。该组已经约化，标准单项式为 \(1,x,y\)，其旧坐标分别为 \(1,y^2,y\)。
\end{solution}

\begin{exercise}\label{b3:ex:08-separating-coordinate}
在 \(\mathbb Q^2\) 中取三点 \((0,0),(1,0),(0,1)\)，令 \(A\) 为它们的坐标代数。证明 \(t=x+2y\) 生成 \(A\)，并写出满足 \(x=q_x(t),y=q_y(t)\) 的二次以下多项式。
\end{exercise}
\begin{solution}
\(t\) 的取值依次为 \(0,1,2\)，互不相同。取
\[
p(T)=T(T-1)(T-2),\quad q_x(T)=T(2-T),\quad
q_y(T)=\dfrac{1}{2}T(T-1).
\]
这两个坐标多项式在三点分别给出 \(x,y\) 的取值。求值同态给出 \(A\simeq\mathbb Q^3\)，而 \(T\) 的三个取值不同，使 \(\mathbb Q[T]/(p)\to A\) 为同构。分母二也说明这份具体公式不能未经重算就搬到特征二。
\end{solution}

\begin{exercise}\label{b3:ex:08-rur-denominator}
设 \(p(T)=T(T-1)(T-2)\in\mathbb Q[T]\)。判断 \(T+1\) 与 \(T\) 在 \(\mathbb Q[T]/(p)\) 中是否可逆；对前者给出逆，对后者说明使其可逆会丢去哪一个点。
\end{exercise}
\begin{solution}
对 \(T+1\) 的逆在 \(0,1,2\) 处应分别取 \(1,\dfrac12,\dfrac13\)。插值得
\[
u(T)=1-\dfrac{2}{3}T+\dfrac{1}{6}T^2,\qquad
(T+1)u(T)-1=\dfrac{1}{6}p(T).
\]
故 \(u\) 是逆。\(T\) 在零处为零，不是单位；把它变为单位只保留根一、二对应的两个因子，丢掉根零。因此带分母的单变量表示必须同时检查分母和 \(p\) 互素，或明确分支范围。
\end{solution}

\begin{exercise}\label{b3:ex:08-nonmonogenic}
设 \(k\) 为域，\(r\) 为不小于二的整数，\(A=k[x_1,\ldots,x_r]/(x_1,\ldots,x_r)^2\)。证明 \(A\) 作为 \(k\) 代数的最少生成元数恰为 \(r\)。再说明有限域上的约化有限代数也未必单生成。
\end{exercise}
\begin{solution}
记 \(\mathfrak m=(\bar x_1,\ldots,\bar x_r)\)，则 \(\mathfrak m^2=0\)，且 \(\dim_k\mathfrak m=r\)。若代数生成元为 \(a_j=c_j+n_j\)、\(1\le j\le s\)，其中 \(n_j\in\mathfrak m\)，由于任何两个 \(n_j\) 的乘积为零，每个生成的多项式都在 \(k+\sum_jkn_j\) 中。因此若它们生成全部 \(A\)，必须使 \(n_j\) 线性生成 \(\mathfrak m\)，从而 \(s\ge r\)。原来的 \(r\) 个变量已经生成，所以最少数为 \(r\)。约化例取 \(B=\mathbb F_2^3\)。一个元素的三个坐标只能取零、一，至少两个坐标相同；其任意多项式在这两坐标也相同，不能生成全部 \(B\)。有限维性本身不保证单变量描述。
\end{solution}


\begin{exercise}\label{b3:ex:08-resultant-quadratic}
设 \(k\) 为域，\(a,b\in k\)。在 \(k[X]\) 中令 \(f=X^2-a\)、\(g=X-b\)。分别用 Sylvester 矩阵和乘法矩阵求结式，判断何时有公共根，并说明 \(a=b=0\) 时公共根是否为单根。
\end{exercise}
\begin{solution}
Sylvester 矩阵为 \(\left(\begin{smallmatrix}1&1&0\\0&-b&1\\-a&0&-b\end{smallmatrix}\right)\)，行列式为 \(b^2-a\)。在基 \((1,X)\) 下，乘 \(X-b\) 的矩阵为 \(\left(\begin{smallmatrix}-b&a\\1&-b\end{smallmatrix}\right)\)，有相同行列式。因此公共根存在恰当且仅当 \(a=b^2\)，此时根为 \(b\)。当 \(a=b=0\)，零是 \(f\) 的二重根、\(g\) 的单根；结式为零只检测公共根存在，不单独给出两个多项式各自的重数。
\end{solution}

\begin{exercise}\label{b3:ex:08-resultant-degree-drop}
在 \(\mathbb Q[t,X]\) 中取 \(f=tX-1\)、\(g=tX^2-X+1\)。计算关于 \(X\) 的结式，并检查其零点是否给出原方程的解。
\end{exercise}
\begin{solution}
在 \(\mathbb Q(t)\) 中由根公式得到 \(\operatorname{Res}_X(f,g)=t^2\cdot g(1/t)=t^2\)；两边均为 \(\mathbb Q[t]\) 中多项式，所以这也是原系数环的行列式。结式只给候选 \(t=0\)，但此时 \(f=-1\)，方程无解。还可直接用 \(g-Xf=1\) 看出原理想为单位理想，因而对任何参数都无解。此例要求同时核查原生成元，不可只返回结式零点。
\end{solution}

\begin{exercise}\label{b3:ex:08-swap-invariants}
设 \(k\) 为任意域，\(C_2\) 在 \(k[x,y]\) 上通过交换 \(x,y\) 作用。证明其不变环为 \(k[x+y,xy]\)，且 \(x+y,xy\) 代数无关。解释为什么这个结论在特征二时仍成立。
\end{exercise}
\begin{solution}
设 \(f(x,y)=f(y,x)\)。先处理非零齐次 \(f\)，在字典序 \(x\succ y\) 下记其首项为 \(cx^ay^b\)。必有 \(a\ge b\)：否则由对称性，\(cx^by^a\) 也出现且更大。多项式
\[
c(x+y)^{a-b}(xy)^b
\]
有相同首项；从 \(f\) 中减去它，所得多项式仍对称、齐次，而且首单项式严格减小。同一次数只有有限多个单项式，故反复减去后得到零，证明每个齐次对称多项式都在 \(k[x+y,xy]\) 中。一般 \(f\) 的各齐次部分仍对称，分别应用此过程即可。两个生成元本身不变，故得到等式。

对非负整数 \(i,j\)，\((x+y)^i(xy)^j\) 的首单项式为 \(x^{i+j}y^j\)，首项系数为一。不同 \((i,j)\) 给出不同指数对。若存在非零关系 \(F(x+y,xy)=0\)，在 \(F\) 的有限多个非零项中，取代入后首单项式最大的那一项；其他项的任何单项式都比它小，不能抵消它，矛盾。因此代入同态 \(k[u,v]\to k[x,y]\) 单射。这些步骤没有除以二，特征二时交换两个变量也仍不是恒等作用，所以结论保持成立。
\end{solution}

\begin{exercise}\label{b3:ex:08-cubic-diagonal-invariants}
设 \(k\) 为特征不为三的域，并含有一个三阶本原单位根 \(\zeta\)。令 \(C_3=\langle s\rangle\) 在 \(k[x,y]\) 上作用为 \(s\cdot x=\zeta x\)、\(s\cdot y=\zeta^{-1}y\)。证明
\[
k[x,y]^{C_3}=k[x^3,xy,y^3]
\cong k[u,v,w]/(uw-v^3),
\]
并把全部关系的核写成一个消去理想。不要只检验 \(uw-v^3\) 确实在核中。
\end{exercise}
\begin{solution}
单项式 \(x^iy^j\) 的作用因子为 \(\zeta^{i-j}\)。因为 \(\zeta\) 恰有三阶，比较系数说明不变多项式只含 \(i\equiv j\pmod3\) 的项。此时唯一地写成 \(i=3a+r,j=3c+r\)，其中 \(r\in\{0,1,2\}\)，从而
\[
x^iy^j=(x^3)^a(xy)^r(y^3)^c.
\]
这证明所列三个元素生成全部不变环。

令 \(\phi:k[u,v,w]\to k[x,y]\) 分别把变量送到 \(x^3,xy,y^3\)。对 \(v^3-uw\) 按 \(v\) 作首一除法，任意余式唯一形如
\[
A_0(u,w)+v\cdot A_1(u,w)+v^2\cdot A_2(u,w).
\]
其中 \(u^av^rw^c\) 的像为 \(x^{3a+r}y^{3c+r}\)。由两个指数模三的余数可恢复 \(r\)，再恢复 \(a,c\)，故不同的这些单项式的像互不相同。若余式映为零，则各系数全为零。于是 \(\ker\phi=(uw-v^3)\)。

最后，\cref{b3:prop:parametric-implicit-ideal}取分母一、参数 \(x,y\)，给出
\[
\ker\phi=(u-x^3,v-xy,w-y^3)\cap k[u,v,w],
\]
取交前的理想在 \(k[x,y,u,v,w]\) 中生成。它等于上述主理想，既有生成关系的验证，也有余式独立性给出的反向包含。
\end{solution}

% !TeX root = ../../Book3.tex
% 纲要标题：## 第9章　Krull 维数与高度定理
% 纲要位置：工作记录/计划纲要.md，第 2200 行。
% 纲要细目（写作范围）：

\chapter{Krull 维数与高度定理}
\label{b3:ch:09}
\BookChapterNumber{b3:ch:09}

\begin{chapterabstract}
本章用素理想链定义 Krull 维数，比较商环、局部化与支集中的链，并证明主理想定理和高度定理。Noether 正规化把有限型代数的维数与超越次数联系起来；局部维数和极大理想的最少生成元数也得到区分。
\end{chapterabstract}

\begin{learninggoals}
\begin{itemize}
\item 按严格包含的次数计算链长，区分高度、商环维数和模的支集维数。
\item 在主理想定理中辨认极小性，并用高度界限制方程和生成元的个数。
\item 利用超越次数、整扩张及正规化计算有限型代数的维数。
\item 构造参数系，计算嵌入维数，区分参数、极大理想生成元和正则性。
\end{itemize}
\end{learninggoals}

设 \(k\) 为域。在 \(k[x,y]\) 中，可沿 \((0)\subsetneq(x)\subsetneq(x,y)\) 走两步；加上关系 \(xy=0\) 后，\((0)\) 不再是素理想，相应的两个不可约分量各只容许一步素理想链。给环添一个方程，维数往往下降，但下降多少不能只看方程个数：若方程已经由原关系推出，就不会再减少任何链。维数应当从真正存在的素理想链，而非从展示中的符号数目读取。

\ProseParagraphBreak

局部化又只保留落在某个素理想以下的链，因此同一环在不同点附近可能显出不同的局部维数。另一方面，一个局部环的极大理想需要几个生成元，问的是局部坐标有多少独立的一阶方向，与最长素理想链不是同一个问题。比较这两个数字，将为后面对正则点与奇点的区分做好准备。

\ProseParagraphBreak

本章使用\chapref{b3:ch:01}至\chapref{b3:ch:06}中的素谱、局部化、Noether 条件、准素理想、Nakayama 引理和整扩张；超越次数仍指一组超越基的元素个数，其域论基础见第一卷第14.3节。高度定理的证明先于 Hilbert 函数和 Artin--Rees 引理，后面的章节可以直接调用它。正则局部环在这里仅建立定义和基本算例，一般结构定理留在\chapref{b3:ch:18}。

\ProseParagraphBreak
有限型整代数的维数公式可参阅 Matsumura 的《Commutative Ring Theory》\cite[Theorem 5.6, pp. 34--35]{Matsumura1989CommutativeRingTheory}；高度定理和逐次商环的维数见\cite[Theorems 13.5--13.6, pp. 100--101]{Matsumura1989CommutativeRingTheory}，参数系与嵌入维数的定义见\cite[\S14, Systems of parameters and multiplicity, pp. 104--105]{Matsumura1989CommutativeRingTheory}。本章从素理想链证明高度定理。

% !TeX root = ../../Book3.tex
% 纲要标题：### 9.1 维数与高度
% 纲要位置：工作记录/计划纲要.md，第 2079 行。
% 纲要细目（写作范围）：
% * 素理想链、环维数和理想高度
% * 模的维数与支集
% * 局部化对维数的影响

\section{维数与高度}
\label{b3:sec:0901}

素谱中的点既可以独立出现，也可以按包含关系排列。一个不可约闭集内还有多少层更小的不可约闭集，是维数要记录的问题。以下的链长只计算严格包含的次数；不假设所有极大链一样长，也不预设任何几何图像中的维数公式。

\subsection{素理想链、环维数与理想高度}
\label{b3:subsec:0901:01}

\begin{definition}\label{b3:def:krull-dimension-height}
设 \(R\) 为交换环。其有限素理想链
\[
\mathfrak p_0\subsetneq\mathfrak p_1\subsetneq\cdots\subsetneq\mathfrak p_r
\]
的长度 \(r\) 的上确界称为 \(R\) 的\term[Krull-weishu]{Krull 维数}[Krull dimension]，记为 \(\dim R\)。对素理想 \(\mathfrak p\subseteq R\)，以 \(\mathfrak p\) 为末项的此类链长的上确界称为它的\term[gaodu]{高度}[height]，记为 \(\operatorname{ht}\mathfrak p\)。真理想 \(I\subseteq R\) 的高度定义为
\[
\operatorname{ht}I=\inf_{\mathfrak p\supseteq I}\operatorname{ht}\mathfrak p.
\]
允许维数或高度为 \(+\infty\)。零环的谱为空，约定其维数为 \(-\infty\)；单位理想的高度不在本章使用。
\end{definition}
\symidx{\(\dim R\)：环的 Krull 维数}
\symidx{\(\operatorname{ht}\mathfrak p\)、\(\operatorname{ht}I\)：素理想与真理想的高度}

域只有一个素理想 \(0\)，所以维数为零。整数环的非零素理想都是极大理想，故
\(\dim\mathbb Z=1\)。一元多项式环 \(k[t]\) 同样为一维：它是主理想整环，每个非零素理想由不可约多项式生成，因而极大。零维并不表示素谱只有一个点；有限个域的直积已经给出多个点。

\begin{proposition}\label{b3:prop:dimension-basic-comparison}
设 \(R\) 为交换环，\(I\subseteq R\) 为理想，\(R_{\mathrm{red}}=R/\sqrt{(0)}\)。则有 \(\dim(R/I)\le\dim R\)，而
\(\dim R_{\mathrm{red}}=\dim R\)。若 \(\mathfrak p\) 为素理想，且有关数均有限，则
\[
\operatorname{ht}\mathfrak p+\dim(R/\mathfrak p)\le\dim R.
\]
若 \(R_1,\ldots,R_s\) 为交换环且至少一个非零，则
\(\dim(R_1\times\cdots\times R_s)\) 等于非零因子 \(R_i\) 的维数的最大值。
\end{proposition}
\begin{proof}
商环的素理想链对应原环中包含 \(I\) 的链；模掉幂零根时，所有素理想都保留，故链长完全不变。对第二式，取一条末项为 \(\mathfrak p\) 的最长链，以及一条首项为 \(\mathfrak p\) 的最长链，接在一起即可。有限整数的上确界能取到，故这些链存在。直积环的素理想分别来自一个因子，不同因子所对应的素理想互不包含，所以任何链都留在同一个因子中。
\end{proof}

上述不等式不能在没有附加条件时写成等式。高度测量点以下的链，商环维数测量点以上的链；两段能连接起来，不表示任何最长链都经过这个指定点。

\subsection{模的维数与支集}
\label{b3:subsec:0901:02}

\begin{definition}\label{b3:def:module-dimension}
设 \(R\) 为交换环，\(M\) 为 \(R\) 模。支集
\(\operatorname{Supp}_R M=\{\,\mathfrak p\in\operatorname{Spec}R\mid M_{\mathfrak p}\ne0\,\}\)
内有限严格素理想链长的上确界称为 \(M\) 的\term[mo-weishu]{维数}[dimension of a module]，记为 \(\dim_R M\)。空支集的维数取 \(-\infty\)。
\end{definition}
\symidx{\(\dim_R M\)：模的支集维数}

若 \(M\) 有限生成，\cref{b3:thm:support-finite}给出
\(\operatorname{Supp}_R M=V(\operatorname{Ann}_R M)\)，所以
\begin{equation}\label{b3:eq:module-dimension-annihilator}
\dim_R M=\dim\bigl(R/\operatorname{Ann}_R M\bigr).
\end{equation}
这里的维数不是生成元个数，也不是向量空间维数。例如 \(R=k[x]\) 上的 \(R^r\)、\(r\ge1\)，都有支集 \(\operatorname{Spec}R\)，故维数均为一；而 \(R/(x^n)\)、\(n\ge1\)，的支集只有一个闭点，维数为零，其 \(k\) 向量空间维数却为 \(n\)。

\begin{proposition}\label{b3:prop:dimension-associated-components}
若 \(R\) 为交换 Noether 环，\(M\ne0\) 为有限 \(R\) 模，则
\[
\dim_R M=\max_{\mathfrak p\in\operatorname{Ass}_R M}\dim(R/\mathfrak p).
\]
若 \(0\rightarrow L\rightarrow M\rightarrow N\rightarrow0\) 为有限生成 \(R\) 模的短正合列，则
\[
\dim_R M=\max\{\dim_R L,\dim_R N\}.
\]
\end{proposition}
\begin{proof}
由\cref{b3:prop:minimal-primes-associated}，支集中每个素理想都包含某个关联素理想，且关联素理想有限。因此支集等于这些 \(V(\mathfrak p)\) 的有限并。任一链的最小项包含某个这样的 \(\mathfrak p\)，整条链便落在 \(V(\mathfrak p)\) 中，得到第一式。
短正合列的支集是两端支集之并，见\cref{b3:prop:support-exact}。因有限模的支集是闭集，链的最小项落在哪一端，整条链就留在该端，故得到最大值公式。
\end{proof}

\subsection{局部化对维数的影响}
\label{b3:subsec:0901:03}

\begin{proposition}\label{b3:prop:dimension-localization}
设 \(R\) 为交换环，\(\mathfrak p\subseteq R\) 为素理想，\(S\subseteq R\) 为乘法集。则有
\[
\dim R_{\mathfrak p}=\operatorname{ht}\mathfrak p,\qquad
\dim S^{-1}R\le\dim R,\qquad
\dim R=\sup_{\mathfrak m\text{ 极大}}\dim R_{\mathfrak m}.
\]
若 \(M\) 为有限 \(R\) 模，则 \(\dim_{S^{-1}R}S^{-1}M\le\dim_R M\)。零环情形的空上确界同样取 \(-\infty\)。
\end{proposition}
\begin{proof}
局部化素理想对应保留包含关系。\(R_{\mathfrak p}\) 的素理想恰好对应包含于
\(\mathfrak p\) 的素理想，每条这样的链都可在末端补上 \(\mathfrak p\)，所以其维数就是高度。一般局部化只保留不遇到 \(S\) 的素理想，故不能增加链长。任一有限素链的末项包含于某个极大理想；在该极大理想处局部化保留整条链，得到第三式。

对模，若 \(\mathfrak q\cap S=\varnothing\)，则先在 \(S\) 处局部化、再在
\(S^{-1}\mathfrak q\) 处局部化，与直接取 \(M_{\mathfrak q}\) 相同。两者均由允许
\(R\setminus\mathfrak q\) 中的分母给出，分式的互逆公式保证这一识别。因此局部模的支集链对应原支集中避开 \(S\) 的链，得到不等式。
\end{proof}

一般点处的局部化可能大幅降低维数。例如整环的零理想处局部环是分式域，维数为零，不论原环有多少维。相反，在多项式环的闭点处，局部化会保留许多不同方向的素理想链；下一节将用理想生成元的个数限制这些链的长度。

% !TeX root = ../../Book3.tex
% 纲要标题：### 9.2 高度定理
% 纲要位置：工作记录/计划纲要.md，第 2085 行。
% 纲要细目（写作范围）：
% * Krull 主理想定理
% * 广义高度定理
% * 生成元数与余维约束

\section{高度定理}
\label{b3:sec:0902}

一个方程不应一次切掉任意多层素理想链，这个直觉在 Noether 环中由主理想定理保证。证明需要小心区分极小素理想与任意包含该方程的素理想。本节先用准素性和 Nakayama 引理建立单个方程的结论，再归纳处理多个方程。

\subsection{Krull 主理想定理}
\label{b3:subsec:0902:01}

\begin{lemma}\label{b3:lem:principal-primary-local-domain}
设 \((R,\mathfrak m)\) 为 Noether 局部整环，\(a\in R\)，且
\(\sqrt{aR}=\mathfrak m\)。则 \(\dim R\le1\)。
\end{lemma}
\begin{proof}
若 \(a=0\)，整环性给出 \(\mathfrak m=0\)，故 \(R\) 是域。以下设 \(a\ne0\)。
反设存在非零素理想 \(0\ne\mathfrak p\subsetneq\mathfrak m\)。因
\(\sqrt{aR}=\mathfrak m\)，有 \(a\notin\mathfrak p\)。

环 \(R/aR\) 是零维 Noether 环，所以由\cref{b3:thm:noether-dimension-zero}为 Artin 环。考虑符号幂
\(Q_n=\mathfrak p^nR_{\mathfrak p}\cap R\)。在 \(R_{\mathfrak p}\) 中，\(\mathfrak p^nR_{\mathfrak p}\) 的根为唯一极大理想，故由\cref{b3:prop:maximal-radical-primary}准素；再由\cref{b3:prop:primary-localization}，收缩后的每个 \(Q_n\) 都是 \(\mathfrak p\) 准素理想。它们在 \(R/aR\) 中的像组成降链，因此存在 \(n\ge1\) 使
\[
Q_n+aR=Q_{n+1}+aR.
\]
对 \(x\in Q_n\)，可写 \(x=y+ar\)，其中 \(y\in Q_{n+1}\)。于是 \(ar\in Q_n\)；
因 \(a\notin\sqrt{Q_n}=\mathfrak p\)，准素性迫使 \(r\in Q_n\)。所以
\[
Q_n=Q_{n+1}+aQ_n.
\]
有限生成模 \(Q_n/Q_{n+1}\) 因而等于自身的 \(a\) 倍。由 \(a\in\mathfrak m\) 及
\cref{b3:thm:nakayama}，它为零，故 \(Q_n=Q_{n+1}\)。

在 \(\mathfrak p\) 处局部化得到
\(\mathfrak p^nR_{\mathfrak p}=\mathfrak p^{n+1}R_{\mathfrak p}\)。
再对 Noether 局部环 \(R_{\mathfrak p}\) 上的有限生成模
\(\mathfrak p^nR_{\mathfrak p}\) 使用 Nakayama 引理，得到它为零。但
\(R_{\mathfrak p}\) 是整环，且非零的 \(\mathfrak p\) 在其中仍非零，其正整数次幂不可能为零，矛盾。于是除零理想外，没有严格小于 \(\mathfrak m\) 的素理想。
\end{proof}

\begin{theorem}[Krull 主理想定理]\label{b3:thm:principal-ideal-theorem}
设 \(R\) 为交换 Noether 环，\(a\in R\)。每个在 \((a)\) 上极小的素理想 \(\mathfrak p\) 都满足
\[
\operatorname{ht}\mathfrak p\le1.
\]
\end{theorem}
\begin{proof}
在 \(\mathfrak p\) 处局部化后，只存在一个包含 \(a\) 的素理想，即唯一极大理想；
因此可设 \(R\) 为局部环且 \(\sqrt{aR}=\mathfrak m\)。若存在长度至少二的素链，可截取并在末端补上 \(\mathfrak m\)，得到
\(\mathfrak q_0\subsetneq\mathfrak q_1\subsetneq\mathfrak m\)。
在局部整环 \(R/\mathfrak q_0\) 中，\(a\) 的像的根仍为极大理想的像，上一引理却排除了非零的中间素理想
\(\mathfrak q_1/\mathfrak q_0\)，矛盾。因此局部维数至多为一，即所求高度界。
\end{proof}

这称为\term[Krull-zhulixiang-dingli]{Krull 主理想定理}[Krull's principal ideal theorem]。
在整环中，若 \(a\) 非零且非单位，其上极小素理想非零，故高度恰好为一。
“极小”不可删除：例如 \((x)\subseteq(x,y)\) 于 \(k[x,y]\)，后者包含主理想
\((x)\)，却有链 \(0\subsetneq(x)\subsetneq(x,y)\)，高度至少为二。

\subsection{广义高度定理}
\label{b3:subsec:0902:02}

\begin{theorem}[广义高度定理]\label{b3:thm:height-theorem}
若 \(R\) 为交换 Noether 环，\(I=(a_1,\ldots,a_r)\subsetneq R\)，其中 \(r\ge0\)，则每个在 \(I\) 上极小的素理想 \(\mathfrak p\) 都满足
\[
\operatorname{ht}\mathfrak p\le r.
\]
\end{theorem}
\begin{proof}
对 \(r\) 归纳。\(r=0\) 时，\(\mathfrak p\) 是极小素理想，高度为零；\(r=1\) 已由主理想定理证明。一般情形在 \(\mathfrak p\) 处局部化，可设
\((R,\mathfrak m)\) 局部且 \(\sqrt I=\mathfrak m\)。

取任意链
\[
\mathfrak p_0\subsetneq\cdots\subsetneq\mathfrak p_d=\mathfrak m.
\]
若 \(d=0\) 无须证明。因 \(I\not\subseteq\mathfrak p_{d-1}\)，重排生成元后可设
\(a_1\notin\mathfrak p_{d-1}\)。令 \(\mathfrak q_d=\mathfrak m\)，从
\(i=d-1\) 向下逐次在 \(\mathfrak q_{i+1}\) 内选一个在
\(\mathfrak p_i+(a_1)\) 上极小的素理想 \(\mathfrak q_i\)。这样的选择可以先在
\(\mathfrak q_{i+1}\) 处局部化、取一个极小素理想，再收缩；任何更小的候选也在
\(\mathfrak q_{i+1}\) 内，故它在原环中仍然极小。

对 \(i\le d-2\)，若 \(\mathfrak q_i=\mathfrak q_{i+1}\)，则
\[
\mathfrak p_i\subsetneq\mathfrak p_{i+1}\subsetneq\mathfrak q_i.
\]
第二个包含严格，因为 \(a_1\in\mathfrak q_i\) 而
\(a_1\notin\mathfrak p_{i+1}\)。但在 \(R/\mathfrak p_i\) 中，
\(\mathfrak q_i/\mathfrak p_i\) 在一个主理想上极小，主理想定理禁止它具有这样的长度二链。故
\(\mathfrak q_0\subsetneq\cdots\subsetneq\mathfrak q_{d-1}\) 严格。

这是一条包含 \(a_1\) 的长度 \(d-1\) 的链。在局部商环 \(R/(a_1)\) 中，
其极大理想是剩余 \(r-1\) 个生成元的根，由归纳假设高度至多 \(r-1\)。
所以 \(d-1\le r-1\)，即 \(d\le r\)。原链任意，结论成立。
\end{proof}

这一定理也称\term[Krull-gaodu-dingli]{Krull 高度定理}[Krull's height theorem]。
证明没有将
\(\operatorname{ht}\mathfrak p\) 擅自拆成两个高度之和；一般环上的素链未必具有允许这种拆分的长度性质。Matsumura 的《Commutative Ring Theory》给出相同高度结论\cite[Theorem 13.5, pp. 100--101]{Matsumura1989CommutativeRingTheory}，但其证明次序先用 Samuel 函数；本卷采用上述独立论证，使后面的增长理论能够调用高度定理。

\subsection{生成元数与余维约束}
\label{b3:subsec:0902:03}

若 \((R,\mathfrak m)\) 为 Noether 局部环，\(\mathfrak m\) 有有限个生成元，故
\cref{b3:thm:height-theorem}说明 \(\dim R<\infty\)。更一般地，若
\(\sqrt{(a_1,\ldots,a_r)}=\mathfrak m\)，就有 \(\dim R\le r\)。
这限制的是使零点集中到闭点所需的方程数；不表示任一主理想的任一零点都只有余维一。

\begin{lemma}\label{b3:lem:dimension-cutting}
设 \((R,\mathfrak m)\ne0\) 为交换 Noether 局部环，\(d=\dim R\)。存在
\(x_1,\ldots,x_d\in\mathfrak m\)，使
\[
\sqrt{(x_1,\ldots,x_d)}=\mathfrak m.
\]
\(d=0\) 时取空组。
\end{lemma}
\begin{proof}
对有限整数 \(d\) 归纳。\(d=0\) 时唯一素理想为 \(\mathfrak m\)，所以幂零根为
\(\mathfrak m\)，空组满足要求。若 \(d>0\)，环的极小素理想只有有限多个，且都严格小于
\(\mathfrak m\)；否则 \(\mathfrak m\) 极小便排除了其他素理想，维数将为零。
由\cref{b3:lem:finite-prime-avoidance}选
\(x_1\in\mathfrak m\) 避开全部极小素理想。

任一包含 \(x_1\) 的素链，其首项不是极小素理想，可以在下面补上一项。
因此 \(\dim R/(x_1)\le d-1\)。对这个局部商环使用归纳假设，提升其至多
\(d-1\) 个元素，得到一个根为 \(\mathfrak m\)、至多由 \(d\) 个元素生成的理想。
若少于 \(d\) 个，则广义高度定理会给出 \(\dim R<d\)，矛盾。所以总数恰为 \(d\)，完成归纳。
\end{proof}

这个引理与高度定理合起来，将局部维数解释为生成一个 \(\mathfrak m\) 准素理想所需的最少元素个数。\secref{b3:sec:0904}将把这样的元素组命名为参数系。它们生成的理想通常小于 \(\mathfrak m\)，因此参数个数和极大理想的最少生成元数仍需区分。

% !TeX root = ../../Book3.tex
% 纲要标题：### 9.3 有限型代数的维数
% 纲要位置：工作记录/计划纲要.md，第 2214 行。
% 纲要细目（写作范围）：
% * 多项式环的维数
% * Noether 正规化和超越次数
% * 整扩张的维数保持

\section{有限型代数的维数}
\label{b3:sec:0903}

多项式环中有明显的素理想链，但仅展示一条长链只能给出下界。上界需要证明：在有限型整代数中，每次模掉一个非零素理想，分式域的超越次数就会下降。Noether 正规化随后把一般有限型整代数与多项式环联系起来，借整性将多项式环中的长链提升到原环。


\subsection{多项式环的维数}
\label{b3:subsec:0903:01}

取 \(A=k[x,y]\)、\(\mathfrak p=(x)\)，商环 \(A/\mathfrak p=k[\bar y]\) 的分式域的一个超越基是 \(\{\bar y\}\)。把 \(\bar y\) 提升为 \(y\)，再添入被取商杀掉的非零元素 \(x\)，就在 \(\operatorname{Frac}(A)\) 中得到代数独立的 \(x,y\)。下面的引理说明，在一般有限型整代数中，从商环的有限代数生成组中选出的超越基提升后，也能这样添入一个非零素理想中的元素；困难在于证明添入后仍然代数独立。

\begin{lemma}\label{b3:lem:transcendence-degree-drop}
设 \(A\) 为域 \(k\) 上的有限型整代数，\(0\ne\mathfrak p\) 为素理想。则
\[
\operatorname{trdeg}_k\operatorname{Frac}(A/\mathfrak p)
<\operatorname{trdeg}_k\operatorname{Frac}(A).
\]
\end{lemma}
\begin{proof}
从 \(A/\mathfrak p\) 的有限代数生成组中取极大代数独立组
\(\bar y_1,\ldots,\bar y_r\)，并在 \(A\) 中取提升 \(y_i\)。其他生成元在
\(k(\bar y_1,\ldots,\bar y_r)\) 上代数，故左侧超越次数是 \(r\)。提升后的 \(y_i\) 也代数独立。

取 \(0\ne a\in\mathfrak p\)。我们证明 \(a\) 在 \(k(y_1,\ldots,y_r)\) 上超越。
若它代数，取它在域 \(K=k(y_1,\ldots,y_r)\) 上的首一最小多项式
\(\mu(T)\in K[T]\)。为了在模掉 \(\mathfrak p\) 后留下关于 \(\bar y_i\) 的非零关系，需要这条关系有一个不含 \(a\) 的非零常数项。
若 \(\mu(0)=0\)，则 \(T\) 整除不可约首一多项式 \(\mu\)，故只能有 \(\mu(T)=T\)，从而 \(a=0\)，与所选元素非零矛盾。因此 \(\mu(0)\ne0\)。
选取 \(\mu\) 各系数的非零公共分母 \(D(y)\in k[y_1,\ldots,y_r]\)，
将 \(\mu(a)=0\) 乘以 \(D(y)\)，得到
\[
c_s(y)\cdot a^s+\cdots+c_1(y)\cdot a+c_0(y)=0,
\qquad 0\ne c_0\in k[Y_1,\ldots,Y_r].
\]
这里 \(D(y)\mu(0)\ne0\)，因为 \(D(y)\) 和 \(\mu(0)\) 均为域 \(K\) 中的非零元；
\(y_1,\ldots,y_r\) 的代数独立性又保证对应的常数项多项式 \(c_0(Y)\) 非零。
模掉 \(\mathfrak p\) 得
\(c_0(\bar y)=0\)，与代数独立性矛盾。因此 \(y_1,\ldots,y_r,a\) 代数独立，右侧超越次数至少是 \(r+1\)。
\end{proof}

\begin{theorem}\label{b3:thm:polynomial-dimension}
任意域 \(k\) 及整数 \(n\ge0\) 满足
\[
\dim k[X_1,\ldots,X_n]=n.
\]
更一般地，有限型整 \(k\) 代数 \(A\) 的维数不超过
\(\operatorname{trdeg}_k\operatorname{Frac}(A)\)。
\end{theorem}
\begin{proof}
对严格素链 \(\mathfrak p_0\subsetneq\cdots\subsetneq\mathfrak p_r\)，逐次将上一引理用于
\(A/\mathfrak p_i\) 及其非零素理想 \(\mathfrak p_{i+1}/\mathfrak p_i\)。严格包含使这个新加入的素关系非零，因此超越次数每一步至少下降一。一条长度 \(r\) 的素链就需要起点至少有 \(r\) 个代数独立参数；而起点的超越次数不超过 \(\operatorname{trdeg}_k\operatorname{Frac}(A)\)，因为商环中的代数独立元素提升后仍然代数独立。这证明上界。对于多项式环，
\[
(0)\subsetneq(X_1)\subsetneq(X_1,X_2)\subsetneq\cdots
\subsetneq(X_1,\ldots,X_n)
\]
的各项都是素理想，给出长度 \(n\) 的链。\(n=0\) 时环为域，结论同样成立。
\end{proof}

例如 \(k[X,Y,Z]/(XY)\) 的两个极小素理想为 \((\bar X)\) 和 \((\bar Y)\)，相应商环均为二元多项式环，故其维数为二。这个计算按不可约分量分别进行，不把有零因子的环擅自当成具有一个分式域的整环。

\subsection{Noether 正规化与超越次数}
\label{b3:subsec:0903:02}

\begin{theorem}\label{b3:thm:affine-dimension}\label{b3:thm:affine-dimension-trdeg}
设 \(k\) 为域。若 \(A\) 是有限型整 \(k\) 代数，则
\[
\dim A=\operatorname{trdeg}_k\operatorname{Frac}(A).
\]
若 \(A\ne0\) 只是有限型交换 \(k\) 代数，则
\[
\dim A=\max_{\mathfrak p\in\operatorname{Min}A}
\operatorname{trdeg}_k\operatorname{Frac}(A/\mathfrak p).
\]
这里 \(\operatorname{Min}A\) 表示极小素理想的集合。
\end{theorem}
\begin{proof}
整环情形由\cref{b3:thm:noether-normalization}选出代数独立元素
\(y_1,\ldots,y_d\)，使 \(A\) 是 \(B=k[y_1,\ldots,y_d]\) 上的有限生成模。
\secref{b3:sec:0504}已经证明，这种有限整包含给出有限域扩张
\(k(y_1,\ldots,y_d)\subseteq\operatorname{Frac}(A)\)。因此这些 \(y_i\) 构成该分式域的超越基，超越次数为 \(d\)。由\cref{b3:thm:polynomial-dimension}，这给出上界 \(\dim A\le d\)。

下界来自 \(B\subseteq A\) 的整性。多项式环 \(B\) 中有长度 \(d\) 的素链，由\cref{b3:thm:lying-over}提升首项，再反复用\cref{b3:thm:going-up}提升其余各项；收缩严格增加，提升链也严格增加。因此 \(\dim A\ge d\)，两界合起来得到维数公式。

一般情形下，\(A\) 为 Noether 环，只有有限多个极小素理想。任意素链的首项都包含某个极小素理想，因此整条链都来自某个 \(A/\mathfrak p\)。反之这些商环中的链也是 \(A\) 中的链。这给出
\(\dim A=\max_{\mathfrak p\in\operatorname{Min}A}\dim(A/\mathfrak p)\)，再应用整环情形。
\end{proof}

一般情形的公式说明，全环维数由维数最大的不可约分量决定，不同分量可以有不同维数。幂零厚度不改变素理想链；嵌入素理想所对应的闭集已包含在某个不可约分量中，也不会增加这个最大值。

尖点环 \(k[t^2,t^3]\) 整包含 \(k[t^2]\)，其分式域为 \(k(t)\)，所以维数为一。
变量之间有关系 \((t^3)^2=(t^2)^3\)，两个代数生成元不意味着二维。
相反，\(k[x,y,x^{-1}]\) 虽添入了一个逆元，分式域仍为 \(k(x,y)\)，所以维数仍为二。
有限型条件在这里保证 Noether 正规化可用；不能把该公式未经说明用于任意子环。

\subsection{整扩张的维数保持}
\label{b3:subsec:0903:03}

\begin{theorem}\label{b3:thm:integral-dimension}
若交换环间的包含 \(A\subseteq B\) 是整扩张，则 \(\dim A=\dim B\)。不要求扩张有限，也不要求两环为整环。
\end{theorem}
\begin{proof}
取 \(B\) 中的严格素链 \(\mathfrak q_0\subsetneq\cdots\subsetneq\mathfrak q_r\)，令 \(\mathfrak p_i=\mathfrak q_i\cap A\)。这些收缩仍然严格：若相邻两项相等，\cref{b3:thm:incomparability}就会迫使相应的 \(\mathfrak q_i\) 相等。因此收缩保留链长，得到 \(\dim B\le\dim A\)。

反之，从 \(A\) 中任意有限严格素链出发，先用\cref{b3:thm:lying-over}提升首项，再用\cref{b3:thm:going-up}逐项提升。每个提升的收缩是原来指定的素理想，不同的收缩保证每次包含严格。所以提升也保留链长，得到 \(\dim B\ge\dim A\)。这比较的是所有有限链的长度，维数无穷时也成立。
\end{proof}

整包含 \(A\subseteq B\) 使两个素谱的整体维数相同，但这不能直接给出一个指定素理想 \(\mathfrak q\) 与其收缩的高度相等。Going-up 提升链时，上端的素理想由构造产生，未必等于预先指定的 \(\mathfrak q\)。若要固定这个上端，逐项向下补出原像，就需要 Going-down；当基环 \(A\) 为整闭整环、扩环 \(B\) 也为整环时，它给出下面的点处结论。
\begin{corollary}\label{b3:cor:integral-height-normal-base}
若 \(A\subseteq B\) 为整环间的整扩张，且 \(A\) 整闭，则对
\(\mathfrak q\in\operatorname{Spec}B\)、\(\mathfrak p=\mathfrak q\cap A\)，有
\(\operatorname{ht}\mathfrak q=\operatorname{ht}\mathfrak p\)。
\end{corollary}
\begin{proof}
以 \(\mathfrak q\) 为末项的链收缩后仍严格，故 \(\operatorname{ht}\mathfrak q\le\operatorname{ht}\mathfrak p\)。反向由\cref{b3:thm:going-down}，从指定的 \(\mathfrak q\) 开始向下提升每条以 \(\mathfrak p\) 为末项的有限链。不同的收缩使提升链严格，于是 \(\operatorname{ht}\mathfrak p\le\operatorname{ht}\mathfrak q\)。
\end{proof}

整同态若不是单射，应先把基环换成其像：此时能推出的是
\[
\dim B=\dim\bigl(A/\ker(A\rightarrow B)\bigr).
\]
例如商同态
\(k[x]\rightarrow k\)、\(x\mapsto0\) 是整同态，却并不保持原基环的维数。

% !TeX root = ../../Book3.tex
% 纲要标题：### 9.4 局部维数、参数系与嵌入维数
% 纲要位置：工作记录/计划纲要.md，第 2097 行。
% 纲要细目（写作范围）：
% * 参数系
% * 嵌入维数与 Krull 维数
% * 正则局部环的定义先行，结构理论留给第18章
% ---


\section{局部维数、参数系与嵌入维数}
\label{b3:sec:0904}

在 Noether 局部环中，维数有两种容易混淆的比较对象：使零点只剩闭点需要多少个方程，以及生成整个极大理想需要多少个元素。前者等于 Krull 维数，后者可能更大。两者相等时得到正则局部环的定义。


\subsection{参数系}
\label{b3:subsec:0904:01}

\begin{definition}\label{b3:def:system-of-parameters}
设 \((R,\mathfrak m)\ne0\) 为交换 Noether 局部环，\(d=\dim R\)。若一组元素
\(x_1,\ldots,x_d\in\mathfrak m\) 满足
\(\sqrt{(x_1,\ldots,x_d)}=\mathfrak m\)，则称它们为 \(R\) 的\term[canshuxi]{参数系}[system of parameters]，称它们生成的理想为\term[canshu-lixiang]{参数理想}[parameter ideal]。
\end{definition}

\begin{theorem}\label{b3:thm:system-of-parameters}
设 \((R,\mathfrak m)\ne0\) 为交换 Noether 局部环，\(d=\dim R\)。则 \(R\) 有参数系，且 \(d\) 等于生成一个根为 \(\mathfrak m\) 的理想所需的最少元素个数。
若 \(x_1,\ldots,x_d\) 是参数系，则对 \(0\le i\le d\)，有
\[
\dim R/(x_1,\ldots,x_i)=d-i.
\]
此外，对任意 \(a_1,\ldots,a_r\in\mathfrak m\)，都有
\(\dim R/(a_1,\ldots,a_r)\ge d-r\)。
\end{theorem}
\begin{proof}
存在性已由\cref{b3:lem:dimension-cutting}证明，高度定理说明少于 \(d\) 个元素不可能生成根为 \(\mathfrak m\) 的理想。
设 \(e=\dim R/(a_1,\ldots,a_r)\)，在此商环中选一个有 \(e\) 个元素的参数系并提升。
它与 \(a_1,\ldots,a_r\) 合起来生成根为 \(\mathfrak m\) 的理想，所以 \(d\le r+e\)。
对参数系的前 \(i\) 项应用这个不等式，得到商维数至少是 \(d-i\)；剩余 \(d-i\) 项在商环中生成极大根理想，高度定理又给出反向不等式。
\end{proof}

\(R/(x_1,\ldots,x_d)\) 是零维 Noether 局部环，故是 Artin 环，且作为 \(R\) 模具有有限长度。
反之，若商模具有有限长度，其支集只含闭点，于是生成理想的根为 \(\mathfrak m\)。所以参数系也可以通过商环的有限长度识别。

\ProseParagraphBreak
例如在 \(k[x,y]_{(x,y)}\) 中，\(x,y\) 和 \(x^2,y^3\) 都是参数系，后者并不生成极大理想。
在尖点局部环 \(k[t^2,t^3]_{(t^2,t^3)}\) 中，\(t^2\) 单独构成参数系：模掉它后
\((t^3)^2=(t^2)^3=0\)，故商环只有一个素理想。参数并不要求能作自由的坐标。

\subsection{嵌入维数与 Krull 维数}
\label{b3:subsec:0904:02}

\begin{definition}\label{b3:def:embedding-dimension}
设 \((R,\mathfrak m)\ne0\) 为交换 Noether 局部环，剩余域为 \(\kappa=R/\mathfrak m\)。向量空间 \(\mathfrak m/\mathfrak m^2\) 在 \(\kappa\) 上的维数
\[
\dim_\kappa(\mathfrak m/\mathfrak m^2)
\]
称为 \(R\) 的\term[qianru-weishu]{嵌入维数}[embedding dimension]，记为 \(\operatorname{edim}R\)。
\end{definition}
\symidx{\(\operatorname{edim}R\)：局部环的嵌入维数}

\(\mathfrak m/\mathfrak m^2\) 上的标量作用来自 \(R\)，并因 \(\mathfrak m\) 零化该商模而降为 \(\kappa\) 作用。
\cref{b3:cor:local-generators}说明，嵌入维数恰好等于极大理想的最少生成元数。
由高度定理立即得到
\begin{equation}\label{b3:eq:dimension-embedding-bound}
\dim R\le\operatorname{edim}R.
\end{equation}

这个向量空间仅保留函数的一次信息；乘积落入 \(\mathfrak m^2\) 而消失。取
\[
R=k[x_1,\ldots,x_n]_{(x_1,\ldots,x_n)},
\]
则变量的像构成
\(\mathfrak m/\mathfrak m^2\) 的基。生成性显然；若线性式属于局部化后的平方理想，则乘一个常数项非零的分母后，其一次部分仍是该线性式乘以非零常数，故各系数必须为零。因此嵌入维数为 \(n\)。坐标素链在此局部化中全部保留，Krull 维数也为 \(n\)。

\ProseParagraphBreak
尖点的情形不同。\(t^2,t^3\) 在 \(\mathfrak m/\mathfrak m^2\) 中线性独立，因为平方理想中每个非零元素的最低次数至少为四，局部化的单位分母不改变这一结论。因此该环的嵌入维数为二，Krull 维数为一。这两个数字之差已经能辨认这个点与直线上的点不同。

\subsection{正则局部环的定义}
\label{b3:subsec:0904:03}

\begin{definition}\label{b3:def:regular-local-ring}
若非零交换 Noether 局部环 \((R,\mathfrak m)\) 满足
\(\dim R=\dim_{R/\mathfrak m}(\mathfrak m/\mathfrak m^2)\)，则称它为
\term[zhengze-jubuhuan]{正则局部环}[regular local ring]。
\end{definition}

域是零维正则局部环；反过来，零维正则局部环的极大理想满足
\(\mathfrak m/\mathfrak m^2=0\)，由 Nakayama 引理有 \(\mathfrak m=0\)，所以它是域。
上小节的多项式局部环是正则的，尖点局部环则不是。
另一个简单的非正则例子为 \(k[\epsilon]/(\epsilon^2)\)：其 Krull 维数为零，而嵌入维数为一。

\ProseParagraphBreak
设 \((R,\mathfrak m)\) 为 \(d\) 维正则局部环。极大理想 \(\mathfrak m\) 的任一极小生成组
都有 \(d\) 项，并且生成的理想就是 \(\mathfrak m\)，因而构成参数系；这样的生成组称为\term[zhengze-canshuxi]{正则参数系}[regular system of parameters]。
正则序列在\chapref{b3:ch:15}引入，正则参数系构成正则序列的证明见\secref{b3:sec:1801}。
正规表示整闭，正则表示上述维数等式，它们是不同定义。下一章将在一维整环情形研究二者与离散赋值环的关系。

\begin{table}[htbp]
\centering
\caption[维数类不变量]{维数类不变量。\(R\) 为交换含幺环，\(M\) 为 \(R\) 模；局部环行中的 \(\mathfrak m\) 表示唯一极大理想，\(\kappa=R/\mathfrak m\)。空支集与零环的维数取 \(-\infty\)。}
\label{b3:tab:dimension-invariants}
\begin{tabularx}{\textwidth}{@{}>{\raggedright\arraybackslash}p{0.19\textwidth}>{\raggedright\arraybackslash}p{0.23\textwidth}>{\raggedright\arraybackslash}X@{}}
\toprule
不变量 & 条件 & 所测量的对象与结论\\
\midrule
\(\operatorname{ht}\mathfrak p\) & \(\mathfrak p\in\operatorname{Spec}R\) & 以 \(\mathfrak p\) 为末项的严格素链长的上确界，等于 \(\dim R_{\mathfrak p}\)\\
\(\dim(R/\mathfrak p)\) & \(\mathfrak p\in\operatorname{Spec}R\) & 以 \(\mathfrak p\) 为首项的严格素链长的上确界，即闭子空间 \(V(\mathfrak p)\) 的维数\\
\(\dim_R M\) & 任意 \(R\) 模 \(M\) & 支集 \(\operatorname{Supp}_R M\) 内的严格素链长的上确界；若 \(M\) 有限生成，等于 \(\dim(R/\operatorname{Ann}_R M)\)\\
\(\operatorname{edim}R\) & \((R,\mathfrak m)\ne0\) 为 Noether 局部环 & \(\dim_\kappa(\mathfrak m/\mathfrak m^2)\)，等于生成 \(\mathfrak m\) 所需的最少元素个数\\
生成 \(\mathfrak m\) 准素理想的最少元素数 & \((R,\mathfrak m)\ne0\) 为 Noether 局部环 & 在所有 \(\mathfrak m\) 准素理想的生成组中取元素个数的最小值，等于 \(\dim R\)；达到最小值的生成组是参数系\\
\bottomrule
\end{tabularx}
\end{table}

\cref{b3:tab:dimension-invariants}中的高度测量指定素理想以下的链，商环维数测量它以上的链。若各量有限，\cref{b3:prop:dimension-basic-comparison}给出
\[
\operatorname{ht}\mathfrak p+\dim(R/\mathfrak p)\le\dim R.
\]
\begin{samepage}
取一条经过 \(\mathfrak p\)、以极大理想 \(\mathfrak m\) 为末项的素链，可以在 \(\mathfrak p\) 处分成两段：
\[
\underbrace{\mathfrak p_0\subsetneq\cdots\subsetneq\mathfrak p}
_{u\le\operatorname{ht}\mathfrak p}
\underbrace{\subsetneq\cdots\subsetneq\mathfrak m}
_{v\le\dim(R/\mathfrak p)}.
\]
这里 \(u,v\) 是两段各自的长度。对两段分别取最长链再拼接，就得到上述不等式；原环的最长链未必经过这个指定素理想。
\end{samepage}


\begin{chaptersummary}
素链给出维数的定义，也给出最直接的下界。上界需要额外论证：Noether 条件下由理想的生成元数控制高度，有限型整代数中则由超越次数沿严格素链的下降控制链长。Lying-over 提供素链首项的提升，Going-up 逐项向上提升，Incomparability 保证收缩后的包含仍然严格，从而整扩张保持整体维数；指定点的高度比较另外用到 Going-down 的假设。

\ProseParagraphBreak
在 Noether 局部环中，参数系存在且长度恰为维数，但它生成的理想通常不是极大理想。极大理想最少需要多少个生成元，由剩余域向量空间 \(\mathfrak m/\mathfrak m^2\) 决定。尖点的 Krull 维数为一，而嵌入维数为二，说明这两种计数回答的是不同问题。
\end{chaptersummary}

\BookExercises

\begin{exercise}\label{b3:ex:09-product-dimension}
设 \(k\) 为域。对任意整数 \(d\ge2\) 和 \(0\le h\le d-2\)，构造一个 \(k\) 上的 Noether 代数 \(R\) 及其素理想 \(\mathfrak p\)，使
\[
\dim R=d,\qquad \operatorname{ht}\mathfrak p=h,\qquad
\dim(R/\mathfrak p)=1.
\]
由此说明高度与商环维数之和可以比全环维数小任意给定的正整数。
\end{exercise}
\begin{solution}
取
\[
C=k[x_1,\ldots,x_h,t],\qquad B=k[y_1,\ldots,y_d],\qquad
R=C\times B,\qquad \mathfrak p=(x_1,\ldots,x_h)C\times B.
\]
两个多项式环都 Noether，故其有限直积也 Noether。由\cref{b3:thm:polynomial-dimension}和直积维数公式，
\(\dim R=\max\{h+1,d\}=d\)。商环 \(R/\mathfrak p\cong k[t]\)，所以 \(\mathfrak p\) 素且商维数为一。
直积中的素链只留在同一个因子内，故 \(\operatorname{ht}_R\mathfrak p\) 等于 \((x_1,\ldots,x_h)\) 在 \(C\) 中的高度。逐个加入 \(x_i\) 得到长度 \(h\) 的素链，而高度定理给出反向上界，因而该高度恰为 \(h\)；\(h=0\) 时这个理想是整环 \(C\) 的零理想，高度为零。
差值为 \(d-h-1\)。给定正整数 \(g\)，取 \(h=0,d=g+1\)，便得到差值 \(g\)。
\end{solution}

\begin{exercise}\label{b3:ex:09-reduced-dimension}
设 \(k\) 为域，\(A=k[x,y,\epsilon]/(\epsilon^3)\)。计算 \(\dim A\) 和 \(\dim_A A/(x,\epsilon)\)，并解释增加幂零元为何没有增加维数。
\end{exercise}
\begin{solution}
每个素理想都含 \(\epsilon\)，故 \(A\) 的谱与 \(k[x,y]\) 有相同素链，维数二。
商模是 \(k[y]\)，其支集对应商环的谱，模维数一。幂零元给出不同的环乘法信息，却不产生新的素理想或链。
\end{solution}

\begin{exercise}\label{b3:ex:09-nonfinite-support}
将 \(\mathbb Q/\mathbb Z\) 视为 \(\mathbb Z\) 模，计算其零化子、支集和模维数，说明有限生成模的零化子公式不能任意推广。
\end{exercise}
\begin{solution}
没有非零整数零化全部有理数类，故零化子为零。在零素理想处局部化，每个元素原本被某个非零整数零化，所以局部化为零。
在 \((p)\) 处，\(1/p+\mathbb Z\) 不能被不属于 \((p)\) 的整数零化，因此局部化非零。
支集恰为全部非零素理想，这些素理想互不可比，模维数为零；而
\(\dim\bigl(\mathbb Z/\operatorname{Ann}(\mathbb Q/\mathbb Z)\bigr)=1\)。
\end{solution}

\begin{exercise}\label{b3:ex:09-support-sum}
设 \(k\) 为域，\(R=k[x,y,z]\)。计算 \(R/(x)\oplus R/(y,z)\) 及子模
\(xR\subseteq R\) 的维数。再由短正合列核对 \(\dim R\)。
\end{exercise}
\begin{solution}
两直和项的维数为二和一，直和维数为二。乘以 \(x\) 给出 \(R\simeq xR\)，故后者维数三。
短正合列 \(0\rightarrow xR\rightarrow R\rightarrow R/(x)\rightarrow0\) 的维数公式为
\(3=\max\{3,2\}\)。非零子模不必比商模低维。
\end{solution}

\begin{exercise}\label{b3:ex:09-local-coordinate-height}
设 \(k\) 为域，\(n\ge0\) 为整数，\(R=k[x_1,\ldots,x_n]\)。令 \(\mathfrak p=(x_1,\ldots,x_r)\)、\(0\le r\le n\)。
计算 \(\dim R_{\mathfrak p}\) 和 \(\dim R/\mathfrak p\)。
\end{exercise}
\begin{solution}
前 \(r\) 个变量逐个加入给出末项为 \(\mathfrak p\) 的长度 \(r\) 素链。
另一方面 \(\mathfrak p\) 在自己生成的 \(r\) 元理想上极小，高度定理给出高度至多 \(r\)。所以
\(\dim R_{\mathfrak p}=\operatorname{ht}\mathfrak p=r\)，而商环为 \(n-r\) 元多项式环，维数为 \(n-r\)。
\end{solution}

\begin{exercise}\label{b3:ex:09-height-components}
设 \(k\) 为域，\(R=k[x,y,z]\)。求 \(I=(xy,xz)\) 的极小素理想及其高度，并计算 \(\dim R/I\)。
\end{exercise}
\begin{solution}
包含 \(I\) 的素理想若不含 \(x\)，必同时含 \(y,z\)，所以极小素理想恰为
\((x)\) 与 \((y,z)\)，高度分别为一和二。这符合两个生成元的高度上界，却不要求所有极小素理想有相同高度。
两个不可约分量的坐标环维数分别为二和一，因此 \(\dim R/I=2\)。
\end{solution}

\begin{exercise}\label{b3:ex:09-primary-generators}
设 \(k\) 为域，\(R=k[x,y]_{(x,y)}\)，\(a,b\) 为正整数。证明任何 \((x,y)R\) 准素理想都不能由一个元素生成。
再求 \(R/(x^a,y^b)\) 的长度。
\end{exercise}
\begin{solution}
该局部环的维数为二。若一个主理想的根为极大理想，主理想定理会给出维数至多一，矛盾。
所给商环中 \(x^iy^j\)、\(0\le i<a,0\le j<b\)，构成 \(k\) 基；常数项非零的分母在这个环中已经可逆，因为其非常数部分幂零。商环为局部 Artin 环且剩余域为 \(k\)，每个简单因子是一维 \(k\) 空间，所以长度为 \(ab\)。
\end{solution}

\begin{exercise}\label{b3:ex:09-hypersurface-dimension}
设 \(k\) 为域，\(n\ge1\) 为整数，\(f\in k[x_1,\ldots,x_n]\) 为非零非单位。证明
\(\dim k[x_1,\ldots,x_n]/(f)=n-1\)。
\end{exercise}
\begin{solution}
先设 \(f\) 不可约，重排变量使其对 \(x_n\) 的次数为正。它不能整除只含前
\(n-1\) 个变量的非零多项式，所以这些变量在商环中仍代数独立。
关系 \(f=0\) 又说明 \(x_n\) 在它们的分式域上代数；首项系数虽未必为单位，在分式域中却可逆。
因此商整环的分式域超越次数为 \(n-1\)，用有限型维数公式得结论。
一般 \(f\) 的极小素理想由其不可约因子生成；有限多个对应商环都有维数 \(n-1\)，故结论不变。
\end{solution}

\begin{exercise}\label{b3:ex:09-open-dimension}
设 \(k\) 为域，\(A\) 为有限型整 \(k\) 代数，\(0\ne f\in A\)。证明 \(\dim A_f=\dim A\)。
解释这为何不与 \(\dim A_{(0)}=0\) 冲突。
\end{exercise}
\begin{solution}
只倒置 \(f\) 可写为 \(A_f\cong A[u]/(fu-1)\)，所以 \(A_f\) 仍是有限型整 \(k\) 代数。它与 \(A\) 具有相同分式域，由\cref{b3:thm:affine-dimension}得到两者维数相等。
零素理想处的局部化则倒置全部非零元素，得到 \(A_{(0)}=\operatorname{Frac}(A)\)。若 \(\dim A>0\)，这个域不可能是有限型 \(k\) 代数：否则由 Zariski 引理\cref{b3:lem:zariski}，它在 \(k\) 上有限代数，从而超越次数为零，与 \(\dim A>0\) 矛盾。因而不能对它继续使用有限型整代数的维数公式。例如 \(k[t]_{(0)}=k(t)\) 的环维数为零，超越次数却为一。
\end{solution}

\begin{exercise}\label{b3:ex:09-monomial-curve-dimension}
设 \(k\) 为域，\(t\) 为不定元，\(a,b\) 为互素正整数。计算 \(k[t^a,t^b]\) 的维数和分式域，并说明哪一步使用互素性。
\end{exercise}
\begin{solution}
整数 Bézout 等式 \(ra+sb=1\) 给出 \(t=(t^a)^r(t^b)^s\)，允许负指数后说明分式域为 \(k(t)\)，所以维数一。
互素性只用于认出分式域等于 \(k(t)\)。若最大公因数为 \(d\)，相同论证给出分式域 \(k(t^d)\)，超越次数仍为一，因此维数仍为一。
\end{solution}

\begin{exercise}\label{b3:ex:09-nonfinite-integral}
设 \(L/k\) 为无限次代数扩张。证明 \(k[t]\subseteq L[t]\) 为非有限的整扩张，并计算两环维数。
\end{exercise}
\begin{solution}
任意一个 \(L[t]\) 元素只有有限多个系数，它们生成有限域扩张 \(k'/k\)。环
\(k'[t]\) 是 \(k[t]\) 上有限自由模，所以其中每个元素都整。这证明整性。
若 \(L[t]\) 是有限生成的 \(k[t]\) 模，模掉 \(t\) 后会使 \(L\) 成为有限维 \(k\) 空间，矛盾。
两环分别是一元多项式环，维数均为一，符合整扩张的维数不变性。
\end{solution}

\begin{exercise}\label{b3:ex:09-local-integral-height}
设 \(k\) 为域，\(A=k[x,y]\)，\(B=k[x,y,z]/(z^2-xy)\)。证明 \(A\) 自然嵌入 \(B\)，且 \(B\) 是整环。计算 \(B\) 在素理想 \((x,z)\) 和 \((x,y,z)\) 处的局部维数，并说明高度比较为什么需要把素链的末项固定。
\end{exercise}
\begin{solution}
关于 \(z\) 的首一关系允许带余除法，每个商类都唯一写成 \(a+bz\)，其中 \(a,b\in A\)。因此 \(B\cong A\oplus Az\) 作为 \(A\) 模，常数商类使 \(A\rightarrow B\) 单射；\(B\) 又是有限生成 \(A\) 模，故该包含为有限整扩张。

整环性另由关系的不可约性得到。多项式 \(z^2-xy\) 对 \(k[x,y]\) 中的素元 \(x\) 满足 Eisenstein 判据\footnote{艾森斯坦（Ferdinand Gotthold Max Eisenstein，1823-1852），德国数学家，研究数论与椭圆函数。}：常数项被 \(x\) 整除但不被 \(x^2\) 整除，其余非首项系数也被 \(x\) 整除。因此它不可约，唯一分解环 \(k[x,y,z]\) 中的不可约元为素元，故商环 \(B\) 是整环。
两理想的商环分别为 \(k[y]\) 和 \(k\)，所以确为素理想；收缩到 \(A\) 后分别是 \((x)\) 和 \((x,y)\)。基环 \(A\) 为整闭整环，\cref{b3:cor:integral-height-normal-base}使两处高度分别等于一和二，这也就是所求局部维数。
Going-up 从已选的底端素理想向上构造链，末项未必是这两个指定素理想之一；Going-down 则能从指定末项向下提升。这解释了所用高度比较中整闭基环的作用。
\end{solution}

\begin{exercise}\label{b3:ex:09-parameter-zero-divisor}
设 \(k\) 为域，\(n\ge2\) 为整数，\(R=\bigl(k[x,y]/(x^2,xy^n)\bigr)_{(x,y)}\)。计算 \(R\) 的 Krull 维数、嵌入维数及 \(\operatorname{Ann}_R(y)\)。再对 \(a,b\in k\)，判定 \(ax+by\) 何时构成参数，并证明满足条件的这些参数全是零因子。
\end{exercise}
\begin{solution}
商环中的每个元素唯一写成 \(f(y)+xg(y)\)，其中 \(f\in k[y]\)，\(g\in k[y]/(y^n)\)。局部化后相应地有
\[
R=k[y]_{(y)}\oplus x\bigl(k[y]/(y^n)\bigr)
\]
作为 \(k[y]_{(y)}\) 模的分解：常数项非零的分母在第二项上也可逆，而第二项的平方为零。幂零根因此为 \((x)\)，约化局部环为 \(k[y]_{(y)}\)，故 \(\dim R=1\)。原关系都没有一次项，\(x,y\) 的像在 \(\mathfrak m/\mathfrak m^2\) 中构成一组基，故 \(\operatorname{edim}R=2\)。

若 \(y(f+xg)=0\)，上面的直和分解使 \(yf=0\) 及 \(yg=0\) 在各自的项中成立。第一项是整环，故 \(f=0\)；第二项中 \(yg=0\) 当且仅当 \(g\in k y^{n-1}\)。所以
\[
\operatorname{Ann}_R(y)=(xy^{n-1}),\qquad xy^{n-1}\ne0.
\]
元素 \(ax+by\) 在约化环中的像是 \(by\)。若 \(b=0\)，该元素幂零，生成理想的根为 \((x)\)，不等于 \(\mathfrak m\)。若 \(b\ne0\)，商环中可令 \(y=-ab^{-1}x\)，得到 \(R/(ax+by)\cong k[x]/(x^2)\)，其维数为零，所以该元素是参数。此时
\((ax+by)xy^{n-1}=ax^2y^{n-1}+bxy^n=0\)，被零化的 \(xy^{n-1}\) 非零，故这个参数是零因子。
\end{solution}

\begin{exercise}\label{b3:ex:09-parameter-powers}
设 \((R,\mathfrak m)\ne0\) 为交换 Noether 局部环，\(d=\dim R\)，\(x_1,\ldots,x_d\) 为 \(R\) 的参数系。证明其中每个元素可独立换为任意正整数次幂，仍得参数系。又证明它们可以任意重排，且每个前 \(i\) 项的商环维数为 \(d-i\)（\(0\le i\le d\)）。
\end{exercise}
\begin{solution}
一个素理想包含元素的正整数次幂当且仅当包含元素本身，因此这些理想有相同根。
重排不改变生成理想。对重排后的参数系应用\cref{b3:thm:system-of-parameters}，便得每个前缀的维数公式。
\end{solution}

\begin{exercise}\label{b3:ex:09-artin-regular}
设 \(R\ne0\) 为交换 Artin 局部环。证明 \(R\) 正则当且仅当它是域。再设 \(k\) 为域、\(n\ge0\) 为整数，计算
\(k[x_1,\ldots,x_n]/(x_1,\ldots,x_n)^2\) 的嵌入维数。
\end{exercise}
\begin{solution}
Artin 环 Noether 且零维。正则性等价于 \(\mathfrak m/\mathfrak m^2=0\)，Nakayama 引理使 \(\mathfrak m=0\)，故为域；域显然正则。
所给环的极大理想平方为零，\(x_1,\ldots,x_n\) 的像构成其 \(k\) 基，所以嵌入维数为 \(n\)。当 \(n>0\) 时它不正则。
\end{solution}

\begin{exercise}\label{b3:ex:09-tangent-linear-relations}
设 \(k\) 为域，\(n\ge1\) 为整数。令 \(\mathfrak n=(x_1,\ldots,x_n)\subseteq k[x_1,\ldots,x_n]\)，并设
\(I=(f_1,\ldots,f_r)\subseteq\mathfrak n\)，\(R=\bigl(k[x_1,\ldots,x_n]/I\bigr)_{\mathfrak n/I}\)。
证明 \(\operatorname{edim}R\) 等于 \(n\) 减去这些 \(f_i\) 的一次部分的秩。
\end{exercise}
\begin{solution}
极大理想的平方商为 \(\mathfrak n/(\mathfrak n^2+I)\)，局部化不改变这个被 \(\mathfrak n\) 零化的 \(k\) 向量空间。
模掉 \(\mathfrak n^2\) 后，每个 \(f_i\) 只留下其一次部分；任意组合
\(\sum_i g_if_i\) 的一次部分为 \(\sum_i g_i(0)(f_i)_{\mathrm{lin}}\)。
所以全部关系恰是这些线性形式的张成空间，商空间维数为题述数值。
\end{solution}

\begin{exercise}\label{b3:ex:09-nonpure-parameter}
设 \(k\) 为域，\(R=\bigl(k[x,y,z]/(xy,xz)\bigr)_{(x,y,z)}\)。证明其维数为二、嵌入维数为三，并验证 \(x+y,z\) 是参数系。
\end{exercise}
\begin{solution}
两个极小素理想对应的商环为局部化后的 \(k[y,z]\) 和 \(k[x]\)，维数分别为二和一，所以总维数二。
所定义关系都为二次式，没有一次关系，嵌入维数三。
令 \(z=0,y=-x\) 后，商环为 \(k[x]_{(x)}/(x^2)\)，是零维局部环，故两个元素构成参数系。
这个例子也说明参数系的存在不要求各个不可约分量等维。
\end{solution}

\begin{exercise}\label{b3:ex:09-module-parameters}
设 \((R,\mathfrak m)\) 为交换 Noether 局部环，\(M\ne0\) 为有限生成的 \(R\) 模，\(d=\dim_R M\)。
证明存在 \(x_1,\ldots,x_d\in\mathfrak m\) 使 \(M/(x_1,\ldots,x_d)M\) 具有有限长度，并证明用 \(\mathfrak m\) 中少于 \(d\) 个元素不能做到这一点。
\end{exercise}
\begin{solution}
由\cref{b3:thm:support-finite}，\(\operatorname{Supp}_R M=V(\operatorname{Ann}_R M)\)，所以 \(A=R/\operatorname{Ann}_R M\) 的素谱恰好记录 \(M\) 的支集，且 \(\dim A=d\)。对 \(A\) 取参数系，相当于在这个支集上选出足够多的方程。将参数提升到 \(R\)，得到理想
\(I=(x_1,\ldots,x_d)\)。有限生成模的支集及局部 Nakayama 引理说明
\[
\operatorname{Supp}(M/IM)=\operatorname{Supp}(M)\cap V(I).
\]
具体地，局部化后若 \(I\not\subseteq\mathfrak p\)，商为零；若 \(I\subseteq\mathfrak p\) 且
\(M_{\mathfrak p}\ne0\)，Nakayama 引理禁止其等于 \(I M_{\mathfrak p}\)，所以商非零。
令 \(Q=M/IM\)。参数条件使 \(\operatorname{Supp}_R Q=\{\mathfrak m\}\)；\(Q\) 有限生成，故由支集的零化子公式得
\[
\sqrt{\operatorname{Ann}_R Q}=\mathfrak m.
\]
这意味着 \(S=R/\operatorname{Ann}_R Q\) 是只有一个素理想的 Noether 环，故由\cref{b3:thm:noether-dimension-zero}为 Artin 环。现在 \(Q\) 是有限生成的 \(S\) 模，\cref{b3:thm:artin-finite-length}给出其有限长度。这里调用的证明正是利用幂零极大理想的有限滤过：\(\mathfrak m^NQ=0\) 对某个 \(N\) 成立，且各商层 \(\mathfrak m^iQ/\mathfrak m^{i+1}Q\) 是有限维剩余域向量空间。\(R\) 与 \(S\) 在 \(Q\) 上有相同的子模，因此这也是它作为 \(R\) 模的长度。

反之，若 \(J=(y_1,\ldots,y_r)\subseteq\mathfrak m\) 使 \(M/JM\) 有有限长度，Nakayama 引理保证这个商非零。其支集因此为 \(\{\mathfrak m\}\)。由上述支集等式，\(J\) 在 \(A\) 中的像的根是 \(\mathfrak m A\)；高度定理给出 \(d=\dim A\le r\)，所以少于 \(d\) 个元素不能做到这一点。
\end{solution}

% !TeX root = ../../Book3.tex
% 纲要标题：## 第10章　赋值环、离散赋值环与 Dedekind 整环
% 纲要位置：工作记录/计划纲要.md，第 2228 行。
% 纲要细目（写作范围）：

\chapter{赋值环、离散赋值环与 Dedekind 整环}
\label{b3:ch:10}
\BookChapterNumber{b3:ch:10}

\begin{chapterabstract}
赋值把整除关系化为有序群中的大小比较。离散赋值环中，每个非零元素由统一化元的幂和一个单位相乘而成；本章从一维、Noether 性与整闭性出发证明这一结构。Dedekind 整环把这些局部整数阶合成为全局理想的唯一分解。数域整数环、理想类群和正规曲线展示这种分解怎样应用于算术与几何。
\end{chapterabstract}

\begin{learninggoals}
\begin{itemize}
\item 从赋值构造局部环，区分赋值的秩、有理秩与离散性。
\item 证明一维 Noether 正规局部整环是 DVR，计算理想的阶数与商模长度。
\item 通过局部化恢复 Dedekind 整环的非零理想分解，并构造分式理想的逆。
\item 区分可逆理想与主理想，说明理想类群怎样由秩一投射模的同构类描述。
\end{itemize}
\end{learninggoals}

整数的素因子分解唯一，换到别的整环，这个熟悉的结论还成立吗？在 \(\mathbb Z[\sqrt{-5}]\) 中，
\[
6=2\cdot3=(1+\sqrt{-5})(1-\sqrt{-5}).
\]
这里的四个因子都不可约，两组因子也不能通过乘单位和重排彼此得到；本章稍后用范数核实这些事实。元素分解的唯一性因此失效。若把元素换成它生成的理想，能否重新找到唯一的分解？先回看整数中的一个例子：
\[
(60)=(2)^2(3)(5),\qquad v_2(60)=2,\quad v_3(60)=v_5(60)=1.
\]
除这三个素数外，其余素数在 \(60\) 中的指数都是零。整数分解已经可以用各素数处的指数完整记录；Dedekind 理论把这些指数推广为非零理想在各非零素理想处的阶数。离散赋值使局部指数仍为整数，而局部化让其他素因子成为单位，从而只留下眼前这个素点的指数。

\ProseParagraphBreak

在一个离散赋值环里，所有非零元素都可按同一统一化元的幂排出层次，理想也随之线性排列。Dedekind 整环要求在各高度一的点处出现这种简单局部结构；要从局部指数拼出全局理想分解，还需证明只涉及有限多个素理想，以及拼出的理想确实唯一。元素分解失败与理想分解成功的差别就在这里。

\ProseParagraphBreak

本章使用第2章的局部化检测、第3章的行列式技巧与 Nakayama 引理、第5—6章的整闭性和正规化有限性，以及第9章的维数语言。数域整数环的自由性还使用第二卷主理想整环上的有限生成模结构定理。Picard 群\footnote{皮卡尔（Charles Émile Picard，1856-1941），法国数学家，研究微分方程与复分析。}部分用到的有限投射模、对偶和张量积，也已在第二卷建立。

\ProseParagraphBreak
一般赋值可参阅 Matsumura 的《Commutative Ring Theory》\cite[\S10, General valuations, pp. 71--77]{Matsumura1989CommutativeRingTheory}。DVR 的等价刻画与可逆理想见\cite[Theorems 11.1--11.3, pp. 78--81]{Matsumura1989CommutativeRingTheory}；Dedekind 整环的分解定理见\cite[Theorem 11.6, pp. 82--84]{Matsumura1989CommutativeRingTheory}。本章保留一般赋值与离散赋值的区别，再在 Noether 条件下开展理想分解。

% !TeX root = ../../Book3.tex
% 纲要标题：### 10.1 赋值与赋值环
% 纲要位置：工作记录/计划纲要.md，第 2230 行。
% 纲要细目（写作范围）：
% * 赋值群及赋值环的整除性质
% * 一般赋值环与离散赋值环的区别
% * 分式域和局部性
% * 赋值的秩与赋值环的 Krull 维数：用凸子群与素理想的反包含对应说明维数，区分秩与有理秩，并使用第9章已建立的素链语言

\section{赋值与赋值环}
\label{b3:sec:1001}

整数中可以比较两个数含有多少个素因子 \(p\)，有理函数中也可以比较它们在一点的零点或极点阶数。这些阶数在乘法中相加，例如 \(p^r p^s=p^{r+s}\)；相加时则可能因最低阶项抵消而升高阶数。赋值把这种计算写成有序群中的加法。允许一般的有序群，才能看清哪些性质来自整除的全序性，哪些性质另外需要离散性或 Noether 性。


\subsection{赋值群与赋值环的整除性质}
\label{b3:subsec:1001:01}

\begin{definition}\label{b3:def:valuation}
设 \(\Gamma\) 为带有全序的加法交换群。若对所有 \(\alpha,\beta,\gamma\in\Gamma\) 有
\(\alpha\le\beta\Rightarrow\alpha+\gamma\le\beta+\gamma\)，则称 \(\Gamma\) 为\term[quanxu-jiaohuanqun]{全序交换群}[totally ordered abelian group]。
设 \(K\) 为域，\(\Gamma\) 为全序交换群。给定满射
\(v:K^*\rightarrow\Gamma\)，并规定 \(v(0)=+\infty\)。若对非零 \(a,b\in K\) 满足第一式，且对任意 \(a,b\in K\) 满足第二式：
\[
v(ab)=v(a)+v(b),\qquad
v(a+b)\ge\min\bigl\{v(a),v(b)\bigr\},
\]
则称 \(v\) 为 \(K\) 上取值于 \(\Gamma\) 的\term[fuzhi]{赋值}[valuation]。其中 \(+\infty\) 大于每个群元素，且
\(\gamma+(+\infty)=+\infty\)。称 \(\Gamma\) 为赋值群。
\end{definition}
\symidx{\(v:K^*\to\Gamma\)：域的赋值及其值群}

本书约定 \(\Gamma\) 就是实际取到的值群，所以要求 \(v\) 满射。若采用允许映射不满射的定义，只须把所给有序群换成子群 \(v(K^*)\)，就得到这里的约定；赋值不等式及相应赋值环都不改变。

\ProseParagraphBreak
由第一条公理有 \(v(1)=0\)、\(v(a^{-1})=-v(a)\)。全序群无挠，故由
\(2v(-1)=0\) 得 \(v(-1)=0\)。如果 \(v(a)<v(b)\)，则
\(v(a+b)=v(a)\)：否则把 \(a=(a+b)-b\) 代入三角不等式便矛盾。
这说明最低阶项只有在至少两项同阶时才可能互相抵消。例如对任意域 \(k\) 及不定元 \(t\)，\(t^2+t^5=t^2(1+t^3)\)，括号内在 \(t=0\) 处的值为一，所以和在零点处的阶仍为二。

\begin{definition}\label{b3:def:valuation-ring}
设 \(K\) 为域，\(v:K^*\rightarrow\Gamma\) 为取值于全序交换群 \(\Gamma\) 的赋值。定义子环
\[
V_v=\{0\}\cup\bigl\{a\in K^*\mid v(a)\ge0\bigr\},
\]
并称 \(V_v\) 为 \(v\) 的\term[fuzhihuan]{赋值环}[valuation ring]。
等价地，若整环 \(V\) 的分式域为 \(K\)，且对每个 \(x\in K^*\) 都有
\(x\in V\) 或 \(x^{-1}\in V\)，则称 \(V\) 为 \(K\) 的赋值环。允许 \(V=K\)，这对应零赋值群，称为平凡赋值。
\end{definition}
\symidx{\(V_v\)：赋值 \(v\) 的赋值环}

说明两种说法确实等价。赋值不等式保证 \(V_v\) 对加法、减法、乘法封闭，且含 \(1\)；全序性给出倒数条件。
反过来，若 \(V\) 满足倒数条件，乘一个 \(V\) 的单位不会改变整除关系。于是把只差一个单位的非零元素看成具有同一个阶，即把满足 \(a/b\in V^*\) 的 \(a,b\) 识别为同一类。所得商群是 \(\Gamma=K^*/V^*\)，它忘掉单位因子，保留整除的比较；把商群运算写成加法，并规定
\[
[a]\le[b]\quad\Longleftrightarrow\quad b/a\in V.
\]
换乘单位不改变此条件；若两个方向同时成立，商为单位，故两类相等。
倒数条件给出全序，乘法给出与平移的相容性。设 \(v(a)=[a]\)，若
\(v(a)\le v(b)\)，则 \((a+b)/a=1+b/a\in V\)，于是赋值不等式成立。

\begin{proposition}\label{b3:prop:valuation-ideal-comparison}
设 \(v:K^*\rightarrow\Gamma\) 为域 \(K\) 上的赋值，\(V_v\) 为其赋值环。则 \(V_v\) 的任意两个理想可以按包含关系比较。因此每个有限生成理想都是主理想。
对非零元素 \(a,b\in V_v\)，有
\(a\mid b\Longleftrightarrow v(a)\le v(b)\)。
\end{proposition}
\begin{proof}
末一等价式就是 \(b/a\in V_v\)。全序性因此先保证两个非零元素生成的主理想可以比较，再由此比较任意理想。若理想 \(I\) 不包含于 \(J\)，取
\(a\in I\setminus J\)。任取 \(0\ne b\in J\)，不能有 \(a\in bV_v\)，故必有
\(b\in aV_v\subseteq I\)，于是 \(J\subseteq I\)。

若 \(I=(a_1,\ldots,a_r)\ne0\)，在非零生成元中取赋值最小的 \(a_j\)。每个非零 \(a_i\) 都满足 \(v(a_j)\le v(a_i)\)，故 \(a_j\mid a_i\)，从而 \(I=a_jV_v\)。零理想本身由零生成，所以每个有限生成理想都是主理想。
\end{proof}

\subsection{一般赋值环与离散赋值环}
\label{b3:subsec:1001:02}

\begin{definition}\label{b3:def:discrete-valuation-ring}
设 \(K\) 为域，\(v:K^*\rightarrow\Gamma\) 为域 \(K\) 上的赋值。若赋值群 \(\Gamma\) 与通常排序的 \(\mathbb Z\) 有序同构，则称相应赋值环 \(V_v\) 为
离散赋值环\termidx[lisan-fuzhi-huan]{离散赋值环}，缩写为 DVR。本书的 DVR 不是域。
\end{definition}

“离散”说的是赋值群由整数阶组成，并不是说极大理想的幂所定义的拓扑为离散拓扑。

\ProseParagraphBreak
对素数 \(p\)，将非零有理数写成 \(p^n a/b\)，其中 \(a,b\) 为整数且 \(p\nmid ab\)。指数 \(n\) 唯一，定义
\(v_p(p^na/b)=n\)，所得赋值环是 \(\mathbb Z_{(p)}\)。

\ProseParagraphBreak
设 \(k\) 为域，\(t\) 为不定元。每个非零 \(f\in k[t]\) 唯一写成 \(f=t^ru\)，其中 \(r\ge0\)、\(u\in k[t]\) 且 \(u(0)\ne0\)。记 \(\nu_t(f)=r\)，即 \(t\) 在 \(f\) 中作为因子的重数。对 \(f,g\in k[t]\setminus\{0\}\)，定义
\[
\operatorname{ord}_0(f/g)=\nu_t(f)-\nu_t(g).
\]
因子重数与多项式次数不同，例如 \(\operatorname{ord}_0(t+1)=0\)。这个赋值的赋值环为 \(k[t]_{(t)}\)。两例中乘法阶数相加、加法阶数不下降，都可直接由提取相应素因子的幂验证。

\ProseParagraphBreak
一般赋值未必离散。在 \(\mathbb Q\) 中，任何正值 \(q\) 下方仍有正值 \(q/2\)，所以正值没有最小元。若一个赋值以 \(\mathbb Q\) 为实际值群，正赋值元素连同零组成的理想就不能有限生成：有限生成会给出一个赋值最小的生成元，其正值记为 \(q\)；满射性却给出赋值为 \(q/2\) 的元素，它属于这个理想而不能被该生成元整除。

\ProseParagraphBreak
这样的赋值可如下构造。设 \(k\) 为域，\(t\) 为不定元，在一个固定扩域中相容选取 \(t^{1/n!}\)，并令
\(K=\bigcup_{n\ge1}k(t^{1/n!})\)；
取 \(v(t^{1/n!})=1/n!\)，在 \(k(t^{1/n!})\) 中把以 \(t^{1/n!}\) 为变量的零点阶除以 \(n!\)。这些赋值在嵌套子域上相容，因而定义了 \(K\) 上以 \(\mathbb Q\) 为值群的赋值。
每个有理数的分母都整除某个 \(n!\)，所以各子域的值群 \(\dfrac1{n!}\mathbb Z\) 的并恰为 \(\mathbb Q\)。
因此所得赋值环不是 Noether 环。

\subsection{分式域与局部性}
\label{b3:subsec:1001:03}

\begin{proposition}\label{b3:prop:valuation-local-normal}
设 \(v:K^*\rightarrow\Gamma\) 为域 \(K\) 上的赋值，\(V_v\) 为其赋值环。则 \(V_v\) 的分式域是 \(K\)，它是局部整环，唯一极大理想为
\[
\mathfrak m_v=\{0\}\cup\bigl\{a\in K^*\mid v(a)>0\bigr\},
\]
单位恰为赋值零的元素。每个赋值环都在 \(K\) 中整闭。
\end{proposition}
\symidx{\(\mathfrak m_v\)：赋值环中正赋值元素组成的极大理想}
\begin{proof}
任意 \(x\in K^*\) 或属于 \(V_v\)，或为 \(V_v\) 中非零元素的倒数，故分式域为 \(K\)。
赋值公理说明 \(\mathfrak m_v\) 是理想；非零 \(a\in V_v\) 可逆恰当 \(v(a)=0\)，所以非单位全体恰是该真理想。这证明局部性。

整闭性要求排除负赋值元素满足 \(V_v\) 系数首一方程的可能。若 \(x\in K\) 满足首一方程 \(x^n+a_1x^{n-1}+\cdots+a_n=0\)、\(a_i\in V_v\)，而 \(v(x)<0\)，则 \(x^n\) 的赋值严格低于其余各个非零项的赋值：系数的赋值非负，较低次的 \(x\) 的幂反而具有较高赋值。这唯一的最低阶项不能被抵消。把这一点写成理想中的等式，除以 \(x^n\) 得
\(1+a_1x^{-1}+\cdots+a_nx^{-n}=0\)。除首项外每项均属于 \(\mathfrak m_v\)，于是 \(1\in\mathfrak m_v\)，矛盾。因此 \(x\in V_v\)。
\end{proof}

局部性不保证是赋值环。例如对域 \(k\)，在 \(R=k[x,y]_{(x,y)}\) 中，主理想 \((x)\) 与 \((y)\) 不能比较：模掉 \((y)\) 后得到 \(k[x]_{(x)}\)，其中 \(x\) 非零，故 \(x\notin(y)\)；交换两变量，同样有 \(y\notin(x)\)。所以 \(R\) 不是赋值环；这两个不包含关系也说明 \(x/y\) 与 \(y/x\) 均不在 \(R\) 中。多项式环 \(k[x,y]\) 整闭，由\cref{b3:cor:integrally-closed-localization}，\(R\) 也整闭，因此整闭局部环仍不必是赋值环。

\subsection{赋值的秩与赋值环的 Krull 维数}
\label{b3:subsec:1001:04}

设 \(\Gamma\) 为全序交换群，\(H\) 为其子群。若对所有 \(\gamma\in\Gamma\) 和 \(h\in H\)，
\(0\le\gamma\le h\) 均蕴含 \(\gamma\in H\)，则称 \(H\) 为\term[tuzhiqun]{凸子群}[convex subgroup]。
凸子群按包含关系组成全序：若 \(H_1\not\subseteq H_2\)，取正的
\(h\in H_1\setminus H_2\)，则 \(H_2\) 的每个正元素均小于 \(h\)，故属于 \(H_1\)。
赋值的\term[fuzhi-zhixu]{秩}[rank of a valuation]定义为其赋值群中有限严格凸子群链的长度上确界。它测量凸子群能分出多少层，所数的不是群生成元的个数。赋值的\term[youli-zhixu]{有理秩}[rational rank]则定义为
\[
\operatorname{rat.rank}(v)
=\dim_{\mathbb Q}(\Gamma\otimes_{\mathbb Z}\mathbb Q),
\]
它测量值群有理化后的向量空间维数。
\symidx{\(\operatorname{rat.rank}(v)\)：赋值 \(v\) 的有理秩}

\begin{samepage}
以按字典序排序的 \(\Gamma=\mathbb Z^2\) 为例，先比较第一坐标，第一坐标相同时才比较第二坐标。因此 \((0,n)\) 对任意整数 \(n\) 都小于 \((1,0)\)。子群 \(H=\{0\}\times\mathbb Z\) 是凸的：夹在零与 \(H\) 的正元素之间的元素，第一坐标只能为零。非零凸子群若只含第一坐标为零的元素，凸性便迫使它包含 \((0,1)\)，故等于 \(H\)；若含第一坐标为正的元素，其正整数倍最终超过任意给定正元素，凸性便使该子群等于全群。

\ProseParagraphBreak
对于值群为 \(\mathbb Z^2\) 的赋值，令 \(\mathfrak p\) 由零元素以及赋值第一坐标为正的元素组成。它只检测第一层阶数，忽略第二坐标。凸子群与素理想之间的对应排列为
\[
\begin{array}{ccccc}
0&\subsetneq&\{0\}\times\mathbb Z&\subsetneq&\mathbb Z^2\\
\updownarrow&&\updownarrow&&\updownarrow\\
\mathfrak m_v&\supsetneq&\mathfrak p&\supsetneq&(0).
\end{array}
\]
保留的凸子群越大，对应素理想越小。一般值群也有这样的反向对应。
\end{samepage}

\begin{proposition}\label{b3:prop:valuation-rank-dimension}
设 \(v:K^*\rightarrow\Gamma\) 为域 \(K\) 上的赋值，\(V_v\) 为其赋值环。则 \(\Gamma\) 的凸子群与 \(V_v\) 的素理想有反包含对应
\[
H\longmapsto\mathfrak p_H
=\{0\}\cup\bigl\{a\in V_v\setminus\{0\}\mid v(a)>h\;\text{对所有}\;h\in H\bigr\}.
\]
特别地，赋值的秩等于赋值环的 Krull 维数。
\end{proposition}
\begin{proof}
对非负群元素，属于 \(H\) 或大于 \(H\) 中所有元素，二者必居其一；这是凸性的直接结果。
所以 \(V_v\setminus\mathfrak p_H\) 恰由赋值属于 \(H\) 的非零元素组成。
这个集合对乘法封闭，而 \(\mathfrak p_H\) 对相加、乘以 \(V_v\) 元素封闭，故为素理想。

反过来，给定素理想 \(\mathfrak p\)，令
\(S=\bigl\{v(a):a\in V_v\setminus\mathfrak p\bigr\}\subseteq\Gamma_{\ge0}\)。
它对加法封闭。若 \(0\le\gamma\le s=v(a)\in S\)，取 \(b\in V_v\) 使 \(v(b)=\gamma\)；
因 \(b\mid a\)，\(b\in\mathfrak p\) 会迫使 \(a\in\mathfrak p\)，所以 \(\gamma\in S\)。
于是 \(S\) 向下封闭；若 \(s,t\in S\) 且 \(s\ge t\)，则 \(0\le s-t\le s\)，故 \(s-t\in S\)。所以 \(H=S-S=S\cup(-S)\) 是凸子群，其非负部分恰为 \(S\)。
因此 \(\mathfrak p=\mathfrak p_H\)，两构造互逆。包含的反转由公式直接看出。
有限严格链由此互相对应，长度不变。
\end{proof}

实数加法群的任意非零有序子群都只有两个凸子群：若凸子群含正元素 \(h\)，则 \(h\) 的正整数倍最终超过任意给定正元素，凸性使它等于全群。因此通常排序的 \(\mathbb Z\)、\(\mathbb Q\) 和 \(\mathbb Z+\sqrt2\mathbb Z\) 的秩均为一。前两个群有理化后都是一维 \(\mathbb Q\) 向量空间，有理秩为一；第三个群中的 \(1,\sqrt2\) 在 \(\mathbb Q\) 上线性独立，有理秩为二。字典序 \(\mathbb Z^2\) 的凸子群链则有两层，秩为二，其有理秩也为二。

\ProseParagraphBreak
这四个值群中，只有通常排序的 \(\mathbb Z\) 对应本书的 DVR。\(\mathbb Q\) 和 \(\mathbb Z+\sqrt2\mathbb Z\) 虽然秩一，却不与 \(\mathbb Z\) 有序同构；字典序 \(\mathbb Z^2\) 虽有最小正值 \((0,1)\)，秩仍为二。因此秩一或有最小正值都不能单独替代离散赋值环的定义。上面的 \(\mathbb Q\) 值群例子是一维整闭局部环，却不是 DVR；下一节将考察一维局部整环再满足 Noether 性和整闭性时的结构。

% !TeX root = ../../Book3.tex
% 纲要标题：### 10.2 离散赋值环
% 纲要位置：工作记录/计划纲要.md，第 2237 行。
% 纲要细目（写作范围）：
% * 一维 Noether 正规局部整环的刻画
% * 极大理想的统一化元
% * 理想与阶函数：由商模长度恢复整数阶，并用尖点说明长度与外部赋值的相等不能单独刻画 DVR

\section{离散赋值环}
\label{b3:sec:1002}

离散赋值环最有用之处在于，所有非零理想都由同一个元素的幂生成。因此它既能按整数测量阶数，也能从局部环的有限性和整闭性识别。


\subsection{一维 Noether 正规局部整环的刻画}
\label{b3:subsec:1002:01}

\begin{theorem}\label{b3:thm:dvr-characterizations}
对非域整环 \(R\)，下列条件等价：
\begin{equivalences}
\item\label{b3:item:dvr-valuation} \(R\) 为离散赋值环；
\item\label{b3:item:dvr-principal-maximal} \(R\) 为 Noether 局部环，且其极大理想由一个元素生成；
\item\label{b3:item:dvr-normal-dimension-one} \(R\) 为一维 Noether 正规局部整环。
\end{equivalences}
它们也等价于 \(R\) 是非域的局部主理想整环。
\end{theorem}
\begin{proof}
先设 \(R=V_v\)、\(v(K^*)=\mathbb Z\)，选 \(\pi\) 满足 \(v(\pi)=1\)。
每个 \(0\ne a\in R\) 唯一写成 \(a=\pi^nu\)，其中 \(n\ge0\)、\(u\) 为单位。
非零理想中非零元素的赋值有最小值 \(n\)，于是该理想恰为 \((\pi^n)\)。所以 \(R\) 是主理想整环，且只有 \((0)\)、\((\pi)\) 两个素理想。由\cref{b3:prop:valuation-local-normal}，赋值环整闭，这同时得到\itemref{b3:item:dvr-principal-maximal}和\itemref{b3:item:dvr-normal-dimension-one}。

设\itemref{b3:item:dvr-principal-maximal}成立，\(\mathfrak m=(\pi)\)。要把非零元素写成统一化元的有限次幂乘以单位，须排除它被 \(\pi\) 的所有幂整除的情形，因此令 \(J=\bigcap_{n\ge0}\pi^nR\)。
若 \(x\in J\)，写 \(x=\pi y\)；由整环中可约去 \(\pi\)，从
\(x\in\pi^{n+1}R\) 得 \(y\in\pi^nR\)，故 \(y\in J\)。于是 \(J=\pi J\)。
Noether 性使理想 \(J\) 有限生成，且 \((\pi)\subseteq J(R)=\mathfrak m\)，所以 Nakayama 引理给出 \(J=0\)。对每个非零 \(a\in R\)，存在 \(N\) 使 \(a\notin\pi^NR\)，而 \(a\in\pi^0R\)。由于这些理想依次缩小，满足 \(a\in\pi^nR\) 的非负整数构成有限初段，故有最大值 \(n(a)\)。于是 \(a=\pi^{n(a)}u\)，且 \(u\notin\mathfrak m\)，否则 \(a\in\pi^{n(a)+1}R\)；所以 \(u\) 为单位。
两个这样的分解相乘，单位部分的乘积仍是单位，故 \(n(ab)=n(a)+n(b)\)。对 \(a,b\in R\setminus\{0\}\) 定义 \(v(a/b)=n(a)-n(b)\)。若 \(a/b=c/d\)，则 \(ad=bc\)，从而 \(n(a)+n(d)=n(b)+n(c)\)，这说明分式表示无关。乘法指数相加；对环中元素相加，提出较小的 \(\pi\) 次幂后剩余部分仍在 \(R\)，故阶数不小于两个输入阶数的最小值。对任意分式通分，便得到同样的赋值不等式。
其值群为 \(\mathbb Z\)，非负值元素恰属于 \(R\)，所以得到\itemref{b3:item:dvr-valuation}。

最后设\itemref{b3:item:dvr-normal-dimension-one}成立，记极大理想为 \(\mathfrak m\)，分式域为 \(K\)。为得到 \(\mathfrak m\) 的一个生成元，可以寻找 \(x\in K\setminus R\) 使 \(x\mathfrak m\subseteq R\)，再用整闭性证明这个理想必须含有单位。
取 \(0\ne a\in\mathfrak m\)。包含 \(a\) 的素理想非零，由维数一和局部性只能是 \(\mathfrak m\)，所以 \(\sqrt{aR}=\mathfrak m\)，故某个
\(\mathfrak m^n\subseteq aR\)。这里用到 \(\mathfrak m\) 有限生成：每个生成元的某个幂属于 \(aR\)，展开充分高次的乘积即可。
取最小的 \(n\ge1\)，再取 \(b\in\mathfrak m^{n-1}\setminus aR\)。这样的 \(b\) 存在：当 \(n=1\) 时，\(R\not\subseteq aR\)；当 \(n>1\) 时，这是 \(n\) 的最小性。
于是 \(x=b/a\notin R\)，而 \(b\mathfrak m\subseteq\mathfrak m^n\subseteq aR\)，正好给出 \(x\mathfrak m\subseteq R\)。
若 \(x\mathfrak m\subseteq\mathfrak m\)，乘以 \(x\) 就定义了有限生成 \(R\) 模 \(\mathfrak m\) 上的 \(R\) 线性自同态 \(u_x:y\mapsto xy\)。对它使用\cref{b3:thm:determinant-trick}，取该定理中的理想为 \(R\)，得到首一多项式 \(P(T)\in R[T]\) 使 \(P(u_x)=0\)。因此对每个 \(y\in\mathfrak m\) 都有 \(P(x)y=0\)。这里的忠实性来自整环中 \(\mathfrak m\ne0\)：取 \(0\ne y\in\mathfrak m\)，在 \(K\) 中约去 \(y\)，得到 \(P(x)=0\)，故 \(x\) 在 \(R\) 上整。
正规性将推出 \(x\in R\)，矛盾。
所以理想 \(x\mathfrak m\) 含有单位，从而 \(x\mathfrak m=R\)。选 \(\pi\in\mathfrak m\) 满足 \(x\pi=1\)。
每个 \(y\in\mathfrak m\) 都满足 \(xy\in R\)，于是 \(y=\pi(xy)\)，得到 \(\mathfrak m=(\pi)\)，即\itemref{b3:item:dvr-principal-maximal}。
局部主理想整环满足\itemref{b3:item:dvr-principal-maximal}，而第一段已证明\itemref{b3:item:dvr-valuation}使每个理想主生成，故最后一个等价条件也成立。
\end{proof}

Noether 性使证明中的 \(J\) 和 \(\mathfrak m\) 都有限生成：前者用于 Nakayama 引理，后者用于从 \(\sqrt{aR}=\mathfrak m\) 得到幂包含关系，并应用行列式技巧。若直接删除有限性，上一节的一维非离散赋值环就给出反例。

\subsection{极大理想的统一化元}
\label{b3:subsec:1002:02}

\begin{definition}\label{b3:def:uniformizer}
设 \((R,\mathfrak m)\) 为 DVR，\(K=\operatorname{Frac}(R)\)。若元素 \(\pi\in R\) 满足 \(\mathfrak m=(\pi)\)，则称其为
\term[tongyihuayuan]{统一化元}[uniformizer]。将相应赋值的赋值群与通常排序的 \(\mathbb Z\) 认同，并令 \(v(\pi)=1\)，所得满射 \(v:K^*\to\mathbb Z\) 称为 \(R\) 的\term[guiyihua-lisan-fuzhi]{归一化离散赋值}[normalized discrete valuation]。
\end{definition}
\symidx{\(\pi\)：离散赋值环的统一化元}

统一化元并不唯一：\(\pi'\) 是另一个统一化元，当且仅当 \(\pi'=u\pi\)、\(u\in R^*\)。
所以元素的展开 \(a=u\pi^n\) 中单位部分依赖选择，整数 \(n\) 则与选择无关。
对 \(\mathbb Z_{(p)}\) 可取 \(p\)，对 \(k[t]_{(t-c)}\) 可取 \(t-c\)。
归一化离散赋值确定了理想的阶数。环与形式幂级数环的关系，还涉及完备性、特征及剩余域中系数的提升。

\begin{proposition}\label{b3:prop:noether-valuation-dvr}
非域赋值环为 Noether 环，当且仅当它为 DVR。
\end{proposition}
\begin{proof}
DVR 已证明为主理想整环。反向，Noether 赋值环的极大理想有限生成，因而由\cref{b3:prop:valuation-ideal-comparison}主生成，再应用\cref{b3:thm:dvr-characterizations}。
\end{proof}

极大理想主生成本身也不能推出 Noether 性。设 \(k\) 为域，取 \(K=k(\!(u)\!)(\!(t)\!)\)，对非零
\(f=t^n\cdot a_n(u)+t^{n+1}\cdot a_{n+1}(u)+\cdots\) 定义
\(v(f)=(n,\operatorname{ord}_u a_n)\)，赋值群按字典序排序的 \(\mathbb Z^2\)。
首项乘法和加法直接验证赋值公理。最小正值是 \((0,1)\)，所以极大理想由 \(u\) 生成；但零元素及赋值第一坐标为正的元素组成非零素理想
\(\mathfrak p\subsetneq\mathfrak m\)。对任意 \(0\ne a\in\mathfrak p\)，都能取非负整数 \(r\) 使 \(v(t/u^r)=(1,-r)\prec v(a)\)，所以 \(\mathfrak p\) 中没有赋值最小的非零元素。
若 \(\mathfrak p\) 有限生成，\cref{b3:prop:valuation-ideal-comparison}会使它主生成，其生成元的赋值就是非零元素的最小赋值，矛盾。因此 \(\mathfrak p\) 不是有限生成理想，这个环不是 Noether 环，也不是 DVR。这正是上一节字典序 \(\mathbb Z^2\) 对应图中的中间素理想。

\subsection{理想与阶函数}
\label{b3:subsec:1002:03}

\begin{proposition}\label{b3:prop:dvr-order-ideals}
在离散赋值环 \((R,\mathfrak m)\) 中，每个非零理想唯一写成 \(\mathfrak m^n\)，其中 \(n\ge0\) 为整数。
对非负整数 \(a,b\)，有
\[
\mathfrak m^a\mathfrak m^b=\mathfrak m^{a+b},\qquad
\mathfrak m^a+\mathfrak m^b=\mathfrak m^{\min(a,b)},\qquad
\mathfrak m^a\cap\mathfrak m^b=\mathfrak m^{\max(a,b)}.
\]
并且 \(\ell_R(R/\mathfrak m^n)=n\)。若 \(v\) 为 \(R\) 的归一化离散赋值，则对每个 \(0\ne a\in R\) 有
\[
v(a)=\ell_R(R/aR).
\]
不过，这一长度等式对某些元素成立并不能反过来刻画 DVR。对任意域 \(k\) 及不定元 \(t\)，尖点局部环 \(C=k[t^2,t^3]_{(t^2,t^3)}\) 满足
\[
\ell_C(C/t^2C)=2=\operatorname{ord}_0(t^2),\qquad
\ell_C(C/t^3C)=3=\operatorname{ord}_0(t^3),
\]
但 \(t\) 在 \(C\) 上整而不属于 \(C\)，所以 \(C\) 不是 DVR。这里 \(\operatorname{ord}_0\) 是 \(k(t)\) 上的零点阶赋值。
\end{proposition}
\begin{proof}
取统一化元 \(\pi\)。理想形式已在刻画定理中证明；\(\pi^n\notin\mathfrak m^{n+1}\) 保证指数唯一。
乘积公式由生成元相乘得到。指数大小与理想包含方向相反：\(a\le b\) 时 \(\mathfrak m^a\supseteq\mathfrak m^b\)，所以和取较大的理想，对应较小的指数；交取较小的理想，对应较大的指数。
对每个 \(i\ge0\)，乘以 \(\pi^i\) 给出
\(R/\mathfrak m\simeq\mathfrak m^i/\mathfrak m^{i+1}\)：满射显然，核为 \(\mathfrak m\)，因为可在整环中约去 \(\pi^i\)。
因此 \(R/\mathfrak m^n\) 的幂滤过有恰好 \(n\) 个简单商，长度为 \(n\)。
对 \(0\ne a\in R\)，写 \(a=u\pi^n\)，其中 \(u\) 为单位。于是 \(aR=\mathfrak m^n\)，从而 \(v(a)=n=\ell_R(R/\mathfrak m^n)=\ell_R(R/aR)\)。

对尖点，记 \(A=k[t^2,t^3]\)。其中出现的幂指数恰为 \(0,2,3,4,\ldots\)，所以 \(A/t^2A\) 的 \(k\) 基为 \(1,t^3\)，\(A/t^3A\) 的 \(k\) 基为 \(1,t^2,t^4\)。这两个商环分别同构于 \(k[z]/(z^2)\) 和 \(k[z]/(z^3)\)，其中 \(z\) 分别对应 \(t^3\) 和 \(t^2\)。常数项非零的元素在两商环中已经可逆，因此局部化不改变它们。沿 \(z\) 的幂所得的链分别有二个和三个简单商，均为剩余域 \(k\)，故作为 \(C\) 模的长度分别为二和三。
另一方面，\(t=t^3/t^2\) 属于 \(\operatorname{Frac}(C)=k(t)\)，并满足首一方程 \(T^2-t^2=0\)，故在 \(C\) 上整。若 \(t\in C\)，可写 \(t=f/g\)，其中 \(f,g\in A\)、\(g(0)\ne0\)。但 \(tg=f\) 的左端一次项系数为 \(g(0)\ne0\)，右端没有一次项，矛盾。因此 \(C\) 不整闭，不能为 DVR。
\end{proof}

在 DVR 上，整数阶因此可以完全由商模的合成层数读取。在一维正规曲线中，这将给出零点重数的一个局部代数描述。

\ProseParagraphBreak
一般赋值环的理想已按包含关系排成全序；Noether 条件则使非域赋值环成为 DVR。下一节的 Dedekind 整环把一维正规条件放在整个环上，在各个非零素理想处得到 DVR，进而合成全局的理想分解。\cref{b3:tab:valuation-dvr-dedekind}并列这些条件与结构。

\begin{center}
\begin{minipage}{\textwidth}
% 此表随正文排放，避免浮到后续定理与证明之间；沿用标准表题及编号。
\makeatletter
\def\@captype{table}
\makeatother
\centering
\caption[赋值环、DVR 与 Dedekind 整环]{赋值环、DVR 与 Dedekind 整环。赋值环行中，\(v:K^*\to\Gamma\) 为域 \(K\) 上的满赋值，\(V_v\) 为其赋值环。Dedekind 整环行限定为非域，其定义及全局理想分解见\secref{b3:sec:1003}。}
\label{b3:tab:valuation-dvr-dedekind}
\begin{tabularx}{\textwidth}{@{}>{\raggedright\arraybackslash}p{0.18\textwidth}>{\raggedright\arraybackslash}p{0.23\textwidth}>{\raggedright\arraybackslash}p{0.14\textwidth}>{\raggedright\arraybackslash}X@{}}
\toprule
对象 & 局部性与有限性 & Krull 维数 & 整除或理想结构\\
\midrule
一般赋值环 \(V_v\) & 局部整闭整环，未必 Noether & \(\operatorname{rank}(v)\) & 理想按包含关系全序；每个有限生成理想主生成\\
秩一赋值环 & 局部整闭整环，未必 Noether & \(1\) & 赋值群未必同构于 \(\mathbb Z\)，极大理想未必主生成\\
DVR & Noether 局部整闭整环，非域 & \(1\) & 值群同构于 \(\mathbb Z\)，非零理想为极大理想的幂；等价于一维 Noether 正规局部整环\\
非域 Dedekind 整环 & Noether 整闭整环，未必局部 & \(1\) & 非零素理想处的局部环为 DVR；非零理想有唯一素理想分解\\
\bottomrule
\end{tabularx}
\end{minipage}
\end{center}

% !TeX root = ../../Book3.tex
% 纲要标题：### 10.3 Dedekind 整环
% 纲要位置：工作记录/计划纲要.md，第 2243 行。
% 纲要细目（写作范围）：
% * 定义及局部刻画
% * 非零理想的素理想分解
% * 分式理想与可逆性

\section{Dedekind 整环}
\label{b3:sec:1003}

一个整数的素因子分解，可以分别在每个素数处读取指数。Dedekind 理论保留了这种办法，不过应分解的对象是非零理想。每个非零素理想处的局部环是 DVR，先在那里读出整数阶，再把有限多个局部数据合成原理想。


\subsection{Dedekind 整环的定义与局部刻画}
\label{b3:subsec:1003:01}

\begin{definition}\label{b3:def:dedekind-domain}
设 \(R\) 为整环。若 \(R\) 为 Noether 环、在其分式域中整闭，且每个非零素理想都是极大理想，则称 \(R\) 为\term[Dedekind-zhenghuan]{Dedekind 整环}[Dedekind domain]。
本书允许域作为零维情形；非域的 Dedekind 整环恰为一维 Noether 正规整环。
\end{definition}

在非域整环中，\((0)\) 是素理想，每个极大理想都非零。若非零素理想都极大，素理想链只能是 \((0)\subsetneq\mathfrak m\)，所以维数为一；反过来，维数一也排除了 \((0)\subsetneq\mathfrak p\subsetneq\mathfrak m\) 这样的链。这说明定义中的素理想条件正是在限制维数。

\begin{theorem}\label{b3:thm:dedekind-local-characterization}
设 \(R\) 为非域 Noether 整环。则 \(R\) 是 Dedekind 整环，当且仅当每个极大理想
\(\mathfrak m\) 处的 \(R_{\mathfrak m}\) 都是 DVR。
\end{theorem}
\begin{proof}
若 \(R\) 为 Dedekind 整环，则 \(R_{\mathfrak m}\) 为一维 Noether 局部整环。
整闭性由\cref{b3:cor:integrally-closed-localization}传至局部化，所以可用\cref{b3:thm:dvr-characterizations}。

反向，任意非零素理想 \(\mathfrak p\) 包含于某个 \(\mathfrak m\)。其局部化是
DVR 中的非零素理想，只能等于 \(\mathfrak mR_{\mathfrak m}\)；收缩回去得
\(\mathfrak p=\mathfrak m\)。再证整闭性。第二章将\cref{b3:cor:submodule-local-detection}应用于 \(R\subseteq K=\operatorname{Frac}(R)\)，已经得到
\[
R=\bigcap_{\mathfrak m\text{ 极大}}R_{\mathfrak m}.
\]
若 \(x\in K\) 在 \(R\) 上整，
它在每个 \(R_{\mathfrak m}\) 上也整，因各 DVR 整闭而属于上述交集，故属于 \(R\)。
\end{proof}

各极大理想处的局部环都是 DVR，并不必然使整个整环成为 Noether 环。例如，令
\(K=\mathbb Q(\zeta_p:p\text{ 为素数})\)，其中 \(\zeta_p\) 是本原 \(p\) 次单位根。
\(K\) 中所有代数整数构成的整环，在每个极大理想处局部化都是 DVR，却不是 Noether 环。这个数论构造的出处与上述性质见所引文献\cite[Section 2, p. 280]{loper2006almostDedekind}。
在已知 Dedekind 的环中，局部化仍为 Dedekind 整环或域：Noether 性、整闭性均保持，而素链只会减少。

\subsection{非零理想的素理想分解}
\label{b3:subsec:1003:02}

对 Dedekind 整环 \(R\) 的非零素理想 \(\mathfrak p\)，将
\(R_{\mathfrak p}\) 的归一化离散赋值记为 \(v_{\mathfrak p}\)。
给定非零理想 \(I\)，定义 \(v_{\mathfrak p}(I)\ge0\) 为唯一满足
\(IR_{\mathfrak p}=(\mathfrak pR_{\mathfrak p})^{v_{\mathfrak p}(I)}\) 的整数。
\symidx{\(v_{\mathfrak p}(I)\)：理想在素理想处的阶数}

例如，在 \(\mathbb Z\) 中取 \(I=(60)\)。局部化到 \(\mathbb Z_{(2)}\) 后，\(15\) 成为单位，所以 \(I\mathbb Z_{(2)}=2^2\mathbb Z_{(2)}\)；在 \(\mathbb Z_{(3)}\) 与 \(\mathbb Z_{(5)}\) 中，\(20\) 与 \(12\) 分别成为单位，故相应理想由 \(3\) 与 \(5\) 生成。对其余素数 \(p\)，\(60\) 本身就是 \(\mathbb Z_{(p)}\) 的单位。因此
\[
v_{(2)}(I)=2,\qquad v_{(3)}(I)=v_{(5)}(I)=1,
\qquad v_{(p)}(I)=0\quad(p\ne2,3,5),
\]
这些局部指数恰好给出通常的理想分解 \((60)=(2)^2(3)(5)\)。下面的定理说明，在 Dedekind 整环中也能从局部指数恢复整个理想。

\begin{theorem}\label{b3:thm:dedekind-factorization}
设 \(R\) 为 Dedekind 整环，\(I\subseteq R\) 为非零理想；对非零素理想 \(\mathfrak p\)，令 \(v_{\mathfrak p}(I)\) 表示 \(IR_{\mathfrak p}\) 在 DVR \(R_{\mathfrak p}\) 中的幂指数。则有唯一分解
\[
I=\prod_{\mathfrak p\ne0}\mathfrak p^{v_{\mathfrak p}(I)},
\]
其中只有有限多个指数非零。单位理想对应空乘积。
\end{theorem}
\begin{proof}
若 \(I=R\)，空乘积就是所求分解；以下设 \(I\subsetneq R\)，此时 \(R\) 不是域。取 \(0\ne a\in I\)。若 \(v_{\mathfrak p}(I)>0\)，则 \(I\subseteq\mathfrak p\)，故 \(a\in\mathfrak p\)。所以，只需证明含 \(a\) 的素理想只有有限多个。它们全是非零极大理想，因而 \(R/(a)\) 是零维 Noether 环，由\cref{b3:thm:noether-dimension-zero}为 Artin 环，只有有限多个素理想。这就得到局部指数的有限支撑。

令右端有限乘积为 \(J\)。在任一非零极大理想 \(\mathfrak q\) 处局部化，所有
\(\mathfrak p\ne\mathfrak q\) 都含有一个不属于 \(\mathfrak q\) 的元素，因而其局部化为单位理想。
只有 \(\mathfrak q\) 因子保留，于是 \(JR_{\mathfrak q}=IR_{\mathfrak q}\)。
由\cref{b3:cor:submodule-local-detection}得到 \(I=J\)：拼出的乘积具有与 \(I\) 相同的全部局部指数，而这些局部理想已经足以检测全局相等。

任一有限素理想乘积在 \(\mathfrak p\) 处局部化时，只保留对应素理想的幂；DVR 中幂指数唯一，所以分解中每个指数都被原理想确定。
\end{proof}

\subsection{分式理想与可逆性}
\label{b3:subsec:1003:03}

\begin{definition}\label{b3:def:fractional-invertible-ideal}
设 \(R\) 为整环，\(K\) 为其分式域，\(I\subseteq K\) 为非零 \(R\) 子模。若存在 \(0\ne d\in R\) 使 \(dI\subseteq R\)，则称 \(I\) 为
\term[fenshi-lixiang]{分式理想}[fractional ideal]。
若 \(I,J\) 均为 \(R\) 的非零分式理想，则其乘积 \(IJ\) 由全部有限和 \(\sum_i a_ib_i\)、\(a_i\in I,b_i\in J\)，组成。
对非零分式理想 \(I\)，定义
\[
I^{-1}\coloneqq\{x\in K:xI\subseteq R\}.
\]
若 \(II^{-1}=R\)，则称 \(I\) 为\term[keni-lixiang]{可逆理想}[invertible ideal]。
\end{definition}
\symidx{\(I^{-1}\)：分式理想的逆理想}

定义中的 \(J=dI\) 是非零通常理想，并且 \(I=d^{-1}J\)，所以分式理想就是将一个通常理想整体除以同一个非零元素。\(I^{-1}\) 则由能把整个 \(I\) 乘回 \(R\) 的分式组成，并非把 \(I\) 的非零元素逐个取倒数所得的集合。例如，对 \(0\ne a\in K\)，有 \((aR)^{-1}=a^{-1}R\)。

在 Noether 整环上，分式理想都是有限生成模，因为乘以 \(d\) 将它同构到通常理想 \(J\)。
\(I^{-1}\) 也是分式理想：它包含非零的 \(dR\)；取 \(0\ne a\in I\)，则
\(aI^{-1}\subseteq R\)，再清除 \(a\) 自身的分母即可得到定义中所需的 \(R\) 元素。

全局上，一个分式理想可能需要若干生成元；在每个局部环中只需一个生成元，则可望构造它的局部逆，再由局部检测证明全局可逆。反过来，等式 \(II^{-1}=R\) 会给出有限个乘积之和为 \(1\)，由这一等式在各局部环中选出生成元。

\begin{proposition}\label{b3:prop:invertible-local-principal}
设 \(R\) 为整环，\(K\) 为其分式域。若分式理想 \(I\subseteq K\) 为有限生成的 \(R\) 模，则它可逆，当且仅当每个极大理想 \(\mathfrak m\subseteq R\) 处的
\(IR_{\mathfrak m}\) 都由一个非零分式生成。可逆分式理想必为有限生成投射模。
\end{proposition}
\begin{proof}
若 \(II^{-1}=R\)，取有限和 \(\sum_i a_ib_i=1\)，其中 \(a_i\in I\)、\(b_i\in I^{-1}\)。
每个 \(x\in I\) 都可写成 \(x=\sum_i (xb_i)a_i\)，说明 \(a_i\) 生成 \(I\)。
映射 \(R^r\rightarrow I\)、\((c_i)\mapsto\sum_i c_ia_i\) 有截面
\(x\mapsto(xb_i)_i\)，故 \(I\) 是有限自由模的直和因子，因而投射。
在 \(R_{\mathfrak m}\) 中，若所有 \(a_ib_i\) 都在极大理想 \(\mathfrak mR_{\mathfrak m}\) 中，它们的和也在其中，不可能等于 \(1\)。所以至少一个 \(a_ib_i\) 不在该极大理想中，因而是单位。
对任意 \(x\in I\)，有 \(x/a_i=(xb_i)/(a_ib_i)\in R_{\mathfrak m}\)，于是 \(I_{\mathfrak m}=a_iR_{\mathfrak m}\)。

反向需要说明逆理想与局部化相容。包含 \((I^{-1})_{\mathfrak m}\subseteq(I_{\mathfrak m})^{-1}\) 直接来自定义。为证明反向包含，设 \(I=(a_1,\ldots,a_r)\)、\(x\in K\)，且
\(xI_{\mathfrak m}\subseteq R_{\mathfrak m}\)，为每个 \(xa_i\) 选一个分母
\(s_i\notin\mathfrak m\) 使 \(s_ix a_i\in R\)。令 \(s=\prod_i s_i\)，便有
\(sx\in I^{-1}\)，因而 \(x=(sx)/s\in(I^{-1})_{\mathfrak m}\)。这里正是有限生成性使全部 \(xa_i\) 能用一个共同分母乘回 \(R\)。因此 \((I^{-1})_{\mathfrak m}=(I_{\mathfrak m})^{-1}\)。
若 \(I_{\mathfrak m}\) 主生成且非零，它与自己的逆相乘为 \(R_{\mathfrak m}\)，故
\((II^{-1})_{\mathfrak m}=R_{\mathfrak m}\)。若 \(II^{-1}\) 为真理想，将其包含在某个极大理想中会与这个等式矛盾。于是 \(II^{-1}=R\)。
\end{proof}

\begin{corollary}\label{b3:cor:dedekind-fractional-group}
设 \(R\) 为 Dedekind 整环。它的全部非零分式理想在乘法下组成交换群。每个分式理想 \(I\) 有唯一表达式
\[
I=\prod_{\mathfrak p\ne0}\mathfrak p^{n_{\mathfrak p}},
\qquad n_{\mathfrak p}\in\mathbb Z,
\]
其中只有有限多个指数非零。因此该群是以非零素理想为基的自由交换群。
\end{corollary}
\begin{proof}
若 \(R\) 为域，则唯一非零分式理想为 \(R\)，结论成立。以下设 \(R\) 非域。分式理想有限生成，每个极大理想非零，相应局部环为 DVR，局部化理想均主生成，故由上一命题可逆。
选 \(0\ne d\in R\) 使 \(dI\) 为通常理想，则 \(I=(dI)(dR)^{-1}\)。
分别分解这两个通常理想便得到整数指数；在每个 DVR 处读出指数又证明唯一性。乘法的结合律、交换律来自有限和与环乘法，单位元是 \(R\)，所定义的 \(I^{-1}\) 即群逆元。
\end{proof}

分式理想的局部指数仍记为 \(v_{\mathfrak p}(I)\)，但现在允许取负整数。
若 \(\pi\) 是 DVR \(R_{\mathfrak p}\) 的统一化元，则
\(IR_{\mathfrak p}=\pi^{v_{\mathfrak p}(I)}R_{\mathfrak p}\)；
负指数表示统一化元的逆元的正整数次幂。
指数的大小与理想的包含方向相反。例如，\(\mathbb Z\) 中 \((12)\subsetneq(4)\)，但在 \((2)\) 处两者的指数同为 \(2\)，在 \((3)\) 处的指数分别为 \(1\) 与 \(0\)，其他素数处同为 \(0\)。所以较小理想的每个局部指数都不小于较大理想的相应指数。
整数指数的相加、取负、取最小值与取最大值，分别对应\cref{b3:tab:dedekind-ideal-exponents}中的理想运算。

\begin{table}[htbp]
\centering
\caption[Dedekind 整环中的局部指数运算]{Dedekind 整环中的局部指数运算。\(R\) 为非域 Dedekind 整环，\(K\) 为其分式域，\(I,J\subseteq K\) 为非零分式理想。对非零素理想 \(\mathfrak p\subseteq R\)，\(v_{\mathfrak p}(I)\in\mathbb Z\) 是满足 \(IR_{\mathfrak p}=(\mathfrak pR_{\mathfrak p})^{v_{\mathfrak p}(I)}\) 的唯一整数。}
\label{b3:tab:dedekind-ideal-exponents}
\begin{tabularx}{\textwidth}{@{}p{0.27\textwidth}X@{}}
\toprule
理想运算或包含关系 & 局部指数\\
\midrule
乘积 \(IJ\) & \(v_{\mathfrak p}(IJ)=v_{\mathfrak p}(I)+v_{\mathfrak p}(J)\)\\
和 \(I+J\) & \(v_{\mathfrak p}(I+J)=\min\{v_{\mathfrak p}(I),v_{\mathfrak p}(J)\}\)\\
交 \(I\cap J\) & \(v_{\mathfrak p}(I\cap J)=\max\{v_{\mathfrak p}(I),v_{\mathfrak p}(J)\}\)\\
逆 \(I^{-1}\) & \(v_{\mathfrak p}(I^{-1})=-v_{\mathfrak p}(I)\)\\
\(I\subseteq J\) & 当且仅当对所有非零素理想 \(\mathfrak p\subseteq R\)，均有 \(v_{\mathfrak p}(I)\ge v_{\mathfrak p}(J)\)\\
\bottomrule
\end{tabularx}
\end{table}

这些运算公式来自 DVR 中的幂理想运算，因为局部化与乘积、和及有限交相容。
Dedekind 整环上的分式理想有限生成，故\cref{b3:prop:invertible-local-principal}的证明也给出了这里逆理想与局部化的相容性。
局部幂理想的包含方向与指数方向相反，再由\cref{b3:cor:submodule-local-detection}得到表中的全局包含判据。

可逆不等于主生成。主分式理想 \(xR\) 总有逆 \(x^{-1}R\)，但一般 Dedekind 整环中还存在非主可逆理想。下一节给出具体算例，并说明它们怎样组成理想类群。

% !TeX root = ../../Book3.tex
% 纲要标题：### 10.4 数域整数环、理想类群与曲线赋值
% 纲要位置：工作记录/计划纲要.md，第 2126 行。
% 纲要细目（写作范围）：
% * 数域整数环的 Dedekind 性
% * 理想类群及 Picard 群的初步关系
% * 正规曲线与离散赋值
% ---


\section{数域整数环、理想类群与曲线赋值}
\label{b3:sec:1004}

整数环和一元多项式环都具有唯一的元素分解，但其有限扩张未必仍有这个性质。Dedekind 理论说明，转向非零理想之后仍能保持唯一分解。曲线上的零点阶数则把同一结构解释为几何数据。


\subsection{数域整数环的 Dedekind 性}
\label{b3:subsec:1004:01}

\begin{definition}\label{b3:def:number-field-integers}
给定有限域扩张 \(K/\mathbb Q\)，称 \(K\) 为\term[shuyu]{数域}[number field]。在 \(K\) 中，把所有满足某个首一整数系数多项式的元素组成的环称为 \(K\) 的\term[zhengshuhuan]{整数环}[ring of integers]，记为 \(\mathcal O_K\)；它也就是 \(\mathbb Z\) 在 \(K\) 中的整闭包。
\end{definition}
\symidx{\(\mathcal O_K\)：数域的整数环}

\begin{theorem}\label{b3:thm:number-integers-dedekind}
设 \(K/\mathbb Q\) 为有限域扩张，\(\mathcal O_K\) 为 \(\mathbb Z\) 在 \(K\) 中的整闭包。则 \(\mathcal O_K\) 是有限自由 \(\mathbb Z\) 模，秩为 \([K:\mathbb Q]\)，并且是 Dedekind 整环。
\end{theorem}
\begin{proof}
特征零保证 \(K/\mathbb Q\) 可分。\(\mathbb Z\) 是整闭 Noether 整环，故\cref{b3:lem:normalization-separable-finite}给出 \(\mathcal O_K\) 的有限生成性。
为直接说明自由性，令 \(n=[K:\mathbb Q]\)，选取有理基 \(\beta_1,\ldots,\beta_n\) 及坐标同构
\[
\theta:K\longrightarrow\mathbb Q^n,\qquad
\sum_{i=1}^n q_i\beta_i\longmapsto(q_1,\ldots,q_n).
\]
为 \(\mathcal O_K\) 的有限个整数模生成元的坐标取正整数公共分母 \(D\)，便有
\(\theta(\mathcal O_K)\subseteq D^{-1}\mathbb Z^n\)。映射
\[
\mathcal O_K\longrightarrow H\coloneqq D\theta(\mathcal O_K)\subseteq\mathbb Z^n,
\qquad a\longmapsto D\theta(a)
\]
是 \(\mathbb Z\) 模同构。每个子群 \(H\subseteq\mathbb Z^n\) 都自由，可对 \(n\) 归纳：第一坐标像若为零，就在 \(\mathbb Z^{n-1}\) 中使用归纳；否则其像为 \(d\mathbb Z\)、\(d>0\)，取 \(h\in H\) 的第一坐标为 \(d\)。每个元素唯一写成 \(mh+h_0\)，其中 \(m\in\mathbb Z\)、\(h_0\) 第一坐标为零，故 \(H=\mathbb Zh\oplus H_0\)，再对 \(H_0\) 归纳。这证明 \(\mathcal O_K\) 为有限自由模。
任何 \(x\in K\) 乘一个非零整数后都对 \(\mathbb Z\) 整：对 \(x\) 的首一有理系数方程清除分母，正如\cref{b3:lem:normalization-separable-finite}证明中的缩放。因此
\(\mathbb Q\otimes_{\mathbb Z}\mathcal O_K=K\)，自由秩为 \([K:\mathbb Q]\)，分式域也为 \(K\)。

\(\mathcal O_K\) 作为 \(\mathbb Z\) 模有限生成，因而是 Noether 环；整闭包对整性传递而整闭。
由\cref{b3:thm:integral-dimension}，它与 \(\mathbb Z\) 同为一维。于是它是 Dedekind 整环。
\end{proof}

这个证明也说明正规化有限性在算术中的作用：它先给出 Noether 性，不能只凭“所有元素都整”就断言整数环 Noether。
对任意 \(0\ne a\in\mathcal O_K\)，乘以 \(a\) 是有限自由整数模上的单射，且在 \(\mathbb Q\) 上为同构，所以余核是有限交换群。
因此 \(\mathcal O_K/(a)\) 有限；非零素理想的剩余域也有限。

\subsection{理想类群与 Picard 群}
\label{b3:subsec:1004:02}

在 Dedekind 整环上，非零分式理想都可逆，并由\cref{b3:prop:invertible-local-principal}成为局部自由秩一的有限生成投射模。若只保留其模结构，就得到一种在各局部环上都像一条自由直线的对象。这样的模在任意交换环上仍可用张量积相乘；Picard 群把它们的同构类组织起来。

\begin{definition}\label{b3:def:ideal-class-picard}
设 \(R\) 为 Dedekind 整环，\(K=\operatorname{Frac}(R)\)。分式理想群对其主分式理想子群的商
\[
\operatorname{Cl}(R)\coloneqq
\{\text{非零分式理想}\}/\{xR:x\in K^*\}
\]
称为 \(R\) 的\term[lixiangleiqun]{理想类群}[ideal class group]。
另设 \(R\) 为任意交换含幺环。考虑有限生成投射 \(R\) 模 \(L\) 的同构类，要求对每个素理想 \(\mathfrak p\subseteq R\)，\(L_{\mathfrak p}\) 都是秩一自由 \(R_{\mathfrak p}\) 模。这些同构类以 \(R\) 上的张量积为乘法所成的群，称为 \(R\) 的\term[Picard-qun]{Picard 群}[Picard group]，记为 \(\operatorname{Pic}(R)\)。
\end{definition}
\symidx{\(\operatorname{Cl}(R)\)：Dedekind 整环的理想类群}
\symidx{\(\operatorname{Pic}(R)\)：秩一有限投射模的 Picard 群}

说明后一构造确实为群。两个有限投射模都是有限自由模的直和因子，张量后仍为有限投射模；局部化后两个秩一自由模的张量积仍秩一。
对偶 \(L^*=\operatorname{Hom}_R(L,R)\) 也为有限投射模：将分解 \(R^r=L\oplus L'\) 取对偶便得到相应的直和分解。有限投射模也有限展示，因为这项分解中 \(L'\) 是由 \(R^r\) 的有限基经投影得到的有限生成核。因此\cref{b3:thm:hom-localization}适用于 \(L\)，它使求值映射
\(L\otimes_RL^*\rightarrow R\) 在每个素理想处成为秩一自由模的求值同构。
由\cref{b3:thm:local-morphism-detection}它本身即为同构。因此 \([L^*]\) 是逆元，\([R]\) 是单位元。

\begin{proposition}\label{b3:prop:dedekind-picard-class}
设 \(R\) 为 Dedekind 整环，\(K=\operatorname{Frac}(R)\)。则有自然群同构
\(\operatorname{Cl}(R)\simeq\operatorname{Pic}(R)\)，由 \([I]\mapsto[I]\) 给出。
而 \(\operatorname{Cl}(R)=0\) 当且仅当 \(R\) 是主理想整环。
\end{proposition}
\begin{proof}
可逆理想由\cref{b3:prop:invertible-local-principal}为有限投射模，并且局部自由秩一。
乘法映射 \(I\otimes_RJ\rightarrow IJ\) 在每个极大理想 \(\mathfrak m\) 处，把两个非零主分式理想 \(I_{\mathfrak m}=aR_{\mathfrak m}\)、\(J_{\mathfrak m}=bR_{\mathfrak m}\) 的张量积的基 \(a\otimes b\) 送到 \((IJ)_{\mathfrak m}\) 的基 \(ab\)，故局部为同构，再由局部检测知它是同构。
主分式理想同构于 \(R\)，所以映射在类群上良定义且保持乘法。

若有 \(R\) 模同构 \(\varphi:I\xrightarrow{\sim}J\)，借助自然乘法同构 \(K\otimes_RI\to K\)、\(\lambda\otimes a\mapsto\lambda a\)，以及 \(J\) 的对应同构，将 \(1\otimes\varphi\) 识别为 \(K\) 的一个线性自同构 \(\varphi_K\)。它必为乘以某个 \(x\in K^*\)。对每个 \(a\in I\)，自然嵌入 \(I,J\subseteq K\) 与张量映射相容，因而在 \(K\) 中有
\[
\varphi(a)=\varphi_K(a)=xa.
\]
所以 \(J=\varphi(I)=xI\)，说明类群映射单射。
给定有限投射且局部秩一的 \(L\)，它是自由模的直和因子，因而无挠；所以
\(L\rightarrow K\otimes_RL\) 单射。后者一维，选择它与 \(K\) 的同构，将 \(L\) 识别为 \(K\) 内有限生成的非零子模。为有限个生成元清除分母后，它成为一个分式理想。因此映射满射。
最后，类群平凡恰好表示每个非零分式理想都主生成，特别是每个非零通常理想主生成；反向由清分母同样成立。
\end{proof}

取 \(K=\mathbb Q(\sqrt{-5})\)，记 \(s=\sqrt{-5}\)。此时 \(R=\mathcal O_K=\mathbb Z[s]\)。
为核实这一步，设 \(\alpha=a+bs\) 对 \(\mathbb Z\) 整，其中 \(a,b\in\mathbb Q\)。它的共轭 \(\bar\alpha=a-bs\) 满足同一个首一整数系数多项式，因此也整。迹 \(\alpha+\bar\alpha=2a\) 和范数 \(\alpha\bar\alpha=a^2+5b^2\) 作为整元素的和与积也对 \(\mathbb Z\) 整，又都属于 \(\mathbb Q\)，由 \(\mathbb Z\) 整闭可知它们是整数。于是
\(-20b^2=(2a)^2-4(a^2+5b^2)\) 为整数。将 \(b\) 写成既约分数，分母的平方整除 \(20\)，所以 \(2b\in\mathbb Z\)。
写 \(a=A/2,b=B/2\)，由 \(A^2+5B^2\equiv0\pmod4\) 得 \(A,B\) 都为偶数，故 \(a,b\in\mathbb Z\)。反向由 \(s\) 整显然成立。

\ProseParagraphBreak
在 \(R=\mathbb Z[s]\) 中，范数 \(N(a+bs)=a^2+5b^2\) 是非负整数并具有乘法性。单位的范数必为 \(1\)，由这个公式可知单位恰为 \(\pm1\)。环中没有范数为 \(2\) 或 \(3\) 的元素：若 \(b\ne0\)，范数至少为 \(5\)；若 \(b=0\)，范数就是整数的平方。而
\[
N(2)=4,\qquad N(3)=9,\qquad N(1+s)=N(1-s)=6.
\]
若这些元素之一能分成两个非单位，其范数就会分成两个大于 \(1\) 的整数。但 \(4,9,6\) 的每种非平凡乘积分解都含因子 \(2\) 或 \(3\)，故 \(2,3,1+s,1-s\) 都不可约。相伴元素有相同范数，而 \(1+s\ne\pm(1-s)\)，所以这四个元素两两不相伴。因此
\[
6=2\cdot3=(1+s)(1-s)
\]
是两种不能由重排与乘单位互相得到的不可约分解，\(R\) 不是唯一分解整环。

\ProseParagraphBreak
理想 \(\mathfrak p=(2,1+s)\) 的商环为 \(\mathbb F_2\)，所以是素理想。
有 \(\mathfrak p^2=(2)\)：平方的生成元均被 \(2\) 整除，而
\(2(1+s)-(1+s)^2=6\)，连同 \(4\in\mathfrak p^2\) 得 \(2\in\mathfrak p^2\)。
它不是主理想。若 \(\mathfrak p=(\alpha)\)，则由 \(\mathfrak p^2=(2)\) 得 \(\alpha^2=2u\)，其中 \(u\in R^*=\{\pm1\}\)。取范数得 \(N(\alpha)^2=4\)，所以 \(N(\alpha)=2\)，与上述范数的取值矛盾。
因此 \([\mathfrak p]\) 是类群中阶二的非平凡元素。这里已经看到元素分解与理想分解的区别，不需要先计算整个类群。

\subsection{正规曲线与离散赋值}
\label{b3:subsec:1004:03}

设 \(k\) 为域，\(A\) 为有限型整 \(k\) 代数。若 \(\dim A=1\)，本节便称 \(A\) 所对应的仿射对象为一条不可约仿射曲线。这里用 \(\operatorname{Spec}A\) 表示该仿射对象，以便包含非代数闭基域的闭点。由\cref{b3:thm:affine-dimension}，其有理函数域在 \(k\) 上的超越次数为一。
若 \(A\) 整闭，则称曲线正规。此时 \(A\) 为 Dedekind 整环，每个闭点
\(\mathfrak p\) 的局部环是 DVR。

\ProseParagraphBreak
对 \(0\ne f\in\operatorname{Frac}(A)\)，整数 \(v_{\mathfrak p}(f)\) 表示在该点的阶数：正值是零点阶数，负值是极点阶数，零值表示在该局部环中为单位。
主分式理想的分解说明只有有限多个非零阶，并给出
\[
fA=\prod_{\mathfrak p\ne0}\mathfrak p^{v_{\mathfrak p}(f)}.
\]
这个分解的指数由 \(\operatorname{Spec}A\) 中的闭点给出。

\ProseParagraphBreak
在 \(A=k[t]\) 上，闭点由首一不可约多项式 \(q(t)\) 表示，统一化元为 \(q\)，阶数即 \(q\) 在分子和分母中的重数之差。
特别地，函数 \(t\) 在闭点 \((t)\) 处的阶数为 \(1\)，在其余闭点 \((q)\)（\(q\ne t\)）处的阶数都为 \(0\)。在同一函数域中令 \(u=t^{-1}\)，倒数坐标图 \(k[u]\) 中点 \((u)\) 的局部环是 \(k[u]_{(u)}\)，统一化元为 \(u\)。于是
\[
v_\infty(t)=v_{(u)}(u^{-1})=-1.
\]
这一点是在 \(t\) 坐标图之外补入的无穷远点。仿射环 \(k[t]\) 的理想分解记录了 \(t\) 在原点的一阶零点；倒数坐标则显示它在无穷远的一阶极点。

\ProseParagraphBreak
在 \(k[t,t^{-1}]\) 上，\((t)\) 已因局部化而消失，\(t\) 是单位，在留下的所有闭点的阶数都为零。
这解释了删去一点会怎样改变可逆函数，而不改变函数域。

\ProseParagraphBreak
尖点 \(k[t^2,t^3]\) 的正规化为 \(k[t]\)，见\cref{b3:prop:cusp-normalization}。
原点上方正规局部环的统一化元是 \(t\)，原来两个坐标的阶数分别为二和三。
但原来的局部环不含 \(t\)，它的极大理想需要两个生成元，因此自身不是 DVR。
正规化使零点阶成为一个整数赋值；这不等于把原局部环已经具有的生成关系抹去。


\begin{chaptersummary}
赋值环的整除关系是全序的，这使有限生成理想成为主理想；但值群可以稠密，也可以具有多个凸子群。因此整闭、局部、一维或极大理想主生成等零散条件，都不能单独代替 DVR 的完整刻画。Noether 性使行列式技巧和 Nakayama 引理能够把一维正规局部结构转化为统一化元的整数幂。

\ProseParagraphBreak
Dedekind 整环的非零理想在各个 DVR 处留下有限多个整数指数，局部检测再恢复唯一的全局分解。分式理想允许负指数，于是组成自由交换群；模掉主分式理想得到类群。类群也可以用秩一有限投射模来描述。在章首的 \(\mathbb Z[\sqrt{-5}]\) 中，元素的非唯一不可约分解并未破坏非零理想的唯一素理想分解。理想 \(\mathfrak p=(2,1+\sqrt{-5})\) 满足 \(\mathfrak p^2=(2)\)，却不是主理想；它在类群中给出阶二的非平凡元素。理想类群由此记录了可逆理想不能由单个元素生成的障碍。
\end{chaptersummary}

\BookExercises

\begin{exercise}\label{b3:ex:10-lowest-valuation}
设 \(K\) 为域，\(v:K^*\twoheadrightarrow\Gamma\) 为取值于全序交换群 \(\Gamma\) 的赋值，且 \(a_1,\ldots,a_r\in K^*\) 满足 \(a_1+\cdots+a_r=0\)。证明最小的
\(v(a_i)\) 至少出现两次。给出两个同阶元素之和阶数严格增大的例子。
\end{exercise}
\begin{solution}
和为零而各项非零，所以 \(r\ge2\)。若最小值只在 \(a_1\) 处出现，其余有限项的最小值便严格大于 \(v(a_1)\)，于是
\[
v(a_2+\cdots+a_r)\ge\min_{i\ge2}v(a_i)>v(a_1).
\]
但左侧为 \(v(-a_1)=v(a_1)\)，矛盾。
在 \(k(t)\) 的零点赋值中，\(1\) 与 \(-1+t\) 都为零阶，其和 \(t\) 为一阶；这在任意特征都成立。
\end{solution}

\begin{exercise}\label{b3:ex:10-rational-orders}
对每个素数 \(p\)，令 \(v_p\) 为 \(\mathbb Q\) 上满足 \(v_p(p)=1\) 的 \(p\)-进赋值。计算 \(v_2(84/75)\)、\(v_3(84/75)\)、\(v_5(84/75)\)，判断该有理数属于哪些 \(\mathbb Z_{(p)}\)，并指出在哪些局部环中为单位。
\end{exercise}
\begin{solution}
\[
\dfrac{84}{75}=\dfrac{2^2\cdot3\cdot7}{3\cdot5^2}
=\dfrac{28}{25}=2^2\cdot7\cdot5^{-2}.
\]
因此
\[
v_2(84/75)=2,\qquad v_3(84/75)=0,\qquad v_5(84/75)=-2.
\]
它恰在 \(p\ne5\) 时属于 \(\mathbb Z_{(p)}\)；在 \(p=2,7\) 处是非单位，在其余 \(p\ne5\) 的素数处都是单位。特别地，它在 \(\mathbb Z_{(3)}\) 中是单位。
\end{solution}

\begin{exercise}\label{b3:ex:10-irrational-weight}
设 \(k\) 为域，取正无理数 \(\alpha\)。对 \(0\ne f=\sum c_{ij}x^iy^j\in k[x,y]\)，令
\(v(f)=\min\{i+\alpha j:c_{ij}\ne0\}\)。证明它延拓为 \(k(x,y)\) 的赋值，并求其值群的秩和有理秩。这一构造见\cite[Exercise 10.11, p. 78]{Matsumura1989CommutativeRingTheory}，这里补入两种秩的比较。
\end{exercise}
\begin{solution}
无理性使不同单项式具有不同权重，所以每个非零多项式有唯一最低权项。两个最低权项的乘积是乘积多项式唯一最低权项，系数非零，故
\(v(fg)=v(f)+v(g)\)。相加只能消去已有项，因而满足赋值不等式。
定义 \(v(f/g)=v(f)-v(g)\)，乘法等式保证与分式表示无关。值群为
\(\mathbb Z+\alpha\mathbb Z\)，作为实数有序子群，其非零凸子群必为全群，秩一。
\(1,\alpha\) 在 \(\mathbb Q\) 上线性独立，所以有理秩二。
\end{solution}

\begin{exercise}\label{b3:ex:10-rational-value-maximal}
设 \(K\) 为域，\(v:K^*\twoheadrightarrow\mathbb Q\) 为赋值，\(\mathbb Q\) 取通常的次序，\(V=\{0\}\cup\{a\in K^*:v(a)\ge0\}\)，\(\mathfrak m=\{0\}\cup\{a\in K^*:v(a)>0\}\)。证明 \(V\) 的极大理想 \(\mathfrak m\) 满足
\(\mathfrak m^2=\mathfrak m\)，并说明为什么 Nakayama 引理不能使它为零。
\end{exercise}
\begin{solution}
取 \(0\ne a\in\mathfrak m\)，设 \(v(a)=q>0\)。由赋值满射可取 \(b\) 使
\(v(b)=q/2\)；于是 \(b,a/b\in\mathfrak m\)，故 \(a\in\mathfrak m^2\)。反向包含总成立。
\(\mathfrak m\) 非零却不有限生成；它没有最小正赋值元素，有限生成会迫使主生成。
这正是 Nakayama 引理的有限性条件失效之处。
\end{solution}

\begin{exercise}\label{b3:ex:10-principal-comparison}
证明整环 \(V\) 为赋值环，当且仅当任意两个主理想可按包含比较。再说明“每个有限生成理想都主生成”单独是否足够。
\end{exercise}
\begin{solution}
赋值环的主理想可比较已在正文证明。反向，对 \(x=a/b\in\operatorname{Frac}(V)^*\)，
若 \(aV\subseteq bV\)，则 \(x\in V\)；若 \(bV\subseteq aV\)，则 \(x^{-1}\in V\)。
仅要求有限生成理想主生成不够，例如 \(\mathbb Z\) 是主理想整环，但 \((2)\) 和 \((3)\) 不可比较，\(2/3\) 及其倒数均不在 \(\mathbb Z\)。
\end{solution}

\begin{exercise}\label{b3:ex:10-overrings}
设 \(V\) 为 \(K\) 的赋值环，且 \(V\subseteq W\subseteq K\) 为子环。证明存在
\(\mathfrak p\in\operatorname{Spec}V\) 使 \(W=V_{\mathfrak p}\)。
\end{exercise}
\begin{solution}
对 \(x\in K^*\)，\(x\) 或 \(x^{-1}\) 已在 \(V\)，当然也在 \(W\)，故 \(W\) 为赋值环。
设 \(\mathfrak n\) 为其极大理想，\(\mathfrak p=\mathfrak n\cap V\)。
\(V\setminus\mathfrak p\) 中的元在 \(W\) 可逆，所以 \(V_{\mathfrak p}\subseteq W\)。
对 \(x\in W\setminus V\)，倒数条件给 \(x^{-1}\in V\)，它在 \(W\) 中的逆为 \(x\)，故
\(x^{-1}\notin\mathfrak p\)，于是 \(x\in V_{\mathfrak p}\)。这证明反向包含。
\end{solution}

\begin{exercise}\label{b3:ex:10-lexicographic-primes}
设 \(k\) 为域，\(K=k(\!(u)\!)(\!(t)\!)\)。对非零 Laurent 级数
\(f=t^n a_n(u)+t^{n+1}a_{n+1}(u)+\cdots\)，其中 \(a_n(u)\ne0\)，令
\(v(f)=(n,\operatorname{ord}_u a_n)\in\mathbb Z^2\)，值群取字典序。列出赋值环 \(V=\{0\}\cup\{f\in K^*:v(f)\ge(0,0)\}\) 的全部素理想，并证明它的极大理想主生成而环不 Noether。
\end{exercise}
\begin{solution}
凸子群恰为 \(0\)、\(\{0\}\times\mathbb Z\) 和 \(\mathbb Z^2\)。若凸子群含某个第一坐标非零的元素，换号后可令第一坐标为正，其正整数倍超过任意给定元素，凸性使子群为全群；其余凸子群归结为 \(\mathbb Z\) 的凸子群。
对应素理想为 \(\mathfrak m\)、\(\mathfrak p=\{0\}\cup\bigl\{f:v(f)_1>0\bigr\}\) 和 \((0)\)。
最小正值是 \((0,1)=v(u)\)，所以 \(\mathfrak m=(u)\)。若 \(\mathfrak p\) 有限生成，它会主生成并有最小赋值，但 \((1,-r)=v(t/u^r)\) 随 \(r\) 无界下降，且第一坐标一是 \(\mathfrak p\) 中最小可能第一坐标，故不存在最小赋值。环不 Noether。
\end{solution}

\begin{exercise}\label{b3:ex:10-cusp-fails-dvr}
设 \(k\) 为域，\(R=k[t^3,t^4,t^5]_{(t^3,t^4,t^5)}\)，\(\mathfrak m\) 为其极大理想。对每个整数 \(n\ge3\)，求 \(R/t^nR\) 的一组 \(k\) 基及其 \(R\) 模长度；再求 \(\dim_k\mathfrak m/\mathfrak m^2\)，说明为什么这些长度仍等于 \(t\)-阶，却不能使 \(R\) 成为 DVR。
\end{exercise}
\begin{solution}
环 \(B=k[t^3,t^4,t^5]\) 中的幂指数集合为 \(H=\{0\}\cup\{3,4,5,\ldots\}\)。理想 \(t^nB\) 的指数集合为 \(n+H\)，所以商环的一组 \(k\) 基是
\[
1,\ t^3,t^4,\ldots,t^{n-1},\ t^{n+1},t^{n+2},
\]
其中 \(n=3\) 时中间一段为空。共计 \(n\) 个基元素。商环中所有正次幂都幂零，故常数项非零的元素已经可逆，局部化不改变商环。其唯一简单模为 \(k\)，长度因此等于 \(k\) 维数，即 \(n=\operatorname{ord}_0(t^n)\)。
\(\mathfrak m\) 由 \(t^3,t^4,t^5\) 生成，而 \(\mathfrak m^2\) 中的元素的 \(t\)-阶至少为六。三者在 \(\mathfrak m/\mathfrak m^2\) 中线性无关：若某个 \(k\) 线性组合属于 \(\mathfrak m^2\)，清除常数项非零的分母后比较三至五次项便得全部系数为零。因此该空间维数为三，极大理想不能由一个元素生成，\(R\) 不是 DVR。这里对所有 \(n\ge3\) 都成立的长度公式，仍不足以替代 DVR 的环结构条件。
\end{solution}

\begin{exercise}\label{b3:ex:10-dvr-colon}
设 \((R,\mathfrak m)\) 为 DVR，\(\pi\) 为统一化元，\(a,b,n,r\) 为非负整数。求 \((\mathfrak m^a:\mathfrak m^b)\)，以及 \(R/\mathfrak m^n\) 中被 \(\pi^r\) 零化的子模的长度。
\end{exercise}
\begin{solution}
条件 \(x\mathfrak m^b\subseteq\mathfrak m^a\) 等价于 \(v(x)+b\ge a\)，且 \(x\in R\)，
所以理想商为 \(\mathfrak m^{\max(a-b,0)}\)。第二个子模为
\(\mathfrak m^{\max(n-r,0)}/\mathfrak m^n\)，其长度为 \(\min(r,n)\)，这里 \(r\ge0\)。
\end{solution}

\begin{exercise}\label{b3:ex:10-dvr-presentation}
设 \(R\) 为 DVR，\(\pi\) 为统一化元。计算矩阵
\(\left(\begin{smallmatrix}\pi^2&\pi\\0&\pi^3\end{smallmatrix}\right)\)
所定义的 \(R^2\rightarrow R^2\) 的余核分解与长度。
\end{exercise}
\begin{solution}
交换两列得到 \(\left(\begin{smallmatrix}\pi&\pi^2\\\pi^3&0\end{smallmatrix}\right)\)。
第二行减去第一行的 \(\pi^2\) 倍，再令第二列减去第一列的 \(\pi\) 倍，得到
\(\operatorname{diag}(\pi,-\pi^4)\)。这些均为可逆基变换，故余核为
\(R/(\pi)\oplus R/(\pi^4)\)，长度五。行列式的阶也为五，在这里与长度计算一致。
\end{solution}

\begin{exercise}\label{b3:ex:10-ramification-basic}
设 \(R\subseteq S\) 为 DVR 间的局部单射，\(\pi,\tau\) 分别为统一化元，\(v_R,v_S\) 分别为满足 \(v_R(\pi)=v_S(\tau)=1\) 的归一化离散赋值。
证明存在整数 \(e\ge1\) 使 \(v_S(x)=e\,v_R(x)\) 对全部 \(x\in\operatorname{Frac}(R)^*\) 成立。
计算 \(k[t]_{(t)}\subseteq k[u]_{(u)}\)、\(t=u^e\)，的这个整数。
\end{exercise}
\begin{solution}
局部同态使 \(\pi\) 落入 \(S\) 的极大理想，又因单射而非零，故
\(v_S(\pi)=e\ge1\)。\(R\) 中的单位在 \(S\) 仍是单位，因为其逆也被映入。
写 \(x=\pi^n a\)、\(a\in R^*\)、\(n\in\mathbb Z\)，便得公式。
例子中 \(v_u(t)=v_u(u^e)=e\)，恰为所求整数；结论不要求特征不整除 \(e\)。
\end{solution}

\begin{exercise}\label{b3:ex:10-fractional-integers}
在 \(\mathbb Z\) 中取分式理想 \(I=(6/35)\)、\(J=(14/15)\)。求它们的和、交、积和各自的逆。
\end{exercise}
\begin{solution}
统一分母为 \(105\)，分子理想分别为 \((18)\) 和 \((98)\)。最大公因数为二，最小公倍数为 \(882\)，故
\[
I+J=(2/105),\qquad I\cap J=(42/5),\qquad IJ=(4/25).
\]
主分式理想的逆由生成元取倒数给出，分别为 \((35/6)\) 和 \((15/14)\)。
\end{solution}

\begin{exercise}\label{b3:ex:10-hom-fractional}
若 \(R\) 为 Dedekind 整环，\(I,J\) 为非零分式理想，证明
\(\operatorname{Hom}_R(I,J)\simeq JI^{-1}\)，右边的元素通过乘法定义映射。
\end{exercise}
\begin{solution}
令 \(K=\operatorname{Frac}(R)\)。通过 \(c\otimes a\mapsto ca\)，将 \(K\otimes_R I\) 和 \(K\otimes_R J\) 都识别为 \(K\)。任一 \(R\) 线性映射 \(\varphi:I\to J\) 张量到 \(K\) 后就是乘以某个 \(x\in K\)。对 \(a\in I\)，自然嵌入 \(J\hookrightarrow K\) 将 \(\varphi(a)\) 送到 \(x a\)，所以原映射本身也是乘以 \(x\) 的限制。
它取 \(I\) 到 \(J\) 当且仅当 \(xI\subseteq J\)。乘以 \(I^{-1}\) 得这等价于
\(xR\subseteq JI^{-1}\)，即 \(x\in JI^{-1}\)。
反向包含也直接由 \(I^{-1}I=R\) 得到。零映射对应 \(x=0\)，所以这是全部 Hom 模的同构。
\end{solution}

\begin{exercise}\label{b3:ex:10-dedekind-quotient-length}
设 \(R\) 为 Dedekind 整环。若 \(I=\prod_{i=1}^r\mathfrak p_i^{n_i}\ne0\)，其中 \(n_i>0\)、各 \(\mathfrak p_i\) 为不同的非零素理想，证明
\[
R/I\simeq\prod_{i=1}^rR/\mathfrak p_i^{n_i},
\qquad \ell_R(R/I)=\sum_{i=1}^r n_i.
\]
\end{exercise}
\begin{solution}
不同极大理想的幂互素：若 \(a+b=1\)、\(a\in\mathfrak p_i,b\in\mathfrak p_j\)，展开
\((a+b)^{n_i+n_j-1}\) 即见每项属于两幂之一。中国剩余定理给出直积同构。
\(R/\mathfrak p^n\) 中每个 \(s\notin\mathfrak p\) 都可逆，因为 \(R/\mathfrak p^n\) 只有一个极大理想。
因此该商环同构于 \(R_{\mathfrak p}/\mathfrak p^nR_{\mathfrak p}\)，它的 DVR 幂滤过有 \(n\) 个简单因子，每个作为 \(R\) 模也同构于 \(R/\mathfrak p\)。长度为 \(n\)，对有限直积相加得到结论。
\end{solution}

\begin{exercise}\label{b3:ex:10-semilocal-dedekind}
证明只有有限多个极大理想的 Dedekind 整环是主理想整环。此题见\cite[Exercise 11.7, p. 86]{Matsumura1989CommutativeRingTheory}。
\end{exercise}
\begin{solution}
域的情形显然。设极大理想为 \(\mathfrak p_1,\ldots,\mathfrak p_r\)，取非零理想
\(I=\prod_i\mathfrak p_i^{n_i}\)。每个
\(\mathfrak p_i^{n_i}/\mathfrak p_i^{n_i+1}\ne0\)，可选
\(a_i\in\mathfrak p_i^{n_i}\setminus\mathfrak p_i^{n_i+1}\)。
中国剩余定理给 \(a\in R\) 使 \(a\equiv a_i\pmod{\mathfrak p_i^{n_i+1}}\) 同时成立，故
\(v_{\mathfrak p_i}(a)=n_i\)。已列出所有非零素理想，所以唯一分解给 \(aR=I\)。
\end{solution}

\begin{exercise}\label{b3:ex:10-localization-class-group}
设 \(R\) 为 Dedekind 整环，\(S\) 为不含零的乘法集。证明
\(\operatorname{Cl}(R)\rightarrow\operatorname{Cl}(S^{-1}R)\)、\([I]\mapsto[S^{-1}I]\)，满射，且核由所有满足 \(\mathfrak p\cap S\ne\varnothing\) 的素理想类生成。
\end{exercise}
\begin{solution}
局部化后，遇到 \(S\) 的素理想成为单位理想，其余非零素理想与局部化环的非零素理想一一对应。因此分式理想群之间的映射删去前一类基元素，是满射，并将主理想映为主理想。
若 \(S^{-1}I=xS^{-1}R\)，则分式理想 \(x^{-1}I\) 在所有避开 \(S\) 的素理想处的指数为零，因而由遇到 \(S\) 的素理想及其逆组成。
在类群中 \([x^{-1}I]=[I]\)，所以核恰由这些类生成。
\end{solution}

\begin{exercise}\label{b3:ex:10-every-ideal-invertible}
设 \(R\) 为整环，且每个非零通常理想都可逆。证明 \(R\) 为 Dedekind 整环。
\end{exercise}
\begin{solution}
可逆理想有限生成，见\cref{b3:prop:invertible-local-principal}证明中的有限单位分解，因此全部理想有限生成，\(R\) 为 Noether 环。
若 \(R\) 非域，每个极大理想 \(\mathfrak m\ne0\) 可逆，在 \(\mathfrak m\) 处局部主生成。
于是 \(R_{\mathfrak m}\) 是非域 Noether 局部整环，极大理想主生成，故为 DVR。
由\cref{b3:thm:dedekind-local-characterization}，\(R\) 是 Dedekind 整环。域也符合定义。
\end{solution}

\begin{exercise}\label{b3:ex:10-quadratic-ideal-three}
在 \(R=\mathbb Z[\sqrt{-5}]\) 中，令 \(s=\sqrt{-5}\)、
\(\mathfrak q=(3,1+s)\)、\(\mathfrak q'=(3,1-s)\)。证明它们为不同素理想，并计算
\(\mathfrak q\mathfrak q'\) 及 \((1+s)\) 的素理想分解。
\end{exercise}
\begin{solution}
两个商环均为 \(\mathbb F_3\)，但 \(s\) 的像分别为 \(-1\) 与 \(1\)，故是不同素理想。
乘积由 \(9,3(1+s),3(1-s),6\) 生成，所有生成元属于 \((3)\)，而 \(9-6=3\) 属于乘积，故乘积为 \((3)\)。
记 \(\mathfrak p=(2,1+s)\)。因 \(\mathfrak p\) 和 \(\mathfrak q\) 都含 \(1+s\)，有
\((1+s)\subseteq\mathfrak p\cap\mathfrak q=\mathfrak p\mathfrak q\)：不同极大理想满足 \(\mathfrak p+\mathfrak q=R\)，互素性才给出交等于积。
在商环 \(R/(1+s)\) 中有 \(s=-1\)，原来的关系 \(s^2+5=0\) 变成 \(6=0\)，故
\(R/(1+s)\cong\mathbb Z/6\mathbb Z\)。中国剩余定理给
\(|R/\mathfrak p\mathfrak q|=2\cdot3=6\)，故两个有限指数子群相等，即
\((1+s)=\mathfrak p\mathfrak q\)。
\end{solution}

\begin{exercise}\label{b3:ex:10-invertible-nonfree}
设 \(R=\mathbb Z[\sqrt{-5}]\)，\(\mathfrak p=(2,1+\sqrt{-5})\subseteq R\)。证明 \(\mathfrak p\) 是秩一有限生成投射 \(R\) 模，但不是自由模；再求
\(\mathfrak p\otimes_R\mathfrak p\) 的同构类型。
\end{exercise}
\begin{solution}
Dedekind 整环中非零理想可逆，故有限投射且局部秩一。若自由，其与分式域张量后维数一，故只能自由秩一，从而主生成，和正文中的范数计算矛盾。
乘法给出 \(\mathfrak p\otimes_R\mathfrak p\simeq\mathfrak p^2=(2)\simeq R\)。因此非自由秩一投射模的张量平方可以自由。
\end{solution}

\begin{exercise}\label{b3:ex:10-number-prime-residue}
设 \(K/\mathbb Q\) 为次数 \(n\) 的数域，\(p\) 为素数。证明
\(|\mathcal O_K/p\mathcal O_K|=p^n\)。若
\(p\mathcal O_K=\prod_i\mathfrak p_i^{e_i}\) 为非零素理想分解，其中各 \(\mathfrak p_i\) 互异、\(e_i\ge1\)，且
\(|\mathcal O_K/\mathfrak p_i|=p^{f_i}\)，证明 \(\sum_i e_if_i=n\)。
\end{exercise}
\begin{solution}
整数环是秩 \(n\) 的自由整数模，模掉 \(p\) 后为 \(n\) 维 \(\mathbb F_p\) 空间，元素数为 \(p^n\)。
中国剩余定理分解商环。每个 \(\mathcal O_K/\mathfrak p_i^{e_i}\) 的幂滤过有
\(e_i\) 个同构于剩余域的因子，因此其元素数为 \(p^{e_if_i}\)。相乘并比较 \(p\) 的指数，得到公式。
\end{solution}

\begin{exercise}\label{b3:ex:10-affine-orders}
设 \(k\) 为域。求 \(f=(t-1)^3/\bigl(t^2(t+1)\bigr)\) 在仿射直线各闭点的阶数，分别讨论
\(\operatorname{char}k=2\) 与不等于二的情形。再将环换成 \(k[t,t^{-1}]\)。
\end{exercise}
\begin{solution}
特征不为二时，\(0,1,-1\) 是互异的 \(k\) 点，阶数分别为 \(-2,3,-1\)，其余闭点为零。
特征二时 \(t+1=t-1\)，约去后 \(f=(t-1)^2/t^2\)，在 \(0,1\) 的阶数分别为 \(-2,2\)，其余为零。
换成 Laurent 多项式环后，点 \((t)\) 被删除；其他点和阶数保持，只需删去原点处的条目。这里统计的是各仿射图所含的闭点。对原仿射直线，若补入射影直线的无穷远点，还须计入 \(v_\infty(f)\)；本题分子与分母同为三次，在任意特征下该阶数均为零。
\end{solution}

\begin{exercise}\label{b3:ex:10-normal-affine-membership}
设 \(k\) 为域，\(A\) 为一维整闭有限型整 \(k\) 代数，\(f\in\operatorname{Frac}(A)^*\)。对每个闭点 \(\mathfrak p\in\operatorname{Spec}A\)，令 \(v_{\mathfrak p}\) 为离散赋值环 \(A_{\mathfrak p}\) 的归一化离散赋值。证明 \(f\in A\) 当且仅当所有闭点处 \(v_{\mathfrak p}(f)\ge0\)；证明 \(f\in A^*\) 当且仅当各处阶数都为零。
\end{exercise}
\begin{solution}
由\cref{b3:cor:submodule-local-detection}在分式域中对 \(A\) 检测成员，\(A\) 等于全部极大理想局部环之交。
每个局部环是 DVR，属于该环恰等价于赋值非负，故得到第一结论。
单位条件要求 \(f\) 和 \(f^{-1}\) 同时属于 \(A\)，等价于各赋值同时非负和非正，即全部为零。
\end{solution}

% !TeX root = ../../Book3.tex
% 纲要标题：## 第11章　分次环、Hilbert 函数与多项式
% 纲要位置：工作记录/计划纲要.md，第 2257 行。
% 纲要细目（写作范围）：

\chapter{分次环、Hilbert 函数与多项式}
\label{b3:ch:11}
\BookChapterNumber{b3:ch:11}

\begin{chapterabstract}
本章研究分次模各次数分量的增长，证明 Hilbert 级数的有理性，并在标准分次时由最终多项式读取支集维数。局部环与有限生成模则以理想幂的逐层商得到 Hilbert–Samuel 多项式；参数理想、曲线奇点及切锥算例说明增长次数和首项系数各自记录什么。
\end{chapterabstract}

\begin{learninggoals}
\begin{itemize}
\item 区分正分次与标准分次，正确处理齐次关系及次数平移。
\item 从分次正合列计算 Hilbert 级数，判断何时存在最终的单个多项式，并在已知分次满射时用级数核验展示。
\item 说明 Hilbert 多项式、累计增长、维数与重数之间的关系。
\item 构造伴随分次环，并用参数理想和切锥计算局部长度。
\item 用标准单项式计数求 Hilbert 数据，并计算有限点集求值码的维数。
\end{itemize}
\end{learninggoals}

设 \(k\) 为域。在 \(k[x,y]\) 中，次数为 \(n\) 的单项式有 \(n+1\) 个，所以分次维数列为 \(1,2,3,\ldots\)。若加上关系 \(x^2=0\)，次数 \(n\ge1\) 时只剩 \(y^n\) 与 \(xy^{n-1}\) 两类，维数稳定为 \(2\)。一条方程改变了增长速度，却没有使所有高次部分消失。把这些数汇成生成函数，可以更清楚地看见分母中的维数信息。

\ProseParagraphBreak

局部环往往没有预先给定的次数，但极大理想的幂提供了逐层逼近：比较 \(M/\mathfrak m^{n+1}M\) 随 \(n\) 增长的长度。这里既要解释为何最终出现多项式，又要分清增长次数和首项系数分别记录什么。与分次环的计算并列，才能理解局部重数与切锥所保留的信息。

\ProseParagraphBreak

预备知识为\chapref{b3:ch:03}的有限性和长度、\chapref{b3:ch:05}的 Noether 正规化，以及\chapref{b3:ch:09}的维数和参数系。局部增长证明需要的诱导滤过估计在\cref{b3:lem:induced-filtration-bound}中由有限分次生成元建立；\chapref{b3:ch:13}将进一步证明 Artin–Rees 等式，并研究完备化。

\ProseParagraphBreak
Hilbert 级数与局部增长可对照松村英之的《Commutative Ring Theory》\cite[\S13, Theorems 13.2--13.4, pp. 94--100]{Matsumura1989CommutativeRingTheory}。本章以 \(h(n)\) 记单层长度、\(\chi(n)\) 记累计长度；两者的多项式次数相差一，零维时还须单独处理。

% !TeX root = ../../Book3.tex
% 纲要标题：### 11.1 分次环与分次模
% 纲要位置：工作记录/计划纲要.md，第 2237 行。
% 纲要细目（写作范围）：
% * 正分次与标准分次
% * 齐次理想和分次展示
% * 次数平移及有限生成

\section{分次环与分次模}
\label{b3:sec:1101}

多项式的全次数把一个无限维空间分成有限维的片。若方程齐次，取商以后仍可逐次计数；若方程不齐次，则通常只能先保留一个次数滤过。分次理论先研究前一种情形。本章的环都交换，分次环从次数零开始，分次模允许整数次数，但有限生成模的非零分次片总有一个共同的下界。

\subsection{正分次与标准分次}
\label{b3:subsec:1101:01}

\begin{definition}\label{b3:def:positive-standard-grading}
设 \(A\) 为交换环。若其加法群具有满足以下条件的直和分解
\[
A=\bigoplus_{n\ge0}A_n,\qquad 1\in A_0,\qquad A_iA_j\subseteq A_{i+j},
\]
则称此分解为 \(A\) 的一个\term[zheng-fenci]{正分次}[positive grading]。属于单个 \(A_n\) 的元素称为次数 \(n\) 的\term[qici-yuansu]{齐次元素}[homogeneous element]。若 \(A=A_0[A_1]\)，且 \(A_1\) 是有限生成 \(A_0\) 模，则称此分次为\term[biaozhun-fenci]{标准分次}[standard grading]。给定 \(A\) 模 \(M\)，若它具有直和分解 \(M=\bigoplus_{n\in\mathbb Z}M_n\)，满足 \(A_iM_j\subseteq M_{i+j}\)，则称 \(M\) 为分次 \(A\) 模。分次同态在本章均指保持次数的模同态。
\end{definition}

正次数部分 \(A_+=\bigoplus_{n>0}A_n\) 是理想，且 \(A/A_+\cong A_0\)。多项式环的通常全次数给出标准分次；若给变量赋正整数权 \(d_i\)，令 \(\deg x_i=d_i\)，也得到正分次，但未必标准。例如 \(k[x]\) 取 \(\deg x=2\)，奇数次分量全为零。

\ProseParagraphBreak
直和意味着每个元素只有有限个非零齐次分量。形式幂级数通常有无限多个分量，因此不能把 \(k[\![x]\!]\) 当作这里的 \(\bigoplus_{n\ge0}kx^n\)。它更适合用\chapref{b3:ch:13}的滤过与完备化描述。

\begin{proposition}\label{b3:prop:graded-noetherian}
设 \(A=\bigoplus_{n\ge0}A_n\) 为交换正分次环。若 \(A_0\) 是 Noether 环，\(A\) 由有限个正次数齐次元素作为 \(A_0\) 代数生成，则 \(A\) 为 Noether 环，每个 \(A_n\) 是有限生成的 \(A_0\) 模。反之，若正分次环 \(A\) 为 Noether 环，则上述两个条件成立。
\end{proposition}
\begin{proof}
正向由\cref{b3:cor:hilbert-module-consequences}得到 Noether 性。若生成元次数为正数 \(d_1,\ldots,d_r\)，给定 \(n\) 时，满足 \(\sum d_i a_i=n\) 的非负整数组只有有限多组，对应单项式生成 \(A_n\)。

反向，\(A_0\cong A/A_+\) 是 Noether 环。对理想 \(A_+\) 取有限组生成元，再把各生成元分成齐次部分，可得到有限组正次数齐次理想生成元 \(x_i\)。若 \(a\in A_n\)、\(n>0\)，在等式 \(a=\sum f_i x_i\) 中只取次数 \(n\) 的部分，得到 \(a=\sum (f_i)_{n-\deg x_i}x_i\)。各系数的次数低于 \(n\)，对 \(n\) 归纳可知 \(a\in A_0[x_1,\ldots,x_r]\)。
\end{proof}

\subsection{齐次理想与分次展示}
\label{b3:subsec:1101:02}

设 \(A\) 为交换正分次环，\(M\) 为分次 \(A\) 模。若子模 \(N\subseteq M\) 包含其中每个元素的所有齐次分量，则称 \(N\) 为\term[qici-zimo]{齐次子模}[graded submodule]。等价地，\(N=\bigoplus_n(N\cap M_n)\)，或 \(N\) 可以由齐次元素生成。后一等价的反向使用了乘积次数相加：把系数也分次后，任意齐次分量仍是原齐次生成元的组合。取 \(M=A\) 就得到齐次理想。

\ProseParagraphBreak
若 \(N\) 齐次，商模以 \((M/N)_n=M_n/(N\cap M_n)\) 分次；分次同态的核与像都是齐次子模，正合性也可以在每个次数上检查。例如 \((x^2,xy)\subset k[x,y]\) 齐次，而 \((x-1)\) 不是齐次理想：它含 \(x-1\)，却不含其常数部分 \(-1\)。

\begin{proposition}\label{b3:prop:graded-presentation}
设 \(A\) 为交换 Noether 正分次环，\(M\) 为有限生成分次 \(A\) 模。则 \(M\) 有有限组齐次生成元，也有由齐次关系组成的有限展示。
\end{proposition}
\begin{proof}
把任意有限组生成元的全部齐次分量收集起来，所得仍为有限组，且生成 \(M\)。设它们为 \(m_i\)，次数为 \(a_i\)。给自由模第 \(i\) 个基向量赋次数 \(a_i\)，便有保持次数的自由满射 \(F\rightarrow M\)。其核 \(K\) 是齐次子模，且由 Noether 性有限生成。再把 \(K\) 的有限组生成元拆为齐次部分，就得到有限组齐次关系及相应的分次展示。
\end{proof}

若两个生成元次数不同，同一个齐次关系中系数的次数也不同。关系的齐次性指“系数次数加上相应生成元次数”相同，而不是要求一列矩阵的所有系数都有相同次数。

\subsection{次数平移与有限生成}
\label{b3:subsec:1101:03}

\begin{definition}\label{b3:def:grading-shift}
设 \(A\) 为交换正分次环，\(M\) 为分次 \(A\) 模，\(a\in\mathbb Z\)。次数平移 \(M(a)\) 定义为
\[
M(a)_n=M_{n+a},
\]
底层模与作用不变。特别地，\(A(-a)\) 的元素 \(1\) 位于次数 \(a\)。
\symidx{\(M(a)\)：分次模的次数平移}
\end{definition}

于是\cref{b3:prop:graded-presentation}的展示可写成
\[
\bigoplus_j A(-b_j)\longrightarrow\bigoplus_i A(-a_i)
\longrightarrow M\longrightarrow0.
\]
相应关系矩阵第 \((i,j)\) 项须齐次且次数为 \(b_j-a_i\)；若此差为负，该项只能为零。这一约定避免在后面级数公式中把平移指数写反。

\ProseParagraphBreak
设 \(A\) 由有限个正次数元素生成，\(M\) 的齐次生成元为 \(m_i\)、\(\deg m_i=a_i\)。逐次取分量有
\[
M_n=\sum_i A_{n-a_i}m_i.
\]
因此当 \(n<\min_i a_i\) 时 \(M_n=0\)，而每个 \(M_n\) 都是有限生成的 \(A_0\) 模。标准分次时，还可选 \(c\) 使 \(M_{n+1}=A_1M_n\) 对 \(n\ge c\) 成立：只须令 \(c\ge\max_i a_i\)，再用 \(A_{j+1}=A_1A_j\)。这就是有限生成在逐次增长中的具体约束。

% !TeX root = ../../Book3.tex
% 纲要标题：### 11.2 Hilbert 函数
% 纲要位置：工作记录/计划纲要.md，第 2265 行。
% 纲要细目（写作范围）：
% * 分次片的维数或长度
% * Hilbert 级数的有理性
% * 从级数读取增长
% * 用初始单项式理想和包含排除实际计算 Hilbert 数据
% * 按原多项式次数计算同时变号不变环的权二分次 Hilbert 级数，用乘以齐次非零因子的正合列与单项式计数互相核验
% * 证明各次有限维的分次代数之间若有保持次数的满射且 Hilbert 级数相同，则该满射为同构；给出无映射时结论失效的反例
% #### 11.2.1 分次片的维数与长度
% #### 11.2.2 Hilbert 级数的有理性
% #### 11.2.3 Hilbert 级数与增长
% #### 11.2.4 从初始理想计算 Hilbert 数据
% #### 11.2.5 不变环的分次与展示核验

\section{Hilbert 函数}
\label{b3:sec:1102}

对每一个次数分别计数，比只说模“无限维”提供了更多信息。例如 \(k[x]\) 的每个分次片只有一维，而 \(k[x,y]\) 的第 \(n\) 片有 \(n+1\) 维。下面把这些数同时装进一个形式级数；级数的运算只使用系数，不涉及分析中的收敛。

\subsection{分次片的维数与长度}
\label{b3:subsec:1102:01}

本节设 \(A\) 为交换正分次环，\(A_0\) 为 Artin 环，\(A\) 由有限个齐次元素 \(x_1,\ldots,x_r\) 作为 \(A_0\) 代数生成，\(d_i=\deg x_i>0\)，\(M\) 为有限生成分次 \(A\) 模。记号 \(A=A_0[x_1,\ldots,x_r]\) 表示这些元素生成了 \(A\)，允许它们满足关系。给不定元 \(X_i\) 赋次数 \(d_i\)，就有分次满同态
\[
A_0[X_1,\ldots,X_r]\twoheadrightarrow A,
\qquad X_i\longmapsto x_i.
\]
由\cref{b3:prop:graded-noetherian}，\(A\) 是 Noether 环；各 \(M_n\) 是有限生成的 \(A_0\) 模，故长度有限。

\begin{definition}\label{b3:def:hilbert-function-series}
设 \(A=\bigoplus_{n\ge0}A_n\) 为交换正分次环，\(A_0\) 为 Artin 环，\(A=A_0[x_1,\ldots,x_r]\)，其中 \(x_i\) 为正次数齐次代数生成元，不要求代数独立。设 \(M\) 为有限生成分次 \(A\) 模。\(M\) 的\term[hilbert-hanshu]{Hilbert 函数}[Hilbert function]及\term[hilbert-jishu]{Hilbert 级数}[Hilbert series]分别为
\[
h_M(n)=\ell_{A_0}(M_n),\qquad
H_M(t)=\sum_{n\in\mathbb Z}h_M(n)\cdot t^n\in\mathbb Z(\!(t)\!).
\]
这里只有有限多个负次幂。若 \(A_0=k\) 为域，长度就是 \(k\) 维数。
\symidx{\(h_M(n)\)、\(H_M(t)\)：分次 Hilbert 函数与级数}
\end{definition}

短正合列 \(0\rightarrow L\rightarrow M\rightarrow N\rightarrow0\) 若保持次数，则 \(h_M(n)=h_L(n)+h_N(n)\)，从而 \(H_M=H_L+H_N\)。长度可加是因为将子模的合成列与商模合成列的逆像连接，就得到中间项的合成列。次数平移则给出
\begin{equation}\label{b3:eq:hilbert-shift}
H_{M(a)}(t)=\sum_n h_M(n+a)t^n=t^{-a}\cdot H_M(t).
\end{equation}

例如 \(k[x_1,\ldots,x_r]\) 取标准分次时，次数 \(n\) 的单项式数为 \(\binom{n+r-1}{r-1}\)、\(r\ge1\)，所以
\[
H_{k[x_1,\ldots,x_r]}(t)=\dfrac1{(1-t)^r}.
\]
这也可直接把每个变量的几何级数相乘得到。若底环为任意 Artin 环 \(A_0\)，同一自由单项式基使分子乘上 \(\ell_{A_0}(A_0)\)。

\subsection{Hilbert 级数的有理性}
\label{b3:subsec:1102:02}

下面的\term[hilbert-serre-dingli]{Hilbert–Serre 定理}[Hilbert--Serre theorem]把有限生成条件转化为级数分母的有限形式。

\begin{theorem}[Hilbert–Serre]\label{b3:thm:hilbert-serre}
设 \(A=\bigoplus_{n\ge0}A_n\) 为交换正分次环，\(A_0\) 为 Artin 环，\(A=A_0[x_1,\ldots,x_r]\)，其中 \(x_i\) 为次数 \(d_i>0\) 的齐次代数生成元，不要求代数独立。设 \(M\) 为有限生成分次 \(A\) 模。则存在整系数 Laurent 多项式 \(Q(t)\)，使
\[
H_M(t)=\dfrac{Q(t)}{\prod_{i=1}^r(1-t^{d_i})}.
\]
若 \(M_n=0\) 对 \(n<0\) 成立，则可取 \(Q(t)\in\mathbb Z[t]\)。
\end{theorem}
\begin{proof}
对代数生成元个数 \(r\) 归纳。\(r=0\) 时，\(M\) 是有限生成的 \(A_0\) 模，有限组齐次生成元的次数之外分量均为零，所以级数为 Laurent 多项式。

为了对生成元数归纳，要让最后一个生成元 \(x_r\) 的作用变成零。商模 \(L=M/x_rM\) 记录模掉它以后剩下的部分；但乘以 \(x_r\) 未必单射，还须用 \(K=(0:_Mx_r)\) 记录被这次乘法杀掉的部分。令 \(d=d_r\)，乘 \(x_r\) 给出分次正合列
\[
0\longrightarrow K(-d)\longrightarrow M(-d)
\xrightarrow{\ x_r\ }M\longrightarrow L\longrightarrow0.
\]
这里 \(m\in M(-d)_n=M_{n-d}\) 时，\(x_rm\in M_n\)，所以乘法保持次数，源中的核正是 \(K(-d)\)。由 \(A\) 的 Noether 性，\(K,L\) 都有限生成；它们被 \(x_r\) 零化，因而是商代数 \(A/(x_r)\) 上的有限生成分次模。这个商代数由前 \(r-1\) 个生成元的像作为 \(A_0\) 代数生成。逐片长度可加与\eqref{b3:eq:hilbert-shift}给出
\begin{equation}\label{b3:eq:hilbert-induction}
(1-t^d)H_M(t)=H_L(t)-t^d\cdot H_K(t).
\end{equation}
对右边应用归纳假设即得结论。若 \(M\) 没有负次分量，则它的级数及所乘的分母都没有负次幂，因此所得 Laurent 多项式实际上没有负次项。
\end{proof}

这里先证明 \(K\) 有限生成再作归纳，是 Noether 性不可省略的一步。此证明可对照 \cite[\S13, Theorem 13.2, pp. 94--95]{Matsumura1989CommutativeRingTheory}。

\subsection{Hilbert 级数与增长}
\label{b3:subsec:1102:03}

标准分次时，分母全由 \(1-t\) 构成。把公共因子约去，可写
\[
H_M(t)=\dfrac{Q(t)}{(1-t)^d},\qquad Q(1)\ne0
\]
或 \(d=0\)、\(H_M\) 本身为 Laurent 多项式。对于非零模，将约分后分母中 \(1-t\) 的指数 \(d\) 称为在 \(t=1\) 处的极点阶数，没有极点时取 \(d=0\)。这是对应有理函数的代数性质，只须检查公因子的约去，不涉及解析邻域或极限。下一节将证明它恰等于模支集的维数。

\ProseParagraphBreak
例如对 \(A=k[x,y]\)、\(M=A/(x^2)\)，乘 \(x^2\) 在 \(A\) 上单射，故
\[
H_M(t)=\dfrac{1-t^2}{(1-t)^2}=\dfrac{1+t}{1-t}.
\]
所以 \(h_M(0)=1\)，而 \(h_M(n)=2\) 对 \(n\ge1\) 成立。若改取 \(A/(x^2,xy)\)，标准单项式为 \(1,x,y,y^2,\ldots\)，因而
\[
H_{A/(x^2,xy)}(t)=\dfrac1{1-t}+t.
\]
额外的嵌入部分只改变有限多个低次值，却不改变最终的常数一。

\ProseParagraphBreak
在 Hilbert–Serre 定理的假设下，一般正分次已经保证级数有理；标准分次进一步把分母统一为 \(1-t\) 的幂，由此可从其系数得到下一节的最终多项式结论。一般正权分次没有这一保证：上节 \(k[x]\)、\(\deg x=2\) 的级数为 \(1/(1-t^2)\)。若其函数最终由一个多项式表示，无穷多个奇数处的零值会迫使这个多项式恒为零，却与偶数处的值为一矛盾。每个同余类可分别呈现多项式行为，但这不是标准分次结论中的同一个多项式。

\subsection{从初始理想计算 Hilbert 数据}
\label{b3:subsec:1102:04}

Hilbert 级数的有理性还可以变成实际的计数方法。齐次理想的 Gröbner 基把每个分次片的基列成标准单项式，计算由此只涉及哪些单项式被排除。

\begin{proposition}\label{b3:prop:hilbert-initial-count}
设 \(k\) 为域，\(S=k[x_1,\ldots,x_n]\) 取标准分次，\(I\subseteq S\) 为齐次理想，并固定全局单项式序。则 \(S/I\) 与 \(S/\operatorname{in}(I)\) 有相同的 Hilbert 函数。若 \(\operatorname{in}(I)=(m_1,\ldots,m_s)\)，则
\[
H_{S/I}(t)=\dfrac{\displaystyle\sum_{J\subseteq\{1,\ldots,s\}}
(-1)^{|J|}t^{\deg\operatorname{lcm}(m_j:j\in J)}}{(1-t)^n},
\]
空集的最小公倍式约定为一。
\end{proposition}
\begin{proof}
从有限齐次生成组开始作 Buchberger 算法，可得到由齐次多项式组成的 Gröbner 基：S 多项式与每一步非零余式仍齐次。对齐次多项式作除法时，每个乘积保持原次数。因此次数为 \(d\) 的标准单项式的剩余类构成 \((S/I)_d\) 的基；\(I\ne S\) 时，独立性由\cref{b3:prop:standard-basis-finite}的首项论证得到，该论证不要求商空间有限维。\(I=S\) 时，两边的商环均为零，标准单项式集合为空，结论也成立。

计数时，若只有两个单项式生成元 \(m_1,m_2\)，被排除的单项式是“被 \(m_1\) 整除”与“被 \(m_2\) 整除”这两个集合的并。两者的交由最小公倍式的倍数构成，所以被排除部分的级数是
\[
\dfrac{t^{\deg m_1}+t^{\deg m_2}-t^{\deg\operatorname{lcm}(m_1,m_2)}}{(1-t)^n}.
\]
一般地，对每个 \(j\)，被 \(m_j\) 整除的单项式集合的计数级数为 \(t^{\deg m_j}/(1-t)^n\)，若同时被若干生成元整除，则等价于被它们的最小公倍式整除。对每个次数中的有限集合应用包含排除，求出这些集合的并，再从全部单项式中扣除这个并，就得到标准单项式的计数。将各次等式组成形式级数，便得到所示公式，其中空集项贡献全部单项式的级数。
\end{proof}

例如在 \(S=k[x,y,z]\) 中取 \(I=(x^2,xy+y^2)\)，按字典序 \(x\succ y\succ z\) 计算。恒等式
\[
y^3=(y-x)(xy+y^2)+yx^2
\]
表明须补入 \(y^3\)。三对 S 多项式中，第一对依次用 \(xy+y^2\)、\(y^3\) 约化为零，另外两对分别为零或 \(y^4\)，也有零余式。因此 \(x^2,xy+y^2,y^3\) 构成 Gröbner 基，其初始理想为 \((x^2,xy,y^3)\)。标准单项式恰为
\[
z^j,\quad xz^j,\quad yz^j,\quad y^2z^j\qquad(j\ge0).
\]
所以
\[
H_{S/I}(t)=\dfrac{1+2t+t^2}{1-t},\qquad
h_{S/I}(d)=\begin{cases}1,&d=0,\\3,&d=1,\\4,&d\ge2.\end{cases}
\]
下一节的维数与重数定理将说明：分母的一阶极点给出维数一，分子在 \(t=1\) 处的值 \(1+2+1=4\) 给出重数四。这里先通过逐项计数得到函数，随后才解释这些数的几何意义。

\subsection{不变环的分次与展示核验}
\label{b3:subsec:1102:invariant-hilbert}

在\cref{b3:prop:sign-invariant-presentation}中，生成性和关系理想都已由单项式证明。现在保留原多项式的全次数，用 Hilbert 级数重新观察这个环。次数的选择不可略去：\(x^2,xy,y^2\) 的次数都是二，所以对应的三变量多项式环应取权二分次。

\begin{example}\label{b3:exm:sign-invariant-hilbert}
设 \(k\) 为特征不为二的域，\(B=k[x,y]^{C_2}\)，其中非单位元把 \(x,y\) 同时变号。给 \(k[x,y]\) 取 \(\deg x=\deg y=1\)，并令 \(B_n=B\cap k[x,y]_n\)。则
\begin{equation}\label{b3:eq:sign-invariant-hilbert}
H_B(t)=\dfrac{1-t^4}{(1-t^2)^3}
=\dfrac{1+t^2}{(1-t^2)^2},
\qquad
\dim_k B_n=
\begin{cases}
n+1,&n\ge0\;\text{为偶数},\\
0,&n\ge0\;\text{为奇数}.
\end{cases}
\end{equation}
\end{example}
\begin{proof}
该作用保持次数，不变多项式的每个齐次部分仍不变，所以所定义的 \(B_n\) 给出分次。由\cref{b3:prop:sign-invariant-presentation}，令 \(S=k[u,v,w]\)、\(\deg u=\deg v=\deg w=2\)，便有分次同构
\[
B\cong S/(q),\qquad q=uw-v^2,\qquad \deg q=4.
\]
\(S\) 是整环且 \(q\ne0\)，因此乘以 \(q\) 在 \(S\) 上单射，得到保持次数的短正合列
\[
0\longrightarrow S(-4)\xrightarrow{\ q\ }S
\longrightarrow B\longrightarrow0.
\]
每个变量的幂贡献级数 \(1/(1-t^2)\)，故 \(H_S(t)=(1-t^2)^{-3}\)。由逐次维数可加和\eqref{b3:eq:hilbert-shift}，便得
\[
H_B(t)=H_S(t)-t^4\cdot H_S(t)=\dfrac{1-t^4}{(1-t^2)^3}.
\]
这里使用的是 \(q\) 在取商前的 \(S\) 上为非零因子；不能仅从“有一条齐次关系”便省略单射性检查。

另一方面，偶数次 \(2m\) 的全部 \(2m+1\) 个单项式都固定，奇数次没有非零不变元。于是也可直接计算
\[
H_B(t)=\sum_{m\ge0}(2m+1)t^{2m}
=\dfrac{2t^2}{(1-t^2)^2}+\dfrac1{1-t^2}
=\dfrac{1+t^2}{(1-t^2)^2}.
\]
中间等式可逐项乘以 \((1-t^2)^2\) 核验，完全是形式级数恒等式。这同时给出所述 Hilbert 函数。
\end{proof}

由于本例所有奇次分量均为零，另定义 \(\widetilde B_m=B_{2m}\) 不会遗漏任何元素。把原次数二改记为次数一后，三个代数生成元都位于一次分量，得到标准分次环及级数
\[
H_{\widetilde B}(z)=\dfrac{1+z}{(1-z)^2},
\qquad H_B(t)=H_{\widetilde B}(t^2).
\]
底层环、乘法和 Krull 维数不变，Hilbert 函数却按新的次数下标计数。原分次的奇次空片使函数不能最终成为一个多项式，新分次则有 \(h_{\widetilde B}(m)=2m+1\)，可应用标准分次的 Hilbert 多项式与重数结论。直接把 \(u,v,w\) 视为次数一而仍声称计算的是 \(B_n\)，会把所有非零分次片移到错误的位置。若一般分次环有非零奇次分量，只取偶次分量会得到一个子环，不能当作同一底层环的这次重编号。

\ProseParagraphBreak
在上述算例中，第二次计数只是核对已经证明的同构。实际计算时，它也能用来证明尚待核验的展示，但必须先有一个保持次数的满射。

\begin{proposition}\label{b3:prop:hilbert-surjection-isomorphism}
设 \(k\) 为域，\(R=\bigoplus_{n\ge0}R_n\)、\(A=\bigoplus_{n\ge0}A_n\) 为各分次片有限维的分次 \(k\) 代数，\(\psi:R\to A\) 为保持次数的 \(k\) 代数满同态。若 \(H_R(t)=H_A(t)\)，则 \(\psi\) 是同构。
\end{proposition}
\begin{proof}
设 \(K=\ker\psi\)。因为 \(\psi\) 保持次数，\(K\) 是齐次理想。任取 \(a\in A_n\) 并取其原像 \(r=\sum_jr_j\)，比较 \(\psi(r)\) 的各次分量可知 \(\psi(r_n)=a\)，所以每个 \(\psi_n:R_n\to A_n\) 都满射。逐次取核得
\[
0\longrightarrow K_n\longrightarrow R_n
\xrightarrow{\ \psi_n\ }A_n\longrightarrow0.
\]
级数相同意味着 \(\dim_kR_n=\dim_kA_n\)，故 \(K_n=0\) 对每个 \(n\) 成立。因 \(K=\bigoplus_nK_n\)，核为零。
\end{proof}

例如在同时变号作用中，若已证明 \(x^2,xy,y^2\) 生成不变环，并验证 \(uw-v^2\) 是关系，就得到 \(S/(uw-v^2)\to B\) 的分次满射。取商前乘以 \(uw-v^2\) 的正合列计算其源的级数，独立的偶数次单项式计数计算其目标的级数；两者相等后，\cref{b3:prop:hilbert-surjection-isomorphism}便给出核恰为这一条关系的另一证明。Hilbert 级数只记录各次分量的维数，没有记录它们之间的乘法；若尚未证明候选元生成目标，或未说明映射保持所用分次，单独写出两个相同的有理函数还不能得到这个结论。

% !TeX root = ../../Book3.tex
% 纲要标题：### 11.3 Hilbert 多项式
% 纲要位置：工作记录/计划纲要.md，第 2285 行。
% 纲要细目（写作范围）：
% * 多项式性的证明
% * 多项式次数与维数的关系
% * 重数和首项系数的归一化

\section{Hilbert 多项式}
\label{b3:sec:1103}

标准分次时，Hilbert 级数的分母只含 \(1-t\)，提取系数便能得到最终的多项式公式。还需解释的是：这个公式的次数为何能测量模支集的维数，而不只反映某种级数写法。证明中选出的齐次参数同时提供可计数的单项式和维数的坐标，两种数量由此联系起来。这里始终要求标准分次，底环 \(A_0\) 为 Artin 环，\(M\ne0\) 为有限生成分次模。

\subsection{Hilbert 函数的多项式性}
\label{b3:subsec:1103:01}

\begin{proposition}\label{b3:prop:hilbert-eventual-polynomial}
设 \(A\) 为交换标准分次环，\(A_0\) 为 Artin 环，\(M\ne0\) 为有限生成分次 \(A\) 模。若其 Hilbert 级数写成 \(H_M(t)=Q(t)/(1-t)^d\)，其中 \(Q(t)\in\mathbb Z[t,t^{-1}]\)、\(d\ge0\)、\(Q(1)\ne0\)，则在 \(d>0\) 时存在次数恰为 \(d-1\) 的多项式 \(P_M(X)\in\mathbb Q[X]\)，使 \(h_M(n)=P_M(n)\) 对充分大的 \(n\) 成立，其首项为
\[
\dfrac{Q(1)}{(d-1)!}X^{d-1}.
\]
若 \(d=0\)，则 \(h_M(n)\) 最终为零，约定 Hilbert 多项式为零。累计函数 \(\sum_{j\le n}h_M(j)\) 在两种情形下均最终是次数为 \(d\)、首项系数为正的多项式。
\end{proposition}
\begin{proof}
先设 \(d>0\)。
写 \(Q(t)=\sum_a q_a t^a\)，只有有限项。公式
\((1-t)^{-d}=\sum_{n\ge0}\binom{n+d-1}{d-1}t^n\)
由 \(d\) 个几何级数相乘得到。于是当 \(n\) 大于所有出现的 \(a\) 时，
\[
h_M(n)=\sum_a q_a\binom{n-a+d-1}{d-1}.
\]
右边是有限个二项式多项式之和，因而是次数至多为 \(d-1\) 的多项式。其 \(n^{d-1}\) 系数为 \(\sum_aq_a/(d-1)!=Q(1)/(d-1)!\ne0\)，所以次数恰为 \(d-1\)。因左边非负，非零首项系数必须为正。

累计函数应计入所有 \(j\le n\) 的分量，包括可能出现的负次分量。有限生成保证这些次数有共同下界，故与几何级数相乘给出
\[
\sum_{n\in\mathbb Z}\left(\sum_{j\le n}h_M(j)\right)t^n
=H_M(t)(1+t+t^2+\cdots)=\frac{H_M(t)}{1-t}.
\]
其 \(t^n\) 系数正好把各个 \(t^j\) 项累计到次数 \(n\)，同一系数计算因此给出次数 \(d\) 及首项系数 \(Q(1)/d!\)。若 \(d=0\)，只有有限片非零，累计函数最终为正整数 \(\sum_j h_M(j)\)。
\end{proof}

这个命题给出的是“大次数以后”的等式。前面的例 \(A/(x^2,xy)\) 的 Hilbert 多项式为一，但第一个次数的额外项不能由这一个多项式恢复。

\subsection{多项式次数与维数}
\label{b3:subsec:1103:02}

\cref{b3:thm:associated-finite}已经把 Noether 环上的有限生成模逐层拆成素商模 \(A/\mathfrak p\)。为了逐次计算 Hilbert 函数，这里的子模还须保持分次；因此要从齐次元素中选择每层生成元，并证明它的零化子是齐次素理想。

\begin{lemma}\label{b3:lem:homogeneous-prime-filtration}
设 \(A\) 为交换 Noether 正分次环，\(M\) 为有限生成分次 \(A\) 模。则 \(M\) 有有限齐次子模滤过，其因子均形如 \((A/\mathfrak p)(-a)\)，其中 \(\mathfrak p\) 为 \(A\) 的齐次素理想、\(a\in\mathbb Z\)。
\end{lemma}
\begin{proof}
若 \(M=0\)，取空滤过即可。以下设 \(M\ne0\)。在非零齐次元素的零化子中选极大元 \(\mathfrak p=\operatorname{Ann}_A(m)\)。因为 \(m\ne0\)，\(\mathfrak p\) 是真理想。若 \(a=\sum_i a_i\in\mathfrak p\)，其中 \(a_i\in A_i\)，则各 \(a_i m\) 的次数 \(i+\deg m\) 互不相同，由 \(am=0\) 及分次的直和性可知每个 \(a_i m=0\)。所以 \(\mathfrak p\) 是齐次理想。

若齐次 \(u,v\) 满足 \(uv\in\mathfrak p\)、\(v\notin\mathfrak p\)，则 \(vm\) 是非零齐次元素，其零化子包含 \(\mathfrak p\)，极大性迫使二者相等，故 \(u\in\mathfrak p\)。于是 \(A/\mathfrak p\) 中两个非零齐次元素的乘积不为零。要把这一性质推广到任意元素，只需观察乘积的最高次数分量。取 \(A/\mathfrak p\) 中任意非零元素 \(a,b\)，令其最高非零齐次项分别为 \(a_r,b_s\)。乘积 \(ab\) 的次数 \(r+s\) 的分量恰为 \(a_rb_s\ne0\)，其余各项次数都较低，不能将它抵消。因此 \(ab\ne0\)，\(\mathfrak p\) 是素理想，并有 \(Am\cong(A/\mathfrak p)(-\deg m)\)。

若此子模不等于 \(M\)，在齐次商模中重复选择，再取原像。每步使齐次子模严格增加，Noether 性保证有限步终止。
\end{proof}

\begin{theorem}\label{b3:thm:hilbert-polynomial-dimension}
设 \(A\) 为交换标准分次环，\(A_0\) 为 Artin 环，\(M\ne0\) 为有限生成分次 \(A\) 模。令
\(\dim_A M=\dim\bigl(A/\operatorname{Ann}_A(M)\bigr)\)。
则 \(H_M(t)\) 在 \(t=1\) 处的极点阶数恰为 \(\dim_A M\)。若此维数为 \(d>0\)，记 \(P_M\) 为满足 \(h_M(n)=P_M(n)\)（对充分大的 \(n\)）的 Hilbert 多项式，则 \(\deg P_M=d-1\)；维数为零时，\(M\) 只有有限多个非零分次片。
\end{theorem}
\symidx{\(P_M\)：分次模 \(M\) 的 Hilbert 多项式}
\begin{proof}
标准分次和 \(A_0\) 的 Artin 性保证 \(A\) 为 Noether 环。由\cref{b3:lem:homogeneous-prime-filtration}，\(H_M\) 是各 \(H_{A/\mathfrak p}\) 的次数平移之和。每个非零因子的累计函数有正首项，次数平移又不改变增长次数，故和的增长次数是各次数的最大值，不会抵消。同时，由\cref{b3:thm:support-finite}及\cref{b3:prop:support-exact}，\(V(\operatorname{Ann}_A M)=\operatorname{Supp}M=\bigcup V(\mathfrak p)\)，支集维数也是各 \(\dim A/\mathfrak p\) 的最大值。只须对标准分次整环 \(B=A/\mathfrak p\) 证明结论。

\(B_0=A_0/(\mathfrak p\cap A_0)\) 是域，因为 Artin 环的素理想都极大；记这个域为 \(k\)。每个 \(B_n\) 的 \(A_0\) 长度等于其 \(k\) 维数。这里选取一次齐次参数，使多项式子环仍取标准分次，便可直接使用前面的计数。\cref{b3:thm:noether-normalization}的一般加权代换未必保持原来的次数，不能直接给出所需的分次比较。

先设 \(k\) 无限。写 \(B=k[z_1,\ldots,z_r]\)、\(\deg z_i=1\)。若这些元素有代数关系，其关系理想齐次，故可选一个非零齐次关系 \(f\)。使用第五章无限域证明中的线性代换，选 \(c_i\in k\) 使 \(f(c_1,\ldots,c_{r-1},1)\ne0\)，以 \(z_i-c_i z_r\) 替换前 \(r-1\) 个生成元。新生成元仍为一次齐次元，而变换后的关系对 \(z_r\) 的最高幂具有非零常数系数；除以这个系数，就得到首一关系。因此 \(B\) 是前 \(r-1\) 个新生成元所生成子环上的有限模。对生成元个数归纳，得到一次代数独立元 \(y_1,\ldots,y_d\)，使 \(B\) 有限于 \(P=k[y_1,\ldots,y_d]\)。

把一组有限 \(P\) 模生成元分成齐次部分，便得到有限组齐次生成元，设其次数为 \(a_i\)。于是有分次单射 \(P\hookrightarrow B\) 和分次满射 \(\bigoplus_iP(-a_i)\twoheadrightarrow B\)。累计维数被 \(P\) 的累计函数与有限个其平移之和夹住；二者都以 \(n^d\) 的正倍数增长，所以累计函数的多项式次数为 \(d\)。另一方面，\cref{b3:thm:integral-dimension}和\cref{b3:thm:polynomial-dimension}给出 \(\dim B=\dim P=d\)。

若 \(k\) 有限，为使用上述线性选择，改用无限域 \(k(u)\)，其中 \(u\) 为不定元。\(B'=B\otimes_k k(u)\) 是 \(B[u]\) 对 \(k[u]\setminus\{0\}\) 的局部化，仍为整环。各分次片满足
\[
B'_n=B_n\otimes_k k(u),\qquad
\dim_{k(u)}B'_n=\dim_k B_n,
\]
所以它们给出同一个 Hilbert 函数。其分式域是 \(\operatorname{Frac}(B)(u)\)。取 \(\operatorname{Frac}(B)/k\) 的超越基 \(T\)，在以 \(u\) 为不定元的多项式中比较系数，可知 \(T\) 在 \(k(u)\) 上仍代数独立；原有的代数方程也仍成立，故它仍是 \(\operatorname{Frac}(B)(u)/k(u)\) 的超越基。由\cref{b3:thm:affine-dimension-trdeg}，\(\dim B'=\dim B\)。这次变基保持了需要比较的增长和维数，应用无限域情形即得结论。
\end{proof}

\subsection{重数与首项系数的归一化}
\label{b3:subsec:1103:03}

\begin{definition}\label{b3:def:graded-multiplicity}
设 \(A\) 为交换标准分次环，\(A_0\) 为 Artin 环，\(M\ne0\) 为有限生成分次 \(A\) 模，\(P_M\) 为 \(M\) 的 Hilbert 多项式。若 \(d=\dim_A M>0\)，把 \(P_M\) 的首项写成
\[
P_M(n)=\dfrac{e(M)}{(d-1)!}n^{d-1}+\text{低次项},
\]
称 \(e(M)\) 为这个分次模的\term[zhongshu]{重数}[multiplicity]。若 \(d=0\)，约定 \(e(M)=\sum_n\ell_{A_0}(M_n)\)。对零模另约定 \(e(0)=0\)。
\end{definition}
\symidx{\(e(M)\)：分次模 \(M\) 的重数}

这里的 \(e(M)\) 按选定的分次计算；只给出底层模而不说明分次，还不能直接使用这一定义。由前两项结果，约去 \(1-t\) 公因子后的分子满足 \(e(M)=Q(1)\)，所以重数是正整数。

对保持次数的短正合列 \(0\to L\to M\to N\to0\)，若中间项维数为 \(d>0\)，则 \(n^{d-1}\) 的系数可加。只将维数为 \(d\) 的非零项计入，就得到重数可加。例如若 \(L\ne0\) 且 \(\dim L<d\)，则 \(P_L\) 为零或次数至多为 \(d-2\)，不贡献 \(n^{d-1}\) 项；此时 \(\dim N=d\)，且 \(e(M)=e(N)\)。若中间项维数为零，则各项都只有有限个非零分次片，逐次利用长度可加再求和，同样得到重数可加。

\ProseParagraphBreak
例如 \(S=k[x_1,\ldots,x_r]\)，若 \(0\ne f\in S\) 齐次且 \(\deg f=a>0\)，则
\[
H_{S/(f)}(t)=\dfrac{1-t^a}{(1-t)^r}
 =\dfrac{1+t+\cdots+t^{a-1}}{(1-t)^{r-1}}.
\]
当 \(r\ge2\) 时维数为 \(r-1\)，重数为 \(a\)；当 \(r=1\) 时商的总维数为 \(a\)，与零维约定一致。由两个次数分别为正整数 \(a,b\) 的齐次元依次作非零因子商，可同样得到分子 \((1-t^a)(1-t^b)\) 与重数 \(ab\)。这里要求第二个元在第一个商环上仍为非零因子，否则乘法正合列的左端核不为零，不能照用该乘积公式。

% !TeX root = ../../Book3.tex
% 纲要标题：### 11.4 局部环中的增长
% 纲要位置：工作记录/计划纲要.md，第 2291 行。
% 纲要细目（写作范围）：
% * 伴随分次环
% * Hilbert–Samuel 函数及多项式
% * 参数理想和重数的第一批例子；切锥的局部标准基计算见附录 F，不用于本章基础结论的证明
% * 固定单层长度 \(\ell(\mathfrak m^nM/\mathfrak m^{n+1}M)\) 与累计长度 \(\ell(M/\mathfrak m^{n+1}M)\) 的记号，分别说明多项式次数和重数的归一化；处理零维例外
% * 定义局部滤过的初始形式和初始理想，证明商环的伴随分次公式；在正文证明给定生成元的最低次项可能遗漏初始关系，并保留切锥坐标环的非约化信息
% * 有限点集上的求值码、次数滤过与维数计数；低次数 Reed–Muller 实例
% #### 11.4.1 伴随分次环
% #### 11.4.2 Hilbert–Samuel 函数与多项式
% #### 11.4.3 参数理想、重数与切锥的例子
% #### 11.4.4 单层长度、累计长度与零维例外
% #### 11.4.5 有限点集上的求值码
% ---

\section{局部环中的增长}
\label{b3:sec:1104}

局部环未必本来分次，但极大理想的幂提供了逐层近似。每层取商后，乘法恰好保存阶数相加，于是得到一个标准分次环。本节把这一构造与前面的计数联系起来，并证明增长次数等于局部维数；所需的子模滤过估计也在这里证明。

\subsection{伴随分次环}
\label{b3:subsec:1104:01}

\begin{definition}\label{b3:def:associated-graded-adic}
设 \(R\) 为交换环、\(I\subseteq R\) 为理想、\(M\) 为 \(R\) 模。定义\term[bansui-fenci-huan]{伴随分次环}[associated graded ring]及相应分次模
\[
\operatorname{gr}_I(R)=\bigoplus_{n\ge0}I^n/I^{n+1},\qquad
\operatorname{gr}_I(M)=\bigoplus_{n\ge0}I^nM/I^{n+1}M.
\]
乘法由 \([a]_i[b]_j=[ab]_{i+j}\) 给出，模作用同样定义。换代表元时多出的项落入高一阶，因此这些运算良定义。
\symidx{\(\operatorname{gr}_I(R)\)、\(\operatorname{gr}_I(M)\)：理想进伴随分次环与模}
\end{definition}

这里的次数来自 \(I\)-进滤过：\(a\in I^n\setminus I^{n+1}\) 的初始类位于次数 \(n\)，未必与原环中的多项式次数相同。例如，对域 \(k\)，在 \(R=k[x,y]_{(x,y)}\)、\(I=(x^2,y^3)R\) 中，\(x^2+I^2\) 与 \(y^3+I^2\) 都是次数一的非零类。

若 \(I=(a_1,\ldots,a_r)\)，则有标准分次满同态
\[
(R/I)[X_1,\ldots,X_r]\twoheadrightarrow\operatorname{gr}_I(R),
\qquad X_i\longmapsto a_i+I^2.
\]
每个 \(X_i\) 的次数都规定为一，代表第一层 \(I/I^2\) 中的类 \(a_i+I^2\)。这些类生成整个伴随分次环，因为 \(I^n\) 由生成元的 \(n\) 次乘积生成。取 \(I=\mathfrak m\) 时，第一层正是第九章定义嵌入维数所用的 \(\mathfrak m/\mathfrak m^2\)。若 \(M\) 有限生成，\(\operatorname{gr}_I(M)\) 由 \(M/IM\) 中的有限组生成元在次数零生成。以下取 \((R,\mathfrak m,k)\) 为非零 Noether 局部环、\(I\) 为 \(\mathfrak m\) 准素理想、\(M\ne0\) 为有限生成模。此时 \(R/I\) 是 Artin 环；又由 \(I\subseteq\mathfrak m=J(R)\) 和\cref{b3:thm:nakayama}，有 \(M/IM\ne0\)，故 \(\operatorname{gr}_I(M)\ne0\)，可以对它使用前面的 Hilbert 理论。

对子模 \(N\subseteq M\)，它自身的 \(I\)-进滤过为 \(I^nN\)，从 \(M\) 继承的滤过则为 \(I^nM\cap N\)。前者总包含于后者，但两者未必相等。下面的估计说明，后者相对于前者的偏移可以由一个与 \(n\) 无关的整数控制，随后便能用它比较商模的长度增长。

\begin{lemma}[诱导滤过的控制]\label{b3:lem:induced-filtration-bound}
设 \(R\) 为交换 Noether 环、\(I\subseteq R\) 为理想、\(M\) 为有限生成 \(R\) 模、\(N\subseteq M\) 为子模。则存在 \(c\ge0\)，使对每个 \(n\ge c\) 有
\[
I^nN\subseteq I^nM\cap N\subseteq I^{n-c}N.
\]
\end{lemma}
\begin{proof}
把各个 \(I^n\) 及 \(I^nM\) 放入同一个分次对象，便可同时控制全部次数。为此引入一个形式变量 \(t\)，用 \(t^n\) 记录滤过次数，考虑正分次代数及分次模
\[
\mathcal R=\bigoplus_{n\ge0}I^nt^n\subseteq R[t],\qquad
\mathcal M=\bigoplus_{n\ge0}I^nMt^n.
\]
若 \(I=(a_1,\ldots,a_r)\)，则 \(\mathcal R=R[a_1t,\ldots,a_rt]\)，由 Hilbert 基定理为 Noether 环。\(M\) 的有限组生成元放在次数零，就生成 \(\mathcal M\) 作为 \(\mathcal R\) 模。于是其齐次子模
\[
\mathcal N=\bigoplus_{n\ge0}(I^nM\cap N)t^n
\]
有限生成。选齐次生成元 \(z_jt^{d_j}\)，并取 \(c\ge\max_jd_j\)；有限个生成元的次数有共同上界，这就给出控制所有层的同一个 \(c\)。比较次数 \(n\ge c\) 的部分，得到
\[
I^nM\cap N=\sum_j I^{n-d_j}z_j\subseteq I^{n-c}N.
\]
另一包含直接来自 \(N\subseteq M\)。零子模时可取 \(c=0\)。
\end{proof}

\chapref{b3:ch:13}将用同一构造证明更精确的 Artin–Rees 等式，并系统研究两种滤过的拓扑。

\subsection{Hilbert–Samuel 函数与多项式}
\label{b3:subsec:1104:02}

\begin{definition}\label{b3:def:hilbert-samuel-function}
设 \((R,\mathfrak m)\) 为交换 Noether 局部环，\(I\subseteq R\) 为 \(\mathfrak m\) 准素理想，\(M\) 为非零有限生成 \(R\) 模。定义单层长度与累计长度
\[
h_{I,M}(n)=\ell_R(I^nM/I^{n+1}M),\qquad
\chi_{I,M}(n)=\ell_R(M/I^{n+1}M)\quad(n\ge0).
\]
累计函数 \(\chi_{I,M}\) 称为\term[hilbert-samuel-hanshu]{Hilbert–Samuel 函数}[Hilbert–Samuel function]；本书始终用下标和 \(\chi\) 区分它与单层函数。
\symidx{\(h_{I,M}(n)\)、\(\chi_{I,M}(n)\)：单层长度与累计 Hilbert--Samuel 函数}
\end{definition}

长度均有限，因为 \(\sqrt I=\mathfrak m\) 使每个 \(R/I^n\)（\(n\ge1\)）是 Artin 局部环。这一条件把相应商模的支集限制在闭点。若在 \(R=k[x,y]_{(x,y)}\) 中取 \(I=(x)\)，则 \(R/I\cong k[y]_{(y)}\)，其中 \((y)\supsetneq(y^2)\supsetneq\cdots\) 严格下降，所以已经在第一次取商时就没有有限长度。对于定义中的准素理想，逐层取商得到
\begin{equation}\label{b3:eq:samuel-cumulative}
\chi_{I,M}(n)=\sum_{j=0}^n h_{I,M}(j).
\end{equation}
\(R\) 长度与 \(R/I\) 长度在各层相同，故\cref{b3:prop:hilbert-eventual-polynomial}已经说明 \(\chi_{I,M}\) 最终是一个首项为正的多项式。现在确定它的次数。

\ProseParagraphBreak
这一维数刻画称为\term[hilbert-samuel-dingli]{Hilbert–Samuel 定理}[Hilbert--Samuel theorem]。

\begin{theorem}[Hilbert–Samuel]\label{b3:thm:hilbert-samuel}
设 \((R,\mathfrak m)\) 为交换 Noether 局部环，\(I\subseteq R\) 为 \(\mathfrak m\) 准素理想，\(M\) 为非零有限生成 \(R\) 模，\(d=\dim_R M=\dim\bigl(R/\operatorname{Ann}_R(M)\bigr)\)。则对充分大的 \(n\)，
\[
\chi_{I,M}(n)=\dfrac{e_I(M)}{d!}n^d+\text{低次项},
\qquad e_I(M)\in\mathbb Z_{>0}.
\]
特别地，\(\dim_{\operatorname{gr}_I(R)}\operatorname{gr}_I(M)=d\)。整数 \(e_I(M)\) 称为 \(M\) 关于 \(I\) 的 Hilbert–Samuel 重数。
\end{theorem}
\begin{proof}
先记最终累计多项式的次数为 \(\delta_I(M)\)。取 \(a\ge1\) 使 \(\mathfrak m^a\subseteq I\subseteq\mathfrak m\)。于是
\[
\chi_{\mathfrak m,M}(n)\le\chi_{I,M}(n)
\le\chi_{\mathfrak m,M}\bigl(a(n+1)-1\bigr).
\]
这里 \(\mathfrak m^{a(n+1)}M\subseteq I^{n+1}M\subseteq\mathfrak m^{n+1}M\)，较小的子模给出较大的商，因而得到所示长度不等式。两个界都是首项为正的多项式增长，故 \(\delta_I(M)=\delta_{\mathfrak m}(M)\)。以下只研究 \(I=\mathfrak m\)，并简记 \(\delta(M)\)。

先取 \(M=R\)。对所有 Noether 局部环中的严格素理想链，按链长 \(r\) 归纳，证明 \(r\le\delta(R)\)。\(r=0\) 时直接成立。若链为 \(\mathfrak p_0\subsetneq\cdots\subsetneq\mathfrak p_r\)、\(r\ge1\)，令 \(B=R/\mathfrak p_0\)、\(\mathfrak m_B=\mathfrak m/\mathfrak p_0\)。它是局部整环，其累计长度不超过 \(R\) 的对应长度，所以 \(\delta(B)\le\delta(R)\)。选 \(0\ne f\in\mathfrak p_1/\mathfrak p_0\)，令 \(C=B/fB\)、\(\mathfrak m_C=\mathfrak m_B/fB\)。这些理想的像在 \(C\) 中给出长度 \(r-1\) 的素链。

现在要让增长次数也至少下降一。长度正合列将从 \(B\) 的增长中扣去 \(fB\) 的贡献；虽然乘以 \(f\) 给出 \(B\cong fB\)，它未必将两边的理想进滤过恰好对齐。Artin--Rees 型估计提供一个固定误差 \(c\)，从而把所需上界化为固定步长差分 \(\chi_B(n)-\chi_B(n-c)\)。无论这个正整数步长多大，多项式差分仍会降一次。为取得该上界，对 \(fB\subseteq B\) 应用\cref{b3:lem:induced-filtration-bound}，可取整数 \(c\ge1\)，使 \(n\ge c\) 时
\[
fB\cap\mathfrak m_B^{n+1}\subseteq\mathfrak m_B^{n+1-c}fB.
\]
因此有满射及同构
\[
\dfrac{fB}{fB\cap\mathfrak m_B^{n+1}}
\twoheadrightarrow\dfrac{fB}{\mathfrak m_B^{n+1-c}fB}
\cong\dfrac{B}{\mathfrak m_B^{n+1-c}}.
\]
分母 \(fB\cap\mathfrak m_B^{n+1}\) 较小，所以它对应的商满射到右侧较大的分母所对应的商，左侧长度至少是右侧长度。最后的同构由乘以 \(f\) 给出，使用了 \(B\) 为整环及 \(f\ne0\)。再取商映射的短正合列
\[
0\longrightarrow\dfrac{fB}{fB\cap\mathfrak m_B^{n+1}}
\longrightarrow\dfrac{B}{\mathfrak m_B^{n+1}}
\longrightarrow\dfrac{C}{\mathfrak m_C^{n+1}}
\longrightarrow0.
\]
其中各模长度有限，\(C\) 模的 \(B\) 长度与 \(C\) 长度相同。从 \(\chi_{\mathfrak m_B,B}(n)\) 中扣去的左端长度至少为 \(\chi_{\mathfrak m_B,B}(n-c)\)，故长度相减给出
\[
\chi_{\mathfrak m_C,C}(n)
\le\chi_{\mathfrak m_B,B}(n)-\chi_{\mathfrak m_B,B}(n-c).
\]
设 \(B\) 的最终累计多项式为 \(P_B(n)=a n^e+\text{低次项}\)，其中 \(a>0\)、\(e=\delta(B)\)。若 \(e=0\)，右边最终为零，与 \(C\ne0\) 使左边为正矛盾。因此 \(e\ge1\)，且
\[
P_B(n)-P_B(n-c)=ace\,n^{e-1}+\text{低次项}.
\]
它的次数恰为 \(e-1\)。归纳假设给出
\(r-1\le\delta(C)\le\delta(B)-1\)，所以 \(r\le\delta(R)\)。对全部链取上确界，得到所需下界。

反向令 \(d=\dim R\)。\cref{b3:thm:system-of-parameters}给出由 \(d\) 个元素生成的 \(\mathfrak m\) 准素理想 \(\mathfrak q\)。\(\operatorname{gr}_{\mathfrak q}(R)\) 是 \((R/\mathfrak q)[X_1,\ldots,X_d]\) 的商。\(d\) 个变量中总次数不超过 \(n\) 的单项式共有 \(\binom{n+d}{d}\) 个，每个单项式对应一份 \(R/\mathfrak q\)，所以前 \(n+1\) 层的累计长度不超过它们的长度之和。由 \(\mathfrak q\subseteq\mathfrak m\)，便得到
\[
\chi_{\mathfrak m,R}(n)\le\chi_{\mathfrak q,R}(n)
\le\ell_R(R/\mathfrak q)\binom{n+d}{d}.
\]
所以 \(\delta(R)\le d\)。\(d=0\) 时参数理想为零，\(R\) 已为 Artin 环，同一界为常数，论证仍成立。

一般有限生成模 \(M\) 令 \(B=R/\operatorname{Ann}_R(M)\)、\(\mathfrak m_B=\mathfrak m/\operatorname{Ann}_R(M)\)，并取生成元 \(m_1,\ldots,m_s\)。满射 \(B^s\twoheadrightarrow M\) 给出 \(M\) 增长的上界；为了得到反向比较，要把 \(B\) 嵌入若干份 \(M\)。作为 \(B\) 模，\(M\) 的零化子为零，故可取单射
\[
\beta:B\longrightarrow M^s,\qquad b\longmapsto(bm_1,\ldots,bm_s).
\]
若 \(\beta(b)=0\)，则 \(b\) 零化每个生成元，因而零化整个 \(M\)，只能有 \(b=0\)，这证明单射性。前述满射取商给出 \(\chi_{\mathfrak m_B,M}(n)\le s\cdot\chi_{\mathfrak m_B,B}(n)\)。令 \(N=\beta(B)\subseteq M^s\)，由\cref{b3:lem:induced-filtration-bound}可取整数 \(c\ge0\)，使充分大的 \(n\) 满足
\[
N\cap\mathfrak m_B^{n+1}M^s\subseteq\mathfrak m_B^{n+1-c}N.
\]
取商后有单射与满射
\[
\begin{gathered}
\dfrac{N}{N\cap\mathfrak m_B^{n+1}M^s}
\hookrightarrow\dfrac{M^s}{\mathfrak m_B^{n+1}M^s},\\[3pt]
\dfrac{N}{N\cap\mathfrak m_B^{n+1}M^s}
\twoheadrightarrow\dfrac{N}{\mathfrak m_B^{n+1-c}N}
\cong\dfrac{B}{\mathfrak m_B^{n+1-c}}.
\end{gathered}
\]
比较长度便得
\[
\chi_{\mathfrak m_B,B}(n-c)\le s\cdot\chi_{\mathfrak m_B,M}(n).
\]
各长度在 \(R\) 与 \(B\) 上相同，且 \(\mathfrak m_B^nM=\mathfrak m^nM\)，所以 \(\delta(M)=\delta(B)=\dim B=d\)。

最后，\(\operatorname{gr}_I(M)\) 的累计函数就是 \(\chi_{I,M}\)。由\cref{b3:thm:hilbert-polynomial-dimension}和\cref{b3:prop:hilbert-eventual-polynomial}，其维数为 \(d\)，将它的 Hilbert 级数约分后可写为
\[
H_{\operatorname{gr}_I(M)}(t)=\frac{Q(t)}{(1-t)^d},
\qquad Q(t)\in\mathbb Z[t],\quad Q(1)>0.
\]
累计首项为 \(Q(1)\cdot n^d/d!\)，故令 \(e_I(M)=Q(1)\) 即得结论；零维情形中，\(Q(t)\) 就是有限层的长度计数多项式，\(Q(1)\) 为全部层长度之和。
\end{proof}

素理想链给出增长次数的下界，参数理想给出上界；这一比较把维数转化为可由有限商 \(M/I^{n+1}M\) 读出的数据。局部增长与维数的经典比较可对照 \cite[\S13, Theorems 13.3--13.4, pp. 98--100]{Matsumura1989CommutativeRingTheory}。

\subsection{参数理想、重数与切锥的例子}
\label{b3:subsec:1104:03}

设 \(k\) 为域。在 \(R=k[x,y]_{(x,y)}\) 中，令 \(\mathfrak q=(x^a,y^b)\)，其中 \(a,b\ge1\) 为整数。模掉 \(\mathfrak q^{n+1}\) 后，剩下的单项式唯一写成
\(x^{au+i}y^{bv+j}\)，其中 \(0\le i<a\)、\(0\le j<b\)、\(u+v\le n\)。因此
\[
\chi_{\mathfrak q,R}(n)=ab\binom{n+2}{2},\qquad e_{\mathfrak q}(R)=ab.
\]
局部化不改变这些商的 \(k\) 维数，因为它们已经是极大理想幂零的局部环。特别地，\(e_{(x,y)}(R)=1\)。参数理想不同，重数可以不同；不能只记一个没有注明理想的数字。

\ProseParagraphBreak
为了计算局部商环的伴随分次，需要保留方程最低次数中的全部项。分式的分母是单位，它的常数项会影响这些项的系数，却不改变阶数。

\begin{definition}\label{b3:def:local-initial-form}
设 \(k\) 为域，\(r\ge1\) 为整数，\(R_0=k[x_1,\ldots,x_r]_{(x_1,\ldots,x_r)}\)，\(\mathfrak m_0=(x_1,\ldots,x_r)R_0\)。对 \(0\ne f\in R_0\)，定义其 \(\mathfrak m_0\) 进阶数及初始形式为
\[
\operatorname{ord}_{\mathfrak m_0}(f)
=\max\{\,n\in\mathbb Z_{\ge0}\mid f\in\mathfrak m_0^n\,\}=a,
\qquad
\operatorname{in}_{\mathfrak m_0}(f)
=f+\mathfrak m_0^{a+1}\in\mathfrak m_0^a/\mathfrak m_0^{a+1}.
\]
若 \(J\subseteq\mathfrak m_0\) 为理想，定义其 \(\mathfrak m_0\) 进初始理想为伴随分次环中的齐次理想
\[
\operatorname{in}_{\mathfrak m_0}(J)
=\bigl(\,\operatorname{in}_{\mathfrak m_0}(g)\mid 0\ne g\in J\,\bigr)
\subseteq\operatorname{gr}_{\mathfrak m_0}(R_0).
\]
当 \(J=0\) 时，此理想为零。
\end{definition}
\symidx{\(\operatorname{ord}_{\mathfrak m_0}(f)\)：元素的 \(\mathfrak m_0\)-进阶数}
\symidx{\(\operatorname{in}_{\mathfrak m_0}(f)\)、\(\operatorname{in}_{\mathfrak m_0}(J)\)：局部初始形式与初始理想}

将 \(f\) 写成 \(g/u\)，其中 \(g,u\in k[x_1,\ldots,x_r]\)、\(g\ne0\)、\(u(0)\ne0\)，并将 \(g\) 的最低非零齐次部分记为 \(g_a\)。局部化后的 \(\mathfrak m_0^n\) 由总次数至少为 \(n\) 的单项式生成。若 \(f\in\mathfrak m_0^n\)，清分母得到某个 \(v(0)\ne0\) 的多项式 \(v\)，使 \(vg\in(x_1,\ldots,x_r)^n\)；乘以 \(v\) 不改变最低次数，故 \(a\ge n\)。反向由 \(g\in(x_1,\ldots,x_r)^n\) 立即成立。因此上述最大值存在且等于 \(a\)。每个分式在次数 \(a\) 的类由 \(g_a/u(0)\) 表示，这给出
\[
\operatorname{gr}_{\mathfrak m_0}(R_0)\cong k[X_1,\ldots,X_r],
\qquad
\operatorname{in}_{\mathfrak m_0}(f)=\dfrac{g_a(X_1,\ldots,X_r)}{u(0)}.
\]
这里 \(X_i\) 对应 \(x_i+\mathfrak m_0^2\)。若另有 \(f=g'/u'\)，则 \(gu'=g'u\)；比较最低非零齐次部分，先得两种表示的最低次数相同，再得 \(u'(0)g_a=u(0)g'_a\)，所以初始形式与表示无关。各个齐次多项式都提供相应次数的类，也只有零齐次多项式给出零类，因而所示映射确为同构。
这一初始形式保留最低次数的全部项，与第八章按单项式序选取的首项不同。例如，\(f=x^2+xy+y^3\) 的局部初始形式是 \(X^2+XY\)，而在字典序 \(x>y\) 下，其 Gröbner 首项是单个单项式 \(x^2\)。前者记录局部滤过中的齐次关系，后者由选定的单项式序决定。

\begin{proposition}\label{b3:prop:quotient-local-initial-ideal}
设 \(k\) 为域，\(r\ge1\) 为整数，\(R_0=k[x_1,\ldots,x_r]_{(x_1,\ldots,x_r)}\)，\(\mathfrak m_0=(x_1,\ldots,x_r)R_0\)，\(J\subseteq\mathfrak m_0\) 为理想。令 \(R=R_0/J\)、\(\mathfrak m=\mathfrak m_0/J\)。则商映射诱导分次环同构
\[
\operatorname{gr}_{\mathfrak m}(R)
\cong\dfrac{\operatorname{gr}_{\mathfrak m_0}(R_0)}
{\operatorname{in}_{\mathfrak m_0}(J)}.
\]
一般不能只用 \(J\) 的给定生成元的初始形式生成 \(\operatorname{in}_{\mathfrak m_0}(J)\)。具体地，在 \(R_0=k[x,y]_{(x,y)}\) 中取 \(J=(x^2+y^3,xy)R_0\)，则有
\[
\operatorname{in}_{\mathfrak m_0}(J)=(X^2,XY,Y^4),\qquad
\operatorname{gr}_{\mathfrak m}(R)\cong k[X,Y]/(X^2,XY,Y^4).
\]
它的 Hilbert 级数为 \(1+2t+t^2+t^3\)，且 \(\ell_R(R)=5\)；只取两个原生成元的初始形式 \(X^2,XY\) 会漏掉 \(Y^4\)。
\end{proposition}
\begin{proof}
因 \(\mathfrak m^n=(\mathfrak m_0^n+J)/J\)，商映射在每个次数都诱导满射
\[
\mathfrak m_0^n/\mathfrak m_0^{n+1}
\longrightarrow\mathfrak m^n/\mathfrak m^{n+1}.
\]
这些映射与乘法相容。在次数 \(n\)，其核为
\[
\dfrac{\mathfrak m_0^n\cap(J+\mathfrak m_0^{n+1})}{\mathfrak m_0^{n+1}}
=\dfrac{(J\cap\mathfrak m_0^n)+\mathfrak m_0^{n+1}}{\mathfrak m_0^{n+1}}.
\]
核中的非零类由某个 \(g\in J\cap\mathfrak m_0^n\) 代表，且 \(g\notin\mathfrak m_0^{n+1}\)，所以恰为 \(\operatorname{in}_{\mathfrak m_0}(g)\)。反过来，每个 \(J\) 中非零元素的初始形式都属于这个分次映射的核。核本身是齐次理想，故正好等于 \(\operatorname{in}_{\mathfrak m_0}(J)\)，得到一般公式。

在具体例子中，记 \(f_1=x^2+y^3\)、\(f_2=xy\)，则
\[
f_3\coloneqq yf_1-xf_2=y^4\in J,
\]
所以 \(X^2,XY,Y^4\) 都属于局部初始理想。为确定这些关系是否已经足够，先在多项式环 \(k[x,y]\) 中使用字典序 \(x>y\)。\(f_1,f_2,f_3\) 的首单项式分别是 \(x^2,xy,y^4\)，三个 \(S\) 多项式为
\[
S(f_1,f_2)=f_3,\qquad
S(f_1,f_3)=y^3f_3,\qquad
S(f_2,f_3)=0.
\]
它们的余式均为零，故\cref{b3:thm:buchberger-criterion}说明这是一组 Gröbner 基。因而多项式商 \(A=k[x,y]/(f_1,f_2)\) 有五个标准单项式组成的基
\[
1,\ \bar x,\ \bar y,\ \bar y^2,\ \bar y^3.
\]
其中 \(\bar x^2=-\bar y^3\)、\(\bar x\bar y=0\)、\(\bar y^4=0\)。故 \((\bar x,\bar y)\) 幂零，每个常数项非零的类都可用有限几何级数取逆，\(A\) 已是剩余域为 \(k\) 的局部环。原点处局部化不再改变它，因而 \(R\cong A\)。

由这些乘法关系，\(\mathfrak m\) 的 \(k\) 基为 \(\bar x,\bar y,\bar y^2,\bar y^3\)，\(\mathfrak m^2\) 的基为 \(\bar y^2,\bar y^3\)，\(\mathfrak m^3\) 的基为 \(\bar y^3\)，且 \(\mathfrak m^4=0\)。所以各层的维数依次为 \(1,2,1,1\)。已得初始关系给出分次满射
\[
k[X,Y]/(X^2,XY,Y^4)\twoheadrightarrow\operatorname{gr}_{\mathfrak m}(R).
\]
左侧的标准单项式也为 \(1,X,Y,Y^2,Y^3\)，总维数为五，与右侧各层之和相等，故满射为同构。这证明初始理想的等式、Hilbert 级数及长度；若仅保留 \((X^2,XY)\)，所有 \(Y\) 的幂都会留下，错误地得到无限维的分次商。
\end{proof}

\begin{proposition}\label{b3:prop:hypersurface-tangent-cone}
设 \(k\) 为域，\(r\ge1\) 为整数，\(R_0=k[x_1,\ldots,x_r]_{(x_1,\ldots,x_r)}\)，\(\mathfrak m_0=(x_1,\ldots,x_r)R_0\)，\(0\ne f\in\mathfrak m_0\)。令 \(a=\operatorname{ord}_{\mathfrak m_0}(f)\)、\(f_a=\operatorname{in}_{\mathfrak m_0}(f)\in k[X_1,\ldots,X_r]\)，以及 \(R=R_0/(f)\)、\(\mathfrak m=\mathfrak m_0/(f)\)。则
\[
\operatorname{gr}_{\mathfrak m}(R)\cong k[X_1,\ldots,X_r]/(f_a).
\]
\end{proposition}
\begin{proof}
取 \(0\ne g\in R_0\)，并令 \(b=\operatorname{ord}_{\mathfrak m_0}(g)\)。伴随分次环 \(k[X_1,\ldots,X_r]\) 为整环，因此 \(\operatorname{in}_{\mathfrak m_0}(g)f_a\ne0\)。乘法的定义便给出
\[
\operatorname{ord}_{\mathfrak m_0}(gf)=a+b,
\qquad
\operatorname{in}_{\mathfrak m_0}(gf)
=\operatorname{in}_{\mathfrak m_0}(g)f_a.
\]
所以 \((f)\) 中每个非零元素的初始形式都是 \(f_a\) 的倍数；而 \(f\) 本身又给出生成元 \(f_a\)。于是 \(\operatorname{in}_{\mathfrak m_0}((f))=(f_a)\)，由\cref{b3:prop:quotient-local-initial-ideal}即得公式。
\end{proof}

此点的\term[qiezhui]{切锥}[tangent cone]以伴随分次环 \(\operatorname{gr}_{\mathfrak m}(R)\) 为坐标环。它的素谱给出点及其拓扑，但非约化信息仍需由伴随分次环本身保存。最低次数可以由缩放看出：引入缩放参数 \(\lambda\)，在尖点方程 \(y^2-x^3=0\) 中代入 \(x=\lambda X,y=\lambda Y\)，先对 \(\lambda\ne0\) 除去共同的 \(\lambda^2\)，得到
\[
Y^2-\lambda X^3=0.
\]
把这一族延伸到 \(\lambda=0\)，便得到 \(Y^2=0\)；在实数的图像解释中，越小的非零 \(\lambda\) 把越靠近原点的坐标放大到固定的 \((X,Y)\) 尺度，而所写的方程在任意域上都成立。对方程 \(y^2-x^2-x^3=0\)，同样得到 \(Y^2-X^2-\lambda X^3=0\)，在 \(\lambda=0\) 留下 \(Y^2-X^2=0\)。所以取最低齐次部分是在保留缩放后的极限方程，包括重数。尖点切锥中的平方不能删去：\(Y^2=0\) 是加厚直线，而 \(Y=0\) 只是它的约化支集。在特征不为二时，节点切锥则是两条不同直线。

尖点方程没有一次项，所以 \(\mathfrak m/\mathfrak m^2\) 中 \(\bar x,\bar y\) 的像独立，其对偶所给出的 Zariski 切空间是二维的整个平面。切锥却保留了最低非零齐次关系 \(Y^2\)，其约化支集只有一条直线。两者记录的信息不同：切空间取一次线性化，切锥还保留更高次的初始齐次关系。

\ProseParagraphBreak
尖点和节点的伴随分次环的 Hilbert 级数都是 \((1+t)/(1-t)\)，因此
\[
h_{\mathfrak m,R}(0)=1,\quad h_{\mathfrak m,R}(n)=2\ (n\ge1),
\qquad \chi_{\mathfrak m,R}(n)=2n+1,
\]
重数都为二；在特征不为二时，前者切锥非约化，后者切锥有两条不同直线。相同重数不意味着相同切锥。本卷附录 F 将进一步讨论计算多方程全部初始关系所用的局部标准基。

\subsection{单层长度、累计长度与零维例外}
\label{b3:subsec:1104:04}

累计长度由各层相加得到，单层长度则是累计长度的差分。维数 \(d>0\) 时，差分把首项中的 \(n^d/d!\) 变成 \(n^{d-1}/(d-1)!\)，所以累计多项式次数为 \(d\)，单层多项式次数为 \(d-1\)，而重数保持相同。零维时，累计长度最终为总长度，单层长度最终为零。\cref{b3:tab:graded-local-hilbert-growth}将分次与局部两种情形并列。

\begin{center}
\begin{minipage}{\textwidth}
\makeatletter
\def\@captype{table}
\makeatother
\centering
\caption[分次与局部 Hilbert 函数的增长]{分次与局部 Hilbert 函数的增长。分次行中，\(A\) 为交换标准分次环，\(A_0\) 为 Artin 环，\(M\ne0\) 为有限生成分次 \(A\) 模，\(d=\dim_A M\)。局部行中，\((R,\mathfrak m)\) 为交换 Noether 局部环，\(I\) 为 \(\mathfrak m\) 准素理想，\(M\ne0\) 为有限生成 \(R\) 模，\(d=\dim_R M\)。所有首项与最终值均对充分大的 \(n\) 而言。}
\label{b3:tab:graded-local-hilbert-growth}
\begin{tabularx}{\textwidth}{@{}>{\raggedright\arraybackslash}p{0.43\textwidth}>{\raggedright\arraybackslash}p{0.25\textwidth}>{\raggedright\arraybackslash}X@{}}
\toprule
函数 & \(d>0\) 时的首项 & \(d=0\) 时的最终值\\
\midrule
分次单层 \(h_M(n)=\ell_{A_0}(M_n)\) & \(\dfrac{e(M)}{(d-1)!}n^{d-1}\) & \(0\)\\
分次累计 \(\sum_{j\le n}\ell_{A_0}(M_j)\) & \(\dfrac{e(M)}{d!}n^d\) & \(\sum_j\ell_{A_0}(M_j)\)\newline\(=e(M)\)\\
局部单层 \(h_{I,M}(n)\)\newline\(=\ell_R(I^nM/I^{n+1}M)\) & \(\dfrac{e_I(M)}{(d-1)!}n^{d-1}\) & \(0\)\\
局部累计 \(\chi_{I,M}(n)\)\newline\(=\ell_R(M/I^{n+1}M)\) & \(\dfrac{e_I(M)}{d!}n^d\) & \(\ell_R(M)=e_I(M)\)\\
\bottomrule
\end{tabularx}
\end{minipage}
\end{center}

在局部情形，若 \(d=0\)，\(M\) 有限长度，某个 \(\mathfrak m^NM\) 为零；因 \(I\subseteq\mathfrak m\)，也有 \(I^NM=0\)，所以各层最终为零，累计长度最终为 \(\ell_R(M)\)。分次情形则只有有限多个非零分次片，其累计值是这些片的 \(A_0\) 长度之和。零维时单层多项式为零，并没有次数为 \(-1\) 的非零多项式。

\ProseParagraphBreak
例如对域 \(k\)，在 \(R=k[x,y]_{(x,y)}\)、\(I=(x,y)R\) 时，单层函数为 \(n+1\)，累计函数为 \(\binom{n+2}{2}\)。在 \(R=k[t]/(t^a)\)、\(I=(t)\) 时，\(a\ge1\) 为整数，单层长度在 \(0\le n<a\) 时为一，以后为零，累计长度为 \(\min(n+1,a)\)。前一个例子显示维数二的两种次数；后一个例子的单层多项式为零，重数却为总长度 \(a\)，不能从零多项式取一个首项系数来定义零维重数。

\subsection{有限点集上的求值码}
\label{b3:subsec:evaluation-codes}
% 保留迁移前的小节标签作为兼容引用。
\label{b3:subsec:1102:05}

对多项式的次数滤过，累计维数还可以衡量一个编码空间的大小。在 \(k=\mathbb F_q\) 上，把全次数不超过 \(d\) 的多项式在有限点集上的值排成向量，就得到码字；若这些向量张成的空间维数为 \(h\)，便有 \(q^h\) 个不同码字。第一卷附录B.3用一元多项式建立 Reed--Solomon 码\footnote{里德（Irving Stoy Reed，1923-2012），美国数学家及工程师，研究纠错码。所罗门（Gustave Solomon，1930-1996），美国数学家及工程师，与 Reed 合作提出代数纠错码。}；现在允许多个变量。仿射点的消失理想通常不齐次，例如 \(x_i^q-x_i\) 把不同全次数的项联系起来，所以商代数通常不能直接按原来的全次数分次。这里保留的是次数上界为 \(d\) 的整个滤过子空间，而不是标准分次意义下的单个次数 \(d\) 分量。

\begin{definition}\label{b3:def:evaluation-code-finite-set}
设 \(q\) 为素数幂，\(k=\mathbb F_q\)，\(X=\{P_1,\ldots,P_N\}\subseteq k^n\) 为非空有限点集，其中各 \(P_i\) 两两不同，\(N=|X|\)，\(d\ge0\) 为整数。记 \(S_{\le d}\subset k[x_1,\ldots,x_n]\) 为全次数不超过 \(d\) 的多项式空间。线性映射
\[
\operatorname{ev}_{X,d}:S_{\le d}\longrightarrow k^N,
\qquad f\longmapsto\bigl(f(P_1),\ldots,f(P_N)\bigr)
\]
的像称为 \(X\) 上次数界为 \(d\) 的\term[qiuzhi-ma]{求值码}[evaluation code]，记为 \(C_X(d)\)。
\end{definition}
\symidx{\(\operatorname{ev}_{X,d}\)：有限点集上次数至多 \(d\) 的求值映射}
\symidx{\(C_X(d)\)：有限点集 \(X\) 的求值码}

\begin{proposition}\label{b3:prop:evaluation-code-hilbert}
设 \(q\) 为素数幂，\(k=\mathbb F_q\)，\(S=k[x_1,\ldots,x_n]\)，\(X\subseteq k^n\) 是非空有限点集，\(I(X)\subseteq S\) 为其消失理想，\(d\ge0\) 为整数，\(C_X(d)\) 为全次数不超过 \(d\) 的多项式在 \(X\) 上的求值码。固定先比较全次数的全局单项式序，则 \(\dim_k C_X(d)\) 等于相对于 \(I(X)\) 的全次数不超过 \(d\) 的标准单项式数。该数也是商代数 \(S/I(X)\) 的全次数滤过第 \(d\) 项的维数。
\end{proposition}
\begin{proof}
求值映射的核正是 \(I(X)\cap S_{\le d}\)，所以像同构于 \(S_{\le d}/(I(X)\cap S_{\le d})\)。因所用次序先比较全次数，每次 Gröbner 除法都不会增大多项式的全次数。故该商由次数不超过 \(d\) 的标准单项式生成；它们的独立性仍由首项判据给出。这个商也同构于 \(S/I(X)\) 中的滤过子空间 \(F_d(S/I(X))=(S_{\le d}+I(X))/I(X)\)，得到最后的断言。
\end{proof}

对整个点集 \(X=\mathbb F_q^n\)，其消失理想为
\[
I(X)=(x_1^q-x_1,\ldots,x_n^q-x_n).
\]
右边的生成元在每一点消失，且其首单项式两两互素，构成 Gröbner 基。若一个余式在所有点上为零，则它关于每个变量的次数都小于 \(q\)。把其他变量固定后用一元根数界，再对变量数归纳，便知余式为零，这证明反向包含。因此求值码的维数为
\[
\#\{\,(a_1,\ldots,a_n)\mid 0\le a_i<q,\ a_1+\cdots+a_n\le d\,\}.
\]
多项式 \((1+t+\cdots+t^{q-1})^n\) 中 \(t^d\) 的系数统计 \(0\le a_i<q\)、\(\sum_i a_i=d\) 的单项式数。乘以 \(1/(1-t)=1+t+t^2+\cdots\) 后，每个系数就把次数不超过 \(d\) 的数目相加，因此
\[
\sum_{d\ge0}\dim_k C_X(d)\,t^d
=\frac{(1+t+\cdots+t^{q-1})^n}{1-t}.
\]
分子按恰好次数计数，分母使计数累计到次数上界；这是滤过的累计维数级数。这类码称为\term[Reed-Muller-ma]{Reed--Muller 码}[Reed--Muller code]；\footnote{穆勒（David Eugene Muller，1924-2008），美国数学家及计算机科学家，研究编码与自动机。}这里采用整个仿射空间和全次数上界的约定。

\ProseParagraphBreak
取 \(q=2,n=2,d=1\)，按 \((0,0),(1,0),(0,1),(1,1)\) 排列求值点。基 \(1,x,y\) 给出生成矩阵
\[
G=\begin{pmatrix}1&1&1&1\\0&1&0&1\\0&0&1&1\end{pmatrix}.
\]
其秩为三，故码有八个码字。非恒定仿射线性函数的零点恰为两点：非零线性函数 \(\mathbb F_2^2\to\mathbb F_2\) 满射，每个纤维有两点。常数一的重量为四，所以最小距离为二。这是一个 \([4,3,2]\) 码，能够检测一个错误，却不能唯一纠正任意一个错误。计算 Hilbert 数据在这里直接求出了可传递信息的维数，距离则还需要另作零点计数。


\begin{chaptersummary}
齐次生成元与齐次关系允许逐次取正合列。乘以一个齐次元所得的核与余核，使 Hilbert 级数按变量数递归，得到有理表达式；标准分次再使分母只含 \(1-t\)，从而产生最终多项式。正权分次虽然仍有有理级数，却可能出现随同余类变化的计数。同时变号不变环的奇数次分量为零，便提醒我们：给生成元重新指定次数，会改变所计算的 Hilbert 函数。在已有分次满射的前提下，逐次维数相等可以证明没有遗漏关系；只有级数相等而没有这种映射，并不能确定环结构。

\ProseParagraphBreak
齐次素理想滤过把增长与维数的比较归结为整环，Noether 正规化则用代数独立参数给出上下界。局部情形改用理想幂的伴随分次：素理想链和参数理想分别给出增长次数的两个方向，最终得到累计长度的次数等于模的维数。

\ProseParagraphBreak
分次重数依赖所选的标准分次，Hilbert–Samuel 重数依赖所选的极大理想准素理想。正维时，两者分别用 \((d-1)!\) 与 \(d!\) 对单层多项式和累计多项式的首项系数作标准归一化。相同维数、相同重数甚至相同 Hilbert 函数，都未必意味着同样的切锥或局部环。标准单项式使分次计数可以实际执行；有限点集的求值码又把滤过维数解释为可编码信息的维数。局部计算则要取理想中全部元素的初始形式，才能保留完整的切锥关系。\chapref{b3:ch:13}将讨论完备化如何保存这些逐层数据，并通过 Artin–Rees 等式比较模及其子模的滤过。
\end{chaptersummary}

\BookExercises

\begin{exercise}\label{b3:ex:11-weighted-count}
设 \(k\) 为域，给 \(k[x,y]\) 赋权 \(\deg x=1\)、\(\deg y=2\)。求 Hilbert 级数与函数，并证明函数最终不是一个多项式。
\end{exercise}
\begin{solution}
次数 \(n\) 的单项式为 \(x^{n-2j}y^j\)，其中 \(0\le j\le\lfloor n/2\rfloor\)。因此
\[
H(t)=\dfrac1{(1-t)(1-t^2)},\qquad h(n)=\lfloor n/2\rfloor+1.
\]
若多项式 \(P\) 最终等于它，则增长速度使 \(P\) 为一次多项式，斜率为 \(1/2\)。在全部充分大的偶数上有 \(P(n)=n/2+1\)，多项式恒等原理迫使 \(P(X)=X/2+1\)，而奇数上的实际值为 \((n+1)/2\)，矛盾。
\end{solution}

\begin{exercise}\label{b3:ex:11-homogeneous-scaling}
设 \(k\) 为无限域，\(I\subseteq k[x_1,\ldots,x_r]\) 为理想。证明 \(I\) 齐次，当且仅当对所有 \(a\in k^\times\) 及 \(f\in I\)，都有 \(f(ax_1,\ldots,ax_r)\in I\)。说明有限域上反向不成立。本题将\cite[\S13, Exercise 13.1, p. 102]{Matsumura1989CommutativeRingTheory}的伸缩判据用于多项式环，并补充有限域反例。
\end{exercise}
\begin{solution}
齐次理想包含每个 \(f\) 的齐次部分 \(f_d\)，所以包含 \(\sum_d a^df_d\)。反向对 \(f=\sum_{d=0}^D f_d\)，选 \(D+1\) 个不同非零标量 \(a_i\)。矩阵 \((a_i^d)_{i,d}\) 的 Vandermonde 行列式\footnote{范德蒙德（Alexandre-Théophile Vandermonde，1735-1796），法国数学家，研究行列式与方程消去。}非零，因而每个 \(f_d\) 是这些伸缩后的多项式的 \(k\) 线性组合，属于 \(I\)。在 \(\mathbb F_q[x]\) 中，\((x^q-x)\) 在所有非零标量伸缩下不变，因为 \(a^q=a\)；但它不包含齐次部分 \(-x\)，所以不是齐次理想。
\end{solution}

\begin{exercise}\label{b3:ex:11-graded-ideal-presentation}
设 \(k\) 为域，\(S=k[x,y]\) 取标准分次。给出 \(I=(x^2,xy)\) 的分次自由展示，计算 \(H_I(t)\) 与 \(H_{S/I}(t)\)。
\end{exercise}
\begin{solution}
关系 \(ax^2+bxy=0\) 可在整环中约去 \(x\)，得到 \(ax+by=0\)。模掉 \(y\) 说明 \(a=yc\)，再得 \(b=-xc\)。因此有正合列
\[
0\to S(-3)\xrightarrow{c\mapsto(yc,-xc)}S(-2)^2\to I\to0.
\]
两个生成元次数为二，关系的全次数为三，故平移如式所示。由\eqref{b3:eq:hilbert-shift}，
\[
H_I(t)=\dfrac{2t^2-t^3}{(1-t)^2},\qquad
H_{S/I}(t)=\dfrac{1-2t^2+t^3}{(1-t)^2}=\dfrac1{1-t}+t.
\]
\end{solution}

\begin{exercise}\label{b3:ex:11-length-not-cardinality}
设 \(p\) 为素数，\(a,r\ge1\) 为整数，令 \(A_0=\mathbb Z/p^a\mathbb Z\)，\(S=A_0[x_1,\ldots,x_r]\) 取标准分次。求 \(S_n\) 的 \(A_0\) 长度，并解释为何不能用它的元素个数代替长度。
\end{exercise}
\begin{solution}
\(A_0\) 的合成列为
\[
0=p^aA_0\subsetneq p^{a-1}A_0\subsetneq\cdots\subsetneq p^0A_0=A_0.
\]
每个因子 \(p^jA_0/p^{j+1}A_0\)（\(0\le j<a\)）均同构于 \(\mathbb F_p\)，故长度为 \(a\)；\(a=1\) 时这条列就是 \(0\subsetneq A_0\)。单项式给出秩 \(\binom{n+r-1}{r-1}\) 的自由基，所以
\[
\ell_{A_0}(S_n)=a\binom{n+r-1}{r-1},\qquad H_S(t)=\dfrac a{(1-t)^r}.
\]
元素个数则为 \(p^{a\binom{n+r-1}{r-1}}\)，短正合列中它相乘而不相加。Hilbert 级数的正合列公式依靠长度可加，故不能这样替换。
\end{solution}

\begin{exercise}\label{b3:ex:11-zero-divisor-series}
设 \(k\) 为域，令 \(A=k[x,y]/(xy)\)，取标准分次。计算 \((0:_A x)\)、\(A/xA\)，并用\eqref{b3:eq:hilbert-induction}求 \(H_A(t)\)。
\end{exercise}
\begin{solution}
\((0:_A x)=(y)\cong k[y](-1)\)，因为任一类可唯一写成 \(c+x\cdot p(x)+y\cdot q(y)\)，乘 \(x\) 为零恰要求前两部分为零。又 \(A/xA\cong k[y]\)。所以
\[
(1-t)H_A(t)=\dfrac1{1-t}-\dfrac{t^2}{1-t},
\qquad H_A(t)=\dfrac{1+t}{1-t}.
\]
若错误地把乘 \(x\) 当作单射，就会漏掉核对应的 \(t^2/(1-t)\)。
\end{solution}

\begin{exercise}\label{b3:ex:11-veronese}
设 \(k\) 为域，\(B=k[x^3,x^2y,xy^2,y^3]\subseteq k[x,y]\)。先保留 \(x,y\) 的通常全次数，再把四个生成元统一改记为次数一。分别计算两种分次的 Hilbert 函数与级数，并在后一分次下求维数和重数。说明前一种 Hilbert 函数为何最终不是单个多项式。
\end{exercise}
\begin{solution}
任意全次数 \(3m\) 的单项式都可拆成 \(m\) 个三次单项式的乘积：若其 \(x\) 指数至少三，取出 \(x^3\)；若 \(x\) 指数为二或一而总次数至少三，分别取出 \(x^2y\) 或 \(xy^2\)；若 \(x\) 指数为零，取出 \(y^3\)。对 \(m\) 归纳即可。故原分次下的 \(B_{3m}\) 等于 \(k[x,y]_{3m}\)，其维数为 \(3m+1\)，其他次数的分量为零。因此
\[
H_B(t)=\sum_{m\ge0}(3m+1)t^{3m}
=\dfrac{1+2t^3}{(1-t^3)^2}.
\]
重新分次得到 \(\widetilde B_m=B_{3m}\)，所以
\[
h_{\widetilde B}(m)=3m+1,\qquad
H_{\widetilde B}(z)=\dfrac{1+2z}{(1-z)^2}.
\]
\(\widetilde B\) 是由四个一次元素生成的标准分次 \(k\) 代数；由\cref{b3:thm:hilbert-polynomial-dimension}，其维数为二，重数为三。两种分次的底层环相同，环的维数也相同；重数三在这里针对 \(\widetilde B\) 的标准分次而言。若原 Hilbert 函数最终等于一个有理系数多项式，该多项式在全部充分大的 \(3m+1\) 处为零，必为零多项式，却在充分大的 \(3m\) 处应取正值 \(3m+1\)，矛盾。
\end{solution}

\begin{exercise}\label{b3:ex:11-complete-intersection-count}
设 \(k\) 为域，\(a,b\ge1\) 为整数。求标准分次环 \(k[x,y,z]/(x^a,y^b)\) 的级数、最终 Hilbert 多项式与重数。
\end{exercise}
\begin{solution}
标准单项式为 \(x^iy^jz^l\)，其中 \(0\le i<a\)、\(0\le j<b\)、\(l\ge0\)。因此
\[
H(t)=\dfrac{(1+t+\cdots+t^{a-1})(1+t+\cdots+t^{b-1})}{1-t}.
\]
当 \(n\ge a+b-2\) 时，每个 \((i,j)\) 都唯一对应 \(l=n-i-j\ge0\)，所以 \(h(n)=ab\)。维数为一，重数为 \(ab\)。这个直接计数也核实了两个非零因子商的级数公式：\(x^a\) 在多项式环中为非零因子，\(y^b\) 在模掉 \(x^a\) 后仍是非零因子。
\end{solution}

\begin{exercise}\label{b3:ex:11-polynomial-forgets-low-degrees}
设 \(k\) 为域，两个商环 \(k[x,y]/(x)\) 与 \(k[x,y]/(x^2,xy)\) 均取标准分次。比较它们的 Hilbert 函数、Hilbert 多项式和累计函数。它们的最终单层多项式是否确定全部累计多项式？
\end{exercise}
\begin{solution}
前者的级数为 \(1/(1-t)\)，后者为 \(1/(1-t)+t\)。前者各次函数为一，后者只有次数一处额外加一。因此二者 Hilbert 多项式都为一；累计函数分别为 \(n+1\) 与 \(n+2\)（后一式对 \(n\ge1\)）。低次分量在累计函数中留下常数差，所以单层最终多项式不确定累计多项式的常数项。
\end{solution}

\begin{exercise}\label{b3:ex:11-zero-dimensional-shift}
设 \(k\) 为域，\(k[x]\) 取标准分次，令 \(M=\bigl(k[x]/(x^a)\bigr)(-b)\)，其中 \(a\ge1\) 为整数、\(b\in\mathbb Z\)。求非零分次片、单层最终多项式及从全部低次开始累计的最终值。
\end{exercise}
\begin{solution}
\(M_n\) 在 \(b\le n\le b+a-1\) 时为一维，其余为零。因此
\(H_M(t)=t^b(1+t+\cdots+t^{a-1})\)。单层最终多项式为零，\(\sum_{j\le n}\dim_kM_j\) 最终等于 \(a\)，且零维重数为 \(a\)。即使 \(b<0\)，也须计入负次的非零片，不能不加说明地从零次开始求和。
\end{solution}

\begin{exercise}\label{b3:ex:11-dvr-adic}
设 \(k\) 为域，\(R=k[t]_{(t)}\)，令 \(I=(t^a)\)，其中 \(a\ge1\) 为整数。构造分次同构 \(\operatorname{gr}_I(R)\cong(R/I)[T]\)，并求 \(\chi_{I,R}(n)\)。
\end{exercise}
\begin{solution}
把 \(T\) 送到 \(t^a+I^2\)。在次数 \(n\)，映射为
\(R/I\to I^n/I^{n+1}\)、\(\bar r\mapsto rt^{an}+I^{n+1}\)。它满射；若像为零，则 \(rt^{an}\in\left(t^{a(n+1)}\right)\)，整环中约去 \(t^{an}\)，得到 \(r\in(t^a)\)，所以单射。各次映射保持乘法，得到同构。每片长度为 \(a\)，故 \(\chi_{I,R}(n)=a(n+1)\)，重数为 \(a\)。
\end{solution}

\begin{exercise}\label{b3:ex:11-three-parameters}
设 \(k\) 为域，令 \(R=k[x,y,z]_{(x,y,z)}\)、\(I=(x^a,y^b,z^c)\)，其中 \(a,b,c\ge1\) 为整数。求 \(R/I^{n+1}\) 的长度及 \(e_I(R)\)。
\end{exercise}
\begin{solution}
每个单项式唯一写成 \(x^{au+i}y^{bv+j}z^{cw+l}\)，其中 \(i<a,j<b,l<c\) 且均非负。它不属于 \(I^{n+1}\) 当且仅当 \(u+v+w\le n\)。这样的三元组有 \(\binom{n+3}{3}\) 个，故长度为 \(abc\binom{n+3}{3}\)，重数为 \(abc\)。商中极大理想幂零，常数项非零的多项式已经可逆，所以局部化不改变该单项式基。
\end{solution}

\begin{exercise}\label{b3:ex:11-higher-cusp}
设 \(k\) 为任意特征的域，\(R=k[x,y]_{(x,y)}/(y^3-x^5)\)，\(\mathfrak m=(\bar x,\bar y)\) 为其极大理想。求 \(\mathfrak m\) 进切锥、Hilbert 级数与累计长度。
\end{exercise}
\begin{solution}
最低齐次部分为 \(y^3\)，故\cref{b3:prop:hypersurface-tangent-cone}给出
\(\operatorname{gr}_{\mathfrak m}R\cong k[X,Y]/(Y^3)\)。级数为 \((1+t+t^2)/(1-t)\)，单层长度依次为 \(1,2,3,3,\ldots\)，累计长度在 \(n=0\) 时为一，在 \(n\ge1\) 时为 \(3n\)。维数为一，重数为三。这个计算只用最低项与整环中乘积的最低项相乘，不使用导数，因而不需排除特征三或五。
\end{solution}

\begin{exercise}\label{b3:ex:11-extra-initial-relation}
设 \(k\) 为任意特征的域，令 \(J=(x^2-y^3,xy-y^2)\subseteq k[x,y]\)、\(A=k[x,y]/J\)，以及 \(R=A_{(\bar x,\bar y)}\)。求 \(R\) 的切锥坐标环、Hilbert 级数与长度，并比较局部化前的 \(\dim_k A\)。说明为什么在局部化以后，\(1-y\) 的可逆性会产生新的关系。
\end{exercise}
\begin{solution}
记 \(f=x^2-y^3\)、\(g=xy-y^2\)。计算
\[
yf-xg=xy^2-y^4,
\qquad yf-xg-yg=y^3-y^4
\]
说明 \(y^3(1-y)\in J\)。在原点局部环中，\(1-y\) 为单位，故 \(y^3\in JR_0\)，其中 \(R_0=k[x,y]_{(x,y)}\)。随即 \(x^2\in JR_0\)，所以
\[
JR_0=(x^2,xy-y^2,y^3)R_0.
\]
右边三个多项式是齐次的。按字典序 \(x\succ y\)，其首单项式为 \(x^2,xy,y^3\)；三个 S 多项式分别约化为零，因为第一对给出 \(xy^2\)，可用 \(xy-y^2\) 将它变为 \(y^3\)，其余两对给出零或 \(-y^4\)。标准单项式为 \(1,x,y,y^2\)，故齐次商有这四元基。其正次数理想的三次幂为零，局部化不改变它。由\cref{b3:prop:quotient-local-initial-ideal}得到
\[
\operatorname{gr}_{\mathfrak m}R\cong k[X,Y]/(X^2,XY-Y^2,Y^3),
\qquad H(t)=1+2t+t^2,
\]
其中 \(\mathfrak m\) 为 \(R\) 的极大理想，且 \(\ell_R(R)=4\)。原来两个生成元的最低项只有 \(X^2,XY-Y^2\)，尚须补上 \(Y^3\)。

局部化前，\(f,g,h=y^4-y^3\) 构成同一次序下的 Gröbner 基。第一对的 S 多项式经 \(g\) 约化为 \(-h\)；另两对分别为
\[
y^4f-x^2h=x^2y^3-y^7,
\qquad y^3g-xh=xy^3-y^5.
\]
它们经 \(f\)、\(g\) 分别约化为 \(-y^3h\)、\(-yh\)，再用 \(h\) 得到零。首单项式为 \(x^2,xy,y^4\)，所以 \(1,x,y,y^2,y^3\) 的类构成 \(A\) 的基，\(\dim_k A=5\)。在 \(A\) 中 \(y^3\ne0\)，因而 \(y^3(1-y)=0\) 说明 \(1-y\) 尚不可逆；点 \((1,1)\) 也满足两条原方程。局部化到原点后，\(1-y\) 成为单位，\(y^3\) 才消失，向量空间维数随之降为四。
\end{solution}

\begin{exercise}\label{b3:ex:11-power-multiplicity}
设 \((R,\mathfrak m)\) 为 Noether 局部环，\(M\) 为非零有限生成 \(R\) 模，\(I\subset R\) 为 \(\mathfrak m\) 准素理想，\(d=\dim_RM\)。证明 \(e_{I^a}(M)=a^d\cdot e_I(M)\)，其中 \(a\ge1\) 为整数。
\end{exercise}
\begin{solution}
由定义
\[
\chi_{I^a,M}(n)=\ell_R\left(M/I^{a(n+1)}M\right)
=\chi_{I,M}\bigl(a(n+1)-1\bigr).
\]
把\cref{b3:thm:hilbert-samuel}的首项 \(e_I(M)\cdot n^d/d!\) 代入，变量的一次替换使首项系数乘 \(a^d\)，即得结论。\(d=0\) 时两边均为 \(\ell_R(M)\)，与 \(a^0=1\) 一致。
\end{solution}

\begin{exercise}\label{b3:ex:11-inclusion-multiplicity}
设 \((R,\mathfrak m)\) 为 Noether 局部环，\(I\subseteq J\subset R\) 为两个 \(\mathfrak m\) 准素理想，\(M\) 为非零有限生成 \(R\) 模。证明 \(e_I(M)\ge e_J(M)\)，并举严格不等的例子。
\end{exercise}
\begin{solution}
\(I^{n+1}M\subseteq J^{n+1}M\) 使前一个商满射到后一个商，故 \(\chi_{I,M}(n)\ge\chi_{J,M}(n)\)。两者由\cref{b3:thm:hilbert-samuel}具有相同次数 \(d=\dim_RM\)，比较首项得到不等式。在 \(k[t]_{(t)}\) 上取 \(M=R\)、\(I=(t^2)\)、\(J=(t)\)，两重数分别为二和一。
\end{solution}

\begin{exercise}\label{b3:ex:11-filtration-shift}
设 \(k\) 为域，取 \(R=k[t]_{(t)}\)、\(I=(t)\)、\(M=R\)、\(N=t^aR\)，其中 \(a\ge1\) 为整数。求 \(I^nM\cap N\)，并求使\cref{b3:lem:induced-filtration-bound}对全部 \(n\ge c\) 成立的最小 \(c\)。
\end{exercise}
\begin{solution}
两个主理想的交为 \(t^{\max(n,a)}R\)，而 \(I^{n-c}N=t^{n-c+a}R\)。取 \(c=a\) 时，对 \(n\ge a\) 两者相等。若 \(c<a\)，取 \(n\ge a\)，则要求 \(t^nR\subseteq t^{n-c+a}R\)，但右边指数严格大于 \(n\)，不成立。因此最小移位为 \(a\)。这也直接展示了诱导滤过与子模自身的理想进滤过通常不逐层相同。
\end{solution}


\begin{exercise}\label{b3:ex:11-monomial-hilbert-count}
设 \(k\) 为域，在标准分次 \(S=k[x,y,z]\) 中令 \(J=(x^2,xy,y^2)\)。从标准单项式求 \(H_{S/J}(t)\)、Hilbert 多项式、维数与重数；再比较 \(S/(x^2,xy,y^3)\)。
\end{exercise}
\begin{solution}
第一商的标准单项式为 \(z^j,xz^j,yz^j\)，故级数为 \((1+2t)/(1-t)\)，Hilbert 函数在零次为一、从一次起为三。多项式为常数三，维数一，重数三。第二商多出全部 \(y^2z^j\)，级数为 \((1+2t+t^2)/(1-t)\)，函数依次为一、三、四、四，维数仍为一而重数为四。增加的这组单项式改变最终增长系数，不能只看两个理想有同样的根理想 \((x,y)\)。
\end{solution}

\begin{exercise}\label{b3:ex:11-reed-muller-plane}
在 \(\mathbb F_3^2\) 的全部九点上求值次数不超过一的多项式。求所得线性码的长度、维数与最小距离，并证明次数不超过四时求值能得到任意九元组。
\end{exercise}
\begin{solution}
一次多项式空间的基为 \(1,x,y\)，而 \(I(\mathbb F_3^2)=(x^3-x,y^3-y)\)，因此求值在该空间上单射，长度九、维数三。非恒定仿射线性函数的零点是一条含三个点的仿射直线，重量六；非零常数的重量九，故最小距离六。标准单项式为 \(x^ay^b\)、\(0\le a,b\le2\)，全部全次数不超过四，共九个，其求值线性无关，故在次数四时像已是 \(\mathbb F_3^9\)。此时可以传递任意九个值，但最小距离降为一，已无冗余用于检错。
\end{solution}

\begin{exercise}\label{b3:ex:11-same-series-different-rings}
设 \(k\) 为域，给 \(R=k[a,b]/(a^2)\)、\(T=k[a,b]/(ab)\) 取标准分次。证明它们的 Hilbert 级数相同，但作为 \(k\) 代数不同构。由此说明\cref{b3:prop:hilbert-surjection-isomorphism}为什么要先给定一个保持次数的满射。
\end{exercise}
\begin{solution}
两条关系 \(a^2,ab\) 都是多项式整环 \(k[a,b]\) 中次数二的非零元。分别使用乘以该元的分次短正合列，得到
\[
H_R(t)=H_T(t)=\dfrac{1-t^2}{(1-t)^2}
=\dfrac{1+t}{1-t}.
\]
\(R\) 中 \(\bar a\ne0\)，因为 \(a\notin(a^2)\)，但 \(\bar a^2=0\)，所以 \(R\) 不约化。另一方面，\((ab)=(a)\cap(b)\)：若 \(af\in(b)\)，模掉 \(b\) 后在整环 \(k[a]\) 中得到 \(a\bar f=0\)，故 \(f\in(b)\)，证明非平凡包含。若 \(h^m\in(ab)\)，则素理想 \((a)\)、\((b)\) 都包含 \(h\)，故 \(h\in(a)\cap(b)=(ab)\)。因此 \(T\) 约化。是否含有非零幂零元由环同构保持，所以两个环不同构。

它们在每个次数只有相同的向量空间维数，这些数不能记录乘法如何区别 \(a^2=0\) 与 \(ab=0\)。若存在保持次数的满射 \(R\to T\) 或 \(T\to R\)，由\cref{b3:prop:hilbert-surjection-isomorphism}就会成为同构，与上述性质矛盾。因此这道题并不违反该命题，反而说明其中映射假设的作用。
\end{solution}

% !TeX root = ../../Book3.tex
% 纲要标题：## 第12章　平坦性、基变换与下降
% 纲要位置：工作记录/计划纲要.md，第 2163 行。


\chapter{平坦性、基变换与下降}
\label{b3:ch:12}
\BookChapterNumber{b3:ch:12}

\begin{chapterabstract}
平坦性控制张量积对关系的影响。本章从理想张量判据出发，证明平坦性的局部性和有限展示平坦模的局部自由性。Fitting 理想将展示矩阵的子式变成内在不变量，用来描述纤维生成元数及局部自由条件；泛自由性则说明整基环上的有限代数数据可在一个稠密开集上同时简化。最后建立忠实平坦基变换的检测作用，并证明带余循环数据的模如何下降。
\end{chapterabstract}

\begin{learninggoals}
\begin{itemize}
\item 用有限生成理想检测平坦性，辨别无挠、平坦及局部自由的区别。
\item 由有限展示计算 Fitting 理想，解释纤维维数所能和不能检测的性质。
\item 用有限组首项和公共分母构造泛自由性定理中的主开集及自由基。
\item 用剩余纤维判断忠实平坦性，并将变基后的正合性、有限性推回原环。
\item 从双重基变换与余循环条件构造模下降的等化子。
\end{itemize}
\end{learninggoals}

预备知识是第二卷的模、有限展示及张量积，以及本卷前两章的局部化和第三章的 Nakayama 引理。平坦性的定义在此重述，所需判据均在本章证明，不以 \(\operatorname{Tor}\) 或后续同调理论为前提。\secref{b3:sec:1206}讨论对象的拼接；第十三章完备化的平坦性只调用\secref{b3:sec:1201}和\secref{b3:sec:1205}的结果。

\ProseParagraphBreak
理想判据、局部环上的有限平坦模及忠实平坦性可对照 Matsumura 的《Commutative Ring Theory》\cite[\S7, Theorems 7.2--7.10, pp. 47--52]{Matsumura1989CommutativeRingTheory}；泛自由性的另一种证明及更一般形式见\cite[\S24, Theorem 24.1, p. 186]{Matsumura1989CommutativeRingTheory}。Fitting 判据用的是理想等式本身，仅保留它们定义的闭点集会丢失必要的信息。

% !TeX root = ../../Book3.tex
% 纲要标题：### 12.1 平坦性的判据
% 纲要位置：工作记录/计划纲要.md，第 2314 行。
% 纲要细目（写作范围）：
% * 理想张量判据
% * 平坦性的局部性
% * 局部化平坦及一般无挠模反例
% * 复用第二卷第7.3节已经证明的理想张量判据，在本节建立其局部检测形式；随后供第13.4节的完备化应用调用

\section{平坦性的判据}\label{b3:sec:1201}

局部化保持正合性，但一般的标量扩张没有这个保证。平坦性正是要求张量积不把单射变坏。判断时无须逐一检查所有模之间的单射：只检查有限生成理想进入原环的包含就够了。这个一般判据的主要证明见第二卷定理7.3.9；本节经自由模子模归纳给出另一证明，再建立素理想处的局部判据。

\subsection{理想张量判据}\label{b3:subsec:1201:01}

\begin{definition}\label{b3:def:flat-module}
设 \(R\) 为交换环，\(M\) 为 \(R\) 模。如果对任意 \(R\) 模之间的单射 \(U\to V\)，诱导映射 \(U\otimes_RM\to V\otimes_RM\) 仍单射，就称 \(M\) 为\term[pingtan-mo]{平坦模}[flat module]。若交换环同态 \(R\to S\) 使 \(S\) 成为平坦 \(R\) 模，就称它为平坦同态。
\end{definition}

对短正合列 \(0\to U\to V\to V/U\to0\)，张量积总是给出右正合列
\[
U\otimes_RM\longrightarrow V\otimes_RM
\longrightarrow(V/U)\otimes_RM\longrightarrow0.
\]
确实，将 \(V/U\) 的生成元和线性关系代入张量积的定义，得到
\((V/U)\otimes_RM\cong(V\otimes_RM)/\operatorname{im}(U\otimes_RM)\)。
唯一未得到保证的是第一个映射仍为单射；平坦性正是排除这里出现新的核。因此此定义等价于张量函子保持短正合列。自由模平坦，因为与自由模张量就是取若干份直和；平坦模的直和及直和项也平坦，可逐分量检查核。

\begin{theorem}[理想张量判据]\label{b3:thm:ideal-flatness-criterion}
设 \(R\) 为交换环，\(M\) 为 \(R\) 模。则 \(M\) 平坦当且仅当对每个有限生成理想 \(I\subseteq R\)，乘法映射
\[
I\otimes_RM\longrightarrow M,\qquad a\otimes m\longmapsto am
\]
单射。此时该映射把 \(I\otimes_RM\) 识别为 \(IM\)。
\end{theorem}
\begin{proof}
必要性直接应用定义于 \(I\subseteq R\)。为证充分性，先把假设推广到任意理想 \(I\)。即使 \(I\) 不有限生成，一个具体张量按定义仍只能涉及有限多个理想元素。因而核中的张量 \(z=\sum_{i=1}^r a_i\otimes m_i\) 来自有限生成子理想 \(I_0=(a_1,\ldots,a_r)\) 的张量积 \(I_0\otimes_RM\) 中的同一个有限和。该有限和映到 \(M\) 中为零，由假设它在 \(I_0\otimes_RM\) 中已经为零，故 \(z=0\)。

理想是秩一自由模的子模；把秩增加一时，可以用最后一个坐标投影将问题再分出一个理想。由此对 \(n\) 归纳，证明对任意子模 \(N\subseteq R^n\)，映射 \(N\otimes_RM\to M^n\) 单射。\(n=1\) 是刚证的理想情形。投影在 \(N\) 上的像为理想 \(I\)，核为 \(N'\subseteq R^{n-1}\)，所以 \(0\to N'\to N\to I\to0\) 正合。张量后的右正合列
\[
N'\otimes_RM\longrightarrow N\otimes_RM\longrightarrow I\otimes_RM\longrightarrow0
\]
说明：若 \(z\in N\otimes_RM\) 映到 \(M^n\) 为零，则其在 \(I\otimes_RM\) 中的像因理想判据为零，故来自某个 \(z'\in N'\otimes_RM\)。这个 \(z'\) 映到 \(M^{n-1}\) 为零，归纳假设给出 \(z'=0\)，于是 \(z=0\)。无限自由模的情形也成立，因为一个核元素只涉及有限多个 \(N\) 元素，它们的坐标支集包含于同一个有限自由直和项。

最后给定任意包含 \(U\subseteq V\)，它未必带有自由坐标，但可以沿自由模满射拉回到已解决的情形。取满射 \(F\to V\)，其中 \(F\) 自由，记核为 \(K\)，并令 \(G\subseteq F\) 为 \(U\) 的原像。于是 \(K\subseteq G\subseteq F\)，且 \(G/K\cong U\)。若 \(z\in U\otimes_RM\) 映到 \(V\otimes_RM\) 为零，先由右正合性将其提升为 \(w\in G\otimes_RM\)。它在 \(F\otimes_RM\) 中的像来自 \(K\otimes_RM\)，所以减去这个来源在 \(G\otimes_RM\) 中的像后，可使 \(w\) 在 \(F\otimes_RM\) 中的像为零。减去的部分映到 \(U\otimes_RM\) 为零，故仍是原来 \(z\) 的提升。刚证的自由模子模情形给出 \(w=0\)，原来的 \(z\) 因而为零。故任意单射都被保持。
\end{proof}

这个证明只用张量积的右正合性、自由模及有限支集，不需要预先引入 \(\operatorname{Tor}\)。同一判据见 Matsumura 的《Commutative Ring Theory》\cite[\S7, Theorem 7.7, p. 50]{Matsumura1989CommutativeRingTheory}。

给定底环系数 \(a_1,\ldots,a_n\)，先在 \(R^n\) 中考察满足 \(\sum_j a_jb_j=0\) 的系数向量。下面的引理说明，若 \(M\) 平坦，\(M\) 中满足同一关系的元组都可用这些底环系数向量作有限组合得到；关系没有因换到 \(M\) 中而多出无法这样表达的解。

\begin{lemma}[平坦模中的有限关系]\label{b3:lem:flat-relations}
设 \(R\) 为交换环，\(M\) 为平坦 \(R\) 模。若 \(a_j\in R\)、\(m_j\in M\)、\(1\le j\le n\)，满足 \(\sum_{j=1}^n a_jm_j=0\)，则存在有限组 \(y_l\in M\)、\(b_{jl}\in R\)，使
\[
m_j=\sum_l b_{jl}y_l,\qquad \sum_j a_jb_{jl}=0\quad\text{对每个}\;l.
\]
\end{lemma}
\begin{proof}
令 \(\varphi:R^n\to R\) 将第 \(j\) 个基向量送到 \(a_j\)，核为 \(K\)。平坦性使
\(K\otimes_RM\to M^n\to M\) 正合，所以元组 \((m_j)_j\) 是有限和 \(\sum_l v_l\otimes y_l\) 的像。写 \(v_l=(b_{jl})_j\in K\)，便同时得到两组等式。
\end{proof}

\subsection{平坦性的局部性}\label{b3:subsec:1201:02}

\begin{theorem}\label{b3:thm:flatness-local}
设 \(R\) 为交换环，\(M\) 为 \(R\) 模。下列条件等价：\(M\) 在 \(R\) 上平坦；对每个素理想 \(\mathfrak p\subseteq R\)，\(M_{\mathfrak p}\) 在 \(R_{\mathfrak p}\) 上平坦；对每个极大理想 \(\mathfrak m\subseteq R\)，\(M_{\mathfrak m}\) 在 \(R_{\mathfrak m}\) 上平坦。
\end{theorem}
\begin{proof}
若 \(M\) 平坦，将 \(R_{\mathfrak p}\) 模之间的单射 \(U\to V\) 看作 \(R\) 模单射，与 \(M\) 在 \(R\) 上张量后仍单射。因为 \(R\setminus\mathfrak p\) 中的元素已在 \(U\) 上可逆，有自然同构
\[
U\otimes_{R_{\mathfrak p}}M_{\mathfrak p}\cong U\otimes_RM,
\qquad u\otimes(m/s)\longmapsto s^{-1}u\otimes m.
\]
其逆把 \(u\otimes m\) 送到 \(u\otimes(m/1)\)，平衡关系保证两种公式相容。对 \(V\) 也作同样识别，得到局部平坦性。

反向，任取包含 \(U\subseteq V\)，设 \(K\) 为 \(U\otimes_RM\to V\otimes_RM\) 的核。由\cref{b3:thm:localization-exact}及张量积局部化，对每个极大理想有
\[
K_{\mathfrak m}
=\ker\left(U_{\mathfrak m}\otimes_{R_{\mathfrak m}}M_{\mathfrak m}
 \longrightarrow V_{\mathfrak m}\otimes_{R_{\mathfrak m}}M_{\mathfrak m}\right)=0.
\]
\cref{b3:thm:local-zero-detection}于是给出 \(K=0\)，即 \(M\) 平坦。
\end{proof}

\begin{proposition}\label{b3:prop:flat-base-change-transitivity}
设 \(R\to S\) 为交换环同态，\(M\) 为 \(R\) 模，\(N\) 为 \(S\) 模。若 \(M\) 在 \(R\) 上平坦，则 \(S\otimes_RM\) 在 \(S\) 上平坦。若 \(S\) 在 \(R\) 上平坦、\(N\) 在 \(S\) 上平坦，则 \(N\) 在 \(R\) 上平坦。
\end{proposition}
\begin{proof}
第一项对 \(S\) 模包含 \(U\subseteq V\) 使用
\(U\otimes_S(S\otimes_RM)\cong U\otimes_RM\)；右侧的单射由 \(M\) 平坦保证。第二项先与 \(S\) 张量，再与 \(N\) 在 \(S\) 上张量，每一步都保持单射，复合由结合律等于与 \(N\) 在 \(R\) 上张量。
\end{proof}

\subsection{局部化的平坦性与无挠模反例}\label{b3:subsec:1201:03}

局部化 \(S^{-1}R\) 平坦，因为
\(M\otimes_RS^{-1}R\cong S^{-1}M\)，而\cref{b3:thm:localization-exact}已证明右侧函子正合。整环上的平坦模必无挠：对非零 \(a\in R\)，映射 \(R\xrightarrow{\cdot a}R\) 单射，张量后得到乘 \(a\) 在 \(M\) 上单射。

\ProseParagraphBreak
反向通常不成立。设 \(k\) 为域，令 \(R=k[x,y]\)、\(M=I=(x,y)\)，它是整环的子模，因而无挠。乘法使两个生成元满足 \(xy-yx=0\)，因此在 \(I\otimes_RI\) 中将两个乘积分别保留为纯张量，就得到一个可能落入乘法核的元素
\[
z=x\otimes y-y\otimes x\in I\otimes_R I
\]
它在乘法映射下为零；要检验平坦性，关键是这个张量是否已经为零。模掉 \(\mathfrak m=(x,y)\) 后，它的像落在
\((I/\mathfrak mI)\otimes_k(I/\mathfrak mI)\)，其中 \(\bar x,\bar y\) 是基，所列两个不同基张量之差非零，任意特征都如此。因此\cref{b3:thm:ideal-flatness-criterion}说明 \(I\) 不平坦。

\subsection{理想张量判据的局部检测形式}\label{b3:subsec:1201:04}

\begin{corollary}\label{b3:cor:local-ideal-flatness}
设 \(R\) 为交换环，\(M\) 为 \(R\) 模。为证明 \(M\) 平坦，只需对每个有限生成理想 \(I\subseteq R\) 及每个极大理想 \(\mathfrak m\subseteq R\)，证明
\[
I_{\mathfrak m}\otimes_{R_{\mathfrak m}}M_{\mathfrak m}
 \longrightarrow M_{\mathfrak m}
\]
单射。在 Noether 环上，全部理想都有限生成。
\end{corollary}
\begin{proof}
该映射正是 \(I\otimes_RM\to M\) 在 \(\mathfrak m\) 处的局部化。由\cref{b3:thm:local-morphism-detection}，所有这些局部映射单射便使原映射单射，再用\cref{b3:thm:ideal-flatness-criterion}即可。
\end{proof}

这里的检验仍在局部环上进行，不能换成只看剩余域中的维数。设 \(k\) 为域，取 \(R=k[\varepsilon]/(\varepsilon^2)\)、\(\mathfrak m=(\varepsilon)\)、\(M=R/(\varepsilon)\)。环 \(R\) 已是局部环且只有这一个素点，所以茎 \(M_{\mathfrak m}=M\) 仍作为 \(R_{\mathfrak m}=R\) 模保留 \(\varepsilon M=0\) 这条关系。纤维 \(M\otimes_R k\cong M/\mathfrak mM\cong k\) 则是一维 \(k\) 向量空间。

在局部环上，\((\varepsilon)\otimes_RM\cong(\varepsilon)/(\varepsilon)^2\cong k\) 非零，但乘法映射 \((\varepsilon)\otimes_RM\to M\) 将每个张量都送到零，因为 \(\varepsilon M=0\)。因此 \(M\) 不平坦，虽然它的纤维有一维；纤维维数没有记录这条包含映射张量后的非零核。\secref{b3:sec:1203}将用 Fitting 理想记录这个差别。

% !TeX root = ../../Book3.tex
% 纲要标题：### 12.2 有限性与局部自由
% 纲要位置：工作记录/计划纲要.md，第 2299 行。
% 纲要细目（写作范围）：
% * 有限展示平坦模
% * 局部环上的有限自由性
% * 有限生成投射模与局部自由模

\section{有限性与局部自由}\label{b3:sec:1202}

平坦性不限制模的大小，例如无穷秩自由模仍平坦。有限展示则同时限制生成元和关系，使一点处的自由性可以扩大到一个邻域，再拼成有限投射性。本节先在局部环中证明自由性，随后解释有限展示为何用于从局部返回整体。

\subsection{有限展示平坦模}\label{b3:subsec:1202:01}

设 \(R\) 为交换环，\(M\) 为 \(R\) 模。如果存在有限非负整数 \(a,b\) 及正合列
\(R^a\to R^b\to M\to0\)，就称 \(M\) 有限展示。如果 \(M\) 同构于某个有限秩自由模的直和项，就称它为有限生成投射模。如果每个素理想都有一个主开邻域 \(D(f)\)，使 \(M_f\cong R_f^r\)、\(r\) 为有限非负整数，就称 \(M\) 局部自由且秩有限；\(r\) 可随邻域改变。

在这些定义下，有限性与平坦性给出以下等价关系：
\[
\boxed{\text{有限展示且平坦}}
\Longleftrightarrow\boxed{\text{有限生成投射}}
\Longleftrightarrow\boxed{\text{有限秩局部自由}}.
\]
最后一项指邻域上的自由性。有限展示使一个局部环上的有限基能够扩展到这样的邻域。

\begin{theorem}\label{b3:thm:finite-presentation-flat-projective}
设 \(R\) 为交换环。下列三类 \(R\) 模相同：
\begin{equivalences}
\item 有限展示平坦模；
\item 有限生成投射模；
\item 有限展示且在每个素理想处自由的模。
\end{equivalences}
它们也恰为局部自由且秩有限的模；这里的局部自由要求在邻域上成立，不能仅把“有限展示”从第三项删去。
\end{theorem}

证明将在下面两小节完成。有限展示的用处不是使张量积右正合，而是使有限个关系可统一处理，使 Hom 与局部化相容。

\subsection{局部环上的有限自由性}\label{b3:subsec:1202:02}

\begin{lemma}\label{b3:lem:finite-flat-local-free}
设 \((R,\mathfrak m)\) 为交换局部环，\(M\) 为平坦 \(R\) 模。如果 \(m_1,\ldots,m_n\in M\) 在 \(M/\mathfrak mM\) 中的像对域 \(R/\mathfrak m\) 线性无关，则它们在 \(M\) 中对 \(R\) 线性无关。特别地，有限生成平坦模在交换局部环上自由。
\end{lemma}
\begin{proof}
对 \(n\) 归纳。若 \(am_1=0\)，由\cref{b3:lem:flat-relations}写
\(m_1=\sum_l b_l y_l\)、\(ab_l=0\)。因为 \(m_1\notin\mathfrak mM\)，至少一个 \(b_l\notin\mathfrak m\)，为单位，于是 \(a=0\)。

一般设 \(\sum_i a_im_i=0\)。有限关系引理给出
\(m_i=\sum_l b_{il}y_l\)、\(\sum_i a_i b_{il}=0\)。
在最后一行，至少有一个 \(b_{nl}\) 是单位，否则 \(m_n\in\mathfrak mM\)。由这一列可写
\(a_n=\sum_{i<n}c_i a_i\)，故
\[
\sum_{i<n}a_i(m_i+c_i m_n)=0.
\]
这 \(n-1\) 个新元素在剩余域中仍线性无关：任何关系都会成为原 \(n\) 个独立像之间的关系。归纳假设给出 \(a_i=0\) 对 \(i<n\)，再得 \(a_n=0\)。

若 \(M\) 有限生成，取 \(M/\mathfrak mM\) 的有限基并提升为 \(m_i\)。\cref{b3:cor:local-generators}即 Nakayama 生成元判据说明它们生成 \(M\)；刚证的独立性说明它们是一组自由基。
\end{proof}

这一步实际上只要求模有限生成。有限展示在下一步扩大邻域时才不可替代。局部环上的结论及其有限关系证明，可对照\cite[\S7, Theorem 7.10, pp. 51--52]{Matsumura1989CommutativeRingTheory}。

\subsection{有限生成投射模与局部自由模}\label{b3:subsec:1202:03}

\begin{lemma}\label{b3:lem:free-stalk-neighborhood}
设 \(R\) 为交换环，\(M\) 为有限展示 \(R\) 模，\(\mathfrak p\subseteq R\) 为素理想。若 \(M_{\mathfrak p}\cong R_{\mathfrak p}^r\)、\(r\) 为非负整数，则存在 \(f\in R\setminus\mathfrak p\)，使 \(M_f\cong R_f^r\)。
\end{lemma}
\begin{proof}
选 \(M\) 中有限个元素，其在 \(\mathfrak p\) 处的像经乘单位后是一组基，定义 \(u:R^r\to M\)。余核有限生成且在 \(\mathfrak p\) 处为零，所以逐个清除生成元的分母，得到 \(f\notin\mathfrak p\)，使 \(u_f\) 满射。

以下在 \(R_f\) 上工作，并暂省去局部化下标，使 \(u\) 确为满射。有限展示保证这个有限自由满射的核 \(K\) 有限生成。具体地，若另有有限展示 \(R^a\xrightarrow D R^b\xrightarrow v M\to0\)，分别将 \(u\) 的基像提升到 \(R^b\)、把 \(v\) 的基像提升到 \(R^r\)，得映射 \(\alpha:R^r\to R^b\)、\(\beta:R^b\to R^r\)，满足 \(v\alpha=u\)、\(u\beta=v\)。对 \(z\in\ker u\)，\(\alpha z\) 属于 \(\operatorname{im}D\)，而
\(z=\beta\alpha z+(1-\beta\alpha)z\)；
因此 \(\ker u\) 由 \(\beta D\) 的有限列及 \(1-\beta\alpha\) 的有限列生成。

核 \(K\) 在 \(\mathfrak p\) 处为零，再清有限个分母。恢复原环记号，可取 \(g\in R\setminus\mathfrak p\)，使 \(u_{fg}\) 同时单射和满射。
\end{proof}

\begin{proof}[\cref{b3:thm:finite-presentation-flat-projective}的证明]
第一项由\cref{b3:lem:finite-flat-local-free}推出第三项；第三项由\cref{b3:thm:flatness-local}推出第一项。若 \(M\) 是有限自由模的直和项，它平坦，且有限展示：设 \(R^n=M\oplus N\)，两投影的像都有限生成，核 \(N\) 因而给出 \(M\) 的有限展示。所以第二项推出第一项。

反向，设 \(M\) 满足第三项，取有限自由满射 \(F\to M\)。考虑
\[
\operatorname{Hom}_R(M,F)\longrightarrow\operatorname{Hom}_R(M,M)
\]
的余核。在任意极大理想处，由\cref{b3:thm:hom-localization}及 \(M\) 有限展示，局部化后的 Hom 正是局部 Hom；\(M_{\mathfrak m}\) 自由使局部满射有截面，故恒等映射在余核中的类处处为零。由\cref{b3:thm:local-zero-detection}，这个类本来就为零，于是 \(F\to M\) 有全局截面，\(M\) 为有限自由模的直和项。

第三项由\cref{b3:lem:free-stalk-neighborhood}给出邻域局部自由。反过来，若在主开邻域上有限自由，由谱拟紧性选有限覆盖 \(D(f_i)\)。\cref{b3:prop:principal-cover-finiteness}先使 \(M\) 有限生成。取 \(R^n\to M\) 满射，其核在每个 \(D(f_i)\) 上是有限自由模的直和项，故有限生成；再次应用同一命题使核整体有限生成。因此 \(M\) 有限展示，并满足第三项。
\end{proof}

局部自由的秩在每个平凡化邻域上恒定，因而为局部常值函数；若素谱连通，秩为一个常数。模仍未必整体自由：从局部的基到一个全局基，还需要在重叠处协调基变换。\secref{b3:sec:1206}将把这种相容性写成具体数据。

% !TeX root = ../../Book3.tex
% 纲要标题：### 12.3 Fitting 理想、纤维生成元数与局部自由
% 纲要位置：工作记录/计划纲要.md，第 2327 行。
% 纲要细目（写作范围）：
% * 从有限展示矩阵的子式定义 Fitting 理想，证明其与展示选择无关，并与任意基变换相容
% * 由剩余域上矩阵的秩解释 \(\dim_{\kappa(\mathfrak p)}(M\otimes_R\kappa(\mathfrak p))\) 的闭条件与上半连续性
% * 在有限展示假设下建立局部自由与常秩的适当 Fitting 判据；区别纤维维数恒定与无附加假设的平坦性断言
% * 连接第二卷第4.2节的子式理想和本章泛自由性；一般环上保留理想不变量，不把它们误写成必然存在的 Smith 对角形
% * 对有限展示模，局部自由秩为 r 的判据写成局部化后的 \(\operatorname{Fitt}_{r-1}(M)=0\)、\(\operatorname{Fitt}_{r}(M)=R\)（约定 \(\operatorname{Fitt}_{-1}=0\)）；用幂零基底上的例子区分这些理想等式与仅在剩余域上看到的维数恒定。

\section{Fitting 理想、纤维生成元数与局部自由}
\label{b3:sec:1203}

有限展示把模写成一个有限矩阵的余核。在域上，矩阵的秩决定余核的维数；在一般环上，子式是否为零、是否生成单位理想，比只看各点的数值保留更多信息。Fitting 理想就是把这些子式组织成与展示无关的不变量。本节的环均为交换环，模均有限展示。


\subsection{Fitting 理想的定义与基变换}
\label{b3:subsec:1203:01}

\begin{definition}\label{b3:def:fitting-ideals}
设 \(R\) 为交换环，\(M\) 为有限展示 \(R\) 模，选取展示
\[
R^m\xrightarrow{D}R^n\longrightarrow M\longrightarrow0,
\]
记 \(I_j(D)\) 为矩阵 \(D\) 的全部 \(j\) 阶子式生成的理想，约定 \(I_0(D)=R\)，超过行数或列数的子式理想为零。对 \(i\ge0\)，定义第 \(i\) 个\term[fitting-lixiang]{Fitting 理想}[Fitting ideal]为
\[
\operatorname{Fitt}_i(M)=I_{n-i}(D),
\]
其中 \(i\ge n\) 时取 \(R\)；另约定 \(\operatorname{Fitt}_{-1}(M)=0\)。
\symidx{\(\operatorname{Fitt}_i(M)\)：模的第 \(i\) 个 Fitting 理想}
\end{definition}

例如，循环模 \(R/(a)\) 的展示 \(R\xrightarrow{a}R\to R/(a)\to0\) 只有一个生成元，故 \(\operatorname{Fitt}_0(R/(a))=(a)\)、\(\operatorname{Fitt}_1(R/(a))=R\)。自由模 \(R^r\) 则有 \(r\) 个生成元和零个关系；零关系矩阵的正阶子式都为零，所以 \(\operatorname{Fitt}_{r-1}(R^r)=0\)、\(\operatorname{Fitt}_r(R^r)=R\)。\(r=0\) 时第一式使用 \(\operatorname{Fitt}_{-1}=0\) 的约定。

这些理想从选定的生成元和关系计算而来，还须证明换用另一组有限展示会得到相同理想。把两组生成元放到同一个自由模中，并各自补上新生成元的定义关系，就能比较它们的关系组。

\begin{lemma}[有限展示的稳定化]\label{b3:lem:presentation-stabilization}
设 \(R\) 为交换环，\(M\) 为 \(R\) 模，给定两个有限展示
\[
R^m\xrightarrow{D}R^n\xrightarrow{\pi}M\longrightarrow0,
\qquad
R^{m'}\xrightarrow{E}R^{n'}\xrightarrow{\rho}M\longrightarrow0,
\]
其中 \(m,n,m',n'\) 为非负整数。存在 \(R\) 线性映射
\(U:R^{n'}\rightarrow R^n\)、\(V:R^n\rightarrow R^{n'}\)，使 \(\pi U=\rho\)、\(\rho V=\pi\)。共同的自由模满射为
\[
\sigma:R^n\oplus R^{n'}\longrightarrow M,
\qquad (x,y)\longmapsto\pi(x)+\rho(y).
\]
其核分别由以下两个矩阵的列生成：
\[
\widehat D=\begin{pmatrix}D&-U\\0&I_{n'}\end{pmatrix},
\qquad
\widehat E=\begin{pmatrix}I_n&0\\-V&E\end{pmatrix}.
\]
在 \(\widehat D\) 中，\(D\) 的旧列加上零分量后仍是原关系；新增的列 \((-Ue_j,e_j)\) 则表示第二组第 \(j\) 个生成元等于第一组生成元的组合 \(Ue_j\)，其中 \(e_j\) 为 \(R^{n'}\) 的第 \(j\) 个标准基向量。\(\widehat E\) 交换两组生成元的角色。将两矩阵的列合并，就得到共同的有限关系组；因此两个展示可以通过添加生成元及其定义关系、添加或删去冗余关系变为同一有限展示。
\end{lemma}
\begin{proof}
为 \(R^{n'}\) 的每个基向量在 \(\pi\) 下选择其 \(\rho\) 像的原像，延拓得到 \(U\)；同样利用 \(\rho\) 的满射性得到 \(V\)。映射 \(\sigma\) 满射，因为它在第一直和项上的限制为 \(\pi\)。

矩阵 \(\widehat D\) 表示映射 \((z,y)\mapsto(Dz-Uy,y)\)。其列都在 \(\ker\sigma\) 中。反过来，若 \((x,y)\in\ker\sigma\)，则 \(\pi(x+Uy)=0\)。由第一展示的正合性，存在 \(z\in R^m\) 使 \(x+Uy=Dz\)，于是
\[
(x,y)=(Dz-Uy,y)=\widehat D(z,y).
\]
所以这些列生成全部核。对第二展示，\(\rho(y+Vx)=0\) 给出 \(z'\in R^{m'}\) 使 \(y+Vx=Ez'\)，故
\[
(x,y)=(x,Ez'-Vx)=\widehat E(x,z').
\]
这证明第二组列也生成同一个核。每组中的列因而都是另一组列的线性组合：从任一组出发添加另一组列，均只添加冗余关系，便得到合并后的共同展示。
\end{proof}

\begin{theorem}\label{b3:thm:fitting-invariance-base-change}
设 \(R\) 为交换环，\(M\) 为有限展示 \(R\) 模。理想 \(\operatorname{Fitt}_i(M)\) 与有限展示的选择无关，并且对任意交换环同态 \(R\to S\) 及 \(i\ge0\) 都有
\[
\operatorname{Fitt}_i(S\otimes_RM)=\operatorname{Fitt}_i(M)S.
\]
此外 \(\operatorname{Fitt}_{i-1}(M)\subseteq\operatorname{Fitt}_i(M)\)。
\end{theorem}
\begin{proof}
\cref{b3:lem:presentation-stabilization}将不同展示的比较化为添加生成元及其定义关系、添加或删去冗余关系；可逆换基则可用于整理这些矩阵。可逆行列变换不改变各阶子式理想：展开变换后子式，它是原同阶子式的线性组合；再作逆变换得到反包含。添加一个已有关系的线性组合也不改变子式理想，因为沿新增列的多线性展开把新子式写成旧子式的组合。

添加一个生成元并用一个关系把它表示为旧生成元的组合，在可逆行变换后将矩阵变成 \(\operatorname{diag}(D,1)\)。对 \(j\ge1\)，新矩阵的 \(j\) 阶子式由 \(D\) 的 \(j\) 阶和 \(j-1\) 阶子式生成；前者包含于后者，故
\(I_j\bigl(\operatorname{diag}(D,1)\bigr)=I_{j-1}(D)\)。同时目标自由模秩增加一，所以 \(I_{n-i}(D)\) 保持不变。约定中的零理想和单位理想也满足这一公式。

由\cref{b3:lem:presentation-stabilization}，两个有限展示添加生成元及其定义关系后，具有同一个自由模满射；两套有限关系列都生成它的核。将两套关系列合并只添加冗余关系。因此，上述三类变换的检查证明了各 Fitting 理想与展示选择无关。

任意标量扩张保持展示的右正合性，故 \(S\otimes_RM\) 由同一矩阵的系数像展示。行列式是整系数多项式，子式因而变成原子式的像，得到基变换公式。最后，将每个 \(j\) 阶子式沿一行展开，得到 \(I_j(D)\subseteq I_{j-1}(D)\)，即所述递增性。
\end{proof}

例如 \(M=R/(a)\oplus R/(b)\) 的对角展示给出
\(\operatorname{Fitt}_0(M)=(ab)\)、\(\operatorname{Fitt}_1(M)=(a,b)\)、\(\operatorname{Fitt}_2(M)=R\)。这组理想记录的不只是一个最大子式。

\subsection{纤维生成元数的上半连续性}
\label{b3:subsec:1203:02}

对 \(\mathfrak p\in\operatorname{Spec}R\)，令
\[
\mu_M(\mathfrak p)=\dim_{\kappa(\mathfrak p)}
 \bigl(M\otimes_R\kappa(\mathfrak p)\bigr).
\]
纤维同构于 \(M_{\mathfrak p}/\mathfrak pM_{\mathfrak p}\)。对局部环 \(R_{\mathfrak p}\) 上的有限生成模应用\cref{b3:cor:local-generators}，提升这个剩余向量空间的一组基便得到生成元；反过来，任何生成组的像都必须生成纤维。因此 \(\mu_M(\mathfrak p)\) 恰是 \(M_{\mathfrak p}\) 的最少生成元数。Fitting 理想把这个数的大小转成子式的消失条件。

\begin{proposition}\label{b3:prop:fitting-fiber-rank}
设 \(R\) 为交换环，\(M\) 为有限展示 \(R\) 模。对素理想 \(\mathfrak p\subseteq R\)，记 \(\kappa(\mathfrak p)=R_{\mathfrak p}/\mathfrak pR_{\mathfrak p}\)，并令 \(\mu_M(\mathfrak p)=\dim_{\kappa(\mathfrak p)}\bigl(M\otimes_R\kappa(\mathfrak p)\bigr)\)。对每个整数 \(r\ge1\)，
\[
\bigl\{\mathfrak p:\mu_M(\mathfrak p)\ge r\bigr\}
 =V\bigl(\operatorname{Fitt}_{r-1}(M)\bigr).
\]
特别地，\(\mu_M\) 是上半连续的整值函数，即其取值不小于每个给定整数的点集为闭集。\(M\) 的支集为 \(V\bigl(\operatorname{Fitt}_0(M)\bigr)\)。
\end{proposition}
\begin{proof}
将定义中的展示与剩余域张量，由右正合性得到
\[
\kappa(\mathfrak p)^m\xrightarrow{D(\mathfrak p)}\kappa(\mathfrak p)^n
\longrightarrow M\otimes_R\kappa(\mathfrak p)\longrightarrow0.
\]
故 \(\mu_M(\mathfrak p)=n-\operatorname{rank}D(\mathfrak p)\)。若 \(r>n\)，左侧点集为空，而 \(\operatorname{Fitt}_{r-1}(M)=R\)，右侧也为空。若 \(1\le r\le n\)，则 \(\mu_M(\mathfrak p)\ge r\) 等价于矩阵的秩至多为 \(n-r\)，又等价于所有 \(n-r+1\) 阶子式在剩余域中为零。一个 \(R\) 中的子式在剩余域中消失，当且仅当它属于 \(\mathfrak p\)，所以此条件就是 \(\operatorname{Fitt}_{r-1}(M)\subseteq\mathfrak p\)。这证明闭集公式。支集公式再用 Nakayama 引理：有限生成模 \(M_{\mathfrak p}\) 为零当且仅当其剩余向量空间为零。
\end{proof}

例如设 \(k\) 为域，则 \(k[t]\) 模 \(M=k[t]/(t)\) 在原点的生成元数为一，在其余点为零。闭集 \(V(t)\) 正好是数值增大的位置。一般地，若 \(\mathfrak p\subseteq\mathfrak q\)，则对每个 \(r\ge1\)，\(\mathfrak p\) 属于闭集 \(V(\operatorname{Fitt}_{r-1}(M))\) 时，\(\mathfrak q\) 也属于它，故 \(\mu_M(\mathfrak p)\le\mu_M(\mathfrak q)\)。这就是向特殊点移动时数值可以向上跳的含义，并不保证各点数值相同。

\subsection{局部自由与常秩的 Fitting 判据}
\label{b3:subsec:1203:03}

\begin{theorem}\label{b3:thm:fitting-local-free}
设 \(R\) 为交换环，\(M\) 为有限展示 \(R\) 模，\(\mathfrak p\in\operatorname{Spec}R\)，\(r\ge0\) 为整数。模 \(M\) 在 \(\mathfrak p\) 的某个主开邻域上自由且秩为 \(r\)，当且仅当
\[
\operatorname{Fitt}_{r-1}(M)_{\mathfrak p}=0,
\qquad \operatorname{Fitt}_{r}(M)_{\mathfrak p}=R_{\mathfrak p}.
\]
因此 \(M\) 在整个素谱上局部自由且处处秩为 \(r\)，当且仅当这两项在 \(R\) 中分别等于零和 \(R\)。
\end{theorem}
\begin{proof}
自由模 \(R^r\) 的零关系展示直接给出所列理想。必要性由\cref{b3:thm:fitting-invariance-base-change}成立。

反向先在局部环 \(R_{\mathfrak p}\) 上讨论，取 \(M_{\mathfrak p}\) 的有限展示矩阵 \(D\)，设其行数为 \(n\)。若 \(r=0\)，\(\operatorname{Fitt}_0=R_{\mathfrak p}\) 意味着 \(D\) 有一个满行数可逆子矩阵，故余核为零。若 \(r>0\)，剩余域秩计算先保证 \(r\le n\)。单位理想条件保证存在一个 \(n-r\) 阶单位子式，因为局部环中一组非单位不能生成单位理想。交换行列并对这个可逆方块作消元，可将矩阵化成
\[
\begin{pmatrix}I_{n-r}&0\\0&E\end{pmatrix}.
\]
单位方块在来源与目标的对应坐标间给出同构，对余核没有贡献；余核只剩下 \(\operatorname{coker}E\)。剩余的 \(E\) 有 \(r\) 行。它的每个元素，差一个符号，都是由单位方块加一行一列组成的 \(n-r+1\) 阶子式，故属于 \(\operatorname{Fitt}_{r-1}=0\)。于是 \(E=0\)，余核为 \(R_{\mathfrak p}^r\)。情形 \(n=r\) 时这句话直接说全部矩阵元素为零。由\cref{b3:lem:free-stalk-neighborhood}，自由性扩大到主开邻域。

若在全部素理想处满足条件，将理想 \(\operatorname{Fitt}_{r-1}(M)\) 视为 \(R\) 模，并取商模 \(R/\operatorname{Fitt}_r(M)\)。这两个 \(R\) 模在全部素理想处的局部化都为零，由\cref{b3:thm:local-zero-detection}得到全局的两项理想等式。反向直接局部化。
\end{proof}

\(\operatorname{Fitt}_r(M)=R\) 排除了纤维生成元数大于 \(r\) 的点。对 \(r\ge1\)，处处有 \(\mu_M\ge r\) 却只能说明 \(\operatorname{Fitt}_{r-1}(M)\) 包含于全部素理想，即包含于 \(\sqrt{(0)}\)。其中仍可能留有非零幂零关系。因此局部自由判据要求这项理想本身为零，比单看纤维上的数值更强。

\subsection{子式理想、泛自由性与 Smith 形的适用条件}
\label{b3:subsec:1203:04}

第二卷第4.2节在主理想整环上用子式理想确定 Smith 不变因子\footnote{史密斯（Henry John Stephen Smith，1826-1883），英国数学家，研究数论与二次型，并提出整数矩阵的标准形。}。若展示矩阵能化为对角形，各 Fitting 理想正是适当个数对角元乘积所生成的理想；一般环上的定义仍然有效，却不保证矩阵可以对角化。

\ProseParagraphBreak
例如设 \(k\) 为域，取 \(R=k[x,y]\)。模 \(R/(x,y)\) 的一行展示为 \(\begin{pmatrix}x&y\end{pmatrix}\)，所以 \(\operatorname{Fitt}_0=(x,y)\)。若这个矩阵能通过可逆行列变换化成 \(\begin{pmatrix}d&0\end{pmatrix}\)，同一理想必须为主理想 \((d)\)，与 \((x,y)\) 非主矛盾。后者可由次数或互素性检查：若 \(d\) 同时整除 \(x,y\)，则 \(d\) 为单位，而 \((x,y)\ne R\)。

\ProseParagraphBreak
若 \(R\) 为整环，任一有限展示模在分式域上当然自由。\cref{b3:lem:free-stalk-neighborhood}把这一事实扩大到某个非空主开集。下一节将处理更强的情形：模只须在一个有限型 \(R\) 代数上有限生成，在 \(R\) 上的自由秩可以无限。

\subsection{幂零基底上的 Fitting 理想与纤维维数}
\label{b3:subsec:1203:05}

设 \(k\) 为域，令 \(R=k[\varepsilon]/(\varepsilon^2)\)，\(M=R/(\varepsilon)\)。素谱只有一点，且该点的纤维维数为一，因此纤维维数是常数。但一行展示 \((\varepsilon)\) 给出
\[
\operatorname{Fitt}_0(M)=(\varepsilon)\ne0,
\qquad \operatorname{Fitt}_1(M)=R.
\]
\cref{b3:thm:fitting-local-free}说明 \(M\) 不是秩一自由模；由\cref{b3:lem:finite-flat-local-free}，它也不平坦。纤维只看见 \(\varepsilon=0\)，而 Fitting 理想仍保存这项非零幂零关系。

\ProseParagraphBreak
有一个可以成立的补充结论。若 \(R\) 约化，且有限展示模的纤维维数处处为固定的 \(r\)，则 \(\operatorname{Fitt}_{r-1}(M)\) 包含于所有素理想，因约化性等于零；\(\operatorname{Fitt}_r(M)\) 不包含于任何素理想，故等于 \(R\)。于是 \(M\) 局部自由。若秩只局部常值，可在相应开邻域分别应用同一论证。不能删除这里的约化性而保留原结论。

% !TeX root = ../../Book3.tex
% 纲要标题：### 12.4 泛自由性与泛平坦性
% 纲要位置：工作记录/计划纲要.md，第 2186 行。
% 纲要细目（写作范围）：
% * 固定正文版本：\(A\) 为 Noether 整环、\(B\) 为有限型 \(A\)-代数、\(M\) 为有限生成 \(B\)-模；证明存在 \(0\ne f\in A\) 使 \(M_f\) 作为 \(A_f\)-模自由，且自由秩未必有限
% * 从 \(\operatorname{Frac}(A)\) 上的有限关系和首项计算出发清除有限个分母；证明中所需的有限自由模项序与 Dickson 型有限性在本节就地建立，不前置第21章的合冲模算法
% * 同时控制代数及其有限生成模，导出稠密主开集上的平坦性；说明基底整性、有限型与有限生成假设的用途
% * 在已经具有概形语言时再解释一般纤维与特殊纤维；此处只用局部化和张量积表述，不要求先读第24章
% * 代数与模均有限展示时的非 Noether 推广作为结论范围的补充说明，不与正文版本混写


\section{泛自由性与泛平坦性}
\label{b3:sec:1204}

设一个方程族的系数属于整环 \(A\)。转到分式域以后，有限关系中的非零系数都可以相除；问题是能否只倒置有限个系数，便在同一个开集上得到一组基。下面用首项与有限分母回答这个问题。证明所需的模单项式在此处建立。


\subsection{Noether 整环上的泛自由性定理}
\label{b3:subsec:1204:01}

\begin{theorem}[泛自由性]\label{b3:thm:generic-freeness}
设 \(A\) 为 Noether 整环，\(B\) 为有限型交换 \(A\) 代数，\(M\) 为有限生成 \(B\) 模。则存在 \(0\ne f\in A\)，使 \(M_f\) 作为 \(A_f\) 模自由。其自由秩允许为零或无限。
\end{theorem}

\term[fan-ziyouxing]{泛自由性}[generic freeness]中的“泛”表示可以从基底去掉一个真闭集。这里没有要求 \(M\) 是有限生成的 \(A\) 模；例如 \(B=M=A[x]\) 已经是自由 \(A\) 模，但基 \(1,x,x^2,\ldots\) 为无限集。定理也允许整个模在局部化后消失，例如 \(M=A/(a)\)、\(a\ne0\)，倒置 \(a\) 后自由秩为零。

\ProseParagraphBreak
证明放在下一小节，核心是为商模选出一组可同时保留的单项式。更一般的泛自由性结果可参阅 Matsumura 的《Commutative Ring Theory》\cite[\S24, Theorem 24.1, p. 186]{Matsumura1989CommutativeRingTheory}；这里给出的证明专门处理本章的有限展示数据。

\subsection{有限自由模项序与分母清除}
\label{b3:subsec:1204:02}

令 \(K=\operatorname{Frac}(A)\)，\(P=A[x_1,\ldots,x_s]\)。选满射 \(P\to B\) 及有限组 \(B\) 模生成元，得到 \(M=P^q/N\)。由 \(A\) 的 Noether 性和 Hilbert 基定理，关系模 \(N\) 有限生成。

\begin{lemma}\label{b3:lem:module-leading-finiteness}
设 \(K\) 为域，\(s,q\) 为非负整数，\(e_1,\ldots,e_q\) 为 \(K[x_1,\ldots,x_s]^q\) 的标准基。给模单项式 \(x^\alpha e_j\) 按全次数、指数的字典序、最后按 \(j\) 排序。这是良序，且同乘一个单项式保持大小。任一 \(K[x_1,\ldots,x_s]\) 子模 \(L\) 的非零元素的首单项式，均被有限多个来自 \(L\) 的首单项式之一整除；整除只在同一基向量 \(e_j\) 上比较。
\end{lemma}
\begin{proof}
给定次数的单项式只有有限多个，所以任一非空集合先选最小次数，再选其中最小项，得到良序；加同一个指数向量不改变比较结果。

所需有限性归结为 \(\mathbb N^s\) 中逐坐标偏序的极小元有限。为证明这一点，先注意任一自然数序列都有无限不减子序列：某个值出现无限次时取常值子序列；否则各有限区间只出现有限次，逐次选更大的项。对 \(s\) 个坐标依次取子序列，可知 \(\mathbb N^s\) 的任意无限序列包含两个按原次序排列且逐坐标不减的项。因此一个集合不可能有无限多个不同的极小元。又对任何 \(\alpha\)，其下方的有限盒子 \(\{\beta:\beta\le\alpha\}\) 中可选极小元，故每个元素都在某个极小元之上。这就是所需的 Dickson 型有限性。

分别对每个位置 \(j\) 的首项指数集合应用它，选出有限组极小首项，并取产生它们的 \(g_i\in L\)。全部其他首项都被这些首项之一整除。
\end{proof}

\begin{proof}[\cref{b3:thm:generic-freeness}的证明]
对 \(L=K\otimes_A N\subseteq K[x_1,\ldots,x_s]^q\) 应用引理，选 \(g_1,\ldots,g_t\)，并把首项系数归一成一。逐次消去可被它们的首项整除的最大项，或把不能整除的最大项移入余式。每一步使尚待处理的首单项式严格下降，故有限步终止。对 \(L\) 中的元素，最后余式若非零，其首项一方面属于 \(L\) 的首项，另一方面不被任何所选首项整除，矛盾。因此这些 \(g_i\) 生成 \(L\)。若 \(L=0\)，下述空列表的论证同样成立。

取 \(N\) 的有限组 \(P\) 模生成元 \(v_1,\ldots,v_h\)。每个 \(g_i\) 是 \(v_j\) 的 \(K[x]\) 线性组合，每个 \(v_j\) 也是 \(g_i\) 的 \(K[x]\) 线性组合。这里只涉及有限多个多项式及有限多个系数，故可选非零 \(f\in A\)，倒置它后使所有系数落在 \(A_f\)。于是 \(g_i\in N_f\)，且它们生成 \(N_f\)，首项系数仍为一。

称不被任何 \(g_i\) 的首单项式整除的模单项式为标准单项式。刚才的首项消去只除以首项系数一，故能直接在 \(A_f[x]^q\) 中进行。这说明标准单项式的类生成 \(P_f^q/N_f=M_f\) 作为 \(A_f\) 模。若这些类之间有一个非零有限线性关系，所得非零向量 \(w\in N_f\) 在 \(K[x]^q\) 中仍非零，因为 \(A_f\subseteq K\)。其首单项式是标准的，却又是 \(L\) 中元素的首单项式，与引理矛盾。因此标准单项式的类线性无关，组成所求自由基。
\end{proof}

有限性在此有两个不同用途：Dickson 型论证限制需要的首项数，关系模的有限生成性限制返回原环时需要清除的分母数。只说“清除分母”而不说明哪一组有限数据，不足以得到同一个 \(f\)。

\subsection{稠密主开集上的平坦性}
\label{b3:subsec:1204:03}

\begin{corollary}[泛平坦性]\label{b3:cor:generic-flatness}
设 \(A\) 为 Noether 整环，\(B\) 为有限型交换 \(A\) 代数，\(M_1,\ldots,M_r\) 为预先给定的有限组有限生成 \(B\) 模。存在同一个 \(0\ne f\in A\)，使 \(B_f\) 与全部 \((M_i)_f\) 都是自由 \(A_f\) 模，因而平坦。
\end{corollary}
\begin{proof}
把有限生成的 \(B\) 模 \(B\oplus\bigoplus_iM_i\) 的每个直和项分别应用泛自由性，取所得非零元素的乘积为 \(f\)。进一步局部化保留自由基，所以各项均自由。自由模平坦已在\cref{b3:def:flat-module}之后直接证明。
\end{proof}

整环 \(A\) 的零素理想属于 \(D(f)\)，且每个非空主开集 \(D(g)\) 都与它相交，因为 \(fg\ne0\)。因此 \(D(f)\) 稠密。这里得到了一个稠密主开集上的平坦性，没有声称整个平坦点集一定由这一个主开集描述。

\subsection{一般纤维与特殊纤维的代数描述}
\label{b3:subsec:1204:04}

对 \(\mathfrak p\in\operatorname{Spec}A\)，代数及模的纤维分别是
\[
B(\mathfrak p)=B\otimes_A\kappa(\mathfrak p),\qquad
M(\mathfrak p)=M\otimes_A\kappa(\mathfrak p).
\]
零素理想处称为一般纤维，使用的是 \(K\)。若 \(M_f\) 已有 \(A_f\) 自由基 \(E\)，则对每个 \(\mathfrak p\in D(f)\)，同一组基的像给出
\(M(\mathfrak p)\cong\kappa(\mathfrak p)^{(E)}\)。因此此开集上的模纤维具有同一基集的大小；当 \(E\) 无限时，不把它误写成某个有限的纤维维数。

\ProseParagraphBreak
例如 \(A=k[t]\)、\(B=A[x]/(tx)\)。在 \(t\ne0\) 处，\(x=0\)，所以 \(B_t\cong A_t\)；在 \(t=0\) 处，纤维却为 \(k[x]\)。作为 \(A\) 模，\(x\ne0\) 而 \(tx=0\)，故 \(B\) 有挠，不平坦。若换成 \(A[x]/(x^2-t)\)，首一除法给出全局基 \(1,x\)，因而处处平坦。即使在 \(t=0\) 处纤维为非约化环 \(k[x]/(x^2)\)，平坦性仍成立；平坦与纤维约化是两种不同的要求。

\subsection{代数与模均有限展示时的非 Noether 推广}
\label{b3:subsec:1204:05}

\begin{proposition}\label{b3:prop:generic-free-finite-presentation}
设 \(A\) 为任意整环，\(B\) 为有限展示交换 \(A\) 代数，\(M\) 为有限展示 \(B\) 模。则存在 \(0\ne f\in A\)，使 \(M_f\) 作为 \(A_f\) 模自由；其自由秩允许为零或无限。
\end{proposition}
\begin{proof}
写 \(B=A[x_1,\ldots,x_s]/(h_1,\ldots,h_u)\)，并取有限展示 \(B^v\to B^q\to M\to0\)。把第一映射的有限列提升到 \(P^q\)，再加上 \(h_a e_j\)（\(1\le a\le u,1\le j\le q\)），这些有限列恰好生成 \(P^q\to M\) 的核：任何核元素先模掉 \((h_a)P^q\)，成为 \(B^q\) 中给定关系的组合，提升该组合后差只属于 \((h_a)P^q\)。因此 \(M\) 在 \(P\) 上有限展示。

前述证明在得到有限组 \(v_j\) 后不再使用 \(A\) 的 Noether 性：分式域仍存在，模首项有限性在域上的多项式模中独立成立，且只清除有限个分母。将该证明逐步应用便得结论。
\end{proof}

非 Noether 版本改变的是假设的表达：不能仅凭“有限型代数、有限生成模”自动断言关系有限。代数和模各自的有限展示条件，分别保证代数关系和模关系都能由有限数据控制。这项推广不改变自由秩可以无限的结论。

% !TeX root = ../../Book3.tex
% 纲要标题：### 12.5 忠实平坦性
% 纲要位置：工作记录/计划纲要.md，第 2321 行。
% 纲要细目（写作范围）：
% * 正合性和消失的反映
% * 忠实平坦基变换下的可检验性质
% * 纤维与下降的基本动机

\section{忠实平坦性}
\label{b3:sec:1205}

平坦基变换保持已有的正合列，却可能把非零模变成零。例如与 \(\mathbb Q\) 张量会使有限挠交换群消失。若希望从变基后的计算推回原来的模，还须保证信息没有这样丢失；忠实平坦性正是这一附加要求。


\subsection{正合性与消失的反映}
\label{b3:subsec:1205:01}

\begin{definition}\label{b3:def:faithfully-flat}
设 \(R\to S\) 为交换环同态。如果 \(S\) 为平坦 \(R\) 模，且对每个 \(R\) 模 \(N\)，从 \(S\otimes_RN=0\) 能推出 \(N=0\)，就称这个同态\term[zhongshi-pingtan]{忠实平坦}[faithfully flat]。
\end{definition}

\begin{theorem}\label{b3:thm:faithfully-flat-criterion}
设 \(R\to S\) 为平坦交换环同态。下列条件等价：
\begin{equivalences}
\item \(R\to S\) 忠实平坦；
\item 对每个极大理想 \(\mathfrak m\subset R\)，有 \(S/\mathfrak mS\ne0\)；
\item \(\operatorname{Spec}S\to\operatorname{Spec}R\) 满射。
\end{equivalences}
特别地，非零局部环之间的平坦局部同态忠实平坦。
\end{theorem}
\begin{proof}
第一项对非零模 \(R/\mathfrak m\) 使用定义，就得到第二项。反过来，设 \(0\ne x\in N\)，则 \(Rx\cong R/I\)，其中 \(I=\operatorname{Ann}_R(x)\) 为真理想。选极大理想 \(\mathfrak m\supseteq I\)，则 \(S/IS\) 满射到非零的 \(S/\mathfrak mS\)，故非零。平坦性使
\(S\otimes_RRx\to S\otimes_RN\) 单射，因而后者非零。这证明第一、二项等价。

若忠实平坦，对每个素理想 \(\mathfrak p\)，非零模 \(\kappa(\mathfrak p)\) 的张量积也非零。非零环 \(S\otimes_R\kappa(\mathfrak p)\) 有素理想，其在 \(S\) 中的收缩恰位于 \(\mathfrak p\) 之上：把纤维写成 \((R\setminus\mathfrak p)^{-1}(S/\mathfrak pS)\)，由局部化素理想对应即可核验。所以谱映射满射。反向，若存在 \(\mathfrak q\subset S\) 收缩为 \(\mathfrak m\)，则 \(\mathfrak mS\subseteq\mathfrak q\ne S\)，于是第二项成立。

局部同态 \((R,\mathfrak m)\to(S,\mathfrak n)\) 满足 \(\mathfrak mS\subseteq\mathfrak n\)，剩余环非零，故最后结论由第二项得到。
\end{proof}

这些判据可对照\cite[\S7, Theorems 7.2--7.3, pp. 47--48]{Matsumura1989CommutativeRingTheory}。检验局部同态时须同时确认平坦和 \(\mathfrak mS\ne S\)；仅仅是局部环之间的一个任意同态并不够。

\subsection{忠实平坦基变换下的性质检测}
\label{b3:subsec:1205:02}

\begin{proposition}\label{b3:prop:faithfully-flat-detection}
设 \(R\to S\) 为忠实平坦交换环同态。下列各项中的模及模同态均在 \(R\) 上，变基后在 \(S\) 上。
\begin{enumerate}
\item 模同态的单射、满射、同构以及模序列的正合性，均可在与 \(S\) 张量后检测。
\item 对每个 \(R\) 模 \(M\)，自然映射 \(M\to S\otimes_RM\)、\(m\mapsto1\otimes m\) 单射；特别地 \(R\to S\) 单射，且对任意理想 \(I\)，有 \(IS\cap R=I\)。
\item \(M\) 在 \(R\) 上平坦，当且仅当 \(S\otimes_RM\) 在 \(S\) 上平坦。
\item 若 \(S\otimes_RM\) 在 \(S\) 上有限生成或有限展示，则 \(M\) 分别在 \(R\) 上有限生成或有限展示。
\end{enumerate}
\end{proposition}
\begin{proof}
平坦性使核、像和余核的短正合列仍然正合。例如
\(S\otimes_R\ker u\cong\ker(1\otimes u)\)，余核也相容，故消失检测给出第一项中的单射、满射与同构。对一个复形 \(L\to M\to N\)，同样把中间同调 \(\ker v/\operatorname{im}u\) 张量，便得正合性检测。若原来还不知道 \(vu=0\)，先对 \(\operatorname{im}(vu)\) 使用同一检测，得到它为零。

若 \(m\ne0\)，平坦性使 \(S\otimes_RRm\) 单射进入 \(S\otimes_RM\)。前者非零，且由 \(1\otimes m\) 作为 \(S\) 模生成，故 \(1\otimes m\ne0\)。这证明第二项的单射；对 \(M=R/I\) 应用它，得到 \(R/I\hookrightarrow S/IS\)，即理想收缩公式。

第三项的正向是\cref{b3:prop:flat-base-change-transitivity}。反向任取单射 \(U\to V\)，先用 \(S\) 的平坦性，继而用 \(S\otimes_RM\) 的 \(S\) 平坦性，得到
\[
S\otimes_R(U\otimes_RM)\longrightarrow
S\otimes_R(V\otimes_RM)
\]
单射。由第一项原映射单射，所以 \(M\) 平坦。

若 \(S\otimes_RM\) 由有限多个张量和生成，收集这些和中出现的有限个 \(m_i\in M\)，令 \(M_0=\sum_iRm_i\)。则 \(S\otimes_R(M/M_0)=0\)，故 \(M=M_0\)，证明有限生成下降。若变基模有限展示，先由刚证结论取 \(R^n\twoheadrightarrow M\)，核为 \(K\)。平坦性给出
\[
S\otimes_RK\cong\ker\bigl(S^n\longrightarrow S\otimes_RM\bigr).
\]
因为 \(S\otimes_RM\) 有限展示，\cref{b3:prop:finite-presentation-kernel}说明任意有限秩自由模满射到它的核均有限生成；故 \(S\otimes_RK\) 有限生成。对 \(K\) 再用刚证的有限生成下降，得到 \(K\) 有限生成，因而 \(R^n\to M\) 给出有限展示。
\end{proof}

结合\cref{b3:thm:finite-presentation-flat-projective}，有限生成投射性也可以忠实平坦地检测。\cref{b3:tab:flat-faithful-properties}集中列出基变换保持和反映的性质。这里“反映”是指变基后满足该性质，就能推出变基前也满足。

\begin{center}
\begin{minipage}{\textwidth}
\makeatletter
\def\@captype{table}
\makeatother
\centering
\caption[平坦与忠实平坦基变换的性质检测]{平坦与忠实平坦基变换的性质检测。设 \(R\to S\) 为交换环同态，表中的变基指 \(M\mapsto S\otimes_RM\) 及相应的模同态；模的性质分别相对于 \(R\) 和 \(S\) 而言。}
\label{b3:tab:flat-faithful-properties}
\begin{tabularx}{\textwidth}{@{}>{\raggedright\arraybackslash}p{0.26\textwidth}>{\raggedright\arraybackslash}p{0.28\textwidth}>{\raggedright\arraybackslash}X@{}}
\toprule
性质 & 平坦基变换 & 忠实平坦基变换\\
\midrule
序列正合 & 保持；未必反映 & 保持且反映\\
模非零 & 未必保持；反映 & 保持且反映\\
单射、满射、同构 & 均保持；未必反映 & 均保持且反映\\
有限生成、有限展示 & 均保持；未必反映 & 均保持且反映\\
平坦、有限生成投射 & 均保持；未必反映 & 均保持且反映\\
有限秩自由 & 保持；未必反映 & 保持；反向只保证有限生成投射，未必全局自由\\
\bottomrule
\end{tabularx}
\end{minipage}
\end{center}

表中有限生成、有限展示和有限秩自由性的保持只用张量积的右正合性及其保持直和；平坦性和有限生成投射性的保持由\cref{b3:prop:flat-base-change-transitivity}及有限自由模直和项的描述得到。若变基模有限秩自由，忠实平坦检测给出原模有限生成投射。取有限主开覆盖使一个有限生成投射模逐块自由，再变基到这些局部化的直积，各块的自由秩相同时就得到自由模；重叠处的基却可能无法拼成原环上的全局基。

\subsection{纤维与下降的动机}
\label{b3:subsec:1205:03}

忠实平坦性在任意基变换和复合下保持。复合的证明是依次使用两个正合函子及两个消失检测。对基变换 \(R\to R'\)，\(S'=R'\otimes_RS\) 在 \(R'\) 上平坦；若 \(N\) 为 \(R'\) 模且 \(S'\otimes_{R'}N=0\)，结合律将它识别为 \(S\otimes_RN=0\)，所以 \(N=0\)。

\ProseParagraphBreak
若 \(D(f_1),\ldots,D(f_r)\) 覆盖 \(\operatorname{Spec}R\)，则
\[
R\longrightarrow S=\prod_{i=1}^rR_{f_i}
\]
忠实平坦。有限直积作为模也是有限直和，故平坦；每个素理想至少属于一个 \(D(f_i)\)，局部化的素理想对应说明谱映射满射。单个真局部化一般不忠实，例如 \(\mathbb Z\to\mathbb Z[1/2]\) 遗漏素理想 \((2)\)，并使 \(\mathbb Z/2\mathbb Z\) 消失。

\ProseParagraphBreak
在这个有限开覆盖的例子中，给定变基模相当于在每个开集上给定一个模。要把它们还原为一个全局模，必须说明在重叠的 \(R_{f_if_j}\) 上如何识别，并要求三重重叠上的识别相容。下一节把这个熟悉的拼接问题写成一般忠实平坦环同态上的代数条件。

% !TeX root = ../../Book3.tex
% 纲要标题：### 12.6* 模的忠实平坦下降
% 纲要位置：工作记录/计划纲要.md，第 2327 行。
% 纲要细目（写作范围）：
% * 双重基变换及余循环条件
% * 忠实平坦下降的模论形式
% * 严格区分带下降数据与只有基变换后对象
% ---

\OptionalSection{模的忠实平坦下降}{b3:sec:1206}

给定忠实平坦 \(R\) 代数 \(S\) 上的模 \(N\)，什么时候能写成 \(S\otimes_RM\)？仅有 \(N\) 还不够；它在两个 \(S\) 因子上的拉回之间必须有指定的相容识别。本节把这些数据写清楚，再用一个等化子实际构造 \(M\)。


\subsection{双重基变换与余循环条件}
\label{b3:subsec:1206:01}

固定忠实平坦环同态 \(R\to S\)，令 \(T=S\otimes_RS\)。同一个 \(S\) 模 \(N\) 有两种变基：
\[
N_1=N\otimes_RS,\qquad N_2=S\otimes_RN.
\]
在前者上，\((a\otimes b)(n\otimes s)=an\otimes bs\)；在后者上，\((a\otimes b)(s\otimes n)=as\otimes bn\)。这两种 \(T\) 模结构一般不能忽略下标而混写。

\begin{definition}\label{b3:def:module-descent-datum}
设 \(R\to S\) 为忠实平坦交换环同态，\(N\) 为 \(S\) 模，\(T=S\otimes_RS\)。赋予 \(N_1=N\otimes_RS\) 及 \(N_2=S\otimes_RN\) 自然的 \(T\) 模结构，即两个 \(S\) 因子分别作用在相应位置。若 \(T\) 线性同构
\(\theta:N_1\to N_2\) 满足下面的\term[yuxunhuan-tiaojian]{余循环条件}[cocycle condition]，就称它为 \(N\) 上的\term[xiajiang-shuju]{下降数据}[descent datum]。具体地，把数据沿三个嵌入
\[
T\longrightarrow S\otimes_RS\otimes_RS,
\quad a\otimes b\longmapsto
a\otimes b\otimes1,\ a\otimes1\otimes b,\ 1\otimes a\otimes b
\]
变基，分别得到从第一个位置移到第二个位置的 \(\theta_{12}\)、从第一个移到第三个的 \(\theta_{13}\)、从第二个移到第三个的 \(\theta_{23}\)。要求
\[
\theta_{23}\theta_{12}=\theta_{13}.
\]
把三重张量积中第 \(i\) 个 \(S\) 因子换成 \(N\) 所得的模记为 \(N^{(i)}\)，其中 \(i=1,2,3\)。余循环条件就是下面三个变基同构构成的图交换：
\[
\begin{tikzcd}[column sep=large,row sep=large]
N^{(1)}\arrow[r,"\theta_{12}"]\arrow[dr,"\theta_{13}"']
&N^{(2)}\arrow[d,"\theta_{23}"]\\
&N^{(3)}
\end{tikzcd}
\]
带下降数据的模之间的态射，是与这些同构相容的 \(S\) 线性映射。
\end{definition}

若 \(N=S\otimes_RM\)，标准下降数据为
\[
\theta\bigl((s\otimes m)\otimes t\bigr)
 =s\otimes(t\otimes m).
\]
它的逆把右边送回左边。在三重张量积上，两种复合都把位于第一个 \(S\) 因子旁的 \(m\) 移到第三个因子旁，系数次序不变，所以满足余循环条件。

\ProseParagraphBreak
为把一般条件用于计算，令
\[
\rho:N\longrightarrow S\otimes_RN,\qquad
\rho(n)=\theta(n\otimes1),\qquad
\mu(s\otimes n)=sn.
\]
\(\rho\) 是 \(R\) 线性映射，满足
\begin{equation}\label{b3:eq:descent-coaction}
\begin{split}
\rho(sn)&=(s\otimes1)\rho(n),\qquad \mu\rho(n)=n,\\
(1\otimes\rho)\rho(n)&=\sum_i s_i\otimes1\otimes n_i
\quad\text{若 }\rho(n)=\sum_i s_i\otimes n_i.
\end{split}
\end{equation}
第一式是 \(T\) 线性。第三式就是把余循环条件用于 \(n\otimes1\otimes1\)：先移动到第二个位置，再移动到第三个，得到左边；直接移到第三个得到右边。第二式来自对角变基：沿 \(S\otimes_RS\to S\)、\(s\otimes t\mapsto st\) 拉回 \(\theta\)，得到 \(N\) 的自同构 \(u=\mu\rho\)。再沿全对角映射
\[
S\otimes_RS\otimes_RS\longrightarrow S,
\qquad s_1\otimes s_2\otimes s_3\longmapsto s_1s_2s_3
\]
变基，三个 \(N^{(i)}\) 都自然认同于 \(N\)，而 \(\theta_{12},\theta_{13},\theta_{23}\) 都变为同一个 \(u\)，因为三个 \(T\) 嵌入与全对角映射的复合都是乘法映射 \(T\to S\)。于是余循环条件成为 \(u^2=u\)。\(u\) 可逆，约去一个 \(u\) 得 \(u=\operatorname{id}_N\)，即 \(\mu\rho=\operatorname{id}_N\)。原同构可由 \(\rho\) 恢复为 \(\theta(n\otimes t)=(1\otimes t)\rho(n)\)。

\subsection{忠实平坦下降的模论形式}
\label{b3:subsec:1206:02}

先解释忠实平坦性为何能够恢复本来就在 \(R\) 上的模。

\begin{lemma}\label{b3:lem:faithful-flat-equalizer}
设 \(R\to S\) 为忠实平坦交换环同态，\(M\) 为任意 \(R\) 模。则下列序列正合：
\[
0\longrightarrow M\xrightarrow{m\mapsto1\otimes m}S\otimes_RM
\xrightarrow{\mathrm{d}}S\otimes_RS\otimes_RM,
\]
其中 \(\mathrm{d}(s\otimes m)=1\otimes s\otimes m-s\otimes1\otimes m\)。
\end{lemma}
\begin{proof}
由\cref{b3:prop:faithfully-flat-detection}，可以先与 \(S\) 张量再检查。新序列的首个映射将 \(a\otimes m\) 送到 \(a\otimes1\otimes m\)，乘合前两个 \(S\) 因子给出左逆，故单射。

若 \(z=\sum a_i\otimes b_i\otimes m_i\) 在后一个映射的核中，则
\[
\sum_i a_i\otimes1\otimes b_i\otimes m_i
=\sum_i a_i\otimes b_i\otimes1\otimes m_i.
\]
乘合前两个 \(S\) 因子，得到
\(z=\sum_i a_ib_i\otimes1\otimes m_i\)，恰好属于前一个映射的像。相反，该像代入差映射显然为零。变基后的正合性因而成立，并反映回原序列。
\end{proof}

\begin{theorem}[忠实平坦模下降]\label{b3:thm:faithfully-flat-module-descent}
设 \(R\to S\) 为忠实平坦交换环同态，\(N\) 为带下降数据 \(\theta:N\otimes_RS\to S\otimes_RN\) 的 \(S\) 模，记 \(\rho(n)=\theta(n\otimes1)\)，并令
\[
M=\bigl\{n\in N:\rho(n)=1\otimes n\bigr\}.
\]
则 \(M\) 是 \(R\) 子模，且乘法映射
\[
\varphi:S\otimes_RM\longrightarrow N,\qquad s\otimes m\longmapsto sm
\]
为保下降数据的同构。这个构造与基变换建立 \(R\) 模范畴和带下降数据的 \(S\) 模范畴之间的等价。若 \(N\) 有限展示，则所得 \(M\) 有限展示。
\end{theorem}
\begin{proof}
记 \(\eta:N\to S\otimes_RN\)、\(\eta(n)=1\otimes n\)。\(\rho\) 与 \(\eta\) 都是 \(R\) 线性映射，且 \(M=\ker(\rho-\eta)\)，所以有左端正合列
\[
0\longrightarrow M\longrightarrow N
\xrightarrow{\ \rho-\eta\ }S\otimes_RN.
\]
与 \(S\) 张量后，由平坦性仍有左端正合列
\[
0\longrightarrow S\otimes_RM\longrightarrow S\otimes_RN
\xrightarrow{\ 1\otimes(\rho-\eta)\ }S\otimes_RS\otimes_RN.
\]
因此可将 \(S\otimes_RM\) 认同于 \(S\otimes_RN\) 中下列两个映射的等化子：
\[
s\otimes n\longmapsto s\otimes\rho(n),\qquad
s\otimes n\longmapsto s\otimes1\otimes n.
\]
\eqref{b3:eq:descent-coaction}的第三式说明 \(\rho(n)\) 落在这个等化子中，故 \(\rho\) 给出 \(N\to S\otimes_RM\) 的映射。第二式说明 \(\varphi\rho=1_N\)。

反过来，若 \(z=\sum_i s_i\otimes n_i\) 属于上述等化子，则
\(\sum_i s_i\otimes\rho(n_i)=\sum_i s_i\otimes1\otimes n_i\)。乘合前两个因子，并用第一式，得到
\[
\rho\left(\sum_i s_in_i\right)=\sum_i s_i\otimes n_i=z.
\]
所以 \(\rho\varphi\) 也是恒等映射，\(\varphi\) 为同构。对 \(m\in M\)，\(\rho(sm)=s\otimes m\)，这正是标准数据在 \(s\otimes m\) 上的公式，故同构与原 \(\theta\) 相容。

若 \(u:N\to N'\) 与数据相容，则 \(u\) 将 \(\rho(n)=1\otimes n\) 的元素送入对应的 \(M'\)，因而限制为 \(R\) 线性映射 \(M\to M'\)。它在 \(S\) 上的扩张通过 \(\varphi\) 恢复原 \(u\)。若从 \(R\) 模 \(M_0\) 出发，\cref{b3:lem:faithful-flat-equalizer}恰说明标准数据的等化子是 \(M_0\) 的自然像，故往返构造为自然同构。这证明范畴等价的对象与态射两部分。最后，\(N\cong S\otimes_RM\)，有限展示由\cref{b3:prop:faithfully-flat-detection}第四项下降。
\end{proof}

等化子用相容性选出能够来自原环的元素，平坦性使这个核的构造与基变换相容；忠实性则保证从原环出发的对象与态射能完整恢复。

\subsection{下降数据与基变换后对象}
\label{b3:subsec:1206:03}

仅给出变基后的对象，还没有指定其中哪些元素应从原环而来。例如 \(R=\mathbb R\)、\(S=\mathbb C\)，在 \(N=\mathbb C\) 上可考虑两种半线性对合
\[
\sigma_+(z)=\bar z,\qquad \sigma_-(z)=-\bar z.
\]
它们都满足 \(\sigma(cz)=\bar c\cdot\sigma(z)\)、\(\sigma^2=1\)，分别选出不动实子空间 \(\mathbb R\) 与 \(i\mathbb R\)。这里
\(\mathbb C\otimes_{\mathbb R}\mathbb C\cong\mathbb C\times\mathbb C\)，映射为 \(z\otimes w\mapsto(zw,z\bar w)\)。这两个分量分别对应恒等嵌入与共轭嵌入。在第一个分量上，\(N_1,N_2\) 都识别为 \(N\)，\(\theta\) 是恒等映射；在第二个分量上，\(N_1\) 为 \(N\)，\(N_2\) 为共轭标量作用的模 \(\overline N\)，其中 \(c\cdot n=\bar c n\)。此处 \(\theta\) 由 \(n\mapsto\sigma(n)\) 给出：半线性等式 \(\sigma(cn)=\bar c\sigma(n)\) 正是从 \(N\) 到 \(\overline N\) 的复线性。因而，两种对合都给出前面定义中的 \(T\) 线性同构。

\ProseParagraphBreak
三重张量积的分量可以用两次相对嵌入 \(g,h\in\{1,\text{共轭}\}\) 标记。\(\theta_{23}\theta_{12}=\theta_{13}\) 在这些分量上要求连续两次相对识别等于复合嵌入的识别；涉及恒等嵌入的条件自动成立，唯一剩下的条件就是 \(\sigma^2=1\)。等化子条件在恒等分量上自动成立，在共轭分量上则为 \(\sigma(n)=n\)，故恢复的实模恰为 \(N^\sigma\)。在每种情况下，乘法映射 \(\mathbb C\otimes_{\mathbb R}N^\sigma\to N\) 都是同构。这两组数据也同构：复线性映射 \(q(z)=iz\) 满足
\[
\sigma_-\bigl(q(z)\bigr)=i\bar z=q\bigl(\sigma_+(z)\bigr),
\]
因而与下降数据相容，并将 \(\mathbb R\) 同构地送到 \(i\mathbb R\)。这个例子说明，\(N\) 内具体的固定子模及其乘法识别，需要由下降数据指定；所得实向量空间的同构类型在这里相同。

\ProseParagraphBreak
开覆盖的例子把下降数据写成换基矩阵。设 \(R\) 为交换环，\(r\) 为非负整数，且 \(D(f_1),\ldots,D(f_q)\) 组成 \(\operatorname{Spec}R\) 的有限主开覆盖。在各开集上取自由模 \(N_i=R_{f_i}^r\)，令
\[
S=\prod_{i=1}^qR_{f_i},\qquad
N=\prod_{i=1}^qN_i\cong S^r.
\]
上一节已证明 \(R\to S\) 忠实平坦，且 \(S\otimes_RS\cong\prod_{i,j}R_{f_if_j}\)。在重叠处指定矩阵
\[
g_{ij}\in\operatorname{GL}_r(R_{f_if_j}),
\qquad g_{jk}g_{ij}=g_{ik}\quad\text{于 }R_{f_if_jf_k}.
\]
这些 \(r\) 阶可逆矩阵给出相容同构 \(\theta_{ij}\)，从而给出 \(N\) 的下降数据。所得等化子 \(M\) 由重叠处相符的局部元素组组成，下降同构在第 \(i\) 个分量上给出 \(M_{f_i}\cong N_i=R_{f_i}^r\)。因此 \(M\) 在主开覆盖上有限局部自由，由\cref{b3:thm:finite-presentation-flat-projective}为有限生成投射模。这些矩阵未必能由各个开集上的换基同时消去，所以不一定得到全局自由模。

\ProseParagraphBreak
余循环条件保证局部识别相容，等化子将它们还原成原环上的模；有限展示性也随之下降，使有限组生成元与关系的描述得以保留。


\begin{chaptersummary}
理想张量判据把全部模之间的单射归结为有限生成理想的包含，局部化又允许在各局部环中检查这些映射。有限平坦模在局部环上自由；有限展示进一步使一点处的基扩展到邻域，从而建立有限展示平坦、有限生成投射及有限秩局部自由之间的等价。

\ProseParagraphBreak
Fitting 理想由有限展示的子式定义，与展示无关且允许任意基变换。纤维生成元数的闭条件只记录子式落入哪些素理想；局部自由还要求相应的较大子式理想真正为零。泛自由性则先在分式域上选首项，再清除有限关系涉及的分母，以标准单项式给出一个可保留到主开集上的自由基。

\ProseParagraphBreak
忠实平坦性除保持正合外，还能检测模是否消失，因此能把正合性、有限性和平坦性推回原环。下降需要更完整的输入：双重基变换上的同构及三重基变换上的余循环条件。等化子选出相容元素，平坦性保证这些元素在变基后恰好恢复给定的模。
\end{chaptersummary}

\BookExercises

\begin{exercise}\label{b3:ex:12-flat-quotient-principal}
设 \(R\) 为交换环、\(a\in R\)。证明 \(R/(a)\) 平坦，当且仅当存在 \(b\in R\) 使 \(a=a^2b\)。并说明此时 \((a)\) 由幂等元生成。
\end{exercise}
\begin{solution}
若平坦，将 \((a)\subseteq R\) 与 \(R/(a)\) 张量。所得映射 \((a)/(a)^2\to R/(a)\) 恒为零，平坦性使其单射，故 \((a)=(a)^2\)，即 \(a=a^2b\)。反向令 \(e=ab\)，则 \(e^2=a^2b^2=ab=e\)，且 \(a=ae\)、\(e\in(a)\)，故 \((a)=(e)\)。分解 \(R=eR\oplus(1-e)R\) 给出 \(R/(a)\cong(1-e)R\)，是自由模的直和项，所以平坦。
\end{solution}

\begin{exercise}\label{b3:ex:12-torsion-flat-distinction}
设 \(R\) 为局部整环，\(I\subseteq R\) 为非零有限生成理想。证明 \(I\) 平坦当且仅当它是主理想。再取 \(R=k[x,y]_{(x,y)}\)，证明 \((x,y)^2R\) 无挠但不平坦。
\end{exercise}
\begin{solution}
若 \(I\) 平坦，\cref{b3:lem:finite-flat-local-free}使它成为有限自由模。令 \(K=\operatorname{Frac}R\)，由于 \(I\ne0\)，有 \(K\otimes_RI=K\)，故这个自由模秩只能为一，于是 \(I\) 为主理想。反向若 \(I=(a)\)、\(a\ne0\)，整环条件使 \(R\xrightarrow{r\mapsto ra}I\) 为同构，故平坦。

在给定的 \(R\) 上，\(\mathfrak m^2/\mathfrak m^3\) 的 \(k\) 基为 \(x^2,xy,y^2\) 的类：一个局部分式模 \(\mathfrak m^3\) 后，只需将分母的非零常数项及其一次修正展开，二次齐次部分仍唯一。因此 \(\mathfrak m^2\) 的最少生成元数为三，不能为主理想，由前半题不平坦。它作为整环的子模仍无挠。
\end{solution}

\begin{exercise}\label{b3:ex:12-finite-local-relation}
设 \((R,\mathfrak m)\) 局部，\(M\) 有限生成且平坦。若 \(m_1,\ldots,m_r\) 为一组极小生成元，证明任何关系 \(\sum a_im_i=0\) 都只有零系数。解释有限生成性在此结论中的具体作用。
\end{exercise}
\begin{solution}
由\cref{b3:cor:local-generators}，极小生成组在剩余域上的像为基。\cref{b3:lem:finite-flat-local-free}的有限关系论证说明其提升线性无关，故所有 \(a_i=0\)。有限生成性用于 Nakayama 引理，以使剩余空间的一组基提升后生成整个模；没有这一步，剩余空间甚至可以为零而原平坦模非零，例如 \(R=\mathbb Z_{(p)}\)、\(M=\mathbb Q\)，有 \(pM=M\)。
\end{solution}

\begin{exercise}\label{b3:ex:12-projective-rank-clopen}
设 \(M\) 为有限生成投射 \(R\) 模。证明其秩函数只取有限多个值，每个定秩集合既开又闭。若 \(\operatorname{Spec}R\) 连通，得到什么结论？
\end{exercise}
\begin{solution}
由\cref{b3:thm:finite-presentation-flat-projective}，每个点有一个使 \(M\) 自由的主开邻域。素谱拟紧性给出有限个这样的邻域覆盖，其自由秩只有有限多个。每个定秩集合是相应邻域的并，故开；其补集是其他定秩集合的并，也开，故它闭。连通时只有一个非空定秩集合，因此秩处处相同；这仍不声称模整体自由。
\end{solution}

\begin{exercise}\label{b3:ex:12-fittings-integer}
计算 \(M=\mathbb Z/12\mathbb Z\oplus\mathbb Z/18\mathbb Z\) 的全部 Fitting 理想，并由它们求每个素理想处纤维的维数。
\end{exercise}
\begin{solution}
对角矩阵 \(\operatorname{diag}(12,18)\) 给出
\[
\operatorname{Fitt}_0(M)=(216),\quad
\operatorname{Fitt}_1(M)=(6),\quad
\operatorname{Fitt}_i(M)=\mathbb Z\ (i\ge2).
\]
在 \((2)\)、\((3)\) 处两个对角元都变成零，纤维维数二；其他素数处均为单位，纤维为零；零素理想处张量到 \(\mathbb Q\) 也为零。这与\cref{b3:prop:fitting-fiber-rank}的闭集公式一致。
\end{solution}

\begin{exercise}\label{b3:ex:12-fittings-three-generators}
令 \(R=k[x,y]\)，\(M\) 为矩阵
\(D=\begin{pmatrix}x&y&0\\0&x&y\end{pmatrix}\)
的余核。求各 Fitting 理想、纤维生成元数与支集。
\end{exercise}
\begin{solution}
二阶子式为 \(x^2,xy,y^2\)，一阶子式生成 \((x,y)\)。故
\(\operatorname{Fitt}_0=(x,y)^2\)、\(\operatorname{Fitt}_1=(x,y)\)、\(\operatorname{Fitt}_2=R\)。原点之外矩阵满行秩，纤维为零；在 \((x,y)\) 处矩阵全零，纤维维数二。支集正是原点。\(M\) 不为零，因为此处的剩余向量空间不为零。
\end{solution}

\begin{exercise}\label{b3:ex:12-fittings-direct-sum}
对有限展示模 \(M,N\)，证明
\[
\operatorname{Fitt}_r(M\oplus N)
=\sum_{i+j=r}\operatorname{Fitt}_i(M)\operatorname{Fitt}_j(N)
\quad(r\ge0).
\]
\end{exercise}
\begin{solution}
取目标自由秩分别为 \(m,n\) 的展示矩阵 \(D,E\)。直和由块对角矩阵展示。一个非零的 \(m+n-r\) 阶子式必须从两个方块分别选择相等的行数与列数，因而是某个 \(a\) 阶 \(D\) 子式与 \(b\) 阶 \(E\) 子式的乘积，\(a+b=m+n-r\)。反向每个这样的乘积也是一个块子式。令 \(i=m-a,j=n-b\)，便有 \(i+j=r\)，得到公式。超出行列数及零阶的情形正由定义中的零理想、单位理想约定覆盖。
\end{solution}

\begin{exercise}\label{b3:ex:12-fitting-nilpotence}
在 \(R=k[\varepsilon]/(\varepsilon^3)\) 上，比较 \(M=R\) 与 \(N=R/(\varepsilon^2)\) 的纤维维数和 Fitting 理想，并判断平坦性。
\end{exercise}
\begin{solution}
两者唯一纤维都是一维。对 \(M\)，\(\operatorname{Fitt}_0(M)=0\)、\(\operatorname{Fitt}_1(M)=R\)；对 \(N\)，对应理想为 \((\varepsilon^2)\ne0\) 和 \(R\)。所以 \(M\) 自由平坦，\(N\) 不自由，由局部环上的有限平坦判据也不平坦。两种第零 Fitting 理想定义相同闭集，却不是相同理想，差别正是被剩余域忽略的幂零关系。
\end{solution}

\begin{exercise}\label{b3:ex:12-reduced-constant-fiber}
设 \(R\) 为约化交换环，\(M\) 仅假定有限生成，且纤维维数函数 \(\mu_M(\mathfrak p)\) 局部常值。证明 \(M\) 有限生成投射。这里没有预先假定有限展示，不能直接调用 Fitting 判据。
\end{exercise}
\begin{solution}
固定 \(\mathfrak p\)，设其纤维维数为 \(r\)。提升一组剩余域基，清除分母后可用 \(M\) 中的 \(r\) 个元素定义 \(u:R^r\to M\)，使 \(u_{\mathfrak p}\) 满射，依据是局部 Nakayama 生成元判据。\(\operatorname{coker}u\) 有限生成，故存在 \(f\notin\mathfrak p\)，使 \(u_f\) 满射；再缩小主开邻域，可令其上全部纤维维数均为 \(r\)。

任取 \(z=(z_1,\ldots,z_r)\in\ker u_f\)。对 \(\mathfrak q\in\operatorname{Spec}R_f\)，张量得到从 \(\kappa(\mathfrak q)^r\) 到一个 \(r\) 维空间的满射，故为同构。因此各 \(z_i\) 在 \(\kappa(\mathfrak q)\) 中为零，即属于 \(\mathfrak q\)。遍历所有 \(\mathfrak q\)，约化性给出 \(z_i=0\)。所以 \(u_f\) 还是单射，\(M_f\cong R_f^r\)。模因而在邻域上局部自由且秩有限，\cref{b3:thm:finite-presentation-flat-projective}遂给出有限生成投射性，并同时导出有限展示性。
\end{solution}

\begin{exercise}\label{b3:ex:12-generic-free-integers}
设 \(B=\mathbb Z[x]/(2x^2+x)\)。在一个明确的主开集上给出 \(B\) 的基，并比较一般纤维与特征二的纤维。
\end{exercise}
\begin{solution}
倒置二后，关系化成首一的 \(x^2+x/2\)，多项式除法给出 \(B[1/2]\) 的 \(\mathbb Z[1/2]\) 基 \(1,x\)。一般纤维 \(\mathbb Q[x]/(2x^2+x)\) 维数二，特征二纤维却为 \(\mathbb F_2[x]/(x)\cong\mathbb F_2\)，维数一。因此不可能在包含素理想 \((2)\) 的主开集上使 \(B\) 自由：若自由，其一般纤维将迫使自由秩为二，特征二纤维也须为二维。
\end{solution}

\begin{exercise}\label{b3:ex:12-flat-empty-fiber}
取 \(A=k[t]\)、\(B=A[x]/(tx-1)\)。求各纤维，判断 \(A\to B\) 的平坦性与忠实平坦性。
\end{exercise}
\begin{solution}
关系说明 \(x=t^{-1}\)，所以 \(B\cong A_t\)，为平坦局部化。在 \(t\notin\mathfrak p\) 的点上，纤维为 \(\kappa(\mathfrak p)\)；在 \(\mathfrak p=(t)\) 上，关系变成 \(-1=0\)，纤维为零环。因此不忠实平坦。这表明平坦并不保证每个点都有非空纤维。
\end{solution}

\begin{exercise}\label{b3:ex:12-infinite-generic-rank}
令 \(A=k[t]\)、\(B=A[x]\)、\(M=(t,x)\subseteq B\)。证明 \(M\) 作为 \(A\) 模自由，写出其基；同时证明它作为 \(B\) 模不是局部自由模。
\end{exercise}
\begin{solution}
元素的常数项必须是 \(t\) 的倍数，其余 \(x\) 的系数任意。因此 \(t,x,x^2,\ldots\) 是 \(A\) 基，独立性来自多项式系数唯一性和 \(A\) 为整环。在 \(B\) 的极大理想 \((t,x)\) 处，模的最少生成元数为二，而张量到 \(k(t,x)\) 后秩为一。若在该点自由，自由秩既要等于二，又要在进一步局部化后等于一，矛盾。这里平坦性所相对的底环不同，两个判断不冲突。
\end{solution}

\begin{exercise}\label{b3:ex:12-standard-module-monomials}
在 \(A[x]^2\)、\(A=k[t]\) 中，令 \(N\) 由 \(txe_1-e_2\)、\(xe_2\) 生成。证明倒置 \(t\) 后商模作为 \(A_t\) 模自由，给出一组基并比较 \(t=0\) 的纤维。
\end{exercise}
\begin{solution}
倒置 \(t\) 后，\(e_2=txe_1\)，第二关系变成 \(tx^2e_1=0\)。因此商同构于 \(A_t[x]/(x^2)\)，基为 \(e_1,xe_1\)。在 \(t=0\) 的纤维，第一关系为 \(e_2=0\)，第二关系随之消失，纤维成为 \(k[x]e_1\)，维数无限。倒置的非零系数恰决定哪一个关系可以消去生成元。
\end{solution}

\begin{exercise}\label{b3:ex:12-nonnoether-finite-presentation}
令 \(A=k[t_1,t_2,\ldots]\)、\(B=A[x]/(t_1x^2+x)\)。验证非 Noether 版本泛自由性的假设，并直接找出所需的非零局部化元。
\end{exercise}
\begin{solution}
\(A\) 是整环，虽不是 Noether 环；\(B\) 有一个代数生成元和一个关系，故有限呈示。把 \(M=B\) 当作 \(B\) 模，它秩一自由，当然有限展示。取 \(f=t_1\)，关系除以单位 \(t_1\) 后首一，得到 \(A_f\) 基 \(1,x\)。这直接核实\cref{b3:prop:generic-free-finite-presentation}，没有把 \(A\) 的有限关系问题交给 Noether 性。
\end{solution}

\begin{exercise}\label{b3:ex:12-faithful-field-extension}
证明任意非零 \(k\) 代数 \(S\) 在域 \(k\) 上忠实平坦，不要求 \(S\) 有限维。若两个 \(k\) 线性映射变基到 \(S\) 后相等，说明它们本来就相等。
\end{exercise}
\begin{solution}
\(S\) 作为向量空间有非空基，所以为非零自由 \(k\) 模，张量是若干份直和，保持正合。唯一极大理想为零，剩余纤维就是非零的 \(S\)，由\cref{b3:thm:faithfully-flat-criterion}忠实平坦。两个映射的差 \(u\) 在变基后为零，其像的张量积为零，由平坦性及消失检测有 \(\operatorname{im}u=0\)，故两映射相同。
\end{solution}

\begin{exercise}\label{b3:ex:12-local-flat-not-local}
说明平坦环同态 \(\mathbb Z_{(p)}\to\mathbb Q\) 为什么不忠实平坦，尽管源与目标都是局部环。这与本章的局部同态判据是否矛盾？
\end{exercise}
\begin{solution}
该映射是局部化，所以平坦；但 \(p\mathbb Q=\mathbb Q\)，唯一闭点纤维为零，不忠实。它不是局部同态：源的极大理想 \((p)\) 没有映入目标的极大理想零，因为 \(p\ne0\) 在 \(\mathbb Q\) 中为单位。判据同时要求映射本身局部，故无矛盾。
\end{solution}

\begin{exercise}\label{b3:ex:12-faithful-ideal-contraction}
设 \(R\to S\) 忠实平坦，\(I,J\subseteq R\)。证明 \(IS\subseteq JS\) 当且仅当 \(I\subseteq J\)。若进一步 \(S\) 为 Noether 环，证明 \(R\) 也为 Noether 环。
\end{exercise}
\begin{solution}
正向收缩并用\cref{b3:prop:faithfully-flat-detection}的 \(IS\cap R=I\)、\(JS\cap R=J\) 即得；反向来自扩张保包含。对 \(R\) 的任意理想升链，扩张后在 Noether 环 \(S\) 中稳定，收缩公式使原链也稳定。因此 \(R\) 满足理想升链条件，为 Noether 环。
\end{solution}

\begin{exercise}\label{b3:ex:12-faithful-free-basis-bound}
设 \(R\to S\) 忠实平坦，\(M\) 为任意 \(R\) 模，且 \(S\otimes_RM\cong S^n\)。证明 \(M\) 有限生成投射，且局部秩处处为 \(n\)。
\end{exercise}
\begin{solution}
由\cref{b3:prop:faithfully-flat-detection}，有限展示和平坦性均下降，故\cref{b3:thm:finite-presentation-flat-projective}给出有限生成投射性。对 \(\mathfrak p\subset R\)，选位于其上的 \(\mathfrak q\subset S\)，存在性由谱满射保证。\(M_{\mathfrak p}\) 自由，设秩为 \(r\)，再变到 \(\kappa(\mathfrak q)\) 后既为 \(r\) 维又为 \(n\) 维，故 \(r=n\)。这里只得到局部自由，不保证一个全局基。
\end{solution}

\begin{exercise}\label{b3:ex:12-two-principal-cover}
证明 \(\mathbb Z\to\mathbb Z[1/2]\times\mathbb Z[1/3]\) 忠实平坦。由此说明一个交换群及其同态的哪些性质可在这两个局部化后同时检测。
\end{exercise}
\begin{solution}
因 \((2,3)=\mathbb Z\)，\(D(2),D(3)\) 覆盖素谱。两个局部化平坦，有限直积仍平坦；每个素理想至少不含二或三，谱映射满射，故忠实平坦。由\cref{b3:prop:faithfully-flat-detection}，群的零性、有限生成性、有限展示性和平坦性可以检测；同态的单射、满射、同构及序列正合性也可以检测。只检查其中一个局部化不够，例如二阶群在倒置二后消失。
\end{solution}

\begin{optionalexercise}\label{b3:ex:12-equalizer-integers}
在 \(\mathbb Q\) 中直接证明 \(\mathbb Z[1/2]\cap\mathbb Z[1/3]=\mathbb Z\)，并据此写出上题开覆盖中标准下降数据的等化子。
\end{optionalexercise}
\begin{solution}
若一个有理数同时可写成分母为二的幂及三的幂的分数，则其既约分母同时整除这两种幂，只能为一。等化子因而是
\[
\bigl\{(a,b)\in\mathbb Z[1/2]\times\mathbb Z[1/3]:
a=b\text{ 在 }\mathbb Z[1/6]\text{ 中}\bigr\}
=\bigl\{(n,n):n\in\mathbb Z\bigr\}.
\]
同一开集与自身的重叠只给出恒等条件，两个不同开集的重叠给出上式中的等式。这是\cref{b3:lem:faithful-flat-equalizer}在这个覆盖上的具体计算。
\end{solution}

\begin{optionalexercise}\label{b3:ex:12-module-without-descent}
令 \(R=k\)、\(S=k\times k\)，环映射取对角嵌入。证明 \(S\) 模 \(N=k\times0\) 不可能带有下降数据。
\end{optionalexercise}
\begin{solution}
\(S\) 是非零自由 \(k\) 模，所以忠实平坦。设存在数据，\cref{b3:thm:faithfully-flat-module-descent}将给出 \(N\cong S\otimes_kM\)。而任一这样的模在两个幂等分量上都是 \(M\)，不可能一个分量为一维、另一个为零。也可直接检查双重基变换：\(S\otimes_kS\cong k^4\)，在第一位置为第一个分量、第二位置为第二个分量的那一因子上，\(N_1\) 维数一而 \(N_2\) 维数零，所以根本不存在 \(\theta\) 同构。
\end{solution}

\begin{optionalexercise}\label{b3:ex:12-real-semilinear}
设复向量空间 \(N\) 有半线性对合 \(\sigma\)，即 \(\sigma(cn)=\bar c\cdot\sigma(n)\)、\(\sigma^2=1\)。令 \(M=N^\sigma\)。不用一般下降定理，直接证明 \(N=M\oplus iM\) 作为实向量空间，并得到 \(\mathbb C\otimes_{\mathbb R}M\cong N\)。
\end{optionalexercise}
\begin{solution}
任取 \(n\)，有分解
\[
n=\frac{n+\sigma(n)}2+i\frac{n-\sigma(n)}{2i}.
\]
两项系数都被 \(\sigma\) 固定。若 \(u=iv\) 且 \(u,v\in M\)，施加 \(\sigma\) 得 \(u=-iv\)，故 \(2iv=0\)，于是 \(u=v=0\)。所以直和成立；复线性映射 \(z\otimes m\mapsto zm\) 在实分解 \(\mathbb C=\mathbb R\oplus i\mathbb R\) 下正是上述直和同构。
\end{solution}

\begin{optionalexercise}\label{b3:ex:12-cocycle-change-bases}
设有限主开覆盖上每个局部模自由秩为 \(r\)，在重叠处以矩阵 \(G_{ij}\in\operatorname{GL}_r(R_{f_if_j})\) 表示从第 \(i\) 个坐标到第 \(j\) 个坐标的识别。写出余循环条件和改变各局部基后的矩阵公式。证明若可将所有矩阵同时变成单位矩阵，则下降模自由。
\end{optionalexercise}
\begin{solution}
坐标约定为 \(v_j=G_{ij}v_i\)，所以三重重叠上必须有 \(G_{jk}G_{ij}=G_{ik}\)。若第 \(i\) 个新基的列在旧基中组成矩阵 \(P_i\)，则 \(v_i=P_iv_i'\)，于是新矩阵为 \(G'_{ij}=P_j^{-1}G_{ij}P_i\)。若它们全部为单位，匹配的局部向量逐坐标相等；每个坐标由\cref{b3:lem:faithful-flat-equalizer}唯一来自 \(R\)。等化子因此为 \(R^r\)，即下降模自由。一般余循环只保证相容，并不保证存在这些 \(P_i\)。
\end{solution}

\begin{optionalexercise}\label{b3:ex:12-descent-morphisms}
设 \(R\to S\) 忠实平坦，\(M,M'\) 为 \(R\) 模。证明一个 \(S\) 线性映射 \(u:S\otimes_RM\to S\otimes_RM'\) 来自 \(R\) 线性映射，当且仅当它与标准下降数据相容。
\end{optionalexercise}
\begin{solution}
若 \(u=1\otimes v\)，把 \((s\otimes m)\otimes t\) 代入标准数据公式，两个次序都得到 \(s\otimes\bigl(t\otimes v(m)\bigr)\)，故相容。反向，相容性使 \(u\) 将标准数据的等化子送入另一个等化子。由\cref{b3:lem:faithful-flat-equalizer}，这些等化子分别是 \(M,M'\) 的自然像，因而得到唯一的 \(R\) 线性映射 \(v:M\to M'\)。\(S\) 线性和 \(1\otimes m\) 的生成性使 \(u=1\otimes v\)。唯一性也可直接由 \(M'\hookrightarrow S\otimes_RM'\) 的单射性验证。
\end{solution}

% !TeX root = ../../Book3.tex
% 纲要标题：## 第13章　滤过、Artin–Rees 与完备化
% 纲要位置：工作记录/计划纲要.md，第 2335 行。

\chapter{滤过、Artin–Rees 与完备化}
\label{b3:ch:13}
\BookChapterNumber{b3:ch:13}

\begin{chapterabstract}
本章把环与模的完备化作为理想幂商的逆极限来构造。Rees 环的有限性导出 Artin–Rees 引理和 Jacobson 根条件下的 Krull 交定理；随后研究有限生成模的完备化何时保持正合，并讨论完备环的 Noether 性、忠实平坦性与局部维数。
\end{chapterabstract}

\begin{learninggoals}
\begin{itemize}
\item 比较子模自身的理想进滤过与环境诱导的滤过，能指出 Artin–Rees 中有限性的作用。
\item 用相容余类构造完备化，解释逆极限一般不保持满射的原因。
\item 由完备化正合性和张量表示证明平坦性，并用剩余域检测忠实性。
\item 计算典型的幂级数完备化，并区分有限阶相等、形式相等和解析收敛。
\end{itemize}
\end{learninggoals}

整数的序列 \(1,1+2,1+2+4,\ldots\) 在通常绝对值下越走越远；模 \(2^n\) 看，第 \(n\) 步以后它却不再改变，且所得相容余数都等于 \(-1\) 的余数。这种收敛取决于理想的幂，而不是实数距离。把所有模 \(I^n\) 的相容答案合在一起，便是 \(I\) 进完备化的基本构造。

\ProseParagraphBreak

对一个模作同样的操作时，逐层商的正合性未必自动传到无限相容族；取子模后所诱导的滤过，也未必显然等于原子模自身的理想幂滤过。Artin–Rees 引理正是控制这两个滤过的差距。弄清这个有限性步骤，才能说明有限生成模的完备化何时保持正合，以及完备局部环为何仍服从 Noether 性。

\ProseParagraphBreak

本章使用\chapref{b3:ch:03}的 Noether 模与 Nakayama 引理、\secref{b3:sec:1104}的 Hilbert–Samuel 理论，以及\secref{b3:sec:1201}、\secref{b3:sec:1205}的平坦和忠实平坦判据。逆极限在本章从相容族定义。除非另有说明，环均交换且有单位。

\ProseParagraphBreak
Artin–Rees、有限生成模的张量表示及完备化性质可对照松村英之的《Commutative Ring Theory》\cite[\S8, Theorems 8.5--8.14, pp. 59--62]{Matsumura1989CommutativeRingTheory}。该书的证明将滤过与有限生成相联系，本章同时展开相容提升和逐阶消去的步骤，以说明完备性在何处真正起作用。


% !TeX root = ../../Book3.tex
% 纲要标题：### 13.1 理想进滤过
% 纲要位置：工作记录/计划纲要.md，第 2359 行。
% 纲要细目（写作范围）：
% * I 进拓扑与分离性
% * 模滤过和伴随分次
% * 相容余类族与滤过完备化；稳定滤过和理想幂滤过之间的自然同构
% * Rees 环与有限生成；理想幂滤过的 Rees 环及模模掉 \(I\) 后得到伴随分次对象，一般稳定滤过只保证在充分高次满足相应等式

\section{理想进滤过}
\label{b3:sec:1301}


多项式只含有限个系数，形式幂级数却要同时保存任意高阶的系数。把一个对象模掉 \(I,I^2,I^3,\ldots\)，就是逐次保留更多信息的代数方法。本节先说明“两个元素相差很小”的含义；这里的“小”由理想的幂衡量，不依赖绝对值。

\subsection{\texorpdfstring{\(I\)}{I} 进拓扑与分离性}
\label{b3:subsec:1301:01}

\begin{definition}\label{b3:def:adic-topology}
设 \(R\) 为交换环、\(I\subseteq R\) 为理想、\(M\) 为 \(R\) 模。以
\[
x+I^nM\qquad(x\in M,\ n\ge0)
\]
为邻域基，得到 \(M\) 的 \term[I-jin-tuopu]{\(I\) 进拓扑}[I-adic topology]：一个集合是开集，是指其中每个点都包含在某个仍位于该集合内的上述邻域中。若
\(\bigcap_{n\ge0}I^nM=0\)，则称这个拓扑\term[fenli-tuopu]{分离}[separated]。对 \(M\) 中的序列 \((x_j)\) 及元素 \(x\in M\)，若对每个 \(n\ge0\)，充分大的 \(j\) 都有 \(x_j-x\in I^nM\)，则称 \(x_j\) 收敛于 \(x\)；若对每个 \(n\ge0\)，充分大的 \(i,j\) 都有 \(x_i-x_j\in I^nM\)，则称 \((x_j)\) 为 Cauchy 序列\footnote{柯西（Augustin-Louis Cauchy，1789-1857），法国数学家，推动极限、收敛与复分析的严格化。}。
\end{definition}

由于 \(I^{n+1}M\subseteq I^nM\)，这些集合确实给出拓扑。商 \(M/I^nM\) 把相差 \(I^nM\) 中元素的两个值看作相同；自然映射 \(M/I^{n+1}M\to M/I^nM\) 则将较高一阶的精度降回第 \(n\) 阶。因此 \(x_j\to x\) 表示：对每个固定的 \(n\)，充分大的 \(j\) 都使 \(x_j\) 与 \(x\) 在 \(M/I^nM\) 中完全相同。加法与取负连续；\(R\) 自身的 \(I\) 进拓扑还使乘法连续，因为 \(I^n\) 是理想。

若拓扑分离，\(x\ne y\) 时就有某个 \(n\) 使 \(x-y\notin I^nM\)。此时两个邻域 \(x+I^nM\)、\(y+I^nM\) 不交，否则交点的两个表示会给出 \(x-y\in I^nM\)。反过来，若 \(0\ne x-y\in\bigcap_n I^nM\)，这两个点在每个商中都相同，不能用开邻域分开。因此分离性等价于 Hausdorff 性，也恰好表示各阶商合起来能够区分原来的不同元素。

\ProseParagraphBreak
设 \(k\) 为域。在 \(k[x]\) 中，\((x)\) 进收敛表示每个固定次数的系数最终不再改变。非零多项式不能被任意高次的 \(x\) 整除，所以此拓扑分离，收敛序列的极限至多有一个。但 \(1+x+\cdots+x^j\) 是 Cauchy 序列，任何极限的每个系数却都必须是一，不可能是多项式。它在 \(k[x]\) 内没有极限，在 \(k[[x]]\) 中才得到极限 \(\sum_{i\ge0}x^i\)。分离性保证已有极限的唯一性，补足这些极限则是另一个问题。另取 \(R=k\times k\)、\(I=k\times0\)，则 \(I^n=I\) 对所有 \(n\ge1\) 成立，拓扑不分离；第二个坐标以外的信息已经在第一个商中被永久遗失。

\begin{proposition}\label{b3:prop:adic-closure}
设 \(R\) 为交换环、\(I\subseteq R\) 为理想、\(M\) 为 \(R\) 模，并在 \(M\) 上取 \(I\) 进拓扑。子模 \(N\subseteq M\) 的闭包为
\[
\overline N=\bigcap_{n\ge0}(N+I^nM).
\]
特别地，\(N\) 闭当且仅当 \(M/N\) 的 \(I\) 进拓扑分离。对任意 \(R\) 模 \(P\)，每个 \(R\) 线性映射 \(f:M\to P\) 在两模的 \(I\) 进拓扑下都连续。
\end{proposition}
\begin{proof}
点 \(x\) 属于闭包，恰指对每个 \(n\)，邻域 \(x+I^nM\) 都与 \(N\) 相交。也就是可以选 \(y_n\in N\)，使 \(x-y_n\in I^nM\)：在任意给定精度下，都能用 \(N\) 中的元素逼近 \(x\)。这等价于 \(x\in N+I^nM\) 对每个 \(n\) 成立，得到闭包公式。又
\(I^n(M/N)=(I^nM+N)/N\)，得到闭性的刻画。若 \(f:M\rightarrow P\) 线性，则 \(f(I^nM)\subseteq I^nP\)，所以原点处连续；平移后处处连续。
\end{proof}

\subsection{模滤过与伴随分次}
\label{b3:subsec:1301:02}

\begin{definition}\label{b3:def:good-adic-filtration}
设 \(R\) 为交换环、\(M\) 为 \(R\) 模。下降的子模列
\(M=F^0M\supseteq F^1M\supseteq\cdots\) 称为一个\term[mo-lvguo]{模滤过}[module filtration]。给定理想 \(I\subseteq R\)，若 \(IF^nM\subseteq F^{n+1}M\) 对每个 \(n\ge0\) 成立，则称它为 \(I\) 滤过；若还有 \(IF^nM=F^{n+1}M\) 对充分大的 \(n\) 成立，则称它为\term[wending-lvguo]{稳定滤过}[stable filtration]。对 \(I\) 滤过，相应伴随分次模为
\[
\operatorname{gr}_F(M)=\bigoplus_{n\ge0}F^nM/F^{n+1}M.
\]
它是 \(\operatorname{gr}_I(R)\) 模，作用由代表元相乘定义。
\end{definition}

良定义需要同时检查两个代表元：若 \(a\in I^r\)、\(x\in F^nM\)，更换 \(a\) 增加的项属于 \(I^{r+1}F^nM\subseteq F^{n+r+1}M\)，更换 \(x\) 增加的项也属于此子模。\cref{b3:def:associated-graded-adic}中的滤过 \(F^nM=I^nM\) 是从第一步就稳定的特例。

\begin{definition}\label{b3:def:filtered-completion}
设 \(R\) 为交换环、\(M\) 为 \(R\) 模，\(M=F^0M\supseteq F^1M\supseteq\cdots\) 为下降滤过。对 \(n\ge1\)，记 \(q^F_n:M/F^{n+1}M\rightarrow M/F^nM\) 为自然商映射。该滤过的相容余类族组成 \(R\) 模
\[
\widehat M_F\coloneqq
\left\{\,(x_n)_{n\ge1}\in\prod_{n\ge1}M/F^nM
\mid q^F_n(x_{n+1})=x_n\text{ 对所有 }n\ge1\,\right\},
\]
加法和标量作用逐坐标进行。这个相容余类模也记为
\(\varprojlim_{n\ge1}M/F^nM\)。对理想 \(I\subseteq R\) 的幂滤过 \(F^nM=I^nM\)，将它记为 \(\widehat M_I\)。自然映射 \(M\rightarrow\widehat M_F\) 把 \(x\) 送到 \((x+F^nM)_n\)。
\end{definition}

伴随分次的第 \(n\) 片 \(F^nM/F^{n+1}M\) 记录这一层相对于下一层留下的信息；相容余类模则保留每个完整的截断 \(M/F^nM\)，以及高阶截断如何映到低阶截断。比较两套滤过时，若它们的精度只相差固定的有限阶数，要求达到任意高阶精度就会给出同一拓扑，并使两个相容余类模自然同构。

\begin{proposition}\label{b3:prop:stable-filtration-equivalent}
设 \(R\) 为交换环、\(I\subseteq R\) 为理想、\(M\) 为 \(R\) 模，\(F^nM\) 为从次数 \(c\ge0\) 起稳定的 \(I\) 滤过。则对 \(n\ge c\) 有
\[
F^nM=I^{n-c}F^cM,\qquad
I^nM\subseteq F^nM\subseteq I^{n-c}M.
\]
因而 \(F^nM\) 与 \(I^nM\) 定义同一拓扑，恒等映射 \(M\rightarrow M\) 还诱导自然同构
\[
\widehat M_I=\varprojlim_{n\ge1}M/I^nM
\xrightarrow{\ \sim\ }\varprojlim_{n\ge1}M/F^nM=\widehat M_F,
\]
它与从 \(M\) 到两个相容余类模的自然映射相容。
\end{proposition}
\begin{proof}
第一个等式由稳定性归纳得到。由 \(F^0M=M\) 及 \(IF^nM\subseteq F^{n+1}M\)，归纳得 \(I^nM\subseteq F^nM\) 对所有 \(n\ge0\) 成立；右侧包含由第一个等式得到。要使差值落入 \(F^nM\)，只需使它落入 \(I^nM\)；要使差值落入 \(I^nM\)，则使它落入 \(F^{n+c}M\) 即可。这就使两个邻域系相互共尾，拓扑相同。

对每个 \(n\ge1\)，包含 \(I^nM\subseteq F^nM\) 及 \(F^{n+c}M\subseteq I^nM\) 给出自然商映射
\[
\alpha_n:M/I^nM\longrightarrow M/F^nM,
\qquad
\beta_n:M/F^{n+c}M\longrightarrow M/I^nM.
\]
它们与各自的过渡商映射交换，所以逐坐标定义
\[
\Phi:\widehat M_I\longrightarrow\widehat M_F,
\quad (x_n)_n\longmapsto\bigl(\alpha_n(x_n)\bigr)_n,
\]
\[
\Psi:\widehat M_F\longrightarrow\widehat M_I,
\quad (y_n)_n\longmapsto\bigl(\beta_n(y_{n+c})\bigr)_n
\]
确实得到相容余类族。复合 \(\beta_n\alpha_{n+c}\) 是从 \(M/I^{n+c}M\) 到 \(M/I^nM\) 的过渡商映射，故 \((\Psi\Phi(x))_n=x_n\)。复合 \(\alpha_n\beta_n\) 是从 \(M/F^{n+c}M\) 到 \(M/F^nM\) 的过渡商映射，故 \((\Phi\Psi(y))_n=y_n\)。因此两映射互逆。

若将 \(c\) 增大，相容性保证 \(\Psi\) 的每个坐标不变。这些商映射都由 \(M\) 的恒等映射诱导，故所得到的同构与自然余类映射相容，也与保持滤过的模同态相容。
\end{proof}

伴随分次未必恢复原环。设 \(p\) 为素数，取 \(A=\mathbb Z/p^2\mathbb Z\)、\(B=\mathbb F_p[t]/(t^2)\)。它们的极大理想分别为 \((p)\) 与 \((t)\)，平方都为零，次数零的片都是 \(\mathbb F_p\)，次数一的片分别由 \(p\)、\(t\) 的类生成。因此两个伴随分次环都为 \(\mathbb F_p[T]/(T^2)\)。但在 \(A\) 中，\(p\) 个单位元之和为 \(p\ne0\)，落入下一层；在 \(B\) 中，\(p\) 个单位元之和已经为零。在伴随分次的次数零部分，这两个和都为零，因而看不到前一种跨到下一层的加法关系，原环的特征仍可不同。

\subsection{Rees 环与有限生成}
\label{b3:subsec:1301:03}

\begin{definition}\label{b3:def:rees-algebra-module}
设 \(R\) 为交换环、\(I\subseteq R\) 为理想、\(M\) 为 \(R\) 模，\(F^nM\) 为 \(I\) 滤过。引入形式变量 \(t\)，定义 \term[rees-huan]{Rees 环}[Rees algebra]与 Rees 模
\[
\mathcal R_I(R)=\bigoplus_{n\ge0}I^nt^n\subseteq R[t],\qquad
\mathcal R_F(M)=\bigoplus_{n\ge0}F^nM\,t^n.
\]
后者的作用为 \((at^r)(xt^n)=(ax)t^{r+n}\)。这只是给各阶加上标记；不要求原模中的元素本来具有次数。对幂滤过 \(F^nM=I^nM\)，将相应 Rees 模记为 \(\mathcal R_I(M)\)。
\symidx{\(\mathcal R_I(R)\)：理想的 Rees 环}
\symidx{\(\mathcal R_F(M)\)：滤过的 Rees 模}
\end{definition}

\cref{b3:lem:induced-filtration-bound}的证明已用这种构造同时处理全部次数。作为 Rees 环上的模，有限生成要求所有层都由有界次数内的有限组齐次元素产生；只知道每一层分别有限生成，还不能得到这样一组共同生成元。

\begin{proposition}\label{b3:prop:rees-finiteness}
设 \(R\) 为交换环、\(I\subseteq R\) 为理想、\(M\) 为 \(R\) 模。幂滤过的 Rees 环及模有自然的分次同构
\[
\mathcal R_I(R)/I\mathcal R_I(R)\cong\operatorname{gr}_I(R),
\qquad
\mathcal R_I(M)/I\mathcal R_I(M)\cong\operatorname{gr}_I(M),
\]
后一个同构的模作用与前一个环同构相容。若 \(R\) 为 Noether 环，则 \(\mathcal R_I(R)\) 为 Noether 环；若还有 \(M\) 有限生成，则对 \(M\) 上的 \(I\) 滤过，\(\mathcal R_F(M)\) 作为 \(\mathcal R_I(R)\) 模有限生成当且仅当滤过稳定。
\end{proposition}
\begin{proof}
理想 \(I\) 放在 Rees 环的次数零部分时，\(I\mathcal R_I(R)\) 的第 \(n\) 片为 \(I^{n+1}t^n\)，故商的第 \(n\) 片为 \(I^n/I^{n+1}\)。逐片把 \(at^n\) 的类送到 \(a+I^{n+1}\)，乘法与代表元相乘相容，就得到第一个同构。同样，\(I\mathcal R_I(M)\) 的第 \(n\) 片为 \(I^{n+1}M\,t^n\)，逐片取商得到第二个同构，且环对模的作用也由相同的代表元乘法给出。

以下设 \(R\) 为 Noether 环。写 \(I=(a_1,\ldots,a_s)\)，则 \(\mathcal R_I(R)=R[a_1t,\ldots,a_st]\)，由 Hilbert 基定理与商环的 Noether 性得到 Rees 环的 Noether 性。

再设 \(M\) 有限生成。若滤过从 \(c\) 起稳定，选 \(F^0M,\ldots,F^cM\) 各自的有限组生成元；这些子模有限生成是因为 \(M\) 为 Noether 模。把这些生成元放入相应次数后，由 \(F^nM=I^{n-c}F^cM\) 可知它们生成全部 Rees 模。

反向可取齐次生成元 \(x_jt^{d_j}\)：将任意有限组生成元分解成齐次部分，所得有限组仍生成。设 \(c=\max_jd_j\)，则
\[
F^nM=\sum_jI^{n-d_j}x_j\quad(n\ge c).
\]
于是 \(F^{n+1}M=IF^nM\) 对 \(n\ge c\) 成立。零模的情形直接成立。
\end{proof}

一般 \(I\) 滤过的 Rees 模模去 \(I\) 后，第 \(n\) 片是 \(F^nM/IF^nM\)，而伴随分次的第 \(n\) 片是 \(F^nM/F^{n+1}M\)。\(I\) 滤过只要求 \(IF^nM\subseteq F^{n+1}M\)，两种商未必相同。即使滤过稳定，也只能保证稳定以后的各次满足等式，不能据此将整个 Rees 商模认同为伴随分次模。

% !TeX root = ../../Book3.tex
% 纲要标题：### 13.2 Artin–Rees 引理
% 纲要位置：工作记录/计划纲要.md，第 2344 行。
% 纲要细目（写作范围）：
% * 子模滤过与诱导滤过
% * Artin–Rees 引理及证明
% * Krull 交定理与子模闭性：Jacobson 根条件及局部环特例

\section{Artin–Rees 引理}
\label{b3:sec:1302}


在子模 \(N\subseteq M\) 内，“\(I\) 的 \(n\) 次幂乘出的元素”组成 \(I^nN\)；“在 \(M\) 中已经达到第 \(n\) 阶的元素”组成 \(N\cap I^nM\)。这两者一般不相等。Artin–Rees 引理说明，在 Noether 条件下，它们的差异能由一个固定的次数偏移控制。

\subsection{子模滤过与诱导滤过}
\label{b3:subsec:1302:01}

设 \(k\) 为域，取 \(R=k[x]\)、\(I=(x)\)、\(M=R\)、\(N=x^rR\)，其中 \(r\ge1\) 为整数。直接计算得到
\[
I^nN=x^{n+r}R,\qquad N\cap I^nM=x^{\max(r,n)}R.
\]
后一滤过到第 \(r\) 阶仍等于整个 \(N\)，前一滤过已经缩小。因此子模的自身理想进滤过和由环境诱导的滤过需要分别处理。

\ProseParagraphBreak
令 \(F^nN=N\cap I^nM\)。它满足 \(F^0N=N\) 及 \(IF^nN\subseteq F^{n+1}N\)，且其 Rees 模是 \(\mathcal R_I(M)\) 的齐次子模。这个观察把滤过比较转化为一个有限生成问题。

\subsection{Artin–Rees 引理及其证明}
\label{b3:subsec:1302:02}

下述\term[artin-rees-yinli]{Artin–Rees 引理}[Artin--Rees lemma]给出诱导滤过最终稳定的精确等式。

\begin{theorem}[Artin–Rees]\label{b3:thm:artin-rees}
设 \(R\) 为交换 Noether 环、\(I\subseteq R\) 为理想、\(M\) 为有限生成 \(R\) 模、\(N\subseteq M\) 为子模。存在整数 \(c\ge0\)，使对所有 \(n\ge c\) 都有
\[
N\cap I^nM=I^{n-c}(N\cap I^cM).
\]
特别地，\(N\) 自身的 \(I\) 进拓扑等于 \(M\) 在 \(N\) 上诱导的拓扑。
\end{theorem}
\begin{proof}
因为 \(R\) 为 Noether 环、\(M\) 有限生成，子模 \(N\) 也是有限生成 \(R\) 模。由\cref{b3:prop:rees-finiteness}，\(\mathcal R_I(R)\) 为 Noether 环，\(\mathcal R_I(M)\) 为有限生成模，所以齐次子模
\[
\mathcal N=\bigoplus_{n\ge0}(N\cap I^nM)t^n
\]
有限生成。由于 \(N\) 有限生成，可再次应用该命题，得到滤过 \(N\cap I^nM\) 从某个次数 \(c\) 起稳定，逐次使用等式便得到所述公式。

具体看齐次生成元，可取 \(\mathcal N\) 的有限组生成元 \(z_jt^{d_j}\) 并增大 \(c\)，使 \(d_j\le c\)。对 \(n\ge c\)，齐次生成性给出
\[
N\cap I^nM=\sum_jI^{n-d_j}z_j.
\]
因为 \(z_j\in N\cap I^{d_j}M\)，有
\[
I^{n-d_j}z_j=I^{n-c}\bigl(I^{c-d_j}z_j\bigr),
\qquad I^{c-d_j}z_j\subseteq N\cap I^cM.
\]
故上面的和包含于 \(I^{n-c}(N\cap I^cM)\)。反向包含则由 \(N\) 为子模及 \(I^{n-c}I^cM\subseteq I^nM\) 得到，这证明所需的精确等式。

最后，
\[
I^nN\subseteq N\cap I^nM\subseteq I^{n-c}N\qquad(n\ge c).
\]
两个邻域系相互包含任意给定的足够小邻域，故拓扑相同。
\end{proof}

\cref{b3:lem:induced-filtration-bound}只需要上述两个包含关系。现在得到的等式还说明诱导滤过何时开始逐次由乘 \(I\) 生成。这里的 \(c\) 可以依赖 \(N\subseteq M\) 和 \(I\)，不能解释为对所有子模统一的常数；前面 \(N=x^rR\) 的例子中，最小的 \(c\) 就是 \(r\)。

\subsection{Krull 交定理与子模闭性}
\label{b3:subsec:1302:03}

\term[krull-jiao-dingli]{Krull 交定理}[Krull intersection theorem]以 Jacobson 根条件保证有限生成模的分离性。

\begin{theorem}[Krull 交定理]\label{b3:thm:krull-intersection}
设 \(R\) 为交换 Noether 环、\(M\) 为有限生成 \(R\) 模、\(I\subseteq R\) 为理想，\(J(R)\) 为 \(R\) 的 Jacobson 根。若 \(I\subseteq J(R)\)，则
\[
\bigcap_{n\ge0}I^nM=0.
\]
此时 \(M\) 的每个子模都闭。特别地，对 Noether 局部环 \((R,\mathfrak m)\)，每个有限生成模在极大理想进拓扑下分离。
\end{theorem}
\begin{proof}
令 \(N=\bigcap_nI^nM\)。它作为 Noether 模的子模有限生成。对 \(N\subseteq M\) 使用\cref{b3:thm:artin-rees}，在 \(n=c+1\) 时得到
\[
N=N\cap I^{c+1}M=I(N\cap I^cM)=IN.
\]
由\cref{b3:thm:nakayama}及 \(I\subseteq J(R)\) 得 \(N=0\)。对任意子模 \(L\)，商模 \(M/L\) 也有限生成，因而分离；\cref{b3:prop:adic-closure}说明 \(L\) 闭。
\end{proof}

有限性不能删去。设 \(p\) 为素数，取 \(R=\mathbb Z_{(p)}\)、\(M=\mathbb Q\)，则 \(p^nM=M\) 对所有 \(n\) 成立，交并不为零。Jacobson 根条件也不能对一般环省略：若 \(k\) 为域，\(k\times k\) 的幂等理想 \(k\times0\) 的各次幂都相同。

\ProseParagraphBreak
本节还给出一种实用判断：设 \((R,\mathfrak m)\) 为 Noether 局部环，\(M\) 为有限生成 \(R\) 模，\(N\subseteq M\) 为子模，\(x\in M\)。若对每个 \(n\ge0\) 都存在 \(y_n\in N\)，使 \(x\equiv y_n\pmod{\mathfrak m^nM}\)，那么 \(x\in N\)，因为 \(N\) 闭。

% !TeX root = ../../Book3.tex
% 纲要标题：### 13.3 完备化
% 纲要位置：工作记录/计划纲要.md，第 2350 行。
% 纲要细目（写作范围）：
% * 逆极限定义
% * 形式幂级数环
% * Noether 情形下有限生成模的完备化正合性
% * 商、有限生成模与张量表示
% * 完备环的 Noether 性；初始形式的有限生成与系数部分和的收敛消去
% * 在完备化局部性质之前证明所需的 Noether 性、完备性与有限生成模的张量表示；把任意逆极限的左正合性与本章有限生成模的完备化正合性区分开

\section{完备化}
\label{b3:sec:1303}


只给出每个精度下的一个答案还不够：高精度答案必须在取商后等于先前的低精度答案。完备化把这种相容性直接写入定义。它首先是一个由商模组成的逆极限；其正合性与 Noether 性则需要另外证明。

\subsection{完备化的逆极限定义}
\label{b3:subsec:1303:01}

\begin{definition}\label{b3:def:inverse-limit-completion}
设 \(R\) 为交换环，\(M_1\xleftarrow{u_1}M_2\xleftarrow{u_2}\cdots\) 为 \(R\) 模及 \(R\) 线性过渡映射组成的逆系统。其\term[niji-xian]{逆极限}[inverse limit]定义为
\[
\varprojlim_nM_n=
\left\{\,(x_n)\in\prod_{n\ge1}M_n\mid u_n(x_{n+1})=x_n\text{ 对所有 }n\,\right\}.
\]
加法与标量作用逐坐标进行。给定理想 \(I\subseteq R\) 与 \(R\) 模 \(M\)，取 \(M_n=M/I^nM\) 和自然商映射，得到 \(M\) 的 \term[wanbeihua]{完备化}[completion]
\[
\widehat M=\varprojlim_n M/I^nM.
\]
自然映射 \(\iota_M:M\rightarrow\widehat M\) 把 \(x\) 送到 \((x+I^nM)_n\)。若此映射为同构，则称 \(M\) 为 \(I\) 进完备模；此处“完备”包含分离性。
\symidx{\(\varprojlim_nM_n\)：模逆系统的逆极限}
\symidx{\(\widehat R\)、\(\widehat M\)：指定理想进拓扑下的完备化}
\end{definition}

逆极限的泛性质可以直接核对：从模 \(P\) 到各 \(M_n\) 的相容映射 \(f_n\)，唯一合成映射 \(p\mapsto\bigl(f_n(p)\bigr)_n\)。另外，\(\ker\iota_M=\bigcap_nI^nM\)。因此自然映射单射所需的正是分离性。

\ProseParagraphBreak
记 \(K_n=\ker(\widehat M\rightarrow M/I^nM)\)。投影 \(\widehat M\rightarrow M/I^nM\) 满射，因为每个余类都可由 \(\iota_M(x)\) 的第 \(n\) 个坐标实现。因此投影诱导的同构
\[
\widehat M/K_n\xrightarrow{\sim}M/I^nM
\]
与过渡商映射相容。在这些同构下，规范映射
\[
\widehat M\longrightarrow\varprojlim_n\widehat M/K_n,
\qquad z\longmapsto(z+K_n)_n
\]
正是 \(\widehat M=\varprojlim_nM/I^nM\) 的恒等映射，故为同构；又 \(\bigcap_nK_n=0\)。

以 \(K_n\) 为邻域基时，这也等价于每个 Cauchy 序列唯一收敛：在每个商中，Cauchy 序列的坐标最终固定，这些固定坐标相容，因而给出唯一极限。反过来，每个相容余类系统逐阶选取代表元便得到 Cauchy 序列，其极限还原这个系统；极限唯一对应分离性。因此这里的逆极限完备性与序列完备性一致。自然映射 \(\iota_M\) 的像稠密，因为任意第 \(n\) 个坐标都可由 \(M\) 中的元素表示。此处的邻域基由投影核 \(K_n\) 构成。

\ProseParagraphBreak
环 \(\widehat R\) 的加法、乘法逐商定义，\(\widehat M\) 是 \(\widehat R\) 模，因为 \((a_n)(m_n)=(a_nm_n)\) 保持相容性。模同态按各阶取商得到完备化同态，并保持恒等映射与复合。

\subsection{形式幂级数环}
\label{b3:subsec:1303:02}

\begin{definition}\label{b3:def:formal-power-series}
设 \(A\) 为交换环，\(r\ge0\) 为整数，\(x_1,\ldots,x_r\) 为形式变量。对 \(\alpha=(\alpha_1,\ldots,\alpha_r)\in\mathbb N^r\)，记 \(x^\alpha=x_1^{\alpha_1}\cdots x_r^{\alpha_r}\)。给定系数族 \((a_\alpha)_{\alpha\in\mathbb N^r}\)，其中 \(a_\alpha\in A\)，把它记为形式和
\(\sum_{\alpha\in\mathbb N^r}a_\alpha x^\alpha\)。全体这样的形式和逐系数相加；对另一个系数族 \((b_\beta)_{\beta\in\mathbb N^r}\)、\(b_\beta\in A\)，按下式相乘：
\[
\left(\sum_\alpha a_\alpha x^\alpha\right)
\left(\sum_\beta b_\beta x^\beta\right)
=\sum_\gamma\left(\sum_{\alpha+\beta=\gamma}a_\alpha b_\beta\right)x^\gamma.
\]
对固定的多重指标 \(\gamma\)，内层和只有有限项，故不需要分析意义的收敛。所得环称为\term[xingshi-miji-shu-huan]{形式幂级数环}[formal power series ring]，记为 \(A[\![x_1,\ldots,x_r]\!]\)。
\end{definition}

交换律、结合律及分配律可在每个固定系数上归结为有限和的同一运算，常数级数 \(1\) 为单位。全次数小于 \(n\) 的截断给出同构
\begin{equation}\label{b3:eq:power-series-limit}
A[\![x_1,\ldots,x_r]\!]\cong
\varprojlim_n A[x_1,\ldots,x_r]/(x_1,\ldots,x_r)^n.
\end{equation}
相容截断唯一指定每个单项式的系数，反过来每个系数族给出相容截断，因此此式既证明双射，也说明环运算为何相同。

\begin{proposition}\label{b3:prop:formal-series-unit}
设 \(A\) 为交换环，\(r\ge0\) 为整数，\(f=\sum_{\alpha\in\mathbb N^r}a_\alpha x^\alpha\in A[\![x_1,\ldots,x_r]\!]\)。则 \(f\) 可逆，当且仅当其常数项 \(a_0\) 在 \(A\) 中可逆。
\end{proposition}
\begin{proof}
若有逆，比较常数项得到必要性。反向写 \(f=a_0(1-h)\)，其中 \(h\) 常数项为零。级数 \(\sum_{j\ge0}h^j\) 有意义，因为全次数小于 \(n\) 的项只受前 \(n\) 项影响。逐阶使用有限等比恒等式得到 \((1-h)\sum_{j\ge0}h^j=1\)，所以 \(f^{-1}=a_0^{-1}\sum_{j\ge0}h^j\)。
\end{proof}

特别地，若 \(k\) 为域，则 \(k[\![x_1,\ldots,x_r]\!]\) 是极大理想为 \((x_1,\ldots,x_r)\) 的局部环。局部多项式环 \(k[x_1,\ldots,x_r]_{(x_1,\ldots,x_r)}\) 也有这个完备化：在每个截断商中，所有原点处不为零的多项式已经可逆，所以先局部化不改变各阶商。若 \(p\) 为素数，则整数环关于 \((p)\) 的完备化为 \(\mathbb Z_p=\varprojlim_n\mathbb Z/p^n\mathbb Z\)。它的进位规则来自这些商环的运算，特征为零，与特征为 \(p\) 的 \(\mathbb F_p[\![t]\!]\) 不同。

\subsection{Noether 环上有限生成模的完备化正合性}
\label{b3:subsec:1303:03}

\begin{lemma}\label{b3:lem:inverse-limit-exact-surjective-kernels}
设 \(R\) 为交换环，给定三个 \(R\) 模逆系统及逐阶短正合列
\(0\rightarrow A_n\rightarrow B_n\rightarrow C_n\rightarrow0\)，其中所有过渡映射与列中的映射交换。若过渡映射 \(A_{n+1}\rightarrow A_n\) 全部满射，则
\[
0\longrightarrow\varprojlim A_n\longrightarrow\varprojlim B_n
\longrightarrow\varprojlim C_n\longrightarrow0
\]
正合。即使不假设满射，左端及中间位置仍正合。
\end{lemma}
\begin{proof}
单射与中间的核可逐坐标检验。要证明右端满射，给定相容的 \((c_n)\)，先选提升 \(b_1\)。若已选 \(b_n\)，任选 \(c_{n+1}\) 的提升 \(b'_{n+1}\)。它投到 \(B_n\) 后与 \(b_n\) 之差落在 \(A_n\)，由满射性可提升成 \(a_{n+1}\in A_{n+1}\)。令 \(b_{n+1}=b'_{n+1}-a_{n+1}\)，就同时保持提升与相容性。递归构造出所需的 \((b_n)\)。
\end{proof}

\begin{theorem}\label{b3:thm:completion-exact}
设 \(R\) 为交换 Noether 环、\(I\subseteq R\) 为理想。对有限生成 \(R\) 模的短正合列
\(0\rightarrow N\rightarrow M\rightarrow P\rightarrow0\)，其 \(I\) 进完备化
\[
0\longrightarrow\widehat N\longrightarrow\widehat M
\longrightarrow\widehat P\longrightarrow0
\]
仍正合。
\end{theorem}
\begin{proof}
把 \(N\) 看作 \(M\) 的子模。逐阶真正正合的列是
\[
0\longrightarrow N/(N\cap I^nM)\longrightarrow M/I^nM
\longrightarrow P/I^nP\longrightarrow0.
\]
其左侧过渡映射为自然商映射，因而满射。使用前一引理取逆极限，右端得到 \(\widehat P\)，中间得到 \(\widehat M\)。左端来自诱导滤过；\cref{b3:thm:artin-rees}给出的
\(I^nN\subseteq N\cap I^nM\subseteq I^{n-c}N\)
说明该滤过与 \(N\) 自身的 \(I\) 幂滤过相互共尾。由\cref{b3:prop:stable-filtration-equivalent}的自然同构，左端识别为 \(\widehat N\)，得到结论。
\end{proof}

这里没有断言任意逆极限右正合。例如短正合列
\(0\rightarrow2^n\mathbb Z\rightarrow\mathbb Z\rightarrow\mathbb Z/2^n\mathbb Z\rightarrow0\)
的逆极限右端为 \(\mathbb Z_2\)，但 \(\mathbb Z\rightarrow\mathbb Z_2\) 不满：每个模 \(2^n\) 的环中 \(3\) 的逆元组成相容系统，若它来自整数 \(a\)，则整数 \(3a-1\) 被所有 \(2^n\) 整除，从而等于零，矛盾。本例的左侧过渡映射为严格包含，不满足引理的满射假设。

\subsection{商、有限生成模与张量表示}
\label{b3:subsec:1303:04}

以下 \(R\) 均为 Noether 环，\(I\) 为固定理想，\(M\) 为有限生成模。

\begin{proposition}\label{b3:prop:completion-quotients}
设 \(R\) 为交换 Noether 环、\(I,J\subseteq R\) 为理想、\(M\) 为有限生成 \(R\) 模，所有完备化均关于 \(I\) 进行。在 \(\widehat M\) 中有
\[
\operatorname{im}(\widehat{JM}\rightarrow\widehat M)=J\widehat M,
\qquad \widehat{M/JM}\cong\widehat M/J\widehat M.
\]
特别地，
\[
\widehat M/I^n\widehat M\cong M/I^nM,
\qquad \ker(\widehat M\rightarrow M/I^nM)=I^n\widehat M.
\]
因此 \(\widehat M\) 对 \(I\widehat R\) 进拓扑完备且分离。
\end{proposition}
\begin{proof}
取 \(J=(a_1,\ldots,a_s)\)。映射
\(M^s\rightarrow JM\)、\((m_j)\mapsto\sum_ja_jm_j\) 满射。由完备化正合性，其完备化仍满；有限直和的完备化逐坐标等于完备化的直和，所以它在 \(\widehat M\) 中的像恰为 \(\sum_ja_j\widehat M=J\widehat M\)。再完备化 \(0\rightarrow JM\rightarrow M\rightarrow M/JM\rightarrow0\)，得到商的同构。

取 \(J=I^n\)。模 \(M/I^nM\) 在第 \(n\) 阶以后各商都不再改变，故其完备化是自身。于是得到第二组公式；逆极限原有的邻域核正是理想的幂，前面已证的完备性和分离性遂成为理想进完备性与分离性。
\end{proof}

\begin{theorem}\label{b3:thm:completion-tensor}
设 \(R\) 为交换 Noether 环、\(I\subseteq R\) 为理想、\(M\) 为有限生成 \(R\) 模，\(\widehat R\) 与 \(\widehat M\) 均表示 \(I\) 进完备化。则自然映射
\[
\widehat R\otimes_RM\longrightarrow\widehat M,
\qquad a\otimes m\longmapsto a\,\iota_M(m)
\]
是同构。特别地，\(\widehat M\) 为有限生成 \(\widehat R\) 模。
\end{theorem}
\begin{proof}
\(R\) 为 Noether 环，故 \(M\) 有有限展示
\(R^u\xrightarrow{A}R^v\rightarrow M\rightarrow0\)。完备化正合性说明 \(\widehat M\) 为矩阵 \(A:\widehat R^u\rightarrow\widehat R^v\) 的余核。张量积的右正合性说明 \(\widehat R\otimes_RM\) 为同一矩阵的余核。自然映射在自由模处就是逐坐标同构，故在余核处也是同构；这里比较余核只需张量积的右正合性。
\end{proof}

\subsection{完备环的 Noether 性}
\label{b3:subsec:1303:noetherian}

商的同构还保留相邻两阶之间的关系，从而识别原环与完备环的伴随分次环。分次环中有限个初始形式能够控制一个理想的全部初始形式；要回到原理想，则需把逐阶消去产生的系数收敛到完备环中。

\begin{theorem}\label{b3:thm:completion-noetherian}
设 \(R\) 为交换 Noether 环、\(I\subseteq R\) 为理想，\(\widehat R\) 为 \(I\) 进完备化。则 \(\widehat R\) 是 Noether 环，并且
\[
\operatorname{gr}_{I\widehat R}(\widehat R)
\cong\operatorname{gr}_I(R).
\]
\end{theorem}
\begin{proof}
记 \(B=\widehat R\)、\(K=IB\)。由\cref{b3:prop:completion-quotients}，\(B\) 对 \(K\) 进拓扑完备且分离。该命题的商同构在相邻两阶之间相容，取核得到 \(K^n/K^{n+1}\cong I^n/I^{n+1}\)，且乘法相容。若 \(I=(a_1,\ldots,a_s)\)，此分次环是 \((R/I)[X_1,\ldots,X_s]\) 的商，故为 Noether 环。

设 \(J\subseteq B\) 为任意理想。对每个非零 \(f\in B\)，分离性使它有唯一阶数 \(d\ge0\)，即 \(f\in K^d\setminus K^{d+1}\)；记其初始形式为 \(\operatorname{in}(f)\)。由所有 \(f\in J\) 的初始形式生成的齐次理想有限生成，故可选有限个 \(f_1,\ldots,f_r\in J\)，使它们的初始形式生成该理想，记阶数为 \(d_i\)。这一步可以选实际初始形式，是因为原生成族中有限个已足以生成同一个理想。

对 \(0\ne f\in J\)，令 \(r_0=f\)。若第 \(\nu\) 次的余项 \(r_\nu\ne0\)，记其阶数为 \(n_\nu\)，将初始形式表示成 \(\sum_i\overline a_{i,\nu}\operatorname{in}(f_i)\)。当 \(n_\nu\ge d_i\) 时，取 \(\overline a_{i,\nu}\) 为次数 \(n_\nu-d_i\) 的齐次元，并提升为 \(a_{i,\nu}\in K^{n_\nu-d_i}\)；当 \(n_\nu<d_i\) 时，取 \(a_{i,\nu}=0\)。于是
\[
r_{\nu+1}=r_\nu-\sum_i a_{i,\nu}f_i
\in J\cap K^{n_\nu+1}.
\]
若余项为零，已经得到有限组合；否则 \(n_{\nu+1}>n_\nu\)。若此过程无限继续，则 \(n_\nu\to\infty\)。设
\[
b_{i,\nu}=\sum_{\mu=0}^{\nu-1}a_{i,\mu},
\qquad r_\nu=f-\sum_i b_{i,\nu}f_i.
\]
对任意 \(N\ge0\)，当 \(\nu\) 充分大时，\(n_\nu-d_i\ge N\) 对每个 \(i\) 成立，所以以后的全部系数增量均属于 \(K^N\)。因此 \((b_{i,\nu})_\nu\) 是 Cauchy 序列，由 \(B\) 的完备性收敛于某个 \(b_i\in B\)。同时 \(r_\nu\in K^{n_\nu}\) 趋于零。对每个 \(N\)，在充分大的 \(\nu\) 处有 \(b_i-b_{i,\nu}\in K^N\) 及 \(r_\nu\in K^N\)，故
\[
f-\sum_i b_if_i
=r_\nu+\sum_i(b_{i,\nu}-b_i)f_i\in K^N.
\]
由 \(\bigcap_NK^N=0\) 得 \(f=\sum_i b_if_i\)。因此 \(J=(f_1,\ldots,f_r)\)；零理想同样有限生成，故 \(B\) 为 Noether 环。
\end{proof}

这个证明同时说明为什么“完备”有实质作用：初始形式只能逐阶消去，系数最后可能成为无穷和。有限阶消去若没有收敛到环中元素的保证，就不能推出原理想由所选元素生成。

% !TeX root = ../../Book3.tex
% 纲要标题：### 13.4 完备化的局部性质
% 纲要位置：工作记录/计划纲要.md，第 2381 行。
% 纲要细目（写作范围）：
% * 完备化与张量积
% * Noether 局部环完备化的忠实平坦性
% * 极大理想、剩余域及维数
% * 用依赖图连接 Rees 有限性、Artin–Rees、完备化正合性、张量表示与平坦性；以有限阶商的长度比较说明维数与重数保持
% * 证明顺序固定为：有限生成模的完备化正合性、\(M\otimes_R\widehat R\cong\widehat M\)、第12.1节的理想张量与局部判据、平坦性、剩余域非零及第12.5节的忠实平坦性判据
% * 极大理想进完备化的维数比较使用第11章的伴随分次与 Hilbert–Samuel 理论；模的忠实平坦下降仍是独立选读，不作为此处证明前提

\section{完备化的局部性质}
\label{b3:sec:1304}


完备化把原环中的计算搬到一个容纳所有相容近似的环中。上一节的正合性与张量表示使我们能够比较两边的模同态：有限生成理想的包含在完备化后仍为单射，因而完备环对原环平坦。在局部环中，剩余域还保持不变；这一点使变基后的计算能够反向检测原来的模。

\subsection{完备化与张量积}
\label{b3:subsec:1304:01}

设 \(R\) 为交换 Noether 环，\(I\subseteq R\) 为理想，各完备化均关于 \(I\) 进行。对有限生成 \(R\) 模 \(M,N\)，\cref{b3:thm:completion-tensor}给出的同构与模同态相容：若 \(f:M\rightarrow N\)，两边都把 \(a\otimes m\) 送到 \(a\,\iota_N\bigl(f(m)\bigr)\)。因此可以把有限生成模之间的完备化映射统一视为沿 \(R\rightarrow\widehat R\) 的标量扩张。

\begin{proposition}\label{b3:prop:completion-finite-tensor-products}
设 \(R\) 为交换 Noether 环、\(I\subseteq R\) 为理想，\(M,N\) 为有限生成 \(R\) 模；记 \(\widehat R\)、\(\widehat M\)、\(\widehat N\) 及 \(\widehat{M\otimes_RN}\) 为各自的 \(I\) 进完备化。则
\[
\widehat{M\otimes_RN}
\cong\widehat M\otimes_{\widehat R}\widehat N.
\]
\end{proposition}
\begin{proof}
左边由张量表示等于 \(\widehat R\otimes_R(M\otimes_RN)\)。右边把两因子写成 \(\widehat R\otimes_RM\)、\(\widehat R\otimes_RN\)，映射
\[
(a\otimes m)\otimes(b\otimes n)\longmapsto ab\otimes(m\otimes n)
\]
保持各个平衡关系，逆映射为
\(c\otimes(m\otimes n)\mapsto(c\otimes m)\otimes(1\otimes n)\)。在纯张量上检查两复合为恒等，便得同构。
\end{proof}

对于不有限生成的模，这个张量表示可能失效。例如设 \(p\) 为素数，取 \(R=\mathbb Z_{(p)}\)、\(I=(p)\)、\(M=\mathbb Q\)。有限阶商 \(R/p^nR\cong\mathbb Z/p^n\mathbb Z\) 给出 \(\widehat R=\mathbb Z_p\)。由于 \(p^nM=M\)，有 \(\widehat M=0\)；但
\[
\widehat R\otimes_RM=\mathbb Z_p[1/p]\ne0.
\]
由局部化的零判据，\(\mathbb Z_p[1/p]=0\) 当且仅当某个 \(p^n\) 在 \(\mathbb Z_p\) 中为零。但对每个 \(n\ge0\)，\(p^n\) 在第 \(n+1\) 阶坐标 \(\mathbb Z/p^{n+1}\mathbb Z\) 中非零，故在逆极限中也非零。因此倒置 \(p\) 后仍为非零环，张量积与完备化在这个模上不同。

\subsection{Noether 局部环完备化的忠实平坦性}
\label{b3:subsec:1304:02}
\label{b3:subsec:1304:04}

\begin{theorem}\label{b3:thm:completion-flat}
若 \(R\) 为交换 Noether 环、\(I\subseteq R\) 为理想，则其 \(I\) 进完备化 \(\widehat R\) 是平坦 \(R\) 模。若 \(R\) 还是局部环，极大理想为 \(\mathfrak m\)，且 \(I\subseteq\mathfrak m\)，则 \(R\rightarrow\widehat R\) 忠实平坦。特别地，极大理想进完备化忠实平坦。
\end{theorem}
\begin{proof}
对任意有限生成理想 \(J\subseteq R\)，映射
\[
J\otimes_R\widehat R\longrightarrow R\otimes_R\widehat R=\widehat R
\]
由自然张量表示等同于 \(\widehat J\rightarrow\widehat R\)。\cref{b3:thm:completion-exact}说明后一个映射单射。于是\cref{b3:thm:ideal-flatness-criterion}给出平坦性。

若 \(R\) 局部且 \(I\subseteq\mathfrak m\)，则 \(R/\mathfrak m\) 的 \(I\) 进滤过第一步就为零。由\cref{b3:prop:completion-quotients}，
\[
\widehat R/\mathfrak m\widehat R\cong R/\mathfrak m\ne0.
\]
这说明完备化没有删去原谱的唯一闭点。\(R\) 只有这一个极大理想，故结合已经证明的平坦性，\cref{b3:thm:faithfully-flat-criterion}给出忠实平坦性，并保证所有素点都有上方点。
\end{proof}

\cref{b3:fig:completion-flat-dependencies}把上述论证所用的结果连在一起。图中“正合”限于有限生成模的完备化；得到环同态的平坦性后，张量积保持正合的结论才适用于任意模。

\begin{center}
\begin{minipage}{\textwidth}
\makeatletter
\def\@captype{figure}
\makeatother
\centering
\[
\begin{tikzcd}[column sep=large,row sep=large]
\boxed{\mathcal R_I(R)\;\text{是 Noether 环}}
  \arrow[r]
&\boxed{\text{Artin--Rees}}
  \arrow[r]
&\boxed{\text{有限生成模的完备化正合}}
  \arrow[d]\\
\boxed{R\rightarrow\widehat R\;\text{忠实平坦}}
&\boxed{R\rightarrow\widehat R\;\text{平坦}}
  \arrow[l,"R\;\text{局部，}I\subseteq\mathfrak m"']
&\boxed{\widehat M\cong\widehat R\otimes_RM}
  \arrow[l,"\text{理想张量判据}"']
\end{tikzcd}
\]
\caption[完备化的正合性、张量表示与平坦性]{完备化的正合性、张量表示与平坦性。设 \(R\) 为交换 Noether 环，\(I\subseteq R\) 为理想，\(M\) 为有限生成 \(R\) 模；局部情形以 \(\mathfrak m\) 表示极大理想。}
\label{b3:fig:completion-flat-dependencies}
\end{minipage}
\end{center}

在 \(R\) 局部且 \(I\subseteq\mathfrak m\) 的情形中，忠实平坦性意味着：一个 \(R\) 模即使不是有限生成的，也不会在张量到 \(\widehat R\) 后从非零变成零；模映射的单射、满射及正合性也能在张量后检测。这里谈的是张量积函子。上一小节的 \(M=\mathbb Q\) 例子表明，对任意模不能把它替换成完备化函子。

\ProseParagraphBreak
全局完备化未必忠实。例如设 \(p,q\) 为不同素数，\(\mathbb Z\rightarrow\mathbb Z_p\) 平坦，但 \(q\) 在 \(\mathbb Z_p\) 中可逆：在各个 \(\mathbb Z/p^n\mathbb Z\) 中，\(q\) 的唯一逆元组成相容系统。因此
\[
(\mathbb Z/q\mathbb Z)\otimes_{\mathbb Z}\mathbb Z_p=0.
\]
几何上，\((q)\) 没有上方素点，因为 \(q\) 可逆，不能属于 \(\mathbb Z_p\) 的任何真理想。先局部化为 \(\mathbb Z_{(p)}\) 再完备化，则由前述定理得到忠实平坦映射；其谱映射满射，而 \(\operatorname{Spec}\mathbb Z_{(p)}\) 在 \(\operatorname{Spec}\mathbb Z\) 中对应 \((0),(p)\)。因此
\[
\operatorname{im}\bigl(\operatorname{Spec}\mathbb Z_p\to\operatorname{Spec}\mathbb Z\bigr)
=\{(0),(p)\}.
\]
关于合数 \(m\) 的 \((m)\) 进完备化，\cref{b3:ex:13-composite-adic}要求通过各有限阶商的中国剩余分解比较不同素数的分量。局部性与忠实性在这种全局完备化中仍须分别检查。

\subsection{极大理想、剩余域与维数}
\label{b3:subsec:1304:03}

\begin{proposition}\label{b3:prop:completed-local-ring}
设 \((R,\mathfrak m)\) 为交换 Noether 局部环，\(k=R/\mathfrak m\) 为剩余域，\(B=\widehat R\) 为极大理想进完备化。则 \(B\) 是 Noether 局部环，其极大理想为 \(\mathfrak n=\mathfrak mB\)，剩余域自然等于 \(k\)。对每个 \(r\ge1\) 有
\[
B/\mathfrak n^r\cong R/\mathfrak m^r,
\qquad \operatorname{gr}_{\mathfrak n}(B)\cong\operatorname{gr}_{\mathfrak m}(R).
\]
\end{proposition}
\begin{proof}
Noether 性和两项同构分别由\cref{b3:thm:completion-noetherian}和\cref{b3:prop:completion-quotients}给出。商 \(B/\mathfrak n=k\) 是域，所以 \(\mathfrak n\) 为极大理想。若 \(b\notin\mathfrak n\)，选 \(a\in R\) 使 \(ab\equiv1\pmod{\mathfrak n}\)，并写 \(ab=1-u\)、\(u\in\mathfrak n\)。因为 \(u^j\in\mathfrak n^j\)，部分和 \(s_N=1+u+\cdots+u^N\) 在模每个固定 \(\mathfrak n^r\) 后最终不变，所以是 \(\mathfrak n\) 进 Cauchy 序列。\(B\) 的完备性给出极限 \(s=\sum_{j\ge0}u^j\)；有限等比恒等式 \((1-u)s_N=1-u^{N+1}\) 在每个有限阶商中给出 \((1-u)s=1\)。于是 \(b(as)=1\)，故 \(b\) 可逆。这与\cref{b3:prop:formal-series-unit}中的逆元构造使用同一个逐阶计算。因此 \(\mathfrak n\) 之外的每个元素均为单位；而 \(\mathfrak n\) 是真理想，其中不含单位，所以它恰由所有非单位组成，是 \(B\) 的唯一极大理想。
\end{proof}

所有有限阶商及剩余域保持不变，并不表示原环与完备环相同。例如设 \(k\) 为域，\(R=k[x]_{(x)}\) 的元素是原点分母非零的有理函数，而 \(\widehat R=k[\![x]\!]\) 中的级数 \(v=\sum_{j\ge0}x^{2^j}\) 不在 \(R\) 中。若 \(v=P/Q\)，其中 \(P,Q\in k[x]\)、\(Q(0)\ne0\)，则对充分大的 \(j\)，\(2^j>\deg P\) 且 \(2^{j-1}>\deg Q\)。在 \(Qv=P\) 中比较 \(x^{2^j}\) 的系数，左侧仅有 \(Q(0)\) 的贡献，右侧为零，矛盾。尽管对每个 \(n\) 都有 \(R/(x^n)\cong\widehat R/(x^n)\)，这个新元素仍需要一整套无限相容的截断才能指定；任一固定阶都可以用原环的多项式表示它。维数保持则由下一小节的长度增长比较给出。

\subsection{Hilbert–Samuel 理论与完备化的维数保持}
\label{b3:subsec:1304:05}

\begin{theorem}\label{b3:thm:completion-dimension}
对交换 Noether 局部环 \((R,\mathfrak m)\) 及其极大理想进完备化 \(B=\widehat R\)、\(\mathfrak n=\mathfrak mB\)，有
\[
\dim B=\dim R,\qquad e_{\mathfrak n}(B)=e_{\mathfrak m}(R).
\]
更一般地，对非零有限生成的 \(R\) 模 \(M\)，
\[
\dim_B\widehat M=\dim_RM,\qquad
 e_{\mathfrak n}(\widehat M)=e_{\mathfrak m}(M).
\]
\end{theorem}
\begin{proof}
\cref{b3:prop:completion-quotients}给出
\(\widehat M/\mathfrak n^{r+1}\widehat M\cong M/\mathfrak m^{r+1}M\)。对每个整数 \(r\ge0\)，两者的标量作用都通过同一个商环 \(B/\mathfrak n^{r+1}\cong R/\mathfrak m^{r+1}\) 分解，所以子模及组成列完全相同，特别是
\[
\ell_B(\widehat M/\mathfrak n^{r+1}\widehat M)
=\ell_R(M/\mathfrak m^{r+1}M).
\]
因此两侧的 Hilbert–Samuel 函数逐项相同。\(B\) 已证为 Noether 局部环，\(\widehat M\) 已证为有限生成 \(B\) 模且由忠实平坦性非零，故两侧都能应用\cref{b3:thm:hilbert-samuel}。相等的累计函数给出相等的最终多项式，比较次数与首项即得维数和重数。取 \(M=R\) 得环的结论。
\end{proof}

因此完备化保留局部长度增长及其决定的维数，却未必保留素理想分解的所有细节。例如曲线在形式邻域中可能分出原环中尚未分开的分支；维数相等并不声称两者素谱同胚。

% !TeX root = ../../Book3.tex
% 纲要标题：### 13.5 有限阶近似与形式邻域
% 纲要位置：工作记录/计划纲要.md，第 2390 行。
% 纲要细目（写作范围）：
% * 商 \(R/\mathfrak m^N\)、截断数据及相容逆系统
% * 有限阶计算、完整形式对象与收敛解析对象的不同信息量
% * Noether 局部环中理想向完备化扩张再收缩的保持；与忠实平坦性的联系
% * 由子模闭性证明线性方程可解性从完备化下降，并以非线性形式根不能下降的例子说明适用范围；习题比较解集、精确解的有限阶逼近与分支数据
% * 算法须注明多项式、代数幂级数或有限截断的输入形式；不能对任意不可计算幂级数声称有限有效求解
% ---

\section{有限阶近似与形式邻域}
\label{b3:sec:1305}


完备化由全部有限阶商组成，但一次计算通常只能得到其中有限多个。本节说明这类计算究竟确定什么，以及何时能够从所有阶的近似恢复原环中的结论。

\subsection{有限阶商、截断数据与相容逆系统}
\label{b3:subsec:1305:01}

\begin{definition}\label{b3:def:finite-order-neighbourhood}
设 \((R,\mathfrak m)\) 为交换 Noether 局部环，\(N\ge0\) 为整数。本书将商环 \(R/\mathfrak m^{N+1}\) 称为该点的 \(N\) 阶邻域的坐标环。
\end{definition}

指数取 \(N+1\)，是为了保留直到第 \(N\) 阶的信息：零阶商 \(R/\mathfrak m\) 只剩剩余域；一阶商 \(R/\mathfrak m^2\) 还保留 \(\mathfrak m/\mathfrak m^2\) 中的一次信息；二阶商再保留 \(\mathfrak m^2/\mathfrak m^3\) 中的二次信息。所有商 \(R/\mathfrak m^q\)、\(q\ge1\)，连同它们之间的自然商映射组成相容逆系统，描述此点的\term[xingshi-linyu]{形式邻域}[formal neighbourhood]。

\ProseParagraphBreak
设 \(k\) 为域。在 \(k[x,y]_{(x,y)}\) 中模掉 \((x,y)^3\)，得到 \(k\) 基
\(1,x,y,x^2,xy,y^2\)。相乘后删去全次数至少三的单项式，即为商环乘法。对尖点局部环
\(R=k[x,y]_{(x,y)}/(y^2-x^3)\)、\(\mathfrak m=(x,y)R\)，二阶邻域的坐标环为
\[
R/\mathfrak m^3\cong k[x,y]/\bigl((x,y)^3,y^2\bigr).
\]
最低关系 \(y^2=0\) 与\cref{b3:prop:hypersurface-tangent-cone}所给切锥的关系 \(Y^2=0\) 相同，但这个有限阶商还截去了全次数至少三的全部单项式。在三阶邻域 \(R/\mathfrak m^4\) 中，右侧的三次项 \(x^3\) 才重新可见。一个低阶商无法单独区分所有高阶方程。

\begin{proposition}\label{b3:prop:compatible-formal-solutions}
设 \((R,\mathfrak m)\) 为交换 Noether 局部环，\(\widehat R\) 为极大理想进完备化，\(f_1,\ldots,f_s\in R[X_1,\ldots,X_r]\)。在 \(\widehat R^r\) 中给定一个共同零点，等价于给定对每个 \(N\ge1\) 的零点
\(a_N\in(R/\mathfrak m^N)^r\)，且 \(a_{N+1}\) 模 \(\mathfrak m^N\) 等于 \(a_N\)。
\end{proposition}
\begin{proof}
零点逐阶取商显然相容。反向由逆极限定义，各坐标的相容余类给出 \(a\in\widehat R^r\)。多项式运算与取商交换，所以每个 \(f_i(a)\) 在所有有限阶商中为零；\(\widehat R\) 分离，故 \(f_i(a)=0\)。
\end{proof}

逐阶任意选取零点 \(a_N\)，只保证每个有限阶的方程成立，未必保证这些选择彼此相容。形式解要求一次选出整个系统，并同时满足
\[
\exists(a_N)_{N\ge1}\quad
\begin{cases}
f_i(a_N)=0\text{ 于 }R/\mathfrak m^N
 &\text{对所有 }i,N,\\
a_{N+1}\bmod\mathfrak m^N=a_N
 &\text{对所有 }N.
\end{cases}
\]
是否能够作出这样的相容选择，还须利用环与方程的具体条件。这与\secref{b3:sec:1303}中逐阶满射未必给出逆极限满射的区别相联系；简单根的逐阶提升将在\chapref{b3:ch:23}系统讨论。

\subsection{有限阶、形式与收敛解析对象}
\label{b3:subsec:1305:02}

在复数域上，多项式、原点附近收敛的幂级数与任意形式幂级数分别给出
\[
\mathbb C[x]\subseteq\mathbb C\{x\}\subseteq\mathbb C[\![x]\!].
\]
这里多项式只允许有限多个非零系数，\(\mathbb C\{x\}\) 允许无限多个系数，但要求正收敛半径；形式幂级数则允许任意系数序列，乘法只作逐阶有限的系数运算。\(\sum_{n\ge0}x^n\) 在 \(|x|<1\) 时收敛，却不是多项式；\(\sum_{n\ge0}n!x^n\) 是形式级数。对任意固定的 \(x\in\mathbb C\setminus\{0\}\)，令 \(a_n=n!x^n\)，则
\[
\dfrac{|a_{n+1}|}{|a_n|}=(n+1)|x|\longrightarrow\infty.
\]
因此从某项起 \(|a_{n+1}|\ge2|a_n|\)，通项 \(n!x^n\) 不趋于零，级数不收敛。

\ProseParagraphBreak
任意给定的形式级数模 \(x^N\) 都由一个次数小于 \(N\) 的多项式实现。相同的截断既可接上全零的尾部，成为一个收敛多项式，也可接上从第 \(N\) 项开始的阶乘系数尾部，成为发散的形式级数。因此有限阶数据不能判定收敛性。完备化是代数上的逐阶补全，不自动赋予解析意义；若研究收敛性，还必须另加系数估计或相应的解析定理。

\subsection{完备化中理想的扩张与收缩}
\label{b3:subsec:1305:03}

\begin{proposition}\label{b3:prop:completion-ideal-contraction}
设 \((R,\mathfrak m)\) 为交换 Noether 局部环，\(B=\widehat R\) 为极大理想进完备化。对每个理想 \(J\subseteq R\)，有
\[
JB\cap R=\bigcap_{n\ge1}(J+\mathfrak m^n)=J,
\qquad B/JB\cong\widehat{R/J}.
\]
后一完备化关于 \((\mathfrak m+J)/J\) 进行。更一般地，设 \(M\) 为有限生成 \(R\) 模，\(N\subseteq M\) 为子模，各模均取 \(\mathfrak m\) 进完备化。自然映射 \(M\to\widehat M\) 及 \(\widehat N\to\widehat M\) 均单射；经这些映射将它们识别为子模后，有
\[
M\cap\widehat N=N\quad\text{于 }\widehat M.
\]
因此，对任意非负整数 \(r,s\)、\(R\) 上的 \(s\times r\) 矩阵 \(A\) 及 \(b\in R^s\)，线性方程 \(Ax=b\) 在 \(B^r\) 中有解，当且仅当在 \(R^r\) 中有解。

一般非线性方程没有这一可解性下降结论。若 \(k\) 为特征不为二的域，\(x\) 为不定元，则 \(k[\![x]\!]\) 中存在唯一常数项为一、满足 \(u^2=1+x\) 的级数；但在 \(k=\mathbb Q\) 时，此方程在原局部环 \(\mathbb Q[x]_{(x)}\) 中无解。
\end{proposition}
\begin{proof}
\(\mathfrak m\) 是 \(R\) 的 Jacobson 根，故\cref{b3:thm:krull-intersection}使每个有限生成模到其完备化的自然映射单射。Noether 性使 \(N\) 与 \(M/N\) 均有限生成；对 \(0\rightarrow N\rightarrow M\rightarrow M/N\rightarrow0\) 应用\cref{b3:thm:completion-exact}，其完备化仍正合。因此 \(\widehat N\) 可识别为 \(\widehat M\) 的子模。一个 \(x\in M\) 的像属于 \(\widehat N\)，当且仅当它在 \(\widehat{M/N}\) 中为零；由 \(M/N\to\widehat{M/N}\) 的单射性，这又等价于 \(x\in N\)，即得交式。

取 \(M=R\)、\(N=J\)，\cref{b3:prop:completion-quotients}将 \(\widehat J\) 的像识别为 \(JB\)，并给出 \(B/JB\cong\widehat{R/J}\)，故得到理想的扩张收缩公式。同一个 Krull 交结论对有限生成 \(R\) 模 \(R/J\) 给出 \(\bigcap_{n\ge1}\mathfrak m^n(R/J)=0\)；取逆像就是 \(\bigcap_{n\ge1}(J+\mathfrak m^n)=J\)。于是，只要原环元素在每个有限精度下都能被 \(J\) 中元素逼近，它本来就属于 \(J\)。

对线性方程，令 \(M=R^s\)、\(N=\operatorname{im}A\)。\(N\) 由 \(A\) 的有限列生成，满射 \(R^r\to N\) 的核也由 Noether 性有限生成。对这个满射及包含 \(N\to R^s\) 分别使用完备化正合性，得到 \(\operatorname{im}(\widehat A)=\widehat N\) 在 \(B^s\) 中的像。因此完备环中有解，恰说明 \(b\) 的像属于 \(\widehat N\)；刚证的交式使它等价于 \(b\in N\)，也就是原环中有解。这是解的存在性结论，选定的形式解仍可在 \(\ker\widehat A\) 中任意改变，并不必来自 \(R^r\)；例如零方程 \(0x=0\) 的每个形式级数都是解。

对于非线性反例，写 \(u=1+\sum_{n\ge1}a_nx^n\)。比较 \(u^2=1+x\) 的系数，得到
\[
2a_1=1,\qquad
2a_n=-\sum_{i=1}^{n-1}a_ia_{n-i}\quad(n\ge2).
\]
因 \(2\) 在 \(k\) 中可逆，这个递归唯一确定所有系数，并使每一阶的方程成立，所以给出唯一常数项为一的形式根。若 \(k=\mathbb Q\) 时有原环中的根，将它写成互素非零多项式之比 \(P/Q\)，就有 \(P^2=(1+x)Q^2\)。不可约因子 \(1+x\) 在左侧的指数为偶数，在右侧的指数为奇数，矛盾。因此连 \(\mathbb Q(x)\) 中也没有根，而 \(\mathbb Q[x]_{(x)}\) 的完备化 \(\mathbb Q[\![x]\!]\) 中有上述形式根。
\end{proof}

Krull 交定理把理想收缩解释为所有阶近似所检验的闭性；另一方面，\cref{b3:thm:completion-flat}已证明 \(R\to B\) 忠实平坦，故同一收缩等式也是\cref{b3:prop:faithfully-flat-detection}第二项的一次应用。

\subsection{幂级数计算的输入形式与有效算法的适用范围}
\label{b3:subsec:1305:04}

设 \(k\) 为域，\(r,N\ge1\) 为整数，记 \(A=k[\![x_1,\ldots,x_r]\!]\)、\(\mathfrak n=(x_1,\ldots,x_r)A\)。当输入为 \(k[x_1,\ldots,x_r]\) 中的多项式并指定精度 \(N\) 时，模 \(\mathfrak n^N\) 的加法、乘法及求单位的逆只涉及有限个系数。若 \(k\) 中的系数具有有限表示，且四则运算与相等判定可有效执行，这些是有限算法。例如对常数项为一的 \(f=1-h\)，有 \(h\in\mathfrak n\)，其模 \(\mathfrak n^N\) 的逆可取
\[
1+h+\cdots+h^{N-1}\pmod{\mathfrak n^N},
\]
因为 \((1-h)(1+h+\cdots+h^{N-1})=1-h^N\)，而 \(h^N\in\mathfrak n^N\)。例如\cref{b3:ex:13-finite-inverse}中的 \(1-x-y^2\) 模 \((x,y)^4\) 的逆，就由这个有限和截去全次数至少四的项算出。

\begin{definition}\label{b3:def:algebraic-series-input}
设 \(k\) 为域，\(r\ge1\) 为整数，\(x_1,\ldots,x_r\) 为不定元，\(f\in k[\![x_1,\ldots,x_r]\!]\)。若 \(f\) 在有理函数域 \(k(x_1,\ldots,x_r)\) 上代数，即存在非零多项式 \(P\in k[x_1,\ldots,x_r,T]\) 使 \(P(x_1,\ldots,x_r,f)=0\)，则称 \(f\) 为\term[algebraicpowerseries]{代数幂级数}[algebraic power series]。
\end{definition}

若输入是代数幂级数，必须同时给出定义方程、所选分支及取得所需系数的方法。例如特征不为二时，\cref{b3:prop:completion-ideal-contraction}中的方程 \(T^2-(1+x)=0\) 有两个形式分支 \(u,-u\)，常数项分别为 \(1,-1\)。只给这个方程尚未指定输入的是哪一个级数；指定常数项后，该例的递归才给出所需有限阶系数。一般定义方程还可能有不能用选定参数表示成幂级数的分支，取得系数的方法须针对所选分支说明。任意形式级数则含有无限数据，未必具有有限编码或可计算的系数序列；它作为逆极限中的元素存在，并不自动给出可供有限算法操作的输入。完备性保证相容近似有极限，对象的有限编码与精确判等还需要另外的可计算性条件。

\ProseParagraphBreak
如果获取一个任意形式级数的信息只靠逐个查询系数，有限次查询不能完成对所有输入的零判定。事实上，设程序在零级数上查询的各单项式全次数至多为 \(N\)，则 \(0\) 与 \(x_1^{N+1}\) 在所有已查询位置给出相同答案，程序的这些查询无法区分二者。因此截断计算给出的是指定精度内的确定结论；要从中推断完整等式，还须有额外的次数界、有限生成关系或已经证明的提升与唯一性结果。局部标准基如何提供这类有限控制，将在本卷附录 F 中进一步讨论。


\begin{chaptersummary}
Artin–Rees 的关键在于把 Noether 环上无限多层子模组成一个有限生成的 Rees 模。它给出诱导滤过的稳定性，使子模自身的拓扑与环境诱导的拓扑一致；配合 Nakayama 引理，便得到 Jacobson 根条件下有限生成模的分离性和子模闭性。

\ProseParagraphBreak
完备化的正合性先在逐阶商中建立，再通过相容提升和滤过比较得到。Noether 性使有限生成模具有有限展示，从而可用同一有限矩阵计算完备化与 \(\widehat R\) 标量扩张，得到 \(\widehat M\cong\widehat R\otimes_RM\)。理想张量判据随后给出平坦性，非零剩余域给出局部忠实性。完备环的 Noether 性由初始形式的有限生成与收敛消去证明，而维数保持由相同的 Hilbert–Samuel 函数证明。

\ProseParagraphBreak
形式对象保存所有相容截断；单个截断只保存有限精度。局部理想的闭性允许从任意高精度的理想成员关系恢复原环中的成员关系，却不能据此把任意形式级数视为可计算或收敛的函数。
\end{chaptersummary}

\BookExercises

\begin{exercise}\label{b3:ex:13-equivalent-ideals}
设 \(R\) 为交换 Noether 环，\(I,J\) 为 \(R\) 的理想且满足 \(\sqrt I=\sqrt J\)。证明任意 \(R\) 模上的 \(I\) 进拓扑与 \(J\) 进拓扑相同，且其完备化自然同构。
\end{exercise}
\begin{solution}
取 \(I=(a_1,\ldots,a_r)\)，各 \(a_i\) 的某个幂属于 \(J\)。共同选择整数 \(c\ge1\) 使 \(a_i^c\in J\)，并令 \(a=r(c-1)+1\)。所有全次数至少 \(a\) 的单项式都有一个指数至少 \(c\)，故 \(I^a\subseteq J\)。交换 \(I,J\) 同理得某个整数 \(b\ge1\) 使 \(J^b\subseteq I\)。于是 \(I^{an}M\subseteq J^nM\)、\(J^{bn}M\subseteq I^nM\)，两组邻域相互共尾。将第 \(an\) 阶相容余类投到 \(M/J^nM\)，以及将第 \(bn\) 阶相容余类投到 \(M/I^nM\)，构成互逆的自然极限同构。
\end{solution}

\begin{exercise}\label{b3:ex:13-nonclosed-submodule}
在 \(\mathbb Z\) 的 \((2)\) 进拓扑中求 \(3\mathbb Z\) 的闭包。比较在 \(\mathbb Z_{(2)}\) 中相应理想的情形。
\end{exercise}
\begin{solution}
因为 \(3\mathbb Z+2^n\mathbb Z=\mathbb Z\)，闭包公式给出 \(\overline{3\mathbb Z}=\mathbb Z\)，故这个真子模不闭。局部化以后 \(3\) 成为单位，\(3\mathbb Z_{(2)}=\mathbb Z_{(2)}\)，当然闭。第一个环中 \((2)\not\subseteq J(\mathbb Z)\)，所以没有违反 Krull 交定理的闭性结论。
\end{solution}

\begin{exercise}\label{b3:ex:13-graded-not-ring}
设 \(k\) 为域，令
\[
R=k[x,y]/(x^2-y^3,xy,y^4),\qquad
S=k[x,y]/(x^2,xy,y^4).
\]
证明它们都是局部环，极大理想均由 \(x,y\) 的像生成，其伴随分次环同构；再证明 \(R,S\) 不同构。可比较平方为零的元素是否全部属于极大理想的平方。
\end{exercise}
\begin{solution}
两环都有 \(k\) 基 \(1,x,y,y^2,y^3\)：在 \(R\) 中用 \(x^2=y^3\) 消去 \(x^2\)，在 \(S\) 中用 \(x^2=0\)，其余混合项和 \(y^4\) 均为零。这些关系确实给出五维代数；结合性只须额外核对 \(x^2y=0=xyx\) 及 \(x^3=0\)。因此上述标准形唯一。两个理想 \(\mathfrak m=(x,y)\) 均满足 \(\mathfrak m^2=(y^2,y^3)\)、\(\mathfrak m^3=(y^3)\)、\(\mathfrak m^4=0\)，且商为 \(k\)，故是唯一极大理想。两环的逐层乘法都给出 \(k[X,Y]/(X^2,XY,Y^4)\)。

若 \(u\in R\) 满足 \(u^2=0\)，先模掉 \(\mathfrak m\) 得到常数项为零。写 \(u=ax+by+cy^2+dy^3\)，则
\[
u^2=b^2y^2+(a^2+2bc)y^3.
\]
由基的独立性先得 \(b=0\)，再得 \(a=0\)，所以 \(u\in\mathfrak m^2\)。但在 \(S\) 中，\(x^2=0\) 且 \(x\notin\mathfrak m^2\)。环同构保持极大理想及其平方，因而这两个环不同构。此例在任意特征下成立，说明即使底域和全部分次片均相同，分次环仍会遗失高阶乘法信息。
\end{solution}

\begin{exercise}\label{b3:ex:13-principal-rees}
设 \(R\) 为交换环，\(a\in R\) 不是零因子，证明 \(\mathcal R_{(a)}(R)\cong R[T]\)，其中 \(T\) 对应 \(at\)。再设 \(k\) 为域，若 \(R=k[\epsilon]/(\epsilon^2)\)、\(a=\epsilon\)，求此满射的核。
\end{exercise}
\begin{solution}
映射 \(R[T]\rightarrow R[t]\) 把 \(\sum r_iT^i\) 送到 \(\sum r_ia^it^i\)，像即 Rees 环。若 \(a\) 不是零因子，则每个 \(a^i\) 也不是零因子，逐系数比较得到核为零。第二种情形中次数至少二全部消失，次数一的系数须满足 \(r_1\epsilon=0\)，即 \(r_1\in(\epsilon)\)，常数项须为零。因此核为 \((\epsilon T,T^2)\)。
\end{solution}

\begin{exercise}\label{b3:ex:13-shifted-filtration}
设 \(R\) 为交换环，\(I\subseteq R\) 为理想，\(M\) 为 \(R\) 模，\(r\) 为非负整数。令 \(F^nM=I^{\max(n-r,0)}M\)。证明这是稳定 \(I\) 滤过；指出稳定开始的一个次数并比较其完备化。
\end{exercise}
\begin{solution}
下降性明显。当 \(n<r\) 时，\(IF^nM=IM\subseteq M=F^{n+1}M\)；当 \(n\ge r\) 时，\(IF^nM=I^{n-r+1}M=F^{n+1}M\)。所以可取稳定次数 \(r\)。滤过只把原滤过向后推迟有限步，与 \(I^nM\) 相互共尾，故给出相同拓扑和自然同构的完备化。
\end{solution}

\begin{exercise}\label{b3:ex:13-artin-rees-number}
设 \(k\) 为域，在 \(R=k[x,y]\) 中取 \(I=(x,y)\)、\(N=(x^a,y^b)\subseteq M=R\)，其中 \(a,b\ge1\)。求 Artin–Rees 等式中最小的 \(c\)。
\end{exercise}
\begin{solution}
令 \(c=\max(a,b)\)。单项式逐项判断给出
\[
N\cap I^n=x^aI^{n-a}+y^bI^{n-b}
=I^{n-c}(x^aI^{c-a}+y^bI^{c-b})\quad(n\ge c).
\]
所以此 \(c\) 可行。不妨 \(b\ge a\)。若 \(c'<b\)，取 \(n=b\)，则 \(y^b\in N\cap I^b\)，但它不属于 \(I^{b-c'}(N\cap I^{c'})\)：后一单项式理想的每个不被 \(x\) 整除的生成单项式，\(y\) 次数至少为 \(b+(b-c')>b\)。故最小值为 \(\max(a,b)\)。
\end{solution}

\begin{exercise}\label{b3:ex:13-krull-domain}
设 \(R\) 为 Noether 整环，\(I\) 为 \(R\) 的真理想。证明 \(\bigcap_nI^n=0\)，即使 \(I\) 不包含于 Jacobson 根也成立。
\end{exercise}
\begin{solution}
令 \(N=\bigcap_nI^n\)。它是 Noether 环 \(R\) 的理想，故有限生成；\cref{b3:thm:artin-rees}给出 \(N=IN\)。对 \(N\) 的恒等自同态应用\cref{b3:thm:determinant-trick}，得到某个 \(a\in I\) 使
\[
(1-a)N=0.
\]
因 \(I\ne R\)，有 \(1-a\ne0\)；整环中的非零元素不是零因子，所以 \(N=0\)。这里用整环的非零因子性质替代了 \(1-a\) 为单位的条件。
\end{solution}

\begin{exercise}\label{b3:ex:13-p-adic-digits}
设 \(p\) 为素数，\(\mathbb Z_p\) 为整数环的 \((p)\) 进完备化。证明每个 \(z\in\mathbb Z_p\) 唯一写成 \(\sum_{i\ge0}a_ip^i\)，其中各 \(a_i\) 为整数且 \(0\le a_i<p\)。证明 \(z\) 可逆当且仅当 \(a_0\ne0\)。
\end{exercise}
\begin{solution}
从第一个商选择唯一 \(a_0\)，再从第二个商中减去 \(a_0\)，剩余是唯一一个 \(a_1p\)；相容性使这种逐位选择持续且唯一。反向任意这样的数位族给出相容截断。如果 \(a_0=0\)，则 \(z\) 在剩余域中为零，不能可逆。若 \(a_0\ne0\)，先选整数 \(b\) 使 \(bz\equiv1\pmod p\)，再用 \((1-u)^{-1}=\sum_{j\ge0}u^j\)、\(u=1-bz\in p\mathbb Z_p\) 构造逆元。
\end{solution}

\begin{exercise}\label{b3:ex:13-composite-adic}
设 \(m=\prod_{i=1}^rp_i^{e_i}>1\)，其中 \(p_i\) 是不同素数。求 \(\mathbb Z\) 的 \((m)\) 进完备化。
\end{exercise}
\begin{solution}
中国剩余定理逐阶给出 \(\mathbb Z/m^n\mathbb Z\cong\prod_i\mathbb Z/p_i^{e_in}\mathbb Z\)，同构与过渡映射相容。逆极限与有限直积交换，因为相容条件逐坐标独立。指数列 \(e_in\) 在正整数中共尾，其极限与全部 \(p_i\) 次幂商的极限相同，所以完备化为 \(\prod_i\mathbb Z_{p_i}\)。
\end{solution}

\begin{exercise}\label{b3:ex:13-complete-finite-module}
设 \(R\) 为交换 Noether 环，\(I\subseteq R\) 为理想，且 \(R\) 对 \(I\) 进拓扑完备。证明每个有限生成的 \(R\) 模 \(M\) 都 \(I\) 进完备。取素数 \(p\)，用 \(R=\mathbb Z_p\)、\(I=(p)\)、\(M=\mathbb Q_p=\mathbb Z_p[1/p]\) 说明有限性不能省略。
\end{exercise}
\begin{solution}
\cref{b3:thm:completion-tensor}使自然映射等同于 \(M\rightarrow R\otimes_RM\)，因此是同构。另一方面，对 \(R=\mathbb Z_p\)、\(I=(p)\)、\(M=\mathbb Q_p=\mathbb Z_p[1/p]\)，有 \(p^nM=M\)，故 \(\widehat M=0\ne M\)。这里失败的首先是分离性。
\end{solution}

\begin{exercise}\label{b3:ex:13-completed-quotient}
设 \(k\) 为域。求 \(R=k[x,y]_{(x,y)}/(y^2,x^3)\) 的极大理想进完备化及其 \(k\) 维数。再求 \(k[x,y]_{(x,y)}/(y^2-x^3)\) 的极大理想进完备化。
\end{exercise}
\begin{solution}
第一环的极大理想四次幂为零，因为任何全次数至少四的单项式都被 \(x^3\) 或 \(y^2\) 整除，故它已完备。基为 \(1,x,x^2,y,xy,x^2y\)，维数六；与 \(k[\![x,y]\!]/(y^2,x^3)\) 相同。第二环由商完备化公式得到 \(k[\![x,y]\!]/(y^2-x^3)\)。这一结论不以把每个形式级数直接代入一个有理函数为依据，而来自全部有限阶商的相容同构。
\end{solution}

\begin{exercise}\label{b3:ex:13-completed-presentation}
设 \(k\) 为域，\(R=k[x]_{(x)}\)、\(A=\operatorname{diag}(x^2,x^4):R^2\rightarrow R^2\)、\(M=\operatorname{coker}A\)。求 \(M\) 的 \((x)\) 进完备化 \(\widehat M\) 的展示及长度。
\end{exercise}
\begin{solution}
在 \(\widehat R=k[\![x]\!]\) 上保留同一个矩阵，得
\(\widehat M\cong k[\![x]\!]/(x^2)\oplus k[\![x]\!]/(x^4)\)。两直和项分别有 \(k\) 基 \(1,x\) 与 \(1,x,x^2,x^3\)，故长度为六。原模已被 \(x^4\) 零化，各阶商最终恒定，所以自然映射 \(M\rightarrow\widehat M\) 本来就是同构。
\end{solution}

\begin{exercise}\label{b3:ex:13-idempotent-completion}
设 \(k\) 为域，取 \(R=k\times k\)、\(I=k\times0\)。求 \(I\) 进完备化映射，并直接证明它平坦而不忠实。
\end{exercise}
\begin{solution}
所有 \(R/I^n\) 都等于第二个 \(k\)，完备化就是投影 \((a,b)\mapsto b\)。作为 \(R\) 模，第二因子为幂等元 \((0,1)\) 生成的直和项，因此投射且平坦。非零模 \(k\times0\) 张量到第二因子后为零，所以不忠实。这也给出自然映射不单射的完备化例子。
\end{solution}

\begin{exercise}\label{b3:ex:13-finite-inverse}
设 \(k\) 为域。在 \(k[\![x,y]\!]/(x,y)^4\) 中求 \(1-x-y^2\) 的逆，写出所有全次数小于四的项。
\end{exercise}
\begin{solution}
令 \(h=x+y^2\)，所求为 \(1+h+h^2+h^3\) 的三阶截断。由于 \(h^2=x^2+2xy^2+y^4\)、\(h^3\equiv x^3\pmod{(x,y)^4}\)，答案是
\(1+x+y^2+x^2+2xy^2+x^3\)。系数 \(2\) 在特征二时为零；等比恒等式说明该式在任意特征都成立。
\end{solution}

\begin{exercise}\label{b3:ex:13-linear-equations-descend}
设 \((R,\mathfrak m)\) 为交换 Noether 局部环，\(\widehat R\) 为极大理想进完备化，\(A\) 是 \(R\) 上的 \(s\times r\) 矩阵，\(b\in R^s\)。证明：若 \(Ax=b\) 在 \(\widehat R^r\) 中有解，则在 \(R^r\) 中也有解。说明这个结论为何不保证选定的形式解本身属于 \(R^r\)。
\end{exercise}
\begin{solution}
令 \(N=\operatorname{im}(A)\subseteq R^s\)。完备化正合性说明 \(\operatorname{im}(\widehat A)=\widehat N\) 在 \(\widehat R^s\) 中的像。形式解存在意味着 \(b\) 的像属于它，由\cref{b3:prop:completion-ideal-contraction}得 \(b\in N\)，即原环中有解。解的选择可能相差核中的任意元素；例如零方程 \(0x=0\) 的任意形式级数都是解，却不必来自 \(R\)。
\end{solution}

\begin{exercise}\label{b3:ex:13-cusp-multiplicity}
设 \(k\) 为域，\(R=k[x,y]_{(x,y)}/(y^2-x^3)\)，\(B=\widehat R\) 为极大理想进完备化。求两环的维数和极大理想进重数。
\end{exercise}
\begin{solution}
\cref{b3:prop:hypersurface-tangent-cone}给出伴随分次环 \(k[X,Y]/(Y^2)\)，其次数零维数一，每个正次数维数二。故累计长度为 \(2n+1\)。由 Hilbert–Samuel 定理，\(\dim R=1\)、\(e(R)=2\)。完备化各阶商相同，得到 \(\dim B=1\)、\(e(B)=2\)。重数是此指定极大理想的重数。
\end{solution}

\begin{exercise}\label{b3:ex:13-approximate-membership}
设 \((R,\mathfrak m)\) 为交换 Noether 局部环，\(J\subseteq R\) 为理想，\(f\in R\)。证明 \(f\in J\) 当且仅当对每个整数 \(N\ge1\) 有 \(f\in J+\mathfrak m^N\)。给出只检查某个固定 \(N\) 不足的例子。
\end{exercise}
\begin{solution}
反向在有限生成模 \(R/J\) 中使用\cref{b3:thm:krull-intersection}，正向显然。设 \(k\) 为域。对任意固定 \(N\)，在 \(R=k[x]_{(x)}\) 中取 \(J=(x^{N+1})\)、\(f=x^N\)。则 \(f\in J+(x)^N\)，却 \(f\notin J\)。所以量词“每个 \(N\)”不能替换成一个事先任意选定的精度。
\end{solution}

\begin{exercise}\label{b3:ex:13-formal-nonlinear-root}
设 \(k\) 为域且 \(\operatorname{char}k\ne2\)。在 \(k[\![x]\!]\) 中构造常数项为一、满足 \(u^2=1+x\) 的唯一级数，并说明在 \(k=\mathbb Q\) 时它不属于 \(\mathbb Q[x]_{(x)}\)。
\end{exercise}
\begin{solution}
写 \(u=1+\sum_{n\ge1}a_nx^n\)。比较系数得 \(2a_1=1\)，且对 \(n\ge2\)，
\(2a_n=-\sum_{i=1}^{n-1}a_ia_{n-i}\)。由于 \(2\) 可逆，递归唯一决定所有系数，并保证每阶方程，故给出唯一形式解。若它在 \(\mathbb Q(x)\) 中，写成互素多项式比 \(P/Q\)，则 \(P^2=(1+x)Q^2\)。不可约因子 \(1+x\) 在左侧出现偶数次，右侧出现奇数次，矛盾。于是形式解的存在并不一般下降为原局部环中非线性方程的解。
\end{solution}

\begin{optionalexercise}\label{b3:ex:13-nonextended-ideal}
设 \(k\) 为可数域。证明存在 \(f\in xk[\![x]\!]\) 在 \(k(x)\) 上超越；再证明 \(k[\![x,y]\!]\) 中的非零理想 \(\bigl(y-f(x)\bigr)\) 与 \(k[x,y]_{(x,y)}\) 的交为零。
\end{optionalexercise}
\begin{solution}
\(k[x,T]\) 可数，每个非零多项式在域 \(k((x))\) 中只有有限个根，因此在 \(k(x)\) 上代数的级数只有可数多个。所有系数仅取两个不同值的级数已不可数，故可选超越 \(f\)，减去其常数项后仍超越。

写 \(F(x,y)=\sum_{i,j\ge0}a_{ij}x^iy^j\)。由于 \(f(x)\in xk[\![x]\!]\)，项 \(x^if(x)^j\) 的阶数至少为 \(i+j\)。在任一固定精度下，代入和因而只涉及有限多个 \((i,j)\)，所以 \(F(x,f(x))\) 良定义，并给出连续满射 \(k[\![x,y]\!]\rightarrow k[\![x]\!]\)。利用有限恒等式
\[
y^j-f(x)^j=(y-f(x))\sum_{\ell=0}^{j-1}y^{j-1-\ell}f(x)^\ell
\quad(j\ge1),
\]
构造
\[
G(x,y)=\sum_{i\ge0,\,j\ge1}a_{ij}x^i
       \sum_{\ell=0}^{j-1}y^{j-1-\ell}f(x)^\ell.
\]
其中每个加项的全次数至少为 \(i+j-1\)，故任一给定总次数只收到有限多个 \((i,j,\ell)\) 的贡献。于是 \(G\) 也是良定义的形式级数，逐有限阶使用上述恒等式得到
\[
F(x,y)-F(x,f(x))=(y-f(x))G(x,y).
\]
这证明代入映射的核恰为 \(\bigl(y-f(x)\bigr)\)。若原局部环的 \(P/Q\) 落入此核，\(Q(0,0)\ne0\) 使 \(Q(x,f(x))\) 可逆，因此 \(P(x,f(x))=0\)。超越性迫使 \(P=0\)。故该非零理想不是收缩理想的扩张，说明并非完备环中的所有理想都来自原环。
\end{solution}

% !TeX root = ../../Book3.tex
% 纲要标题：## 第14章　Kähler 微分与无穷小结构
% 纲要位置：工作记录/计划纲要.md，第 2400 行。
% 纲要细目（写作范围）：
% > 本章研究普遍导子及环映射的无穷小结构；带指定导子的微分环、微分理想与微分方程另见附录 E；两种主题不混称。

\chapter{Kähler 微分与无穷小结构}
\label{b3:ch:14}
\BookChapterNumber{b3:ch:14}

\begin{chapterabstract}
本章构造 Kähler 微分模与普遍导子，证明商环、复合映射、局部化及基变换公式，并用 Jacobian 计算切空间。微分还检测域扩张的可分性和平方零提升的唯一性；平展与光滑分别需要进一步的平坦性及提升存在条件。
\end{chapterabstract}

\begin{learninggoals}
\begin{itemize}
\item 从生成元与关系构造微分模，并把导子解释为平方零扩张中的截面。
\item 正确使用两条微分正合列，保留左端未必单射这一限制。
\item 由 Jacobian 矩阵计算有理点切空间，区分局部环余切空间与相对微分纤维。
\item 证明有限域扩张的微分消失判据，计算纯不可分扩张的微分。
\item 分清提升的存在性、唯一性与平坦性，并在标准例子中实际构造提升。
\end{itemize}
\end{learninggoals}

设 \(k\) 为域。抛物线的坐标环 \(B=k[x,y]/(y-x^2)\) 也可写成 \(k[t]\)，其中 \(x\) 对应 \(t\)，\(y\) 对应 \(t^2\)。对平面方程作形式微分，得到 \(\mathrm{d}y=2x\cdot\mathrm{d}x\)；换成后一种展示，则只需以 \(\mathrm{d}t\) 为基。这两种计算应当描述同一个对象。因此，对选定方程求导还不足以作为定义：我们需要由结构同态 \(k\to B\) 本身决定的微分模，使所有函数的一阶变化及其关系都能在其中表达，而不依赖选用了哪些生成元和方程。

\ProseParagraphBreak

一阶变化还可以用平方为零的无穷小量测试：把 \(x\) 送到 \(x+\varepsilon a\)，忽略 \(\varepsilon^2\) 后，乘法相容性恰好给出导子法则。于是微分不仅计算切方向，也能探测一个映射是否允许唯一提升、是否存在提升。比较这些条件时，平坦性和有限展示又承担不同作用，不能仅凭 Jacobian 的外观互相替代。

\ProseParagraphBreak

预备知识是多项式商、模与张量积、局部化及有限域扩张。\chapref{b3:ch:12}的平坦性判据用于比较非分歧与平展，\chapref{b3:ch:07}的坐标环和双数例子则提供几何背景。光滑性所需的一阶定义和计算在本章给出。

\ProseParagraphBreak
微分模及两条正合列可对照松村英之的《Commutative Ring Theory》\cite[\S25, pp. 191--195]{Matsumura1989CommutativeRingTheory}。具体多项式商的微分可以直接计算，再经基变换和局部化核对；普遍性质解释了这些计算为何相容。


% !TeX root = ../../Book3.tex
% 纲要标题：### 14.1 导子与微分模
% 纲要位置：工作记录/计划纲要.md，第 2404 行。
% 纲要细目（写作范围）：
% * 基环上线性导子
% * 微分模的泛性质
% * 多项式环的微分模

\section{导子与微分模}
\label{b3:sec:1401}


把方程中的变量作一个平方为零的改动，乘法展开时只剩常数项与一次项。导子记录的正是这个一次变化。微分模把所有可能的导子合并成一个普遍对象，使后面的商环、局部化和切空间计算都能从同一构造出发。这里研究的是代数本身允许的导子，不是在环上预先指定某一个微分算子。

\subsection{基环上的线性导子}
\label{b3:subsec:1401:01}

\begin{definition}\label{b3:def:relative-derivation}
设 \(\varphi:A\rightarrow B\) 为交换环同态，\(M\) 为 \(B\) 模。如果映射 \(D:B\rightarrow M\) 对所有 \(b,c\in B\)、\(a\in A\) 都满足
\[
D(b+c)=D(b)+D(c),\qquad D(bc)=b\cdot D(c)+c\cdot D(b),\qquad D\bigl(\varphi(a)\bigr)=0,
\]
就称 \(D\) 为 \(A\) 上的\term[daozi]{导子}[derivation]。全体这样的映射记为 \(\operatorname{Der}_A(B,M)\)。点态相加及 \((b\cdot D)(c)=b\cdot D(c)\) 使它成为 \(B\) 模。
\symidx{\(\operatorname{Der}_A(B,M)\)：从代数到模的相对导子}
\end{definition}

Leibniz 法则给出 \(D(1)=2D(1)\)，故 \(D(1)=0\)，并由归纳得到 \(D(b^n)=nb^{n-1}D(b)\)。又
\(D\bigl(\varphi(a)\cdot b\bigr)=\varphi(a)\cdot D(b)\)，所以导子确为 \(A\) 线性映射；它通常不是 \(B\) 线性的，否则 \(D(b)=b\cdot D(1)=0\)。例如 \(\mathrm{d}/\mathrm{d}x:k[x]\rightarrow k[x]\) 是 \(k\) 导子，而 \((\mathrm{d}/\mathrm{d}x)(x)=1\)，并不满足 \(k[x]\) 线性。

指定基环，就是指定哪些元素必须视为常数。设 \(k\) 为域，同一个环 \(B=k(t)\) 相对于 \(k\) 可以有非零导子：按商法则令
\[
D(f/g)=\frac{f'g-fg'}{g^2}\qquad(f,g\in k[t],\ g\ne0),
\]
就得到 \(D(t)=1\) 的 \(k\) 导子。若改为相对于 \(k(t)\) 本身，则 \(B\) 的每个元素都必须被零化，导子只能为零。

\ProseParagraphBreak
若 \(\operatorname{char}B=p>0\)，则任何导子都满足 \(D(b^p)=0\)。特别地，对特征 \(p\) 的域 \(k\)，在 \(k[x]\) 上仍有 \(D(x)=1\) 的形式求导，但所有 \(k\) 导子都零化 \(x^p\)，尽管 \(x^p\notin k\)。因此被导子零化并不意味着来自基域；导子无法检测 \(p\) 次幂方向，这也预示了不可分扩张中的特殊现象。

对交换环 \(B\)，平方零的变化可以先在双数环 \(B[\varepsilon]/(\varepsilon^2)\) 中计算。每个元素唯一写成 \(b+m\varepsilon\)，而
\[
(b+m\varepsilon)(c+n\varepsilon)=bc+(bn+cm)\varepsilon.
\]
对映射 \(D:B\to B\)，\(b\mapsto b+D(b)\varepsilon\) 保持乘法，恰好要求 \(D\) 满足 Leibniz 法则。若变化量取值于一般 \(B\) 模 \(M\)，同样令两项变化量的乘积为零，就得到下面的构造。

\begin{proposition}\label{b3:prop:derivation-square-zero}
设 \(\varphi:A\to B\) 为交换环同态，\(M\) 为 \(B\) 模，\(D:B\to M\) 为映射。在加法群 \(B\oplus M\) 上定义
\[
(b,m)(c,n)=(bc,bn+cm).
\]
这给出一个单位为 \((1,0)\) 的交换环，并通过结构同态
\[
A\longrightarrow B\oplus M,\qquad a\longmapsto\bigl(\varphi(a),0\bigr)
\]
成为 \(A\) 代数。理想 \(0\oplus M\) 的平方为零。映射 \(b\mapsto\bigl(b,D(b)\bigr)\) 是投影 \(B\oplus M\to B\) 的 \(A\) 代数截面，当且仅当 \(D\in\operatorname{Der}_A(B,M)\)。
\end{proposition}
\begin{proof}
三因子乘积的第二坐标都是 \(bc\,u+b d\,n+c d\,m\)，其中第三因子为 \((d,u)\)，故结合律成立；交换律、分配律和单位可直接由公式检验。所列结构映射是环同态，投影到 \(B\) 也保持这一 \(A\) 代数结构。平方为零来自两因子的第一坐标同时为零。截面必须有 \(\bigl(b,D(b)\bigr)\) 的形式，保持加法、乘法以及 \(A\) 的像分别就是导子的三个条件。
\end{proof}

这个环称为平方零扩张 \(B\oplus M\)。它把“无穷小变化”变成确切的环运算；这里的无穷小并非一个任意小的实数，而是乘积被规定为零的理想元素。

\subsection{微分模的泛性质}
\label{b3:subsec:1401:02}

目标模 \(M\) 改变时，\(\operatorname{Der}_A(B,M)\) 也随之改变。我们希望先用一个固定的 \(B\) 模记录导子的关系，再通过从这个模到 \(M\) 的线性映射，给每种一阶变化指定实际的值。

\begin{theorem}\label{b3:thm:universal-differentials}
设 \(A\) 为交换环，\(B\) 为交换 \(A\) 代数。存在 \(B\) 模 \(\Omega_{B/A}\) 与 \(A\) 导子 \(\mathrm{d}:B\rightarrow\Omega_{B/A}\)，使对任意 \(B\) 模 \(M\)，映射
\[
\operatorname{Hom}_B(\Omega_{B/A},M)\longrightarrow
\operatorname{Der}_A(B,M),\qquad h\longmapsto h\circ \mathrm{d}
\]
为自然双射。二者组成的对象在唯一的保持 \(\mathrm{d}\) 的同构意义下唯一，称为 \term[kahler-weifen-mo]{Kähler 微分模}[module of Kähler differentials]；\(\mathrm{d}\) 称为\term[pubian-daozi]{普遍导子}[universal derivation]。
因此，对任意 \(D\in\operatorname{Der}_A(B,M)\)，都有下列交换图；\(\mathrm{d}\) 和 \(D\) 是 \(A\) 导子，虚线 \(h\) 是唯一的 \(B\) 线性映射：
\[
\begin{tikzcd}[column sep=large,row sep=large]
B\arrow[r,"\mathrm{d}"]\arrow[dr,"D"']
&\Omega_{B/A}\arrow[d,dashed,"h"]\\
&M
\end{tikzcd}
\]
\symidx{\(\Omega_{B/A}\)、\(\mathrm{d}_{B/A}\)：Kähler 微分模与普遍导子}
\end{theorem}
\begin{proof}
记结构同态为 \(\varphi:A\to B\)。取以集合 \(B\) 为指标的自由 \(B\) 模 \(F=\bigoplus_{b\in B}Be_b\)，模掉下列元素生成的子模：
\[
e_{b+c}-e_b-e_c,\qquad
 e_{bc}-be_c-ce_b,\qquad e_{\varphi(a)}
\quad(b,c\in B,\ a\in A).
\]
记所得商为 \(\Omega_{B/A}\)，令 \(\mathrm{d}b\) 为 \(e_b\) 的像。三个关系使 \(\mathrm{d}\) 满足导子的三个公理。对任意导子 \(D\)，从自由模出发的唯一线性映射 \(e_b\mapsto D(b)\) 恰好零化这些关系，故下降为 \(h:\Omega_{B/A}\rightarrow M\)。因各 \(\mathrm{d}b\) 生成微分模，这个 \(h\) 唯一。

若另一对 \((\Omega',\mathrm{d}')\) 也有此性质，分别应用泛性质得到 \(u:\Omega\rightarrow\Omega'\)、\(v:\Omega'\rightarrow\Omega\)，满足 \(u\mathrm{d}=\mathrm{d}'\)、\(v\mathrm{d}'=\mathrm{d}\)。于是 \(vu\mathrm{d}=\mathrm{d}\)，唯一性给出 \(vu=1\)，同理 \(uv=1\)。对 \(M\) 的同态作复合即说明双射的自然性。
\end{proof}

构造中的 \(\mathrm{d}b\) 先作为形式符号记录变化，并只施加导子公理要求的关系；选定 \(B\) 线性映射 \(h:\Omega_{B/A}\to M\) 后，这些符号才取得值 \(h(\mathrm{d}b)=D(b)\)。若 \(B\) 由 \(b_i\) 作为 \(A\) 代数生成，则 \(\Omega_{B/A}\) 由 \(\mathrm{d}b_i\) 作为 \(B\) 模生成，因为乘法法则能把任意多项式的微分化成这些元素的线性组合。生成元之间的代数关系则会带来微分关系。

\subsection{多项式环的微分模}
\label{b3:subsec:1401:03}

\begin{proposition}\label{b3:prop:polynomial-differentials}
设 \(A\) 为交换环，\(n\ge0\) 为整数，\(P=A[X_1,\ldots,X_n]\)。则
\[
\Omega_{P/A}=\bigoplus_{i=1}^nP\,\mathrm{d}X_i,
\qquad \mathrm{d}f=\sum_{i=1}^n\dfrac{\partial f}{\partial X_i}\,\mathrm{d}X_i.
\]
对任意 \(P\) 模 \(M\) 和元素 \(m_1,\ldots,m_n\in M\)，存在唯一 \(A\) 导子 \(D:P\rightarrow M\) 满足 \(D(X_i)=m_i\)。
\end{proposition}
\begin{proof}
在单项式上反复使用 Leibniz 法则，任何导子都必须满足所列公式。反向用该公式定义 \(D(f)=\sum_i(\partial f/\partial X_i)m_i\)，多项式形式偏导的乘法公式直接给出 Leibniz 法则，且它零化 \(A\)，所以存在性也成立。于是自由模 \(\bigoplus_iP\,\mathrm{d}X_i\) 具有普遍导子的泛性质，等同于 \(\Omega_{P/A}\)。也可用导子 \(\partial/\partial X_j\) 检查线性无关：它将假定关系 \(\sum_if_i\,\mathrm{d}X_i=0\) 送到 \(f_j=0\)。
\end{proof}

例如在 \(k[x,y]\) 中，\(\mathrm{d}(x^2y)=2xy\,\mathrm{d}x+x^2\,\mathrm{d}y\)。若再施加关系 \(y^2=x^3\)，则必有 \(2y\,\mathrm{d}y-3x^2\,\mathrm{d}x=0\)；下一节将证明怎样获得全部这样的关系。对形式幂级数环则需要额外区分连续导子与任意代数导子，本章的多项式结论不未经证明地推广到无限级数。

% !TeX root = ../../Book3.tex
% 纲要标题：### 14.2 微分模的正合列与基变换
% 纲要位置：工作记录/计划纲要.md，第 2410 行。
% 纲要细目（写作范围）：
% * 商环的余法正合列
% * 方程微分给出的展示；尖点微分模的非零挠元及其零化理想
% * 复合环映射的微分序列
% * 局部化及基变换

\section{微分模的正合列与基变换}
\label{b3:sec:1402}


给定代数的生成元以后，导子的值不能任意选择：它必须保持方程的微分关系。两条基本正合列分别处理“添加关系”和“改变基环”，局部化与基变换公式则使这些计算能够转到所需的局部环或基域。

\subsection{商环的余法正合列}
\label{b3:subsec:1402:01}

商映射 \(B\to B/J\) 把 \(f\in J\) 宣布为零，因而微分中也必须添加 \(\mathrm{d}f=0\) 的关系。两条方程乘积的微分仍含方程因子，转到 \(B/J\) 后就为零，所以 \(J^2\) 中的元素不会提供新的这一阶关系；记录这些方程类的是 \(J/J^2\)。

\begin{theorem}[余法正合列]\label{b3:thm:conormal-sequence}
设 \(A\) 为交换环，\(B\) 为交换 \(A\) 代数、\(J\subseteq B\) 为理想、\(C=B/J\)。有 \(C\) 模正合列
\[
J/J^2\xrightarrow{\delta}C\otimes_B\Omega_{B/A}
\longrightarrow\Omega_{C/A}\longrightarrow0,
\qquad \delta(f+J^2)=1\otimes \mathrm{d}f.
\]
模 \(J/J^2\) 称为相应商映射的\term[yufa-mo]{余法模}[conormal module]。映射 \(\delta\) 一般不单射。
\end{theorem}
\begin{proof}
若 \(f\in J^2\)，将它写成有限和 \(\sum_i u_iv_i\)、\(u_i,v_i\in J\)，则 \(\mathrm{d}f=\sum_i(u_i\,\mathrm{d}v_i+v_i\,\mathrm{d}u_i)\) 在张量到 \(C\) 后为零，故 \(\delta\) 良定义。又 \(\mathrm{d}(bf)=b\,\mathrm{d}f+f\,\mathrm{d}b\)，第二项在 \(C\) 上为零，所以 \(\delta\) 是 \(C\) 线性的。

将 \(C\otimes_B\Omega_{B/A}\) 模掉全部 \(1\otimes \mathrm{d}f\)、\(f\in J\)，记商为 \(Q\)。映射 \(b+J\mapsto[1\otimes \mathrm{d}b]\) 良定义，因为改变代表元只增加一个被模掉的微分。它是 \(A\) 导子。若 \(D:C\rightarrow M\) 是任意 \(A\) 导子，复合 \(B\rightarrow C\) 后的导子对应一个 \(C\otimes_B\Omega_{B/A}\rightarrow M\)，且零化各 \(\mathrm{d}f\)，故唯一通过 \(Q\)。于是 \(Q\) 有 \(\Omega_{C/A}\) 的泛性质，得到所需正合列。
\end{proof}

\begin{corollary}\label{b3:cor:differential-presentation}
设 \(A\) 为交换环，\(f_1,\ldots,f_s\in A[X_1,\ldots,X_n]\)，\(B=A[X_1,\ldots,X_n]/(f_1,\ldots,f_s)\)，\(x_i\) 为变量的像。则有展示
\[
B^s\xrightarrow{J_f^{\mathsf t}} B^n\longrightarrow\Omega_{B/A}\longrightarrow0,
\qquad (J_f)_{ij}=\dfrac{\partial f_i}{\partial X_j}(x).
\]
因此有限展示交换代数的微分模为有限展示模。
特别地，若 \(k\) 为特征不等于 \(2,3\) 的域，取 \(B=k[x,y]/(y^2-x^3)\)，并以 \(x,y\) 兼记变量在商中的像，则
\[
\omega=2x\,\mathrm{d}y-3y\,\mathrm{d}x
\]
是 \(\Omega_{B/k}\) 中的非零挠元，且 \(y\omega=x^2\omega=0\)。
\end{corollary}
\begin{proof}
对 \(P=A[X]\) 的商应用余法正合列，\(\Omega_{P/A}\) 为以 \(\mathrm{d}X_i\) 为基的自由模。若 \(f=\sum_i g_if_i\)，在 \(\Omega_{P/A}\) 张量到 \(B\) 后有 \(\mathrm{d}f=\sum_i g_i\,\mathrm{d}f_i\)，因为含 \(f_i\,\mathrm{d}g_i\) 的项消失。因此全部微分关系由各 \(\mathrm{d}f_i\) 生成。映射 \(B^s\to B^n\) 把第 \(i\) 个标准基向量送到 \(\mathrm{d}f_i\) 的坐标列 \(\bigl(\frac{\partial f_i}{\partial X_1}(x),\ldots,\frac{\partial f_i}{\partial X_n}(x)\bigr)^{\mathsf t}\)。通常 Jacobian 矩阵将第 \(i\) 个方程的偏导放在第 \(i\) 行，这里将它们放在第 \(i\) 列，所以所需矩阵是其转置。

在尖点的情形中，上述展示给出
\[
\Omega_{B/k}=(B\,\mathrm{d}x\oplus B\,\mathrm{d}y)/B(2y\,\mathrm{d}y-3x^2\,\mathrm{d}x).
\]
由 \(y^2=x^3\)，有
\[
y\omega=x(2y\,\mathrm{d}y-3x^2\,\mathrm{d}x)=0,
\qquad x^2\omega=y(2y\,\mathrm{d}y-3x^2\,\mathrm{d}x)=0.
\]
为证 \(\omega\ne0\)，由\cref{b3:prop:cusp-normalization}将 \(B\) 识别为 \(k[t^2,t^3]\subseteq k[t]\)，其中 \(x=t^2\)、\(y=t^3\)。若 \(\omega=0\)，它在自由模中的代表必须是关系生成元的一个 \(B\) 倍数，因此存在 \(h\in B\) 满足
\[
-3y=-3hx^2,\qquad 2x=2hy.
\]
由于 \(2\) 在 \(k\) 中可逆，后一等式在分式域 \(k(t)\) 中迫使 \(h=t^{-1}\)，但 \(B\subseteq k[t]\) 不含此元素，矛盾。因此 \(\omega\) 非零，而 \(y\ne0\) 是整环 \(B\) 的非零元素，所以上面的零化等式也证明它是挠元。
\end{proof}

余法映射不单射，表示非零的方程类也可能没有留下非零的微分关系。例如在特征 \(p\) 的域上取 \(B=k[t]\)、\(J=(t^p)\)。由于 \(t^p\notin(t^{2p})\)，它在 \(J/J^2\) 中的类非零，却满足 \(\delta(t^p)=p t^{p-1}\mathrm{d}t=0\)。这一类生成整个余法模，所以此例的余法映射甚至为零。

\subsection{复合环映射的微分序列}
\label{b3:subsec:1402:02}

对 \(A\to B\to C\)，从相对于 \(A\) 求导改为相对于 \(B\) 求导，意味着将 \(B\) 的像也视为常数。因此要在 \(\Omega_{C/A}\) 中进一步模掉这些元素的微分，末端微分模就成为中间微分模的商。

\begin{theorem}\label{b3:thm:differential-transitivity}
设 \(A\rightarrow B\xrightarrow{g}C\) 为交换环同态，则有 \(C\) 模正合列
\[
C\otimes_B\Omega_{B/A}\xrightarrow{\alpha}\Omega_{C/A}
\xrightarrow{\beta}\Omega_{C/B}\longrightarrow0,
\]
其中 \(\alpha(c\otimes \mathrm{d}b)=c\,\mathrm{d}\bigl(g(b)\bigr)\)、\(\beta(\mathrm{d}_{C/A}c)=\mathrm{d}_{C/B}c\)。
\end{theorem}
\begin{proof}
复合 \(B\rightarrow C\xrightarrow{\mathrm{d}_{C/A}}\Omega_{C/A}\) 是 \(A\) 导子，故给出 \(\alpha\)。在 \(\Omega_{C/A}\) 中模掉各 \(\mathrm{d}\bigl(g(b)\bigr)\) 生成的子模后，普遍导子就零化 \(B\) 的像，因此成为 \(B\) 导子。任意 \(B\) 导子首先是 \(A\) 导子，并零化这些指定元素，所以唯一经过此商。该商即 \(\Omega_{C/B}\)，证明正合性。
\end{proof}

首箭头可能不单射。例如 \(A=k\)、\(B=k[t^p]\subseteq C=k[t]\)、\(\operatorname{char}k=p\) 时，\(t^p\) 是 \(B\) 的一个多项式生成元，所以 \(\mathrm{d}_{B/k}(t^p)\ne0\)；但它映到 \(\Omega_{C/k}\) 中时变成 \(\mathrm{d}_{C/k}(t^p)=0\)。

\subsection{微分模的局部化与基变换}
\label{b3:subsec:1402:03}

只在上环 \(B\) 中倒置乘法集 \(S\) 时，基环仍为 \(A\)。分母 \(s\in S\) 的微分一般不为零，所以分式的微分需要包含 \(\mathrm{d}s\) 项。若分母来自基环中的乘法集 \(T\)，它们本来就必须被零化；同步倒置基环与上环后，这些分母的逆也成为常数，此时相对基环是 \(T^{-1}A\)。

\begin{proposition}\label{b3:prop:differential-localization}
设 \(\varphi:A\to B\) 为交换环同态，\(S\subseteq B\) 为乘法集。则有 \(S^{-1}B\) 模自然同构
\[
\Omega_{S^{-1}B/A}\cong S^{-1}\Omega_{B/A},
\qquad \mathrm{d}(b/s)=\dfrac{\mathrm{d}b}{s}-\dfrac{b\,\mathrm{d}s}{s^2}.
\]
其中 \(b\in B\)、\(s\in S\)。
若 \(T\subseteq A\) 为乘法集，记 \(T^{-1}B\) 为 \(B\) 对 \(\varphi(T)\) 的局部化，则还有 \(T^{-1}B\) 模自然同构
\[
\Omega_{T^{-1}B/T^{-1}A}\cong T^{-1}\Omega_{B/A}.
\]
右侧按 \(A\) 在 \(\Omega_{B/A}\) 上的作用局部化。
\end{proposition}
\begin{proof}
对取值于任意 \(S^{-1}B\) 模 \(M\) 的 \(A\) 导子 \(D:B\rightarrow M\)，若它能延拓为 \(\widetilde D:S^{-1}B\to M\)，则等式 \(\widetilde D(s\cdot s^{-1})=0\) 迫使其满足上述分式公式。按此公式定义延拓，检查良定义：若 \(b/s=b'/s'\)，则存在 \(t\in S\) 使 \(t(s'b-sb')=0\)。取导数并在 \(M\) 中将 \(t,s,s'\) 视为单位；由于 \(s'b-sb'=0\) 在局部化后成立，得到
\[
s'Db+bDs'-sDb'-b'Ds=0.
\]
连同 \(s'b=sb'\)，整理即得两个分式公式相等。对分式加法和乘法展开可得加法性与 Leibniz 法则：乘法中的导数项恰分成对第一个和第二个因子的两部分。故延拓存在且唯一，由微分模的泛性质得到第一个同构。

对于同步局部化，取任意 \(T^{-1}B\) 模 \(M\)。每个 \(A\) 导子 \(D:B\to M\) 同样唯一延拓为 \(\widetilde D:T^{-1}B\to M\)。对 \(t\in T\)，有
\[
\widetilde D\bigl(\varphi(t)^{-1}\bigr)
=-\varphi(t)^{-2}\cdot D\bigl(\varphi(t)\bigr)=0.
\]
它又零化 \(\varphi(A)\)，所以零化 \(T^{-1}A\) 的像，因而是 \(T^{-1}A\) 导子。反过来，限制一个 \(T^{-1}A\) 导子便得到 \(A\) 导子。两种操作互逆，再用微分模的泛性质得到第二个同构。
\end{proof}

\begin{proposition}\label{b3:prop:differential-base-change}
设 \(B\) 为交换 \(A\) 代数，\(A\to A'\) 为交换环同态，令 \(B'=A'\otimes_AB\)。有自然同构
\[
\Omega_{B'/A'}\cong B'\otimes_B\Omega_{B/A}.
\]
不需要 \(A'\) 在 \(A\) 上平坦。
\end{proposition}
\begin{proof}
任取 \(B'\) 模 \(M\)，沿环同态 \(B\to B'\)、\(b\mapsto1\otimes b\) 将它看作 \(B\) 模。将 \(A'\) 导子 \(\Delta:B'\to M\) 与这个环同态复合，便给出 \(A\) 导子 \(D:B\to M\)。反过来，给定这样的 \(D\)，在纯张量上定义
\[
\Delta(a'\otimes b)=(a'\otimes1)\cdot D(b).
\]
由 \(D\) 的 \(A\) 线性及张量积中 \(A\) 的两个像相等，此公式保持平衡关系，因而对有限和良定义。它零化 \(A'\) 的像；对纯张量的乘积使用 \(D\) 的 Leibniz 法则，再由加法展开，就得到 \(\Delta\) 的 Leibniz 法则。复合与延拓互逆，因为任何 \(A'\) 导子都满足所列纯张量公式。因此两种导子自然一一对应。

又由普遍导子的泛性质，\(D\) 唯一写成 \(h\circ\mathrm{d}_{B/A}\)，其中 \(h:\Omega_{B/A}\to M\) 为 \(B\) 线性映射。这样的 \(h\) 又唯一延拓为 \(B'\) 线性映射
\[
B'\otimes_B\Omega_{B/A}\longrightarrow M,
\qquad b'\otimes\eta\longmapsto b'\cdot h(\eta).
\]
所以右侧张量积和 \(\Omega_{B'/A'}\) 对每个目标模都表示同一个导子函子，普遍对象的唯一性给出所需同构。这个同构将 \(1\otimes\mathrm{d}b\) 送到 \(\mathrm{d}(1\otimes b)\)。
\end{proof}

这两个公式允许先在多项式商上写出 Jacobian 展示，再局部化到指定点，或改变基域后重新考察方程。代数作为环的约化性和正则性可能在基变换后改变；光滑性则是相对于基环的性质，并在任意基变换下保持，具体论证见\secref{b3:sec:1405}。

% !TeX root = ../../Book3.tex
% 纲要标题：### 14.3 切空间与 Jacobian 矩阵
% 纲要位置：工作记录/计划纲要.md，第 2417 行。
% 纲要细目（写作范围）：
% * 极大理想模平方与余切空间
% * 仿射方程的切空间计算
% * 一般素点的剩余域微分方向；不可分闭点与基变换后的幂零结构

\section{切空间与 Jacobian 矩阵}
\label{b3:sec:1403}

一条曲线在一个点附近可能有几条分支，也可能含有重根；只看点集不能分辨这些现象。把点的坐标改成 \(a_i+v_i\varepsilon\)、其中 \(\varepsilon^2=0\)，方程留下的一次关系便限制了允许的方向。先从有理点说明这个计算，再讨论剩余域不是基域的情形。


\subsection{极大理想模平方与余切空间}
\label{b3:subsec:1403:01}

\begin{definition}\label{b3:def:zariski-tangent-space}
设 \(k\) 为域，\(B\) 为有限型交换 \(k\) 代数，\(a:B\to k\) 为一个 \(k\) 代数同态，\(\mathfrak m=\ker a\)。以 \(a\) 赋予 \(k\) 一个 \(B\) 模结构。向量空间
\[
T_a(B)=\operatorname{Der}_k(B,k)
\]
称为 \(B\) 在这个 \(k\) 有理点处的\term[zariski-qiekongjian]{Zariski 切空间}[Zariski tangent space]；\(\mathfrak m/\mathfrak m^2\) 是该点的余切空间。这里的 \(B\) 可以含有非零幂零元，不要求它是某个点集的约化坐标环。
\symidx{\(T_a(B)\)：有理点处的 Zariski 切空间}
\end{definition}

\begin{proposition}\label{b3:prop:tangent-cotangent-rational}
设 \(k\) 为域，\(B\) 为交换 \(k\) 代数，\(a:B\to k\) 为 \(k\) 代数同态，\(\mathfrak m=\ker a\)。则有自然同构
\[
\mathfrak m/\mathfrak m^2\xrightarrow{\sim}
k\otimes_B\Omega_{B/k},\qquad
\bar b\longmapsto1\otimes \mathrm{d}b,
\]
以及
\[
\operatorname{Der}_k(B,k)
\cong\operatorname{Hom}_k(\mathfrak m/\mathfrak m^2,k).
\]
若以局部环 \(B_{\mathfrak m}\) 代替 \(B\)，所得余切空间及切空间相同。
\end{proposition}
\begin{proof}
对 \(b,c\in\mathfrak m\)，\(1\otimes \mathrm{d}(bc)=0\)，故第一个映射良定义。每个 \(b\in B\) 都写成 \(a(b)+\bigl(b-a(b)\bigr)\)，其中后项属于 \(\mathfrak m\)，所以这个映射满射。反向定义
\[
D:B\longrightarrow\mathfrak m/\mathfrak m^2,
\qquad D(b)=b-a(b)\pmod{\mathfrak m^2}.
\]
将 \(b=a(b)+b_0\)、\(c=a(c)+c_0\) 相乘，模掉 \(\mathfrak m^2\) 后得到
\(D(bc)=a(b)\cdot D(c)+a(c)\cdot D(b)\)。故 \(D\) 是取值于剩余域模的 \(k\) 导子，由微分模泛性质诱导逆映射。第二个同构由同一泛性质及第一个同构得到。

若 \(s\notin\mathfrak m\)，它在 \(\mathfrak m/\mathfrak m^2\) 上的乘法就是非零标量 \(a(s)\)，已经可逆。因此局部化不改变该商；由局部化正合性，它又等于
\(\mathfrak mB_{\mathfrak m}/(\mathfrak mB_{\mathfrak m})^2\)。导子的分式延拓公式也给出相同的切空间。
\end{proof}

例如，在 \(k[x_1,\ldots,x_n]_{(x_1,\ldots,x_n)}\) 的原点，记极大理想为 \(\mathfrak m\)。消失于原点的函数模 \(\mathfrak m^2\) 后，只留下 \(\sum_i c_ix_i\) 这样的一次部分；\(\mathfrak m/\mathfrak m^2\) 保存的是函数的一阶线性信息。切向量 \(v=(v_i)\) 则是作用在这些一次部分上的线性泛函，将 \(\sum_i c_ix_i\) 送到 \(\sum_i c_iv_i\)。这就是切空间与余切空间互为对偶的具体含义。

上述导子还有直接的点的解释。给定 \(D\in T_a(B)\)，映射
\[
B\longrightarrow k[\varepsilon]/(\varepsilon^2),\qquad
b\longmapsto a(b)+D(b)\cdot\varepsilon
\]
是一个模 \(\varepsilon\) 后等于 \(a\) 的 \(k\) 代数同态；乘法相容性恰是导子的公式。反过来，每个这样的同态都唯一给出一个 \(D\)。因此，同一个切向量既可视为导子，也可视为余切空间上的线性泛函，或点向双数环中的一阶延伸。

\subsection{仿射方程的切空间计算}
\label{b3:subsec:1403:02}

\begin{theorem}\label{b3:thm:jacobian-tangent-space}
设 \(k\) 为域，\(B=k[X_1,\ldots,X_n]/(f_1,\ldots,f_s)\)，且 \(a=(a_1,\ldots,a_n)\in k^n\) 满足所有 \(f_i(a)=0\)。以 \(X_j\mapsto a_j\) 的求值同态 \(B\to k\) 定义该点的切空间 \(T_a(B)=\operatorname{Der}_k(B,k)\)。令
\[
J_f(a)=\left(\dfrac{\partial f_i}{\partial X_j}(a)\right)_{
1\le i\le s,\,1\le j\le n}.
\]
则在变量坐标给出的识别下，
\[
T_a(B)=\ker\bigl(J_f(a):k^n\to k^s\bigr),
\qquad \dim_kT_a(B)=n-\operatorname{rank}J_f(a).
\]
这个矩阵称为所给方程组的\term[jacobian-juzhen]{Jacobian 矩阵}[Jacobian matrix]，亦称雅可比矩阵\termidx[yakebi-juzhen]{雅可比矩阵}。\footnote{“雅可比矩阵”见菲赫金哥尔茨《微积分学教程》第一卷\cite[第六章“函数矩阵(雅可比矩阵)的乘法”]{Fichtenholz2005CalculusChineseI}。}
\end{theorem}
\begin{proof}
由双数值点的描述，一个方向 \(v=(v_j)\) 对应 \(X_j\mapsto a_j+v_j\varepsilon\)。它能从多项式环下降到 \(B\)，当且仅当每个 \(f_i\) 的像为零。因为 \(\varepsilon^2=0\)，二次以上的改变量全部消失，故
\[
f_i(a_1+v_1\varepsilon,\ldots,a_n+v_n\varepsilon)
=f_i(a)+\varepsilon\sum_j
\dfrac{\partial f_i}{\partial X_j}(a)v_j.
\]
由 \(f_i(a)=0\)，这一条件恰为 \(J_f(a)v=0\)，得到核的描述及维数公式。

\cref{b3:cor:differential-presentation}还将这一计算与环本身的微分模联系起来：展示 \(B^s\to B^n\to\Omega_{B/k}\to0\) 的关系矩阵为 \(J_f^{\mathsf t}\)。张量到 \(k\) 后仍右正合，再取 \(k\) 线性对偶，得到
\[
0\longrightarrow\operatorname{Hom}_k(k\otimes_B\Omega_{B/k},k)
\longrightarrow k^n\xrightarrow{J_f(a)}k^s.
\]
前一命题把左端自然识别为切空间。因此坐标与方程展示虽会改变矩阵，所得导子空间始终由 \(B\) 及这个有理点决定。
\end{proof}

例如设 \(k\) 为域，在 \(B=k[x,y]/(xy-1)\) 中，每个有理点 \((a,b)\) 都有 \(ab=1\)，Jacobian 行为 \((b,a)\)，秩为一，所以切空间由 \(bv_x+av_y=0\) 给出，是一维的。在尖点曲线 \(y^2=x^3\) 的原点，Jacobian 行 \((-3x^2,2y)\) 为零，故切空间是整个 \(k^2\)，任意特征都如此。曲线 \(xy=0\) 的原点也有二维切空间，虽然它有两条分支，而尖点只有一条。这说明切空间维数能发现某些异常，却不能单独数出分支。

\ProseParagraphBreak
尖点的参数化 \(t\mapsto(t^2,t^3)\) 把源直线原点的一阶点 \(t=v\varepsilon\) 送到 \(\bigl((v\varepsilon)^2,(v\varepsilon)^3\bigr)=(0,0)\)。因此它在原点诱导的切映射为零，源点的一阶变化在像中全部消失。

竖直方向则显示 Zariski 一阶切方向与切锥约化方向之间的差别。映射 \(x\mapsto0,y\mapsto\varepsilon\) 在 \(k[\varepsilon]/(\varepsilon^2)\) 中满足尖点方程，因此给出一个切向量。若提升到 \(k[\varepsilon]/(\varepsilon^3)\)，保持这一阶值就必须写成 \(x=a\varepsilon^2\)、\(y=\varepsilon+b\varepsilon^2\)。此时 \(x^3=0\)，而 \(y^2=\varepsilon^2\ne0\)，所以没有这样的提升。Zariski 切空间只检查一次关系；切锥 \(Y^2=0\) 则已保留最低二次关系，其约化支集只有水平直线。因而一阶允许的方向未必能延续成更高阶曲线。

\ProseParagraphBreak
还须保留方程中的幂零结构。\(k[t]/(t^2)\) 只有一个域值点，但该点切空间为一维；约化后的坐标环 \(k\) 则有零切空间。两者的点集相同，一阶延伸却不同。若已经把理想替换成根理想再计算，便无法恢复丢掉的方向。

\ProseParagraphBreak
Jacobian 计算还可以寻找参数族中切空间异常增大的位置。以 \(t\) 为参数，在特征零域 \(k\) 上考虑曲线族
\[
F_t(x,y)=y^2-x^3-t=0.
\]
固定参数后，切空间维数为二恰好要求 \(F_t=0\)、\(\partial F_t/\partial x=-3x^2=0\)、\(\partial F_t/\partial y=2y=0\)。在 \(k[t,x,y]\) 中，相应理想为
\[
(y^2-x^3-t,-3x^2,2y)=(y,x^2,t).
\]
右边的包含关系可直接由 \(t=y^2-x^3-F_t\) 核验，反向也由各生成元的表达式得到。这个 Jacobian 理想记录异常位置的代数结构，其商中还保留 \(x^2=0\) 的加厚信息；若只求底层异常点集，则取根理想得到 \((x,y,t)\)，只剩原点。要判断哪些参数会产生异常，还须消去 \(x,y\)，收缩到参数环得到
\[
(y,x^2,t)\cap k[t]=(t).
\]
这是因为 \(t\) 已属于该理想，而属于它的参数多项式必在原点取值为零。于是异常只在 \(t=0\) 出现；其余参数的每个几何点都有一维切空间。

\subsection{非完美基域上的切空间}
\label{b3:subsec:1403:03}

这里先比较一般素点处的两个一阶对象，差别并不限于不完美基域。设 \(k\) 为域，\(B\) 为交换 \(k\) 代数，\(\mathfrak p\in\operatorname{Spec}B\)，记 \(\kappa=\kappa(\mathfrak p)\)。剩余域 \(\kappa\) 未必等于 \(k\)。局部环自身的余切空间
\(\mathfrak pB_{\mathfrak p}/(\mathfrak pB_{\mathfrak p})^2\)
与相对微分的纤维 \(\kappa\otimes_B\Omega_{B/k}\) 出自不同构造。对 \(B_{\mathfrak p}\to\kappa\) 应用余法正合列及微分局部化，得到下列 \(\kappa\) 向量空间的右正合列：
\begin{equation}\label{b3:eq:cotangent-residue-sequence}
\begin{tikzcd}[column sep=small,row sep=0.8em]
\dfrac{\mathfrak pB_{\mathfrak p}}{(\mathfrak pB_{\mathfrak p})^2}
  \arrow[r]
&\kappa\otimes_B\Omega_{B/k}
  \arrow[r]
&\Omega_{\kappa/k}\arrow[r]&0\\
{\begin{gathered}\text{局部环自身}\\\text{的一阶结构}\end{gathered}}
&{\begin{gathered}\text{相对于基域 }k\\\text{的一阶结构}\end{gathered}}
&{\begin{gathered}\text{剩余域扩张}\\\text{的一阶结构}\end{gathered}}&
\end{tikzcd}
\end{equation}
在 \(k\) 有理点处，\cref{b3:prop:tangent-cotangent-rational}证明左箭头为同构。一般点处的左箭头不必单射，右端也可能非零；它记录了剩余域扩张本身贡献的微分。记左箭头的像为 \(W\)，则中间空间包含子空间 \(W\)，商空间为 \(\Omega_{\kappa/k}\)。向量空间的短正合列可以选基分裂，所以中间空间可非典范地写成 \(W\oplus\Omega_{\kappa/k}\)；但没有自然指定的补空间，而且 \(W\) 只是左端的一个商，不能直接把整个局部余切空间与 \(\Omega_{\kappa/k}\) 作直和。

即使 \(k\) 完美，在 \(B=k[t]\) 的一般点 \(\mathfrak p=(0)\) 处也有 \(\kappa=k(t)\)；由多项式微分与局部化公式，\(\Omega_{\kappa/k}=k(t)\,\mathrm{d}t\ne0\)。此时局部环是域，局部余切空间为零，相对微分纤维却仍有一维。非完美基域上的额外现象是：连闭点的有限剩余域扩张也可能贡献这样的非零微分。

\begin{example}\label{b3:exm:inseparable-tangent}
设 \(p\) 为素数、\(k=\mathbb F_p(s)\)，其中 \(s\) 为超越元。令
\[
L=k[u]/(u^p-s).
\]
由于 \(s\) 在 \(k\) 中不是 \(p\) 次幂，\(u^p-s\) 不可约，故 \(L\) 为域。其唯一局部环的极大理想为零，局部环余切空间也为零。但微分展示给出
\[
\Omega_{L/k}=L\,\mathrm{d}u,
\]
因为 \(s\) 是基域中的常数，且 \(k\) 的特征为 \(p\)，故 \(\mathrm{d}(u^p-s)=0\)。该点的剩余域就是 \(L\)，所以相对微分纤维 \(L\otimes_L\Omega_{L/k}=\Omega_{L/k}\) 非零，完全来自剩余域扩张的不可分性，并非来自局部极大理想。若扩域到 \(K=k(s^{1/p})\)，则
\[
K\otimes_kL\cong K[v]/(v^p),\qquad v=u-s^{1/p},
\]
其唯一 \(K\) 有理点由 \(v\mapsto0\) 给出，局部余切空间 \((v)/(v^2)\) 以 \(v\) 的类为基，故该点的切空间为一维。这里 \(v\) 的像是非零幂零元：原来的环 \(L\) 是域，改变基域后却出现了这个加厚点。
\end{example}

不可约性的一个直接检验如下。在代数闭包中，\(T^p-s=(T-\alpha)^p\)。若 \(\alpha\) 的最小多项式次数 \(r<p\)，它只能为 \((T-\alpha)^r\)，其 \(T^{r-1}\) 系数使 \(r\alpha\in k\)，从而 \(\alpha\in k\)，与 \(s\) 非 \(p\) 次幂矛盾。最后一事实可由有理函数在 \(s=0\) 处的零点阶数核验：\(p\) 次幂的阶数是 \(p\) 的倍数，\(s\) 的阶数则为一。

\ProseParagraphBreak
这一区别在不完美基域上尤其必要。局部环是域，因而在环论意义下最简单，却不能据此断言它相对于给定基域光滑。后面的正则局部环与光滑性判据将分别处理环自身的性质和扩张基域后的性质；本章的微分计算已经显示两者为何不能合并。

% !TeX root = ../../Book3.tex
% 纲要标题：### 14.4 域扩张、可分性与微分
% 纲要位置：工作记录/计划纲要.md，第 2399 行。
% 纲要细目（写作范围）：
% * 单个代数元的微分与有限域扩张的可分性判据
% * 指数一纯不可分扩张；子域 \(K L^p\) 及微分基的构造
% * 纯不可分扩张的次数与微分维数；基变换后的幂零结构

\section{域扩张、可分性与微分}
\label{b3:sec:1404}

代数元的关系 \(f(\alpha)=0\) 对导子给出 \(f'(\alpha)\mathrm{d}\alpha=0\)。若导数在域中非零，这个方向被迫消失；若关系是 \(p\) 次幂，求导却不能看见它。本节从单个代数元的计算出发，说明这种差别如何检测一般有限域扩张的可分性，再比较纯不可分扩张的次数与微分维数。

\subsection{有限代数域扩张的微分模}
\label{b3:subsec:1404:01}

先看一个代数元。设 \(K\) 为域，\(\alpha\) 为某个扩域中对 \(K\) 代数的元素，\(L=K(\alpha)\)，\(f\in K[T]\) 为 \(\alpha\) 的首一不可约多项式。由于 \(L=K[T]/(f)\)，\cref{b3:cor:differential-presentation}给出
\begin{equation}\label{b3:eq:simple-field-differentials}
\Omega_{L/K}\cong L\,\mathrm{d}\alpha/\bigl(f'(\alpha)\mathrm{d}\alpha\bigr).
\end{equation}
若 \(f\) 可分，\(f'(\alpha)\ne0\) 在域中为单位，故微分模为零。若 \(f\) 不可分，其不可约性迫使 \(f'=0\)，于是微分模为一维。一般有限扩张不一定由一个元素生成；不能把这个单生成元计算直接当作一般情形的证明。

\begin{lemma}\label{b3:lem:exponent-one-differentials}
设 \(K\subseteq L\) 为特征 \(p>0\) 的有限域扩张，令 \(F=K L^p\) 为由 \(K\) 及所有 \(p\) 次幂生成的子域。存在 \(\alpha_1,\ldots,\alpha_r\in L\)，使
\[
[L:F]=p^r,\qquad
\Omega_{L/K}\cong\bigoplus_{i=1}^r L\,\mathrm{d}\alpha_i.
\]
其中 \(r=0\) 恰好表示 \(L=K L^p\)。
\end{lemma}
\begin{proof}
每个 \(K\) 导子都零化 \(L^p\)，再由四则运算法则零化 \(F\)，所以 \(K\) 导子与 \(F\) 导子相同，微分模也相同。逐次选取尚未落入已生成子域的 \(\alpha_i\)，直至生成 \(L/F\)。记前一步生成的子域为 \(F_{i-1}=F(\alpha_1,\ldots,\alpha_{i-1})\)。因为 \(\alpha_i^p\in F\)，其在 \(F_{i-1}\) 上的最小多项式整除 \(T^p-\alpha_i^p\)。若次数为 \(j<p\)，则在代数闭包中必为
\[
(T-\alpha_i)^j=T^j-j\alpha_iT^{j-1}+\cdots,
\qquad 1\le j<p.
\]
系数 \(-j\alpha_i\) 属于 \(F_{i-1}\)，而 \(j\) 在特征 \(p\) 的域中非零，故 \(\alpha_i\in F_{i-1}\)，与选取方式矛盾。因此最小多项式恰为 \(T^p-\alpha_i^p\)，每步次数为 \(p\)，而有限次数保证过程终止。

于是 \(\prod_i\alpha_i^{e_i}\)、\(0\le e_i<p\)，组成 \(F\) 基。自然满射
\[
F[X_1,\ldots,X_r]/(X_1^p-\alpha_1^p,\ldots,X_r^p-\alpha_r^p)
\longrightarrow L
\]
两侧维数都是 \(p^r\)，故为同构。关系的微分全部为零，\cref{b3:cor:differential-presentation}遂给出所列自由基。
\end{proof}

\begin{theorem}\label{b3:thm:finite-field-differentials-separable}
设 \(L/K\) 为有限域扩张。则
\[
\Omega_{L/K}=0\quad\Longleftrightarrow\quad L/K\text{ 可分}.
\]
若 \(L/K\) 为特征 \(p>0\) 的非平凡有限纯不可分扩张，则 \(\Omega_{L/K}\ne0\)。
\end{theorem}
\begin{proof}
若扩张可分，对任意 \(\alpha\in L\)，其最小多项式 \(f\) 满足 \(f'(\alpha)\ne0\)。微分等式 \(f'(\alpha)\mathrm{d}\alpha=0\) 因而给出 \(\mathrm{d}\alpha=0\)。这些元素生成微分模，所以该模为零。

先证非平凡纯不可分情形。设 \(L/K\) 的指数不超过 \(e\)，即 \(L^{p^e}\subseteq K\)；有限生成性保证存在同一个 \(e\)。若 \(L=K L^p\)，反复取 \(p\) 次幂并代回便得
\(L=K L^{p^j}\) 对所有 \(j\ge1\) 成立，特别地 \(L=K\)，矛盾。因此 \(K L^p\subsetneq L\)，前一引理给出至少一个非零微分。

最后设 \(L/K\) 不可分。令 \(E\) 为 \(L\) 中对 \(K\) 可分的元素组成的子域；可分代数扩张对复合及子扩张封闭，保证这些元素对四则运算封闭。\(E/K\) 有限可分。对任意 \(\alpha\in L\)，其不可约多项式在特征 \(p\) 下唯一写成 \(g(T^{p^a})\)，其中 \(g'\ne0\)。多项式 \(g\) 仍不可约，否则将其分解代入 \(T^{p^a}\) 就得到原多项式的分解；故 \(g\) 可分，\(\alpha^{p^a}\) 对 \(K\) 可分并属于 \(E\)。于是 \(L/E\) 纯不可分；由于原扩张不可分，它非平凡。复合微分正合列给出
\[
L\otimes_E\Omega_{E/K}\longrightarrow\Omega_{L/K}
\longrightarrow\Omega_{L/E}\longrightarrow0.
\]
左端由已证的可分情形为零，右端由刚证的纯不可分情形非零，所以中间项非零。特征零没有不可分代数扩张，不需此分情形。
\end{proof}

\subsection{纯不可分扩张的微分计算}
\label{b3:subsec:1404:05}

微分维数记录独立的一阶方向，却未必记录纯不可分扩张的全部幂次深度。设 \(p\) 为素数，\(s\) 为超越元，\(k=\mathbb F_p(s)\)、\(L=k(\alpha)\)，其中 \(\alpha^{p^e}=s\)、\(e\ge1\) 为整数。则
\[
\Omega_{L/k}=L\,\mathrm{d}\alpha.
\]
最小多项式 \(T^{p^e}-s\) 的不可约性可用 \(s\) 在 \(\mathbb F_p[s]\) 中的 Eisenstein 判据检验。扩张次数为 \(p^e\)，微分维数却一直为一，因此微分维数不是不可分次数的对数。它测量的是独立的一阶方向，未记录每个方向的全部幂次深度。

\ProseParagraphBreak

多方向的例子可由两个代数独立的元 \(s,t\) 构造。设 \(p\) 为素数，取
\[
k=\mathbb F_p(s,t),\qquad
L=k(\alpha,\beta),\qquad \alpha^p=s,\quad\beta^p=t.
\]
有 \([L:k]=p^2\) 及
\(\Omega_{L/k}=L\,\mathrm{d}\alpha\oplus L\,\mathrm{d}\beta\)。这里可把 \(L\) 识别为有理函数域 \(\mathbb F_p(\alpha,\beta)\)：不同模 \(p\) 的指数类在多项式恒等式中不会相消，所以 \(\alpha^i\beta^j\)、\(0\le i,j<p\)，对 \(k=\mathbb F_p(\alpha^p,\beta^p)\) 线性无关，并显然生成。对应的两个关系微分都为零，给出所列基。

\ProseParagraphBreak
在这一例子中，原来的 \(L\) 是域，扩到 \(L\) 本身后却有
\[
L\otimes_kL\cong L[\varepsilon,\eta]/(\varepsilon^p,\eta^p).
\]
把右边变量分别送到第二因子的 \(\alpha,\beta\) 减去第一因子的相应元素，便得到满射；两侧 \(L\) 维数都是 \(p^2\)，所以为同构。与\cref{b3:exm:inseparable-tangent}一样，这个计算把不可分性表现为基变换后出现的一阶方向和幂零结构。

% !TeX root = ../../Book3.tex
% 纲要标题：### 14.5 形式光滑、非分歧与平展
% 纲要位置：工作记录/计划纲要.md，第 2405 行。
% 纲要细目（写作范围）：
% * 平方零提升的存在性、至多唯一性与存在唯一性
% * 六种无穷小性质的条件对照；非分歧采用有限型，光滑和平展采用有限展示
% * 微分模消失与形式非分歧的等价
% * 标准平展的 Newton 修正与 Jacobian 子式可逆时的光滑提升
% * 形式及普通性质的任意基变换；有限域扩张平展与可分性的等价
% * 一般有限展示形式平展的等价刻画及光滑蕴含平坦在第23章证明
% ---

\section{形式光滑、非分歧与平展}
\label{b3:sec:1405}

导子比较的是两个同态之间的一阶差别；给定商环中的一个同态，能否把它提升到扩张环中，却是另一个问题。有限域扩张的微分计算还显示，环本身是域也不保证它相对于基域具有良好的提升性质。本节将存在性与唯一性分别表述，再用导数可逆的方程构造提升，并解释有限性和平坦性在光滑、非分歧及平展条件中的作用。

\subsection{平方零提升与基本概念}
\label{b3:subsec:1405:01}
\label{b3:subsec:1404:02}

平方零扩张已经把导子解释为同态的一阶变化。现在还要问：商环中的一个同态，是否来自扩张环中的同态？若 \(B=A[X_1,\ldots,X_n]/(f_1,\ldots,f_s)\)，任取变量像的代表后，各 \(f_j\) 的值都落在被商去的理想中。导子控制不同修正之间的差，但关系中的误差未必能够消掉。因此，提升的存在性和唯一性是两个需要分别考虑的问题。

\begin{definition}\label{b3:def:formal-infinitesimal-properties}
设 \(A\to B\) 为交换环同态。对任意交换 \(A\) 代数 \(C\)、满足 \(J^2=0\) 的理想 \(J\subseteq C\)，以及 \(A\) 代数同态 \(u:B\to C/J\)，考虑使下式复合等于 \(u\) 的 \(A\) 代数同态 \(\widetilde u\)：
\[
B\xrightarrow{\widetilde u}C\longrightarrow C/J.
\]
若每组这样的数据都有提升，则称 \(A\to B\) \term[xingshi-guanghua]{形式光滑}[formally smooth]；若每组数据至多有一个提升，则称其\term[xingshi-feifenqi]{形式非分歧}[formally unramified]；若每组数据恰有一个提升，则称其\term[xingshi-pingzhan]{形式平展}[formally étale]。
\end{definition}

“至多一个”允许没有提升，“至少一个”允许有多个，必须分开判断。只用平方零理想已足以检验任意幂零理想：若 \(J^N=0\)，可沿 \(C/J\)、\(C/J^2\)、\ldots、\(C/J^N=C\) 逐次提升，因为 \(J^i/J^{i+1}\) 的平方为零；存在性、唯一性分别逐次传递。

\ProseParagraphBreak
例如设 \(k\) 为域，取
\[
B=k[z]/(z^2),\qquad C=k[\varepsilon]/(\varepsilon^3),
\qquad J=(\varepsilon^2).
\]
在 \(C/J\) 中，\(z\mapsto\varepsilon\) 保持关系，因而给出 \(k\) 代数同态。若它能够提升，则 \(z\) 的像必为 \(\varepsilon+a\varepsilon^2\)，其中 \(a\in k\)。但在 \(C\) 中
\[
(\varepsilon+a\varepsilon^2)^2=\varepsilon^2\ne0,
\]
无论怎样选 \(a\)，误差都不能消去。因此 \(B\) 在 \(k\) 上不形式光滑。若去掉关系，改用 \(k[z]\)，同样的近似同态却有全部 \(z\mapsto\varepsilon+a\varepsilon^2\) 这些提升；自由变量使提升存在，也留下不唯一的变化方向。

\begin{definition}\label{b3:def:finite-presentation-infinitesimal-properties}
设 \(A\to B\) 为交换环同态。若存在有限个 \(b_1,\ldots,b_n\in B\)，使 \(B=A[b_1,\ldots,b_n]\)，则称 \(B\) 为有限型 \(A\) 代数。若 \(B\) 还能写成
\(A[X_1,\ldots,X_n]/(f_1,\ldots,f_s)\)，其中 \(n,s\) 为非负整数，则称 \(B\) 为有限展示 \(A\) 代数（也称有限呈示 \(A\) 代数）。对同态 \(A\to B\)，分别作如下定义：若 \(B\) 有限展示且该同态形式光滑，则称其\term[guanghua]{光滑}[smooth]；若 \(B\) 有限型且 \(\Omega_{B/A}=0\)，则称其\term[feifenqi]{非分歧}[unramified]；若 \(B\) 有限展示、作为 \(A\) 模平坦且 \(\Omega_{B/A}=0\)，则称其\term[pingzhan]{平展}[étale]。
\end{definition}

非分歧只要求有限型；光滑和平展的定义则要求有限展示。由下文的\cref{b3:thm:formally-unramified-differentials}，非分歧也等价于有限型且形式非分歧。\cref{b3:thm:etale-equivalent-characterizations}将证明有限展示且形式平展等价于平展，\cref{b3:thm:smooth-flat}将证明光滑蕴含平坦。六种性质的条件列于\cref{b3:tab:infinitesimal-properties}。

\begin{center}
\begin{minipage}{\textwidth}
\makeatletter
\def\@captype{table}
\makeatother
\centering
\caption[无穷小提升性质的条件]{无穷小提升性质的条件。固定交换环同态 \(A\to B\)；提升栏对任意交换 \(A\) 代数 \(C\)、平方零理想 \(J\subseteq C\) 及 \(A\) 代数同态 \(B\to C/J\) 而言。平展行的提升刻画及光滑的平坦性由第23章的相应定理给出。}
\label{b3:tab:infinitesimal-properties}
\begin{tabularx}{\textwidth}{@{}>{\raggedright\arraybackslash}p{0.17\textwidth}>{\raggedright\arraybackslash}p{0.18\textwidth}>{\raggedright\arraybackslash}p{0.27\textwidth}>{\raggedright\arraybackslash}X@{}}
\toprule
性质 & 平方零提升 & 有限性条件 & 平坦性的地位\\
\midrule
形式光滑 & 至少一个 & 不要求有限性 & 定义不要求\\
形式非分歧 & 至多一个 & 不要求有限性 & 定义不要求\\
形式平展 & 恰好一个 & 不要求有限性 & 定义不要求\\
光滑 & 至少一个 & 有限展示 & 可由定理推出\\
非分歧 & 至多一个 & 有限型 & 不要求\\
平展 & 恰好一个 & 有限展示 & 定义要求\\
\bottomrule
\end{tabularx}
\end{minipage}
\end{center}

\ProseParagraphBreak
例如多项式代数 \(A[X_1,\ldots,X_n]\) 光滑：把每个 \(u(X_i)\) 任意提升到 \(C\)，即可唯一按这些已选的像延拓为代数同态。若 \(n>0\)，选择本身通常不唯一，所以不能据此称形式非分歧。局部化 \(A_f\) 则形式平展：\(f\) 在 \(C/J\) 中的像为单位时，它在 \(C\) 中也是单位；若 \(ab=1+j\)、\(j^2=0\)，则 \(b(1-j)\) 就是 \(a\) 的逆。局部化的泛性质因而给出唯一提升。再用 \(A_f=A[T]/(fT-1)\)、局部化平坦性及 \(\Omega_{A_f/A}=0\)，也可直接核验其平展性。

\subsection{微分模消失与提升唯一性}
\label{b3:subsec:1405:02}
\label{b3:subsec:1404:03}

\begin{theorem}\label{b3:thm:formally-unramified-differentials}
设 \(A\to B\) 为任意交换环同态，不加有限性条件。该映射形式非分歧，当且仅当 \(\Omega_{B/A}=0\)。
\end{theorem}
\begin{proof}
设 \(J^2=0\)，且 \(v,w:B\to C\) 为同一个 \(u:B\to C/J\) 的两个提升。差 \(D=v-w\) 取值于 \(J\)。由于 \(J^2=0\)，\(u\) 使 \(J\) 成为一个 \(B\) 模：取 \(u(b)\) 的任意提升乘 \(J\)，结果不依赖选择。展开
\[
v(b)v(c)-w(b)w(c)
=w(b)D(c)+w(c)D(b)+D(b)D(c)
\]
可见最后一项为零，故 \(D\) 是 \(A\) 导子。如果 \(\Omega_{B/A}=0\)，微分模泛性质使 \(D=0\)，两个提升相同。

反向在平方零扩张 \(C=B\oplus\Omega_{B/A}\) 中比较
\(b\mapsto(b,0)\) 与 \(b\mapsto(b,\mathrm{d}b)\)。\cref{b3:prop:derivation-square-zero}说明两者都是恒等同态 \(B\to B\) 的提升。形式非分歧使它们相同，故所有 \(\mathrm{d}b=0\)。这些元素生成 \(\Omega_{B/A}\)，所以该模为零。
\end{proof}

这项等价也解释了微分模的作用：它的消失控制两个已存在的提升之间的差；提升的存在性与平坦性仍须分别检验。

\begin{example}\label{b3:exm:unramified-not-etale}
设 \(k\) 为域，取商映射 \(A=k[t]\to B=k=A/(t)\)。任意 \(A\) 导子 \(B\to M\) 都为零，因为 \(B\) 的每个元素都来自基环，因此 \(\Omega_{B/A}=0\)，且该映射有限展示、形式非分歧。但它不平坦：单射 \(A\xrightarrow{t}A\) 张量到 \(B\) 后变成非零模 \(B\) 上的零映射。

\ProseParagraphBreak
它也不是形式光滑。把 \(C=k[t]/(t^2)\) 看成 \(A\) 代数，令 \(J=(t)/(t^2)\)。恒等映射 \(B=k\to C/J=k\) 不可能提升为 \(A\) 代数映射 \(B\to C\)，因为 \(t\) 在 \(B\) 中为零，在 \(C\) 中却非零。故它是非分歧而非平展的具体例子。
\end{example}

\subsection{标准平展与 Jacobian 光滑代数}
\label{b3:subsec:1405:03}
\label{b3:subsec:1404:04}

微分判据说明非分歧控制提升的唯一性。要得到提升，还需消去关系中的误差；标准平展代数中，导数可逆使这一步成为一次线性修正。

\begin{proposition}\label{b3:prop:standard-etale-lift}
设 \(A\) 为交换环，\(f\in A[T]\) 为正次数首一多项式，令
\(B=\bigl(A[T]/(f)\bigr)_g\)，其中 \(g\) 是商环中的一个元素，\(t\) 为 \(T\) 的像。若 \(f'(t)\) 在 \(B\) 中可逆，则 \(A\to B\) 平展，并且形式平展。
\end{proposition}
\begin{proof}
商环由首一除法具有有限 \(A\) 基 \(1,t,\ldots,t^{\deg f-1}\)，因而平坦；再局部化仍对 \(A\) 平坦，由\cref{b3:prop:flat-base-change-transitivity}成立。添加一个生成元及关系 \(gU-1\) 表明 \(B\) 有限展示，而
\(\Omega_{B/A}=B\,\mathrm{d}t/\bigl(f'(t)\mathrm{d}t\bigr)=0\)。所以映射平展。

直接检查提升。设 \(J^2=0\)，给定 \(u:B\to C/J\)。取 \(x\in C\) 提升 \(u(t)\)，则 \(f(x)\in J\)，且 \(f'(x)\) 为单位，因为它模 \(J\) 后为单位。作一次 Newton 修正\footnote{牛顿（Isaac Newton，1643-1727），英国数学家及物理学家，创立微积分并提出经典力学定律。}，令
\[
x'=x-f'(x)^{-1}f(x).
\]
平方零条件使 Taylor 展开\footnote{泰勒（Brook Taylor，1685-1731），英国数学家，以函数的级数展开研究微积分。}恰好为
\(f(x')=f(x)+f'(x)(x'-x)=0\)。选取代表 \(g\) 的多项式 \(G\in A[T]\)。因为 \(G(x')\) 模 \(J\) 为单位，它本身为单位；而 \(f(x')=0\) 保证这个值不依赖代表的选择。于是 \(T\mapsto x'\) 唯一延拓为 \(B\to C\)。若另一根 \(x''\) 给出同一模 \(J\) 的像，则
\(0=f(x'')-f(x')=f'(x')(x''-x')\)，故 \(x''=x'\)。由生成元及局部化泛性质，整个提升唯一。
\end{proof}

这种一元计算通常称为标准平展代数的基本情形。条件必须在实际局部化的环内检查。例如
\(A=k[s]\)、\(B=A[t]/(t^2-s)\) 总是秩二自由，但微分模为 \(\bigl(B/(2t)\bigr)\,\mathrm{d}t\)。若 \(\operatorname{char}k\ne2\)，倒置 \(t\) 后导数为单位，得到平展映射；在 \(t=0\) 处导数不能相除。若 \(\operatorname{char}k=2\)，导数处处为零，非零的局部化也不满足这个非分歧条件。

光滑的类似计算允许自由方向：只要选出足够多的变量消去关系误差，其余变量的像仍可自由选择。

\begin{proposition}\label{b3:prop:jacobian-smooth-lifting}
设 \(A\) 为交换环，\(1\le r\le n\) 为整数，\(f_1,\ldots,f_r\in A[X_1,\ldots,X_n]\)，\(g\) 为商环 \(A[X_1,\ldots,X_n]/(f_1,\ldots,f_r)\) 的元素，记
\[
B=\bigl(A[X_1,\ldots,X_n]/(f_1,\ldots,f_r)\bigr)_g.
\]
若 Jacobian 矩阵 \(\bigl(\partial f_i/\partial X_j\bigr)\) 的某个 \(r\times r\) 子式在 \(B\) 中可逆，则 \(A\to B\) 光滑。
\end{proposition}
\begin{proof}
重新排列变量，可设前 \(r\) 列的子式可逆。给定交换 \(A\) 代数 \(C\)、平方零理想 \(J\subseteq C\) 及 \(A\) 代数同态 \(u:B\to C/J\)，先任取所有变量像的提升 \(x_j\in C\)。误差向量 \(\bigl(f_i(x)\bigr)_{i=1}^r\) 的各分量落入 \(J\)。所选子式在 \(C/J\) 中为单位，因而在 \(C\) 中仍为单位。保持后 \(n-r\) 个变量不变，用可逆的 \(r\times r\) 子矩阵解线性方程
\[
\left(\dfrac{\partial f_i}{\partial X_j}(x)\right)_{1\le i,j\le r}
\begin{pmatrix}h_1\\ \vdots\\h_r\end{pmatrix}
=-\begin{pmatrix}f_1(x)\\ \vdots\\f_r(x)\end{pmatrix}.
\]
所得 \(h_j\) 都属于 \(J\)。平方零条件给出准确的一阶等式
\[
\begin{aligned}
&f_i(x_1+h_1,\ldots,x_r+h_r,x_{r+1},\ldots,x_n)\\
&\qquad=f_i(x)+\sum_{j=1}^r
  \dfrac{\partial f_i}{\partial X_j}(x)h_j=0.
\end{aligned}
\]
因为 \(J^2=0\)，展开中含有至少两个 \(h_j\) 的项全部为零；最后一个等式正是所解的线性方程。\(g\) 的值模 \(J\) 后仍为单位，故修正后的值也可逆。于是修正后的变量像满足全部关系，并允许 \(g\) 的逆元，从而给出提升 \(B\to C\)。这证明 \(A\to B\) 形式光滑；添加一个生成元及关系 \(gU-1\) 可知 \(B\) 有限展示，故该同态光滑。
\end{proof}

\subsection{基变换与有限域扩张}
\label{b3:subsec:1405:04}

平方零提升的条件也能直接随基环扩张传递。设 \(A\to B\) 为交换环同态，\(A'\) 为交换 \(A\) 代数，记 \(B'=A'\otimes_AB\)。对任意 \(A'\) 代数 \(C\) 及平方零理想 \(J\)，映射 \(B'\to C/J\) 等价于一个 \(A\) 代数映射 \(B\to C/J\)。它的每个提升 \(B\to C\)，与固定的结构映射 \(A'\to C\) 在 \(A\) 上相同，故由张量积的泛性质唯一合成 \(B'\to C\)；反过来，限制到 \(B\) 还原该提升。因此两边的提升集合一一对应，存在性和至多唯一性分别保持。有限型、有限展示及平坦性也保持基变换，所以光滑、非分歧和平展都保持任意基变换。

\begin{corollary}\label{b3:cor:finite-field-extension-etale}
设 \(L/K\) 为有限域扩张。则 \(K\to L\) 非分歧，当且仅当它平展，也当且仅当 \(L/K\) 可分；此时 \(K\to L\) 还形式平展。
\end{corollary}
\begin{proof}
向量空间自由，故 \(L\) 作为 \(K\) 模平坦。有限扩张也是有限展示 \(K\) 代数：选有限组代数生成元，得到多项式环到 \(L\) 的满射，关系理想由 Hilbert 基定理有限生成。于是非分歧和平展都等价于 \(\Omega_{L/K}=0\)，再用\cref{b3:thm:finite-field-differentials-separable}得到可分性判据。若扩张可分，本原元定理给出 \(L=K(\alpha)\)；其最小多项式的导数在 \(L\) 中可逆，\cref{b3:prop:standard-etale-lift}便给出形式平展性。
\end{proof}

可分扩张的提升性质也可对照松村英之的《Commutative Ring Theory》\cite[\S25, Theorem 25.3, pp. 195--196]{Matsumura1989CommutativeRingTheory}。纯不可分扩张则可能使提升连存在性都失去。设 \(p\) 为素数，\(k=\mathbb F_p(s)\)，其中 \(s\) 为超越元；取 \(L=k[u]/(u^p-s)\)。置 \(f=T^p-s\)、\(C=k[T]/(f^2)\)、\(J=(f)/(f^2)\)。有 \(J^2=0\)、\(C/J=L\)。若恒等同态 \(L\to C/J\) 有提升，\(u\) 的像必为 \(T+h\)、\(h\in J\)。但
\[
(T+h)^p-s=T^p-s+h^p=f\ne0\quad\text{在 }C\text{ 中},
\]
因为 \(p\ge2\) 使 \(h^p=0\)，而 \(f\notin(f^2)\)。这与 \(u^p=s\) 矛盾，所以 \(k\to L\) 不形式光滑。


\begin{chaptersummary}
导子对乘积的一阶变化满足 Leibniz 法则\footnote{莱布尼茨（Gottfried Wilhelm Leibniz，1646-1716），德国数学家及哲学家，创立微积分，并研究符号演算。}，微分模把这些变化统一为一个泛性质。多项式环的微分自由，添加关系后则以关系的偏导数形成展示矩阵。余法正合列与复合微分正合列都只保证右正合，局部化和基变换公式也不替代其他性质所需的条件。

\ProseParagraphBreak
在有理点处，微分模的剩余纤维是极大理想模平方，切空间是它的对偶，也就是 Jacobian 矩阵的核。一般剩余域还会贡献自身的相对微分，不可分闭点表明两种余切对象可能不同。对有限域扩张，微分为零恰好等价于可分；纯不可分扩张中的非零微分可用指数为一的子扩张明确计算。

\ProseParagraphBreak
微分模 \(\Omega_{B/A}\) 消失，恰好表示每个平方零提升问题至多有一个提升；形式光滑则要求每个问题至少有一个提升。非分歧在前者之外要求有限型，光滑在后者之外要求有限展示。平展把有限展示、平坦性与微分消失合在一起；有限展示且形式平展与这些条件的等价将在\chapref{b3:ch:23}证明。标准平展代数的一次 Newton 修正、Jacobian 子式可逆时的线性修正，以及纯不可分扩张的提升障碍，分别展示了这些条件怎样作用于实际方程。
\end{chaptersummary}

\BookExercises

\begin{exercise}\label{b3:ex:14-derivation-constants}
设 \(k\) 为域、\(D:k[x]\to k[x]\) 为形式求导。求 \(\ker D\)，分别讨论特征零与特征 \(p>0\)。再求 \(k[x,y]\) 上同时被两个偏导子零化的多项式。
\end{exercise}
\begin{solution}
写 \(f=\sum_na_nx^n\)，则 \(Df=\sum_{n\ge1}na_nx^{n-1}\)。特征零时核为 \(k\)，特征 \(p\) 时只有指数被 \(p\) 整除的项可以保留，故核为 \(k[x^p]\)。两个变量的情形逐项比较系数，得到特征零时为 \(k\)，特征 \(p\) 时为 \(k[x^p,y^p]\)。这里系数可来自不完美域，不能把 \(k[x^p]\) 无条件写成 \(k[x]\) 中所有 \(p\) 次幂的集合。
\end{solution}

\begin{exercise}\label{b3:ex:14-euler-derivation}
设 \(k\) 为域，\(n\ge1\) 为整数。在 \(P=k[x_1,\ldots,x_n]\) 上令 \(E=\sum_i x_i\partial/\partial x_i\)。证明每个 \(d\) 次齐次多项式满足 \(E(f)=df\)。说明正特征时 \(E(f)=0\) 为什么不能推出 \(f\) 为常数。
\end{exercise}
\begin{solution}
对 \(x^\alpha\)，有 \(E(x^\alpha)=(\sum_i\alpha_i)x^\alpha\)。同次齐次项都有 \(\sum_i\alpha_i=d\)，线性相加得到公式。在特征 \(p\) 时，任意正次数为 \(p\) 倍数的齐次多项式都被零化，例如非常数的 \(x_1^p\)。当 \(n\ge2\) 时，\(x_1^{p-1}x_2\) 也被零化，且它并不是 \(p\) 次幂。
\end{solution}

\begin{exercise}\label{b3:ex:14-lifts-affine-space}
设 \(k\) 为域，\(B=k[x,y]/(xy)\)、\(C=k[\varepsilon]/(\varepsilon^3)\)、\(J=(\varepsilon^2)\subset C\)。对 \(a,b\in k\)，令 \(u:B\to C/J\) 将 \(x,y\) 分别送到 \(a\varepsilon,b\varepsilon\)。判定哪些 \(a,b\) 使 \(u\) 有提升，并求出全部提升。计算以 \(u\) 赋予 \(J\) 模结构后的 \(\operatorname{Der}_k(B,J)\)，说明为什么这个导子模的非零性不能保证提升存在。
\end{exercise}
\begin{solution}
两个生成元像的任意提升都写成
\[
x\longmapsto a\varepsilon+c\varepsilon^2,\qquad
y\longmapsto b\varepsilon+d\varepsilon^2,
\quad c,d\in k.
\]
它们的乘积为 \(ab\varepsilon^2\)，不受 \(c,d\) 影响。因此 \(ab=0\) 时每一对 \(c,d\) 都给出提升；\(ab\ne0\) 时没有提升。

\(x,y\) 对 \(J=k\varepsilon^2\) 的作用都为零，所以关系的微分 \(yD(x)+xD(y)=0\) 自动成立。两个导子值可在 \(J\) 中独立选择，得到 \(\operatorname{Der}_k(B,J)\cong J\oplus J\cong k^2\)，无论 \(a,b\) 如何都是如此。有初始提升时，这两个参数描述全部修正；没有初始提升时，它们却不能消去 \(ab\varepsilon^2\) 这个关系误差。
\end{solution}

\begin{exercise}\label{b3:ex:14-truncated-differentials}
设 \(k\) 为域、\(n\ge2\)、\(B=k[t]/(t^n)\)。计算 \(\Omega_{B/k}\) 的 \(k\) 维数，并计算余法映射 \((t^n)/(t^{2n})\to B\,\mathrm{d}t\) 的核。
\end{exercise}
\begin{solution}
微分展示为 \(\Omega_{B/k}=B\,\mathrm{d}t/(nt^{n-1}\cdot\mathrm{d}t)\)。若 \(\operatorname{char}k=0\)，或 \(\operatorname{char}k=p>0\) 且 \(p\nmid n\)，则 \(n\) 在 \(k\) 中非零，维数为 \(n-1\)；若 \(\operatorname{char}k=p>0\) 且 \(p\mid n\)，则维数为 \(n\)。余法模以 \(t^n\) 的类为一个 \(B\) 自由基，映射等同于乘 \(nt^{n-1}\)。在前一种情形中，其核由 \(t\) 生成；在后一种情形中，整个余法模都是核。两种情形中左箭头都不单射。
\end{solution}

\begin{exercise}\label{b3:ex:14-laurent-differentials}
设 \(k\) 为域，\(B=k[x^{\pm1},y^{\pm1}]\)。证明 \(\mathrm{d}x/x,\mathrm{d}y/y\) 为 \(\Omega_{B/k}\) 的自由基，计算任意整数 \(a,b\) 对应的 \(\mathrm{d}(x^ay^b)\)。
\end{exercise}
\begin{solution}
多项式微分局部化后，\(\mathrm{d}x,\mathrm{d}y\) 仍为自由基；分别乘单位 \(x^{-1},y^{-1}\) 得到所述新基。由 \(\mathrm{d}(x^{-1})=-x^{-2}\cdot\mathrm{d}x\) 及整数幂的归纳，
\[
\mathrm{d}(x^ay^b)=x^ay^b\left(a\dfrac{\mathrm{d}x}{x}+b\dfrac{\mathrm{d}y}{y}\right).
\]
公式中的整数解释为其在 \(k\) 中的像，因此正特征时仍有效。
\end{solution}

\begin{exercise}\label{b3:ex:14-tensor-product-differentials}
设 \(B,C\) 为交换 \(A\) 代数，\(D=B\otimes_A C\)。证明
\[
\Omega_{D/A}\cong
(D\otimes_B\Omega_{B/A})\oplus(D\otimes_C\Omega_{C/A}).
\]
\end{exercise}
\begin{solution}
取任意 \(D\) 模 \(M\)。一个 \(A\) 导子 \(\delta:D\to M\) 限制到 \(B\otimes1\)、\(1\otimes C\)，给出两个导子 \(\delta_B,\delta_C\)。它被这两者决定，因为
\[
\delta(b\otimes c)=(1\otimes c)\delta_B(b)+(b\otimes1)\delta_C(c).
\]
反过来，此公式保持 \(A\) 平衡关系，且乘法展开验证 Leibniz 法则，所以每一对导子都唯一延拓。右侧直和代表这两个导子的自由选择，因而具有 \(\Omega_{D/A}\) 的泛性质，得到同构。
\end{solution}

\begin{exercise}\label{b3:ex:14-power-map-differentials}
设 \(k\) 为域，\(m\ge2\)，环同态 \(k[u]\to k[t]\) 由 \(u\mapsto t^m\) 给定。计算 \(\Omega_{k[t]/k[u]}\)，并判断在何处局部化后微分模为零。
\end{exercise}
\begin{solution}
把 \(k[t]\) 写成 \(k[u][T]/(T^m-u)\)，得到
\(\Omega_{k[t]/k[u]}=k[t]\,\mathrm{d}t/(m t^{m-1}\cdot\mathrm{d}t)\)。若 \(\operatorname{char}k=0\)，或 \(\operatorname{char}k=p>0\) 且 \(p\nmid m\)，则微分模在素理想 \(\mathfrak p\) 处为零当且仅当 \(t\notin\mathfrak p\)：此时系数为单位；若 \(t\in\mathfrak p\)，系数落入局部极大理想，余模非零。若 \(\operatorname{char}k=p>0\) 且 \(p\mid m\)，则微分模自由秩一，在每个素理想处均非零。
\end{solution}

\begin{exercise}\label{b3:ex:14-cusp-torsion}
设 \(k\) 为域且 \(\operatorname{char}k\ne2,3\)，\(B=k[x,y]/(y^2-x^3)\)。正文已证明 \(\omega=2x\cdot\mathrm{d}y-3y\cdot\mathrm{d}x\) 是 \(\Omega_{B/k}\) 中的非零挠元。求它的完整零化理想 \(\operatorname{Ann}_B(\omega)\)，并求循环子模 \(B\omega\) 的 \(k\) 维数。
\end{exercise}
\begin{solution}
由微分展示，\(h\omega=0\) 当且仅当存在 \(q\in B\) 使
\[
(-3hy,2hx)=q(-3x^2,2y).
\]
将 \(B\) 识别为 \(k[t^2,t^3]\)，在分式域中比较第二坐标得 \(q=h/t\)；此值也满足第一坐标。因此 \(h\) 属于零化理想，恰好是 \(h/t\in B\)。\(B\) 中的多项式不含一次项，除以 \(t\) 后仍属于 \(B\)，当且仅当 \(h\) 的常数项与二次项都为零。这些多项式组成理想 \((t^3,t^4)=(y,x^2)\)，故
\[
\operatorname{Ann}_B(\omega)=(y,x^2),\qquad
B\omega\cong B/(y,x^2)\cong k[x]/(x^2).
\]
所以 \(B\omega\) 的 \(k\) 维数为二，基可取 \(\omega,x\omega\)。
\end{solution}

\begin{exercise}\label{b3:ex:14-parabola-characteristic-two}
设 \(k\) 为特征二的域，\(B=k[x,y]/(y-x^2)\)。计算每个 \(k\) 有理点的切空间，并解释 \(\partial(x^2)/\partial x=0\) 为什么不使该曲线具有二维切空间。
\end{exercise}
\begin{solution}
点为 \((a,a^2)\)，Jacobian 行为 \((0,1)\)，所以切空间是 \(\bigl\{(v,0):v\in k\bigr\}\)，维数一。方程对 \(y\) 的偏导仍为单位，给出非零约束。事实上 \(B\cong k[x]\)，微分关系只是 \(\mathrm{d}y=0\)，余下 \(\mathrm{d}x\) 自由；不能只检查某一个偏导数。
\end{solution}

\begin{exercise}\label{b3:ex:14-three-planes-tangent}
设 \(k\) 为域，\(B=k[x,y,z]/(xyz)\)。求每个 \(k\) 有理点的切空间维数，按恰有一个坐标为零和至少两个坐标为零分类。
\end{exercise}
\begin{solution}
Jacobian 行为 \((yz,xz,xy)\)。若恰有一个坐标为零，另外两个非零，三项中恰有一项非零，矩阵秩一，切空间维数二；若至少两个坐标为零，整行全零，切空间维数三。满足方程的点至少有一个零坐标，这两类穷尽所有有理点。计算不需要对 \(k\) 的特征加限制。
\end{solution}

\begin{exercise}\label{b3:ex:14-fat-point-tangent}
设 \(k\) 为域，\(r\ge2\)，\(B=k[x,y]/(x,y)^r\)。求唯一有理点处的切空间维数及 \(\dim_kB\)。说明前者能否恢复后者。
\end{exercise}
\begin{solution}
唯一极大理想为 \(\mathfrak m=(x,y)\)，而 \(\mathfrak m/\mathfrak m^2\) 的基为 \(\bar x,\bar y\)，所以切空间维数二。\(B\) 的基由全次数小于 \(r\) 的单项式组成，故
\(\dim_kB=\sum_{d=0}^{r-1}(d+1)=r(r+1)/2\)。不同 \(r\) 给出相同切空间维数而不同的代数维数，说明切空间只保存一阶信息。
\end{solution}

\begin{exercise}\label{b3:ex:14-cusp-tangent-map}
设 \(k\) 为域，考虑尖点参数化 \(\nu:t\mapsto(t^2,t^3)\)。直接求 \(t=a\) 处诱导的切空间映射。证明它在 \(a=0\) 处为零，在 \(a\ne0\) 处单射。
\end{exercise}
\begin{solution}
把 \(a+v\varepsilon\) 代入得到
\((a^2+2av\varepsilon,a^3+3a^2v\varepsilon)\)，故切映射为
\(v\mapsto(2av,3a^2v)\)。它满足尖点的切方程，因为
\(-3a^4(2av)+2a^3(3a^2v)=0\)。在零点两坐标均为零；若 \(a\ne0\)，整数二和三不可能同时在域中为零，至少一项系数非零，故映射单射。参数化在点集上单射不保证其每一点的切映射单射。
\end{solution}

\begin{exercise}\label{b3:ex:14-characteristic-two-circle}
设 \(k\) 为域。比较 \(B=k[x,y]/(x^2+y^2-1)\) 在特征不为二和特征二时的微分模及有理点切空间。在特征二时求该环的约化商。
\end{exercise}
\begin{solution}
总有 \(\Omega_{B/k}=(B\,\mathrm{d}x\oplus B\,\mathrm{d}y)/(2x\cdot\mathrm{d}x+2y\cdot\mathrm{d}y)\)。特征不为二时，在每个有理点 \((a,b)\) 上，\(a,b\) 不同时为零，故切空间维数一。特征二时，关系多项式为 \((x+y-1)^2\)，微分关系为零，所以微分模自由秩二，所有有理点的切空间维数二。约化商是 \(k[x,y]/(x+y-1)\cong k[x]\)，它在对应点的切空间维数一。这说明改变特征可能把一个方程变成重方程。
\end{solution}

\begin{exercise}\label{b3:ex:14-separable-derivation-extension}
设 \(k\subseteq K\)、\(L=K(\alpha)\)，其中 \(\alpha\) 对 \(K\) 可分，最小多项式为 \(f(T)=\sum_i a_iT^i\)。给定 \(k\) 导子 \(D:K\to L\)，证明它唯一延拓为 \(\widetilde D:L\to L\)，并写出 \(\widetilde D(\alpha)\)。
\end{exercise}
\begin{solution}
求导 \(f(\alpha)=0\) 迫使
\[
\widetilde D(\alpha)=-f'(\alpha)^{-1}\sum_iD(a_i)\alpha^i.
\]
右边有意义，因为可分性使 \(f'(\alpha)\ne0\)。在 \(K[T]\) 上，把系数按 \(D\) 求导、把 \(T\) 的导数指定为该值，得到取值于 \(L\) 的导子，它零化 \(f\)。对任意 \(hf\)，Leibniz 法则使导数在代入 \(\alpha\) 后也为零，所以它下降到 \(K[T]/(f)=L\)。其在系数及生成元上的值均已被迫确定，故唯一。
\end{solution}

\begin{exercise}\label{b3:ex:14-purely-inseparable-tower}
设 \(p\) 为素数，\(s\) 为超越元，令 \(k=\mathbb F_p(s)\)，\(L=k(\alpha)\)、\(E=k(\beta)\)，其中 \(\alpha^p=s\)、\(\beta^p=\alpha\)。计算复合微分正合列的三个非零候选项和两条箭头。
\end{exercise}
\begin{solution}
有 \(\Omega_{L/k}=L\,\mathrm{d}\alpha\)、\(\Omega_{E/k}=E\,\mathrm{d}\beta\)、\(\Omega_{E/L}=E\,\mathrm{d}\beta\)。第一箭头将 \(1\otimes \mathrm{d}\alpha\) 送到 \(\mathrm{d}(\beta^p)=0\)，故为零；第二箭头将 \(\mathrm{d}\beta\) 送到 \(\mathrm{d}\beta\)，为同构。因此该列为
\(E\xrightarrow{0}E\xrightarrow{1}E\to0\)。即使各域扩张次数都有限，首箭头仍可完全丢失一个微分方向。
\end{solution}

\begin{exercise}\label{b3:ex:14-perfect-closure}
设 \(p\) 为素数，\(t_1,t_2,\ldots\) 在 \(\mathbb F_p\) 上代数独立，令
\[
L=\mathbb F_p(t_1,t_2,\ldots),\qquad
K=\mathbb F_p(t_1^p,t_2^p,\ldots).
\]
证明 \(L/K\) 是无限次、指数为一的纯不可分扩张，并证明 \(\Omega_{L/K}\) 是以 \(\mathrm{d}t_i\)、\(i\ge1\)，为基的自由 \(L\) 模。比较此例与正文无限根塔的微分消失现象，说明无限次纯不可分扩张的微分模不必为零。
\end{exercise}
\begin{solution}
每个有理函数只涉及有限多个 \(t_i\)，其 \(p\) 次幂属于 \(K\)，所以扩张纯不可分、指数至多一；\(t_1\notin K\)，故指数恰为一。对任意 \(r\ge1\)，单项式 \(t_1^{e_1}\cdots t_r^{e_r}\)、\(0\le e_i<p\)，对 \(K\) 线性无关：一个假定关系只涉及有限多个变量和分母，清分母后，各项在前 \(r\) 个变量上的指数模 \(p\) 属于不同余数类，不会相消。它们也生成 \(K(t_1,\ldots,t_r)\)，故这个子扩张的次数为 \(p^r\)，从而 \([L:K]\) 无限。

所有 \(\mathrm{d}t_i\) 生成 \(\Omega_{L/K}\)，因为每个元素的微分都能按分式法则写成有限组合。对固定 \(j\)，形式偏导 \(\partial/\partial t_j\) 唯一延拓到 \(L\)，并零化 \(K\)。将它对应的线性映射作用于任何有限关系 \(\sum_i c_i\cdot\mathrm{d}t_i=0\)，便得 \(c_j=0\)。因此
\[
\Omega_{L/K}\cong\bigoplus_{i\ge1}L\,\mathrm{d}t_i.
\]
正文根塔使每个元素都在下一层成为 \(p\) 次幂；本题则不断加入新的独立方向，却不再对它们取 \(p\) 次根，两种无限扩张保留的微分因而不同。
\end{solution}

\begin{exercise}\label{b3:ex:14-arithmetic-etale}
令 \(B=\mathbb Z[t]/(t^2-t-1)\)。证明 \(B[1/5]\) 在 \(\mathbb Z[1/5]\) 上平展，并证明特征五纤维不是非分歧代数。
\end{exercise}
\begin{solution}
首一除法给出基 \(1,t\)，故自由平坦。微分关系为 \((2t-1)\cdot\mathrm{d}t=0\)，而在 \(B\) 中
\((2t-1)^2=4t^2-4t+1=5\)。倒置五后 \(2t-1\) 为单位，所以微分模为零，有限展示条件也成立。模五后，\(t^2-t-1=(t-3)^2\)，纤维为 \(\mathbb F_5[\varepsilon]/(\varepsilon^2)\)，其微分模为 \(\bigl(\mathbb F_5[\varepsilon]/(\varepsilon)\bigr)\mathrm{d}\varepsilon\ne0\)，故不非分歧。
\end{solution}

\begin{exercise}\label{b3:ex:14-square-zero-newton}
在 \(C=\mathbb Q[\varepsilon]/(\varepsilon^2)\) 中，求满足 \(x^2=1+\varepsilon\) 且 \(x\equiv1\pmod\varepsilon\) 的全部元素。若把 \(\mathbb Q\) 改成特征二的域，会发生什么？
\end{exercise}
\begin{solution}
写 \(x=1+a\varepsilon\)，则 \(x^2=1+2a\varepsilon\)，唯一解为 \(1+\varepsilon/2\)。这正是对误差 \(-\varepsilon\) 用导数二作一次修正。特征二时，所有这样的平方都等于一，因此没有解；不是有多个解，而是给定提升根本不存在。
\end{solution}

\begin{exercise}\label{b3:ex:14-free-differentials-not-smooth}
设 \(k\) 为特征 \(p>0\) 的域，\(A=k[s]\)、\(B=A[u]/(u^p-s)\)。证明 \(B\) 是自由秩 \(p\) 的 \(A\) 模，\(\Omega_{B/A}\) 是自由秩一的 \(B\) 模，但 \(A\to B\) 不形式光滑。说明这比“微分模自由但不平坦”的反例多检验了哪个条件。
\end{exercise}
\begin{solution}
首一除法给出 \(A\) 基 \(1,u,\ldots,u^{p-1}\)，所以 \(B\) 对 \(A\) 平坦。相对微分中 \(\mathrm{d}s=0\)，而 \(\mathrm{d}(u^p)=0\)，故关系不增加微分约束，得到 \(\Omega_{B/A}=B\,\mathrm{d}u\)。

置 \(f=T^p-s\)、\(C=A[T]/(f^2)\)、\(J=(f)/(f^2)\)，则 \(J^2=0\)、\(C/J=B\)。若恒等同态 \(B\to C/J\) 有 \(A\) 代数提升，\(u\) 的像必为 \(T+h\)，其中 \(h\in J\)。但
\[
(T+h)^p-s=f+h^p=f\ne0\quad\text{于}\;C,
\]
因为 \(p\ge2\) 使 \(h^p=0\)，而 \(f\notin(f^2)\)。故提升不存在。这里有限展示、平坦性与微分模自由都已成立，仍不能替代形式光滑所要求的提升存在性。
\end{solution}

\begin{exercise}\label{b3:ex:14-smooth-hyperbola-lifting}
设 \(A\) 为交换环，\(B=A[x,y]/(xy-1)\)。直接证明 \(A\to B\) 形式光滑；并在一个明确的平方零扩张中构造同一映射的两个不同提升，说明它一般不形式非分歧。
\end{exercise}
\begin{solution}
给定 \(B\to C/J\)，其中 \(J^2=0\)，\(x\) 的像为单位。任取它的提升 \(u\in C\)，仍为单位，令 \(x\mapsto u\)、\(y\mapsto u^{-1}\)，即得提升。为展示不唯一，取 \(A=k\) 为域，\(C=k[\varepsilon]/(\varepsilon^2)\)，基点为 \(x,y\mapsto1\)。一组提升为 \((1,1)\)，另一组为 \((1+\varepsilon,1-\varepsilon)\)，乘积均为一，且两组不同。
\end{solution}

\begin{exercise}\label{b3:ex:14-formal-base-change}
设交换环同态 \(A\to B\) 光滑、非分歧或平展。对任意交换 \(A\) 代数 \(A'\)，记 \(B'=A'\otimes_AB\)。利用\secref{b3:sec:1405}的形式提升性质，分别验证 \(A'\to B'\) 的同一性质；写出有限型、有限展示及平坦性在验证中的作用。
\end{exercise}
\begin{solution}
若 \(B=A[b_1,\ldots,b_n]\)，则 \(B'\) 由 \(1\otimes b_i\) 作为 \(A'\) 代数生成，故有限型保持。若 \(B=A[X_1,\ldots,X_n]/(f_1,\ldots,f_s)\)，则
\[
B'\cong A'[X_1,\ldots,X_n]/(f'_1,\ldots,f'_s),
\]
其中 \(f'_i\) 由结构同态 \(A\to A'\) 变换系数，故有限展示保持。光滑情形再用正文已证的形式光滑基变换性质。非分歧情形用有限型及
\(\Omega_{B'/A'}\cong A'\otimes_A\Omega_{B/A}=0\)。平展情形还需有限展示以及\cref{b3:prop:flat-base-change-transitivity}给出的平坦性保持。三种性质所需的条件都分别保留下来。
\end{solution}

\begin{exercise}\label{b3:ex:14-formal-composition}
设交换环同态 \(A\to B\to D\) 中的两个映射均形式光滑，或均形式非分歧，或均形式平展。证明复合映射分别具有同一性质。
\end{exercise}
\begin{solution}
给定 \(A\) 代数平方零提升问题 \(D\to C/J\)，先限制到 \(B\)。形式光滑时先提升 \(B\to C\)，把 \(C\) 视为这个结构下的 \(B\) 代数，再利用 \(B\to D\) 的形式光滑性提升整个 \(D\) 映射。形式非分歧时，两个 \(D\) 提升限制到 \(B\) 必相同，因此是同一 \(B\) 代数结构下的两个提升，再由第二映射的唯一性相等。两种性质同时成立便得到形式平展的结论。
\end{solution}

